\hypertarget{uxfcbersicht}{%
\subsection{Übersicht}\label{uxfcbersicht}}

\begin{itemize}
\item
  \textbf{ADR-005:} Project Workspace Architecture
\item
  \textbf{ADR-009:} Textual Source Processing Boundary
\item
  \textbf{ADR-010:} Project Source Registry Architecture
\item
  \textbf{ADR-011:} Semantic Information Unit and Ontology Boundary
\item
  \textbf{ADR-012:} Processing State and Artifact Organization
\item
  \textbf{ADR-013:} Preliminary Coverage and Potential Model Support Assessment
\item
  \textbf{ADR-014:} Project Dashboard Architecture and Evidence Navigation
\item
  \textbf{ADR-015:} Project-bound Agentic Ingestion Integration Architecture
\item
  \textbf{ADR-016:} Human Review Workspace and Approved Input Promotion Architecture
\item
  \textbf{ADR-017:} Simple-by-Default Interaction and Progressive Disclosure
\item
  \textbf{ADR-018:} Model Candidate Layer and Structural Comparability
\item
  \textbf{ADR-019:} Internal Engineering Model Assembly Architecture
\item
  \textbf{ADR-020:} Hybrid Target Projection and Coverage Architecture
\item
  \textbf{ADR-021:} SYSIDE-Compatible SysML v2 Generation Architecture
\item
  \textbf{ADR-022:} SysML v2 Validation Layer Architecture
\item
  \textbf{ADR-023:} Final Model Review and Output Publication Architecture
\item
  \textbf{ADR-024:} Guided Engineering Workflow and UX Projection Architecture
\item
  \textbf{ADR-025:} Semantic Proposal Consolidation and Persona-Aware Consensus
\item
  \textbf{ADR-026:} Source-Anchored Multi-Persona Interpretation and Cross-Unit Semantic Synthesis
\item
  \textbf{ADR-027:} Source-Grounded Evidence Detection and Persona Interpretation Architecture
\item
  \textbf{ADR-028a:} Model Derivation Mode and Review Escalation Architecture
\item
  \textbf{ADR-028b:} Context-Preserving Canonical Engineering Subject Discovery
\item
  \textbf{ADR-029:} Human-Reviewed Model Placement Before Model Assembly
\item
  \textbf{ADR-030:} Semantic Interpretation and Controlled Classification Alignment
\item
  \textbf{ADR-031:} Semantic Field Consistency Alignment
\item
  \textbf{ADR-032:} Project-Level Multi-Source Reconciliation and Controlled Engineering Evolution
\item
  \textbf{ADR-033:} Concern-Centric Project Reconciliation and Coherent Model Handoff
\item
  \textbf{ADR-034:} Source Provenance Does Not Constrain Concern Grouping
\item
  \textbf{ADR-035:} Safe-Demo Review Compression and Final Review Routing
\end{itemize}

\textbackslash newpage

\hypertarget{adr-005--project-workspace-architecture}{%
\subsection{ADR-005 --- Project Workspace Architecture}\label{adr-005--project-workspace-architecture}}

Project Workspace Architecture

Status

Accepted

Date

2026-07-21

Context

Phase P introduces a persistent Project Workspace around the completed Phase F ingestion pipeline.

P1 established the versioned framework template \texttt{TURING\_RFLP\_FRAMEWORK} version \texttt{1.0.0}.

P2 requires a deterministic architecture for:

\begin{itemize}
\tightlist
\item
  project identity
\item
  project metadata
\item
  workspace discovery
\item
  persistence and reopening
\item
  project isolation
\item
  future source, information-unit, run and coverage storage
\item
  validation and recovery behavior
\end{itemize}

The Project Workspace contains software-operational project metadata. It shall not redefine engineering knowledge from the authoritative CATIA SysML v2 model.

External repositories, including the Apollo 11 repository, remain non-normative references and do not define this architecture.

Decision

\hypertarget{project-identity}{%
\subsubsection{Project Identity}\label{project-identity}}

Every project receives an immutable technical identifier named \texttt{project\_id}.

The \texttt{project\_id}:

\begin{itemize}
\tightlist
\item
  is stored as a JSON string
\item
  consists of exactly six decimal digits
\item
  matches \texttt{\^{}{[}0-9{]}\{6\}\$}
\item
  permits leading zeros
\item
  is unique within the configured workspace
\item
  remains unchanged for the lifetime of the project
\item
  is used as the project directory name
\end{itemize}

Valid examples include:

\begin{itemize}
\tightlist
\item
  \texttt{000042}
\item
  \texttt{318604}
\item
  \texttt{999999}
\end{itemize}

The identifier is generated with:

\begin{Shaded}
\begin{Highlighting}[]
\SpecialStringTok{f"}\SpecialCharTok{\{}\NormalTok{secrets}\SpecialCharTok{.}\NormalTok{randbelow(}\DecValTok{1\_000\_000}\NormalTok{)}\SpecialCharTok{:06d\}}\SpecialStringTok{"}
\end{Highlighting}
\end{Shaded}

Before an identifier is accepted, the workspace checks whether the corresponding project directory already exists. A collision causes generation to be repeated.

Failure to produce an available identifier raises \texttt{ProjectIdGenerationError}.

The six-digit identifier is a technical identifier only. It is not a password, authentication credential or security secret.

\hypertarget{project-display-name}{%
\subsubsection{Project Display Name}\label{project-display-name}}

Every project has a human-readable \texttt{display\_name}.

The display name:

\begin{itemize}
\tightlist
\item
  is required
\item
  contains between 1 and 120 characters after trimming
\item
  is mutable
\item
  must be unique across the complete configured workspace
\end{itemize}

Display-name uniqueness is checked during:

\begin{itemize}
\tightlist
\item
  project creation
\item
  project rename
\item
  workspace validation
\end{itemize}

Uniqueness comparison uses one deterministic normalization function:

\begin{enumerate}
\def\labelenumi{\arabic{enumi}.}
\tightlist
\item
  trim leading and trailing whitespace
\item
  collapse consecutive whitespace
\item
  apply Unicode normalization
\item
  apply case-insensitive \texttt{casefold()} comparison
\end{enumerate}

The original validated display text is stored in the manifest. The normalized comparison value is derived and is not persisted.

Two projects must therefore not use names that differ only through casing, equivalent Unicode representation or whitespace formatting.

The immutable \texttt{project\_id} remains the authoritative identity even when the display name changes.

\hypertarget{workspace-root}{%
\subsubsection{Workspace Root}\label{workspace-root}}

The product workspace root is:

\texttt{data/projects}

The root is injectable so that tests and future integrations can use an isolated location.

A persisted project is located at:

\texttt{data/projects/\textless{}project\_id\textgreater{}/}

The P2 project directory initially contains exactly one required file:

\texttt{data/projects/\textless{}project\_id\textgreater{}/project\_manifest.json}

P2 does not pre-create empty directories.

Later phases create their own storage only when required:

\begin{itemize}
\tightlist
\item
  P3 owns \texttt{sources/}
\item
  P4 owns \texttt{information\_units/}
\item
  P5 owns \texttt{runs/}
\item
  P6 owns \texttt{coverage/}
\end{itemize}

This avoids placeholder directories whose contracts have not yet been implemented.

\hypertarget{project-manifest-contract}{%
\subsubsection{Project Manifest Contract}\label{project-manifest-contract}}

The manifest schema version is:

\texttt{1.0.0}

The manifest contains exactly these top-level fields:

\begin{Shaded}
\begin{Highlighting}[]
\FunctionTok{\{}
  \DataTypeTok{"schema\_version"}\FunctionTok{:} \StringTok{"1.0.0"}\FunctionTok{,}
  \DataTypeTok{"project\_id"}\FunctionTok{:} \StringTok{"318604"}\FunctionTok{,}
  \DataTypeTok{"display\_name"}\FunctionTok{:} \StringTok{"Example Project"}\FunctionTok{,}
  \DataTypeTok{"description"}\FunctionTok{:} \StringTok{""}\FunctionTok{,}
  \DataTypeTok{"framework\_template"}\FunctionTok{:} \FunctionTok{\{}
    \DataTypeTok{"template\_id"}\FunctionTok{:} \StringTok{"TURING\_RFLP\_FRAMEWORK"}\FunctionTok{,}
    \DataTypeTok{"template\_version"}\FunctionTok{:} \StringTok{"1.0.0"}
  \FunctionTok{\},}
  \DataTypeTok{"created\_at"}\FunctionTok{:} \StringTok{"2026{-}07{-}21T12:00:00Z"}\FunctionTok{,}
  \DataTypeTok{"updated\_at"}\FunctionTok{:} \StringTok{"2026{-}07{-}21T12:00:00Z"}
\FunctionTok{\}}
\end{Highlighting}
\end{Shaded}

The required rules are:

\begin{itemize}
\tightlist
\item
  \texttt{schema\_version} is exactly \texttt{1.0.0}
\item
  \texttt{project\_id} matches the six-digit identifier contract
\item
  \texttt{project\_id} matches the containing directory name
\item
  \texttt{display\_name} follows the display-name contract
\item
  \texttt{description} is required
\item
  \texttt{description} may be empty
\item
  \texttt{description} contains at most 2000 characters
\item
  \texttt{framework\_template.template\_id} is \texttt{TURING\_RFLP\_FRAMEWORK}
\item
  \texttt{framework\_template.template\_version} is \texttt{1.0.0}
\item
  \texttt{created\_at} is a UTC ISO-8601 timestamp ending in \texttt{Z}
\item
  \texttt{updated\_at} is a UTC ISO-8601 timestamp ending in \texttt{Z}
\item
  \texttt{created\_at} remains unchanged after project creation
\item
  \texttt{updated\_at} is refreshed when mutable metadata changes
\item
  unknown fields are rejected
\end{itemize}

The framework reference is pinned. A project is never silently migrated to a new framework-template version.

A future framework upgrade requires an explicit migration decision and implementation.

The manifest does not contain:

\begin{itemize}
\tightlist
\item
  source lists
\item
  ingestion runs
\item
  reports
\item
  information units
\item
  coverage results
\item
  approval decisions
\item
  generated-model data
\item
  derived counters
\item
  duplicated filesystem paths
\end{itemize}

Those concerns belong to their responsible Phase P or later components.

\hypertarget{project-discovery}{%
\subsubsection{Project Discovery}\label{project-discovery}}

The workspace has no central \texttt{projects.json} index.

Projects are discovered by scanning:

\texttt{data/projects/*/project\_manifest.json}

The directory structure and each validated manifest are the source of truth for project discovery.

Avoiding a central index prevents the index and individual manifests from diverging.

A scan returns a \texttt{WorkspaceScanResult} containing:

\begin{itemize}
\tightlist
\item
  \texttt{valid\_projects}
\item
  \texttt{workspace\_issues}
\end{itemize}

A malformed project must not make valid projects unavailable.

The scanner reports issues for conditions including:

\begin{itemize}
\tightlist
\item
  visible directories whose names are not valid six-digit project identifiers
\item
  missing manifests
\item
  invalid JSON
\item
  missing required fields
\item
  unknown fields
\item
  unsupported schema versions
\item
  invalid framework references
\item
  mismatch between \texttt{project\_id} and directory name
\item
  duplicate normalized display names
\item
  unsafe paths
\item
  symbolic-link project directories
\end{itemize}

When two manifests use the same normalized display name, both projects are reported as conflicting. They remain identifiable through their immutable project IDs so that one can be renamed explicitly.

Hidden filesystem metadata such as \texttt{.DS\_Store} is ignored.

Temporary creation directories matching \texttt{.create-\textless{}project\_id\textgreater{}.tmp} are ignored during discovery but are never deleted automatically.

The workspace does not silently skip validation errors and does not automatically repair or delete malformed data.

\hypertarget{safe-persistence}{%
\subsubsection{Safe Persistence}\label{safe-persistence}}

Project creation uses a temporary sibling directory:

\texttt{data/projects/.create-\textless{}project\_id\textgreater{}.tmp/}

Creation follows this sequence:

\begin{enumerate}
\def\labelenumi{\arabic{enumi}.}
\tightlist
\item
  generate an available project ID
\item
  validate the requested metadata
\item
  create the temporary sibling directory
\item
  write the manifest into the temporary directory
\item
  read and validate the written manifest
\item
  atomically rename the temporary directory to the final project directory
\end{enumerate}

Manifest updates use:

\texttt{project\_manifest.json.tmp}

An update follows this sequence:

\begin{enumerate}
\def\labelenumi{\arabic{enumi}.}
\tightlist
\item
  load and validate the existing manifest
\item
  apply the permitted metadata change
\item
  preserve \texttt{project\_id}, \texttt{created\_at} and the pinned framework reference
\item
  write the updated manifest to the temporary file
\item
  validate the temporary manifest
\item
  atomically replace the existing manifest with \texttt{os.replace}
\end{enumerate}

Interrupted or invalid writes must not replace the last valid manifest.

Temporary data is retained for explicit diagnosis and is not automatically deleted.

\hypertarget{project-isolation-and-path-safety}{%
\subsubsection{Project Isolation and Path Safety}\label{project-isolation-and-path-safety}}

All project operations are resolved below the configured workspace root.

The implementation rejects:

\begin{itemize}
\tightlist
\item
  absolute external project paths
\item
  path traversal
\item
  paths escaping the configured workspace root
\item
  symbolic-link project directories
\item
  project IDs that do not match the six-digit identifier contract
\end{itemize}

The implementation does not follow symbolic links while discovering or loading projects.

Sources, runs, information units and artifacts introduced in later steps must reference the immutable \texttt{project\_id} and remain within that project's directory.

Cross-project data mixing is prohibited.

\hypertarget{module-boundaries}{%
\subsubsection{Module Boundaries}\label{module-boundaries}}

P2 introduces:

\begin{verbatim}
modules/project_workspace/
├── __init__.py
├── errors.py
├── identifiers.py
├── manifest.py
├── types.py
└── workspace.py
\end{verbatim}

\hypertarget{errorspy}{%
\paragraph{\texorpdfstring{\texttt{errors.py}}{errors.py}}\label{errorspy}}

Defines the Project Workspace exception hierarchy:

\begin{itemize}
\tightlist
\item
  \texttt{ProjectWorkspaceError}
\item
  \texttt{ProjectManifestError}
\item
  \texttt{ProjectNotFoundError}
\item
  \texttt{DuplicateProjectNameError}
\item
  \texttt{ProjectIdGenerationError}
\item
  \texttt{UnsafeProjectPathError}
\end{itemize}

\hypertarget{identifierspy}{%
\paragraph{\texorpdfstring{\texttt{identifiers.py}}{identifiers.py}}\label{identifierspy}}

Owns:

\begin{itemize}
\tightlist
\item
  six-digit project-ID generation
\item
  project-ID validation
\item
  deterministic display-name normalization
\end{itemize}

\hypertarget{typespy}{%
\paragraph{\texorpdfstring{\texttt{types.py}}{types.py}}\label{typespy}}

Defines immutable data types:

\begin{itemize}
\tightlist
\item
  \texttt{FrameworkTemplateReference}
\item
  \texttt{ProjectManifest}
\item
  \texttt{WorkspaceIssue}
\item
  \texttt{WorkspaceScanResult}
\end{itemize}

\hypertarget{manifestpy}{%
\paragraph{\texorpdfstring{\texttt{manifest.py}}{manifest.py}}\label{manifestpy}}

Owns:

\begin{itemize}
\tightlist
\item
  manifest parsing
\item
  manifest serialization
\item
  manifest validation
\item
  schema-version validation
\item
  timestamp validation
\item
  framework-reference validation
\item
  rejection of unknown fields
\end{itemize}

It does not access project directories or scan the workspace.

\hypertarget{workspacepy}{%
\paragraph{\texorpdfstring{\texttt{workspace.py}}{workspace.py}}\label{workspacepy}}

Owns:

\begin{itemize}
\tightlist
\item
  workspace scanning
\item
  project creation
\item
  project loading
\item
  project metadata updates
\item
  project-ID collision checks
\item
  display-name uniqueness checks
\item
  path and symbolic-link safety
\item
  atomic filesystem operations
\item
  workspace issue collection
\end{itemize}

The public \texttt{ProjectWorkspace} interface is:

\begin{Shaded}
\begin{Highlighting}[]
\NormalTok{create\_project(display\_name, description}\OperatorTok{=}\StringTok{""}\NormalTok{)}
\NormalTok{load\_project(project\_id)}
\NormalTok{update\_project(}
\NormalTok{    project\_id,}
    \OperatorTok{*}\NormalTok{,}
\NormalTok{    display\_name}\OperatorTok{=}\VariableTok{None}\NormalTok{,}
\NormalTok{    description}\OperatorTok{=}\VariableTok{None}\NormalTok{,}
\NormalTok{)}
\NormalTok{scan\_projects()}
\end{Highlighting}
\end{Shaded}

The workspace root, ID generator and clock are injectable for deterministic tests.

\hypertarget{dependency-direction}{%
\subsubsection{Dependency Direction}\label{dependency-direction}}

The dependency direction is:

\begin{verbatim}
workspace
    -> manifest
    -> identifiers
    -> types
    -> errors
    -> framework template validation
\end{verbatim}

\texttt{modules/framework} does not depend on \texttt{modules/project\_workspace}.

No reverse dependency from the framework-template module to the Project Workspace is permitted.

\hypertarget{p2-scope-boundary}{%
\subsubsection{P2 Scope Boundary}\label{p2-scope-boundary}}

P2 includes:

\begin{itemize}
\tightlist
\item
  project identity
\item
  project metadata
\item
  manifest validation
\item
  safe persistence
\item
  project discovery
\item
  project reopening
\item
  project isolation foundations
\end{itemize}

P2 does not include:

\begin{itemize}
\tightlist
\item
  project deletion or archiving
\item
  source registration
\item
  source upload
\item
  information-unit persistence
\item
  ingestion-run organization
\item
  coverage calculation
\item
  approval decisions
\item
  model-candidate creation
\item
  SysML v2 generation
\item
  Project Dashboard integration
\end{itemize}

Those capabilities remain assigned to later roadmap steps.

Consequences

Positive consequences:

\begin{itemize}
\tightlist
\item
  Projects have compact, immutable and easily comparable identifiers.
\item
  Display names remain human-readable without becoming technical identities.
\item
  Duplicate or confusing display names are prevented.
\item
  Project metadata can be validated independently of the UI.
\item
  Projects can be reopened without maintaining a second index.
\item
  Atomic writes reduce the risk of corrupted manifests.
\item
  Invalid projects do not prevent access to valid projects.
\item
  Future Phase P components receive a stable isolation boundary.
\item
  Framework-template versions remain explicit and reproducible.
\end{itemize}

Trade-offs:

\begin{itemize}
\tightlist
\item
  The identifier space is limited to one million values.
\item
  Identifier generation requires collision detection.
\item
  Workspace scanning performs filesystem validation instead of reading a central index.
\item
  Strict schema validation rejects manifests with unexpected fields.
\item
  Incomplete temporary data requires explicit operator review.
\item
  Framework-template upgrades require explicit migration.
\end{itemize}

Alternatives Considered

UUID project identifiers were rejected because they are unnecessarily difficult to compare manually for the expected workspace size.

Mutable display names as project identities were rejected because renaming would break references and equal or similar names would create ambiguity.

A central project index was rejected because it would duplicate manifest state and introduce synchronization risk.

Pre-creating all future Phase P directories was rejected because their storage contracts belong to later implementation steps.

Permissive manifest parsing and silent issue skipping were rejected because they would hide inconsistent project state.

Automatic repair or deletion of malformed and temporary data was rejected because recovery must remain explicit and auditable.

Affected Components

\begin{itemize}
\tightlist
\item
  \texttt{modules/project\_workspace/}
\item
  \texttt{tests/test\_project\_manifest.py}
\item
  \texttt{tests/test\_project\_workspace.py}
\item
  \texttt{data/projects/}
\item
  later P3--P7 components that reference \texttt{project\_id}
\end{itemize}

Supersedes

None

Related Roadmap Phase

P2 --- Project Manifest and Workspace Structure

Related Implementation

Not yet implemented.

\textbackslash newpage

\hypertarget{adr-009--textual-source-processing-boundary}{%
\subsection{ADR-009 --- Textual Source Processing Boundary}\label{adr-009--textual-source-processing-boundary}}

Textual Source Processing Boundary

Status

Accepted

Date

2026-07-22

Context

The Turing Generator shall ingest heterogeneous customer documents and derive traceable engineering information through the agentic ingestion pipeline.

The original Phase F implementation reads UTF-8 text directly from a source path. Phase P introduces registered project sources and heterogeneous, source-traceable information units.

Modern LLM APIs can accept document files such as PDF, DOCX, PPTX and spreadsheets directly. Some providers and models may additionally interpret page images, diagrams and other visual content.

Supporting arbitrary multimodal engineering sources would require architectural decisions for:

\begin{itemize}
\tightlist
\item
  optical character recognition
\item
  drawing and diagram interpretation
\item
  symbol recognition
\item
  geometric and spatial relationships
\item
  visual provenance
\item
  multimodal uncertainty
\item
  conflict resolution between textual and visual information
\item
  human validation of visually extracted statements
\end{itemize}

These concerns exceed the accepted MVP and thesis scope.

Decision

\hypertarget{supported-information-modality}{%
\subsubsection{Supported Information Modality}\label{supported-information-modality}}

The only supported source-information modality in the MVP is textual information.

The original source file does not have to be a plain-text file.

A binary container format may be used when its relevant textual content can be converted deterministically into a traceable textual representation before the LLM-based engineering extraction begins.

Examples include:

\begin{itemize}
\tightlist
\item
  a PDF containing a machine-readable text layer
\item
  a future DOCX adapter that extracts paragraphs and tables
\item
  other document containers for which a deterministic textual adapter is explicitly implemented and validated
\end{itemize}

File-container support and information-modality support are separate concerns.

Supporting a file extension does not imply support for every type of content that the file may contain.

\hypertarget{immutable-original-source}{%
\subsubsection{Immutable Original Source}\label{immutable-original-source}}

The registered original source remains immutable.

It is stored with:

\begin{itemize}
\tightlist
\item
  \texttt{project\_id}
\item
  \texttt{source\_id}
\item
  original filename
\item
  media type
\item
  byte size
\item
  SHA-256 hash
\item
  explicit source role
\item
  registration timestamp
\end{itemize}

Normalization never replaces or modifies the registered original source.

\hypertarget{deterministic-text-normalization}{%
\subsubsection{Deterministic Text Normalization}\label{deterministic-text-normalization}}

Before an LLM processes a registered source, the source must produce a normalized textual artifact through a deterministic, versioned source adapter.

The normalized artifact must remain traceable to:

\begin{itemize}
\tightlist
\item
  \texttt{project\_id}
\item
  \texttt{source\_id}
\item
  original source SHA-256
\item
  adapter identifier
\item
  adapter version
\item
  source locations used to produce the normalized text
\end{itemize}

Source-location references depend on the input format.

Examples include:

\begin{itemize}
\tightlist
\item
  text and Markdown: line ranges
\item
  PDF: page references
\item
  JSON: JSON Pointer locations
\item
  CSV: row and column references
\item
  future DOCX support: sections, paragraphs or table cells
\end{itemize}

The normalized textual artifact is a derived processing artifact. It is not an authoritative engineering source and does not replace the original file.

\hypertarget{pdf-boundary}{%
\subsubsection{PDF Boundary}\label{pdf-boundary}}

The MVP may support PDF files only when they contain a usable machine-readable text layer.

PDF normalization extracts textual content and retains page boundaries for traceability.

The MVP does not interpret:

\begin{itemize}
\tightlist
\item
  technical drawings
\item
  diagrams
\item
  photographs
\item
  graphical relationships
\item
  dimensions represented only visually
\item
  symbols represented only visually
\item
  scanned pages without a machine-readable text layer
\end{itemize}

A PDF must not be sent as unrestricted multimodal input when doing so would allow the selected model to derive engineering information from visual content outside the accepted MVP scope.

\hypertarget{explicit-failure-states}{%
\subsubsection{Explicit Failure States}\label{explicit-failure-states}}

A source that cannot produce sufficient traceable textual content must not be silently treated as an empty or irrelevant source.

Processing must end with an explicit diagnostic state such as:

\begin{itemize}
\tightlist
\item
  \texttt{unsupported\_non\_textual\_content}
\item
  \texttt{text\_extraction\_insufficient}
\item
  \texttt{unsupported\_source\_format}
\item
  \texttt{source\_normalization\_failed}
\end{itemize}

The original source remains registered and available for diagnosis.

No failed source is silently discarded or promoted.

\hypertarget{llm-processing-boundary}{%
\subsubsection{LLM Processing Boundary}\label{llm-processing-boundary}}

The LLM receives the normalized textual artifact together with explicit provenance information.

The LLM may semantically interpret, classify and derive candidate engineering information from that textual input.

The LLM must not silently invent information that is absent from the normalized source.

The LLM output remains unreviewed and non-authoritative.

\hypertarget{canonical-extraction-output}{%
\subsubsection{Canonical Extraction Output}\label{canonical-extraction-output}}

The canonical result of LLM extraction is validated structured data.

The canonical structured output shall use a versioned JSON contract for source-traceable candidate information units.

Markdown is not the canonical extraction format.

A Markdown review report is generated deterministically from the validated structured extraction data.

The raw LLM output is retained as a technical traceability artifact but is not used directly as approved engineering information.

\hypertarget{human-in-the-loop-boundary}{%
\subsubsection{Human-in-the-Loop Boundary}\label{human-in-the-loop-boundary}}

Extracted information units remain unreviewed until a human decision has been recorded.

The review workflow must allow information units to be:

\begin{itemize}
\tightlist
\item
  accepted
\item
  corrected
\item
  rejected
\item
  marked as requiring clarification
\end{itemize}

Generating or reading a Markdown report does not approve its contents.

Phase P may persist preliminary and review-oriented information, but it does not promote information into generation-ready input.

Approved Input Promotion remains assigned to Phase G.

Only explicitly human-approved structured engineering information may become input to later model-candidate and model-generation stages.

\hypertarget{provider-and-model-independence}{%
\subsubsection{Provider and Model Independence}\label{provider-and-model-independence}}

The Project Workspace does not assume that every provider, endpoint or model supports the same file types.

Provider-specific capabilities are evaluated separately from source registration.

The deterministic normalized textual artifact forms the provider-independent boundary for MVP ingestion.

Direct provider file input may be introduced later only when it preserves the accepted modality boundary, provenance requirements and review behavior.

\hypertarget{excluded-mvp-scope}{%
\subsubsection{Excluded MVP Scope}\label{excluded-mvp-scope}}

The following capabilities are explicitly excluded from the MVP:

\begin{itemize}
\tightlist
\item
  optical character recognition
\item
  scanned-document interpretation
\item
  technical-drawing interpretation
\item
  diagram interpretation
\item
  image-based engineering evidence
\item
  handwritten-content recognition
\item
  audio ingestion
\item
  video ingestion
\item
  spatial or geometric relation extraction from visual media
\item
  general multimodal engineering information extraction
\end{itemize}

Consequences

Positive consequences:

\begin{itemize}
\tightlist
\item
  The MVP scope remains technically and academically manageable.
\item
  Existing text-oriented Phase F pipeline concepts remain reusable.
\item
  Different source containers can use a common ingestion boundary.
\item
  LLM inputs remain reproducible and provider-independent.
\item
  Engineering statements can reference deterministic source locations.
\item
  Human review operates on structured and traceable information units.
\item
  Visual information cannot silently become engineering evidence.
\item
  Future multimodal extensions have a clear architectural starting point.
\end{itemize}

Trade-offs:

\begin{itemize}
\tightlist
\item
  Scanned PDFs cannot be processed without a future OCR capability.
\item
  Relevant information contained only in drawings or diagrams is not extracted.
\item
  Binary document containers require a deterministic text adapter.
\item
  Source normalization introduces an additional persisted artifact and validation step.
\item
  Some direct multimodal capabilities of selected LLMs remain intentionally unused in the MVP.
\end{itemize}

Alternatives Considered

Accepting every file type supported by the selected LLM was rejected because provider support does not establish traceability, reproducibility or a stable human-review boundary.

Sending PDFs directly as unrestricted multimodal model input was rejected for the MVP because it could introduce engineering information derived from drawings, images or spatial relationships without an accepted visual-evidence contract.

Restricting registration to plain-text file extensions was rejected because common customer documents may use binary containers while still containing deterministically extractable textual information.

Using Markdown as the canonical extracted-information store was rejected because Markdown is intended for human-readable reporting rather than strict schema validation and downstream processing.

Replacing the original file with extracted text was rejected because it would destroy source integrity and weaken traceability.

Affected Components

\begin{itemize}
\tightlist
\item
  Project Source Registry introduced in P3
\item
  deterministic source-normalization adapters
\item
  Phase F ingestion integration
\item
  framework-mapped information units introduced in P4
\item
  processing state and artifact organization introduced in P5
\item
  human review and Approved Input Promotion introduced in Phase G
\item
  future provider and model adapters
\item
  future SysML v2 representation of the Turing Generator
\end{itemize}

Model Impact

When the authoritative SysML v2 model is next updated, it shall represent the accepted textual-source-processing constraint.

The model update shall distinguish:

\begin{itemize}
\tightlist
\item
  supported textual information
\item
  text-bearing document containers
\item
  excluded non-textual and multimodal information
\item
  deterministic source normalization
\item
  human review before Approved Input Promotion
\item
  multimodal processing as a future extension
\end{itemize}

Stable model element identifiers and their relationships shall be assigned in the authoritative model. This ADR does not create or redefine those engineering model elements.

SSOT Impact

The next scheduled SSOT UPDATE after completion of Phase P shall document:

\begin{itemize}
\tightlist
\item
  the textual-only MVP information-modality constraint
\item
  deterministic text normalization
\item
  the exclusion of OCR, drawings and multimodal extraction
\item
  structured JSON as the canonical unreviewed extraction format
\item
  Markdown as the derived review representation
\item
  the required future SysML v2 model synchronization
\end{itemize}

Supersedes

None

Related Roadmap Phase

\begin{itemize}
\tightlist
\item
  P3 --- Source Registry and Mandatory Project Assignment
\item
  P4 --- Framework-mapped Heterogeneous Information Units
\item
  P5 --- Processing State and Artifact Organization
\item
  G --- Approved Input Promotion
\end{itemize}

Related Implementation

Not yet implemented.

\textbackslash newpage

\hypertarget{adr-010--project-source-registry-architecture}{%
\subsection{ADR-010 --- Project Source Registry Architecture}\label{adr-010--project-source-registry-architecture}}

Project Source Registry Architecture

Status

Accepted

Date

2026-07-22

Context

Phase P introduces project-oriented processing around the completed Phase F agentic ingestion pipeline.

P2 implemented the persistent Project Workspace and established immutable six-digit project identifiers.

P3 requires a deterministic architecture for:

\begin{itemize}
\tightlist
\item
  mandatory assignment of every new source to exactly one project
\item
  stable source identity
\item
  immutable storage of original source files
\item
  explicit source roles
\item
  duplicate detection
\item
  source discovery and integrity validation
\item
  project isolation
\item
  later traceability from information units, runs, reports and artifacts back to their originating source
\end{itemize}

ADR-009 establishes that the MVP supports textual information only.

The Source Registry does not perform text projection, semantic interpretation, ontological mapping or LLM processing. Those responsibilities begin in P4.

Decision

\hypertarget{source-identity}{%
\subsubsection{Source Identity}\label{source-identity}}

Every registered source receives an immutable technical identifier named \texttt{source\_id}.

The \texttt{source\_id}:

\begin{itemize}
\tightlist
\item
  is stored as a JSON string
\item
  matches \texttt{\^{}SRC-{[}0-9{]}\{6\}\$}
\item
  is unique within its containing project
\item
  remains unchanged for the lifetime of the source
\item
  is used as the source directory name
\end{itemize}

Valid examples include:

\begin{itemize}
\tightlist
\item
  \texttt{SRC-000001}
\item
  \texttt{SRC-004281}
\item
  \texttt{SRC-999999}
\end{itemize}

The \texttt{project\_id} and \texttt{source\_id} are separate identifiers.

The project identifier identifies the project container. The source identifier identifies one individual source within that project.

The globally unambiguous source reference is the pair:

\begin{verbatim}
<project_id>/<source_id>
\end{verbatim}

Example:

\begin{verbatim}
318604/SRC-000001
\end{verbatim}

A source must not be identified only by:

\begin{itemize}
\tightlist
\item
  original filename
\item
  stored filename
\item
  display label
\item
  content hash
\item
  source role
\end{itemize}

\hypertarget{source-id-allocation}{%
\subsubsection{Source-ID Allocation}\label{source-id-allocation}}

Source IDs are allocated sequentially within one project.

The first source ID is:

\texttt{SRC-000001}

The next source ID is greater than every existing or temporarily occupied source ID in that project.

Allocation considers:

\begin{itemize}
\tightlist
\item
  final source directories
\item
  temporary registration directories
\end{itemize}

Gaps are not reused automatically.

Source IDs remain reserved when a temporary registration directory exists.

If \texttt{SRC-999999} has already been reached, registration fails with \texttt{SourceIdExhaustedError}.

If a concurrent operation occupies the selected final directory, the registry rescans the project and attempts registration with the next available sequential identifier.

The six-digit source number is a technical identifier only. It is not a password, authentication credential or security secret.

\hypertarget{mandatory-project-assignment}{%
\subsubsection{Mandatory Project Assignment}\label{mandatory-project-assignment}}

Every source registration requires a valid \texttt{project\_id}.

The referenced project must:

\begin{itemize}
\tightlist
\item
  exist
\item
  contain a valid project manifest
\item
  use the expected Project Workspace schema
\item
  resolve safely below the configured Project Workspace root
\end{itemize}

There is no API for registering a permanent unassigned source.

There is no global source pool.

A source belongs to exactly one project.

The same original file may be registered in different projects because project storage and evidence remain isolated.

\hypertarget{source-storage-structure}{%
\subsubsection{Source Storage Structure}\label{source-storage-structure}}

The source root is created only when the first source is registered.

A registered source is stored under:

\begin{verbatim}
data/projects/<project_id>/sources/<source_id>/
\end{verbatim}

The final source directory contains:

\begin{verbatim}
source_manifest.json
content.<suffix>
\end{verbatim}

Example:

\begin{verbatim}
data/projects/318604/sources/SRC-000001/
├── source_manifest.json
└── content.pdf
\end{verbatim}

The original filename is stored only as metadata.

The original filename is not used directly as a filesystem path.

The stored filename is generated by the registry.

A safe, recognized lowercase suffix may be retained.

Examples include:

\begin{itemize}
\tightlist
\item
  \texttt{content.md}
\item
  \texttt{content.txt}
\item
  \texttt{content.json}
\item
  \texttt{content.csv}
\item
  \texttt{content.pdf}
\end{itemize}

An absent, unsupported or unsafe suffix produces:

\texttt{content.bin}

Unknown media types are represented as:

\texttt{application/octet-stream}

The registry never executes a registered source file.

\hypertarget{format-neutral-registration}{%
\subsubsection{Format-Neutral Registration}\label{format-neutral-registration}}

The Source Registry stores original bytes and technical metadata.

It does not decide whether a registered source can be processed by the selected LLM, provider or future source adapter.

Registration and processability are separate concerns.

A source may therefore be successfully registered but later receive a processing state such as:

\texttt{registered\_but\_not\_processable}

The textual-source-processing constraint from ADR-009 is evaluated before semantic ingestion begins.

A source must not be rejected solely because the current LLM does not support its format.

The registry still rejects sources that violate storage, safety or integrity rules.

\hypertarget{immutable-original-source-1}{%
\subsubsection{Immutable Original Source}\label{immutable-original-source-1}}

The source content is immutable after successful registration.

The registry records:

\begin{itemize}
\tightlist
\item
  byte size
\item
  SHA-256 hash
\end{itemize}

Both values are calculated from the bytes stored in the final source directory.

Loading or scanning a source validates its stored content against the manifest.

A mismatch is an integrity error.

The registry never silently:

\begin{itemize}
\tightlist
\item
  replaces source content
\item
  overwrites an existing source
\item
  updates a hash to match changed content
\item
  repairs a damaged source
\item
  substitutes a different file
\item
  deletes inconsistent source data
\end{itemize}

A changed document must be registered as a new source with a new \texttt{source\_id}.

\hypertarget{source-roles}{%
\subsubsection{Source Roles}\label{source-roles}}

Every source has exactly one explicit source role.

Permitted roles are:

\begin{itemize}
\tightlist
\item
  \texttt{engineering\_source}
\item
  \texttt{context\_only}
\end{itemize}

There is no default source role.

The caller must select the source role during registration.

An \texttt{engineering\_source} may later produce source-traceable engineering information units and contribute to clearly marked preliminary coverage.

A \texttt{context\_only} source may support terminology and interpretation but shall not:

\begin{itemize}
\tightlist
\item
  create engineering evidence
\item
  create framework assignments
\item
  satisfy preliminary coverage
\item
  satisfy approved readiness
\item
  contribute to model generation
\end{itemize}

Source-role behavior remains governed by the accepted Phase P rules.

\hypertarget{source-manifest-contract}{%
\subsubsection{Source Manifest Contract}\label{source-manifest-contract}}

The Source Manifest schema version is:

\texttt{1.0.0}

The manifest contains exactly these top-level fields:

\begin{Shaded}
\begin{Highlighting}[]
\FunctionTok{\{}
  \DataTypeTok{"schema\_version"}\FunctionTok{:} \StringTok{"1.0.0"}\FunctionTok{,}
  \DataTypeTok{"project\_id"}\FunctionTok{:} \StringTok{"318604"}\FunctionTok{,}
  \DataTypeTok{"source\_id"}\FunctionTok{:} \StringTok{"SRC{-}000001"}\FunctionTok{,}
  \DataTypeTok{"source\_role"}\FunctionTok{:} \StringTok{"engineering\_source"}\FunctionTok{,}
  \DataTypeTok{"original\_filename"}\FunctionTok{:} \StringTok{"System Requirements.pdf"}\FunctionTok{,}
  \DataTypeTok{"stored\_filename"}\FunctionTok{:} \StringTok{"content.pdf"}\FunctionTok{,}
  \DataTypeTok{"media\_type"}\FunctionTok{:} \StringTok{"application/pdf"}\FunctionTok{,}
  \DataTypeTok{"size\_bytes"}\FunctionTok{:} \DecValTok{1845203}\FunctionTok{,}
  \DataTypeTok{"sha256"}\FunctionTok{:} \StringTok{"full{-}lowercase{-}sha256{-}value"}\FunctionTok{,}
  \DataTypeTok{"registered\_at"}\FunctionTok{:} \StringTok{"2026{-}07{-}22T10:30:00Z"}\FunctionTok{,}
  \DataTypeTok{"updated\_at"}\FunctionTok{:} \StringTok{"2026{-}07{-}22T10:30:00Z"}
\FunctionTok{\}}
\end{Highlighting}
\end{Shaded}

Required validation rules are:

\begin{itemize}
\tightlist
\item
  \texttt{schema\_version} is exactly \texttt{1.0.0}
\item
  \texttt{project\_id} follows the Project Workspace identifier contract
\item
  \texttt{project\_id} matches the containing project directory
\item
  \texttt{source\_id} follows the Source Registry identifier contract
\item
  \texttt{source\_id} matches the containing source directory
\item
  \texttt{source\_role} is an explicitly permitted source role
\item
  \texttt{original\_filename} is present and non-empty
\item
  \texttt{original\_filename} contains no path semantics
\item
  \texttt{stored\_filename} is a registry-generated basename
\item
  \texttt{stored\_filename} contains no absolute path or traversal
\item
  \texttt{media\_type} is present and non-empty
\item
  \texttt{size\_bytes} is a positive integer
\item
  \texttt{sha256} is a lowercase 64-character hexadecimal SHA-256 value
\item
  \texttt{registered\_at} is a UTC ISO-8601 timestamp ending in \texttt{Z}
\item
  \texttt{updated\_at} is a UTC ISO-8601 timestamp ending in \texttt{Z}
\item
  \texttt{updated\_at} is not earlier than \texttt{registered\_at}
\item
  unknown fields are rejected
\end{itemize}

The manifest does not contain:

\begin{itemize}
\tightlist
\item
  normalized text
\item
  LLM output
\item
  information units
\item
  ontology mappings
\item
  framework assignments
\item
  processing runs
\item
  reports
\item
  coverage calculations
\item
  review decisions
\item
  Approved Input
\item
  model candidates
\item
  generated model content
\end{itemize}

Those concerns belong to later components.

\hypertarget{mutable-and-immutable-metadata}{%
\subsubsection{Mutable and Immutable Metadata}\label{mutable-and-immutable-metadata}}

The following values are immutable:

\begin{itemize}
\tightlist
\item
  \texttt{schema\_version}
\item
  \texttt{project\_id}
\item
  \texttt{source\_id}
\item
  original source bytes
\item
  \texttt{original\_filename}
\item
  \texttt{stored\_filename}
\item
  \texttt{media\_type}
\item
  \texttt{size\_bytes}
\item
  \texttt{sha256}
\item
  \texttt{registered\_at}
\end{itemize}

The following values may be updated under the defined role-change rules:

\begin{itemize}
\tightlist
\item
  \texttt{source\_role}
\item
  \texttt{updated\_at}
\end{itemize}

During P3, a source role may be corrected because no Project Workspace processing runs exist yet.

When P5 introduces persisted processing runs, role changes for sources with dependent runs or artifacts must be blocked until an explicit invalidation and reprocessing contract has been accepted.

Changing a role must never silently reinterpret existing downstream artifacts.

\hypertarget{duplicate-content}{%
\subsubsection{Duplicate Content}\label{duplicate-content}}

Before registration, the registry calculates the input SHA-256 hash and compares it with existing valid sources in the selected project.

If the same SHA-256 already exists in that project:

\begin{itemize}
\tightlist
\item
  registration is rejected
\item
  \texttt{DuplicateSourceContentError} is raised
\item
  the error identifies the existing \texttt{source\_id}
\end{itemize}

This prevents the same source content from being counted repeatedly as apparently independent engineering evidence.

The same SHA-256 may exist in different projects.

Equal original filenames with different content are permitted.

The filename is never used for duplicate detection.

\hypertarget{atomic-registration}{%
\subsubsection{Atomic Registration}\label{atomic-registration}}

Registration accepts an existing local regular file.

The input path must:

\begin{itemize}
\tightlist
\item
  exist
\item
  identify a regular file
\item
  not be a symbolic link
\item
  not be empty
\end{itemize}

Registration follows this sequence:

\begin{enumerate}
\def\labelenumi{\arabic{enumi}.}
\tightlist
\item
  load and validate the selected project
\item
  validate the source role
\item
  validate the input file
\item
  calculate the input size and SHA-256
\item
  scan existing project sources for duplicate content
\item
  allocate the next sequential source ID
\item
  create a temporary sibling directory
\item
  copy the source bytes blockwise into the temporary directory
\item
  calculate size and SHA-256 again while copying
\item
  verify that the source did not change during registration
\item
  create the Source Manifest
\item
  validate the manifest and temporary stored content
\item
  atomically rename the temporary directory to the final source directory
\item
  load and return the persisted Source Manifest
\end{enumerate}

The temporary directory is:

\begin{verbatim}
data/projects/<project_id>/sources/.register-<source_id>.tmp/
\end{verbatim}

Example:

\begin{verbatim}
data/projects/318604/sources/.register-SRC-000001.tmp/
\end{verbatim}

An interrupted registration must not create a final source directory.

Temporary registration directories:

\begin{itemize}
\tightlist
\item
  are ignored during normal source discovery
\item
  reserve their contained source number
\item
  are never automatically deleted
\item
  remain available for explicit diagnosis
\end{itemize}

\hypertarget{atomic-source-role-update}{%
\subsubsection{Atomic Source-Role Update}\label{atomic-source-role-update}}

A source-role update:

\begin{enumerate}
\def\labelenumi{\arabic{enumi}.}
\tightlist
\item
  loads and validates the existing source
\item
  verifies that the new role is permitted
\item
  preserves every immutable manifest value
\item
  changes only \texttt{source\_role} and \texttt{updated\_at}
\item
  writes \texttt{source\_manifest.json.tmp}
\item
  validates the temporary manifest
\item
  atomically replaces the existing manifest with \texttt{os.replace}
\end{enumerate}

An invalid update must not replace the last valid manifest.

\hypertarget{source-discovery}{%
\subsubsection{Source Discovery}\label{source-discovery}}

Sources are discovered by scanning:

\begin{verbatim}
data/projects/<project_id>/sources/*/source_manifest.json
\end{verbatim}

There is no central source index.

The project directory and validated source manifests are the source of truth.

A scan returns a \texttt{SourceScanResult} containing:

\begin{itemize}
\tightlist
\item
  \texttt{valid\_sources}
\item
  \texttt{source\_issues}
\end{itemize}

A malformed source must not make valid sources unavailable.

\hypertarget{source-scan-issues}{%
\subsubsection{Source Scan Issues}\label{source-scan-issues}}

The scanner reports issues including:

\begin{itemize}
\tightlist
\item
  visible directories whose names are not valid Source IDs
\item
  missing Source Manifests
\item
  invalid JSON
\item
  missing required fields
\item
  unknown fields
\item
  unsupported schema versions
\item
  invalid source roles
\item
  mismatch between manifest and containing \texttt{project\_id}
\item
  mismatch between manifest and containing \texttt{source\_id}
\item
  invalid stored filenames
\item
  missing stored content
\item
  unexpected additional content files
\item
  symbolic-link source directories
\item
  symbolic-link content files
\item
  invalid byte sizes
\item
  stored-size mismatches
\item
  invalid SHA-256 values
\item
  stored-content hash mismatches
\item
  duplicate content within one project
\item
  unsafe paths
\end{itemize}

Sources with blocking conflicts are excluded from \texttt{valid\_sources}.

They remain individually identifiable through \texttt{project\_id} and \texttt{source\_id} for diagnosis.

When duplicate content is discovered in existing persisted data, every conflicting source is reported and excluded from \texttt{valid\_sources}.

Hidden operating-system metadata such as \texttt{.DS\_Store} is ignored.

Temporary registration directories are ignored by normal discovery but are considered during Source-ID allocation.

No malformed or temporary data is automatically repaired or deleted.

\hypertarget{path-and-symlink-safety}{%
\subsubsection{Path and Symlink Safety}\label{path-and-symlink-safety}}

All persisted source paths are generated below the validated project directory.

The implementation rejects:

\begin{itemize}
\tightlist
\item
  absolute stored filenames
\item
  stored path traversal
\item
  paths escaping the selected project
\item
  symbolic-link project source roots
\item
  symbolic-link source directories
\item
  symbolic-link source content
\item
  Source IDs that do not satisfy the identifier contract
\item
  filenames that resolve outside their source directory
\end{itemize}

The implementation does not follow source-directory or source-content symbolic links.

The input file may originate outside the Project Workspace, but it is only read and copied into a generated safe destination.

The input file itself must not be a symbolic link.

\hypertarget{module-boundaries-1}{%
\subsubsection{Module Boundaries}\label{module-boundaries-1}}

P3 introduces:

\begin{verbatim}
modules/project_sources/
├── __init__.py
├── errors.py
├── identifiers.py
├── types.py
├── manifest.py
└── registry.py
\end{verbatim}

\hypertarget{errorspy-1}{%
\paragraph{\texorpdfstring{\texttt{errors.py}}{errors.py}}\label{errorspy-1}}

Defines the Source Registry exception hierarchy:

\begin{itemize}
\tightlist
\item
  \texttt{ProjectSourceError}
\item
  \texttt{SourceManifestError}
\item
  \texttt{SourceNotFoundError}
\item
  \texttt{DuplicateSourceContentError}
\item
  \texttt{SourceIdExhaustedError}
\item
  \texttt{UnsupportedSourceRoleError}
\item
  \texttt{UnsafeSourcePathError}
\item
  \texttt{SourceIntegrityError}
\end{itemize}

\hypertarget{identifierspy-1}{%
\paragraph{\texorpdfstring{\texttt{identifiers.py}}{identifiers.py}}\label{identifierspy-1}}

Owns:

\begin{itemize}
\tightlist
\item
  Source-ID validation
\item
  Source-ID formatting
\item
  sequential Source-ID allocation
\item
  Source-ID exhaustion detection
\end{itemize}

It does not access source content or parse manifests.

\hypertarget{typespy-1}{%
\paragraph{\texorpdfstring{\texttt{types.py}}{types.py}}\label{typespy-1}}

Defines immutable data types:

\begin{itemize}
\tightlist
\item
  \texttt{SourceManifest}
\item
  \texttt{SourceIssue}
\item
  \texttt{SourceScanResult}
\end{itemize}

\hypertarget{manifestpy-1}{%
\paragraph{\texorpdfstring{\texttt{manifest.py}}{manifest.py}}\label{manifestpy-1}}

Owns:

\begin{itemize}
\tightlist
\item
  Source Manifest parsing
\item
  Source Manifest serialization
\item
  Source Manifest validation
\item
  schema-version validation
\item
  source-role validation
\item
  timestamp validation
\item
  stored-filename validation
\item
  rejection of unknown fields
\end{itemize}

It does not scan project directories or modify source files.

\hypertarget{registrypy}{%
\paragraph{\texorpdfstring{\texttt{registry.py}}{registry.py}}\label{registrypy}}

Owns:

\begin{itemize}
\tightlist
\item
  source registration
\item
  source loading
\item
  source-content path resolution
\item
  source discovery
\item
  duplicate detection
\item
  source-integrity verification
\item
  source-role updates
\item
  safe path handling
\item
  symbolic-link rejection
\item
  atomic filesystem operations
\item
  Source Issue collection
\end{itemize}

\hypertarget{__init__py}{%
\paragraph{\texorpdfstring{\texttt{\_\_init\_\_.py}}{\_\_init\_\_.py}}\label{__init__py}}

Exposes only the intentional public Source Registry API.

\hypertarget{public-interface}{%
\subsubsection{Public Interface}\label{public-interface}}

The public service is:

\texttt{ProjectSourceRegistry}

Its initial public operations are:

\begin{Shaded}
\begin{Highlighting}[]
\NormalTok{register\_source(}
\NormalTok{    project\_id,}
\NormalTok{    source\_path,}
    \OperatorTok{*}\NormalTok{,}
\NormalTok{    source\_role,}
\NormalTok{)}

\NormalTok{load\_source(}
\NormalTok{    project\_id,}
\NormalTok{    source\_id,}
\NormalTok{)}

\NormalTok{source\_content\_path(}
\NormalTok{    project\_id,}
\NormalTok{    source\_id,}
\NormalTok{)}

\NormalTok{update\_source\_role(}
\NormalTok{    project\_id,}
\NormalTok{    source\_id,}
    \OperatorTok{*}\NormalTok{,}
\NormalTok{    source\_role,}
\NormalTok{)}

\NormalTok{scan\_sources(}
\NormalTok{    project\_id,}
\NormalTok{)}
\end{Highlighting}
\end{Shaded}

The Project Workspace root and clock are injectable for deterministic tests.

Filesystem content is processed blockwise so registration and integrity validation do not require loading complete source files into memory.

\hypertarget{dependency-direction-1}{%
\subsubsection{Dependency Direction}\label{dependency-direction-1}}

The dependency direction is:

\begin{verbatim}
project_sources
    -> project_workspace
    -> Python standard library
\end{verbatim}

\texttt{modules/project\_workspace} does not depend on \texttt{modules/project\_sources}.

The completed P2 module remains independent and stable.

\texttt{modules/project\_sources} does not depend on:

\begin{itemize}
\tightlist
\item
  ingestion agents
\item
  LLM providers
\item
  source-normalization adapters
\item
  information-unit modules
\item
  coverage modules
\item
  UI modules
\item
  model-generation modules
\end{itemize}

\hypertarget{p3-scope-boundary}{%
\subsubsection{P3 Scope Boundary}\label{p3-scope-boundary}}

P3 includes:

\begin{itemize}
\tightlist
\item
  mandatory project assignment
\item
  source identity
\item
  Source Manifest validation
\item
  immutable original-source storage
\item
  duplicate detection
\item
  explicit source roles
\item
  source-role correction before dependent processing
\item
  source discovery
\item
  source integrity validation
\item
  project isolation
\item
  safe and atomic persistence
\end{itemize}

P3 does not include:

\begin{itemize}
\tightlist
\item
  deterministic textual source projection
\item
  PDF text extraction
\item
  semantic interpretation
\item
  terminology or ontology lookup
\item
  LLM execution
\item
  information-unit extraction
\item
  framework mapping
\item
  processing-run persistence
\item
  reports
\item
  human approval
\item
  Approved Input Promotion
\item
  Project Dashboard integration
\item
  model-candidate creation
\item
  SysML v2 generation
\end{itemize}

Deterministic textual source projection and semantic candidate extraction begin in P4 under the boundary established by ADR-009.

Consequences

Positive consequences:

\begin{itemize}
\tightlist
\item
  Every source is assigned to exactly one project.
\item
  Source identity is independent of filenames.
\item
  Short Source IDs remain easy to compare manually.
\item
  Original customer files remain unchanged.
\item
  Duplicate evidence within a project is prevented.
\item
  Source integrity can be verified deterministically.
\item
  Registration remains independent of providers and LLM capabilities.
\item
  Valid sources remain usable when other source directories are malformed.
\item
  P4 receives a stable, traceable source boundary.
\item
  The completed P2 module remains independent.
\end{itemize}

Trade-offs:

\begin{itemize}
\tightlist
\item
  Sequential Source IDs require scanning existing and temporary directories.
\item
  Gaps are not automatically reused.
\item
  Temporary failed registrations require explicit operator review.
\item
  A registered source may not yet be processable.
\item
  Source-role correction becomes more constrained after processing runs exist.
\item
  Strict integrity validation exposes manual filesystem changes as errors.
\item
  There is no automatic cleanup, repair or content replacement.
\end{itemize}

Alternatives Considered

Using the project ID as the source ID was rejected because one project may contain multiple sources.

Using filenames as source identity was rejected because filenames are mutable, repeatable and unsafe as stable references.

Using globally unique Source IDs was rejected because project-scoped identity is sufficient when every external reference contains both \texttt{project\_id} and \texttt{source\_id}.

Using random Source IDs was rejected because sequential project-local IDs are easier to compare and do not require collision generation in the expected project scale.

Maintaining a central source index was rejected because it would duplicate manifest state and introduce synchronization risk.

Restricting registration to the file types currently supported by the Phase F pipeline was rejected because registration and processability are separate responsibilities.

Allowing duplicate content within one project was rejected because duplicate sources could incorrectly multiply apparent engineering evidence.

Automatically deleting incomplete registration directories was rejected because recovery and cleanup must remain explicit and auditable.

Allowing source-content replacement was rejected because it would break traceability from downstream artifacts to the original evidence.

Adding text extraction or semantic interpretation to the Source Registry was rejected because it would violate single responsibility and hide semantic decisions inside persistence logic.

Affected Components

\begin{itemize}
\tightlist
\item
  \texttt{modules/project\_sources/}
\item
  \texttt{tests/test\_source\_manifest.py}
\item
  \texttt{tests/test\_project\_source\_registry.py}
\item
  \texttt{data/projects/\textless{}project\_id\textgreater{}/sources/}
\item
  future P4 information-unit components
\item
  future P5 run and artifact components
\item
  future Project Dashboard integration
\end{itemize}

Model Impact

This ADR defines software persistence and traceability architecture.

The textual information-modality constraint and later semantic-processing boundary are defined by ADR-009 and shall be represented when the authoritative SysML v2 model is updated.

Stable model identifiers remain the responsibility of the authoritative model and are not created by this ADR.

SSOT Impact

The next scheduled SSOT UPDATE after completion of Phase P shall document:

\begin{itemize}
\tightlist
\item
  completion of P3
\item
  project-local Source IDs
\item
  mandatory project assignment
\item
  immutable original-source storage
\item
  explicit source roles
\item
  duplicate detection
\item
  Source Registry integrity behavior
\item
  the separation of source registration from processability
\item
  the dependency from P3 to the P4 textual-source-processing boundary
\end{itemize}

Supersedes

None

Related Roadmap Phase

P3 --- Source Registry and Mandatory Project Assignment

Related Decisions

\begin{itemize}
\tightlist
\item
  ADR-005 --- Project Workspace Architecture
\item
  ADR-009 --- Textual Source Processing Boundary
\end{itemize}

Related Implementation

Not yet implemented.

\textbackslash newpage

\hypertarget{adr-011--semantic-information-unit-and-ontology-boundary}{%
\subsection{ADR-011 --- Semantic Information Unit and Ontology Boundary}\label{adr-011--semantic-information-unit-and-ontology-boundary}}

Semantic Information Unit and Ontology Boundary

Status

Accepted

Date

2026-07-22

Context

Phase P introduces project-oriented processing around the completed Phase F agentic ingestion pipeline.

P1 implemented the accepted Stakeholder, System and Subsystem framework.

P2 implemented the persistent Project Workspace.

P3 implemented the Project Source Registry with mandatory project assignment, immutable original sources and explicit source roles.

P4 shall transform registered textual engineering sources into heterogeneous, source-traceable Information Units and map them to valid nodes of the accepted framework.

ADR-009 limits the MVP to textual information. Binary document containers may be supported only when their relevant textual content can be projected deterministically into a traceable textual representation.

ADR-010 establishes that the Source Registry remains format-neutral and does not perform text projection, semantic interpretation, ontology mapping or LLM processing.

P4 therefore requires explicit architecture decisions for:

\begin{itemize}
\tightlist
\item
  deterministic source projection
\item
  semantic extraction
\item
  Information Unit identity and atomicity
\item
  information classification
\item
  terminology and ontology references
\item
  project-specific terminology
\item
  framework mapping
\item
  multiagent consensus and confidence
\item
  immutable semantic persistence
\item
  boundaries to P5, P6 and Phase G
\end{itemize}

No universal ontology is treated as an implicit or complete MBSE gold standard.

The selected reference architecture combines:

\begin{itemize}
\tightlist
\item
  SysML v2 and KerML semantics
\item
  BFO 2020
\item
  IOF Core
\item
  a curated Turing Core Vocabulary
\item
  a human-reviewed Project Glossary
\end{itemize}

External reference systems shall support semantic consistency without replacing the accepted project framework or the authoritative engineering model.

Decision

\hypertarget{semantic-processing-stages}{%
\subsubsection{Semantic Processing Stages}\label{semantic-processing-stages}}

P4 separates the semantic workflow into four stages.

\hypertarget{p41--deterministic-source-projection}{%
\paragraph{P4.1 --- Deterministic Source Projection}\label{p41--deterministic-source-projection}}

A registered source is converted into a deterministic, traceable textual representation.

This stage performs no semantic interpretation.

\hypertarget{p42--llm-semantic-extraction}{%
\paragraph{P4.2 --- LLM Semantic Extraction}\label{p42--llm-semantic-extraction}}

Independent persona agents identify candidate Information Units from the deterministic Source Projection.

\hypertarget{p43--terminology-and-ontology-candidate-mapping}{%
\paragraph{P4.3 --- Terminology and Ontology Candidate Mapping}\label{p43--terminology-and-ontology-candidate-mapping}}

Extracted terminology is compared with:

\begin{itemize}
\tightlist
\item
  accepted Project Glossary concepts
\item
  the Turing Core Vocabulary
\item
  selected IOF Core and BFO concepts
\end{itemize}

All mappings remain explicit and traceable.

\hypertarget{p44--framework-mapping}{%
\paragraph{P4.4 --- Framework Mapping}\label{p44--framework-mapping}}

Information Units are mapped to zero, one or multiple valid nodes of the accepted framework template.

Framework mapping produces candidates. It does not perform Engineering Approval.

\hypertarget{artifact-separation}{%
\subsubsection{Artifact Separation}\label{artifact-separation}}

P4 distinguishes the following artifact types.

\hypertarget{source-projection}{%
\paragraph{Source Projection}\label{source-projection}}

A Source Projection is a deterministic textual representation of one registered source.

It does not contain semantic interpretations.

\hypertarget{information-unit}{%
\paragraph{Information Unit}\label{information-unit}}

An Information Unit is the smallest independently understandable, source-traceable professional statement that:

\begin{itemize}
\tightlist
\item
  was semantically extracted from an engineering source
\item
  expresses one independently reviewable main claim
\item
  can be independently classified
\item
  can be independently reviewed
\item
  can be independently mapped to the framework
\end{itemize}

\hypertarget{terminology-concept}{%
\paragraph{Terminology Concept}\label{terminology-concept}}

A Terminology Concept represents a controlled project or reference-system meaning.

It is not an engineering statement and does not create framework coverage.

\hypertarget{framework-assignment}{%
\paragraph{Framework Assignment}\label{framework-assignment}}

A Framework Assignment connects exactly one Information Unit to exactly one valid mapping target of a specific framework-template version.

These artifact types shall not be merged into one mutable record.

\hypertarget{source-role-boundary}{%
\subsubsection{Source-Role Boundary}\label{source-role-boundary}}

Only a source with the role:

\texttt{engineering\_source}

may create:

\begin{itemize}
\tightlist
\item
  engineering Information Units
\item
  Framework Assignments
\item
  preliminary engineering coverage
\end{itemize}

A source with the role:

\texttt{context\_only}

may create:

\begin{itemize}
\tightlist
\item
  a Source Projection
\item
  terminology candidates
\item
  candidate definitions
\item
  candidate explanations
\item
  Project Glossary provenance
\end{itemize}

A context-only source shall not create:

\begin{itemize}
\tightlist
\item
  engineering Information Units
\item
  Framework Assignments
\item
  preliminary coverage
\item
  approved readiness
\item
  model-generation input
\end{itemize}

Terminology extracted from a context-only source remains visibly traceable to that source.

\hypertarget{information-unit-identity}{%
\subsubsection{Information Unit Identity}\label{information-unit-identity}}

Every Information Unit receives an immutable project-local identifier named \texttt{information\_unit\_id}.

The identifier:

\begin{itemize}
\tightlist
\item
  is stored as a JSON string
\item
  matches \texttt{\^{}IU-{[}0-9{]}\{6\}\$}
\item
  is unique within one project
\item
  is allocated sequentially
\item
  is not reused
\item
  remains unchanged for the lifetime of the Information Unit
\end{itemize}

Valid examples include:

\begin{itemize}
\tightlist
\item
  \texttt{IU-000001}
\item
  \texttt{IU-004281}
\item
  \texttt{IU-999999}
\end{itemize}

The globally unambiguous reference is the pair:

\begin{verbatim}
<project_id>/<information_unit_id>
\end{verbatim}

Example:

\begin{verbatim}
318604/IU-000001
\end{verbatim}

The two identifiers remain separate stored fields.

The display form does not replace the individual identifiers.

\hypertarget{information-unit-immutability}{%
\subsubsection{Information Unit Immutability}\label{information-unit-immutability}}

The following Information Unit content is immutable after persistence:

\begin{itemize}
\tightlist
\item
  Information Unit identity
\item
  project reference
\item
  source reference
\item
  Source Projection reference
\item
  source anchors
\item
  source excerpt
\item
  interpreted statement
\item
  information classification
\item
  extraction provenance
\end{itemize}

If later processing produces a materially changed interpretation, it creates a new Information Unit with a new identifier.

An existing Information Unit shall not be silently rewritten.

P5 may later record that one Information Unit supersedes or invalidates another. Such lifecycle information remains separate from the immutable Information Unit.

\hypertarget{information-unit-atomicity}{%
\subsubsection{Information Unit Atomicity}\label{information-unit-atomicity}}

An Information Unit contains exactly one independently reviewable main claim.

Independent obligations shall be split into separate Information Units.

Conditions, qualifiers and limits that are necessary to understand one claim remain attached to that claim.

An Information Unit belongs to exactly one engineering source.

It may reference multiple anchors within that source.

P4 shall not synthesize one Information Unit from multiple sources.

Equivalent or overlapping claims from different sources remain separate Information Units.

They may be linked later, but they shall not be automatically merged.

Content fingerprints may support:

\begin{itemize}
\tightlist
\item
  duplicate detection
\item
  idempotence
\item
  comparison
\item
  reproducibility
\end{itemize}

A content fingerprint never replaces the Information Unit identifier.

\hypertarget{source-excerpt-and-interpreted-statement}{%
\subsubsection{Source Excerpt and Interpreted Statement}\label{source-excerpt-and-interpreted-statement}}

An Information Unit distinguishes:

\begin{verbatim}
source_excerpt
interpreted_statement
\end{verbatim}

\texttt{source\_excerpt} contains the unchanged relevant text from the deterministic Source Projection.

\texttt{interpreted\_statement} contains the semantic candidate interpretation.

The field name \texttt{normalized\_statement} shall not be used for the semantic interpretation because deterministic normalization and semantic interpretation are separate operations.

\hypertarget{source-projection-identity}{%
\subsubsection{Source Projection Identity}\label{source-projection-identity}}

Every Source Projection receives an immutable project-local identifier named \texttt{source\_projection\_id}.

The identifier:

\begin{itemize}
\tightlist
\item
  matches \texttt{\^{}SP-{[}0-9{]}\{6\}\$}
\item
  is unique within one project
\item
  is allocated sequentially
\item
  is not reused
\end{itemize}

Valid examples include:

\begin{itemize}
\tightlist
\item
  \texttt{SP-000001}
\item
  \texttt{SP-004281}
\item
  \texttt{SP-999999}
\end{itemize}

A Source Projection belongs to exactly one registered source.

\hypertarget{source-projection-segment-identity}{%
\subsubsection{Source Projection Segment Identity}\label{source-projection-segment-identity}}

Every segment within one Source Projection receives a projection-local identifier named \texttt{segment\_id}.

The identifier:

\begin{itemize}
\tightlist
\item
  matches \texttt{\^{}SEG-{[}0-9{]}\{6\}\$}
\item
  is unique within one Source Projection
\item
  preserves deterministic source order
\end{itemize}

An Information Unit shall reference at least one segment.

All segments referenced by one Information Unit shall belong to the same Source Projection and therefore to the same registered source.

\hypertarget{deterministic-source-projection}{%
\subsubsection{Deterministic Source Projection}\label{deterministic-source-projection}}

The LLM shall not receive an original registered file directly.

Every semantic process begins with a deterministic, persistable Source Projection.

Permitted deterministic operations include:

\begin{itemize}
\tightlist
\item
  deterministic UTF-8 decoding
\item
  removal of a UTF-8 byte-order mark
\item
  normalization of line endings
\item
  deterministic textual extraction from a supported container
\item
  preservation of source order
\item
  deterministic structural segmentation
\item
  generation of source locators
\item
  generation of projection issues
\item
  calculation of hashes and fingerprints
\end{itemize}

The deterministic stage shall not perform:

\begin{itemize}
\tightlist
\item
  spelling correction
\item
  synonym replacement
\item
  terminology harmonization
\item
  unit conversion
\item
  semantic rewriting
\item
  summarization
\item
  requirement rewriting
\item
  ontology mapping
\item
  framework mapping
\item
  domain inference
\end{itemize}

Deterministic source projection is syntactic and structural, not semantic.

\hypertarget{supported-p4-source-adapters}{%
\subsubsection{Supported P4 Source Adapters}\label{supported-p4-source-adapters}}

The P4 MVP provides explicit adapters for:

\begin{verbatim}
.txt
.md
.json
.csv
.tsv
.pdf
\end{verbatim}

An extension is supported only when a deterministic adapter has been implemented and validated.

The P4 MVP does not support:

\begin{verbatim}
.doc
.docx
.odt
.ppt
.pptx
.xls
.xlsx
.rtf
\end{verbatim}

These formats may be added later through explicit, versioned adapters.

The P4 MVP also excludes:

\begin{itemize}
\tightlist
\item
  scanned PDFs
\item
  OCR
\item
  image interpretation
\item
  drawing interpretation
\item
  diagram interpretation
\item
  geometric reasoning
\item
  audio
\item
  video
\item
  handwritten information
\end{itemize}

\hypertarget{text-encoding}{%
\subsubsection{Text Encoding}\label{text-encoding}}

Textual source adapters accept:

\begin{itemize}
\tightlist
\item
  UTF-8
\item
  UTF-8 with byte-order mark
\end{itemize}

Encoding guessing and silent fallback decoding are prohibited.

Unsupported or invalid encodings shall produce a visible projection issue.

\hypertarget{json-projection}{%
\subsubsection{JSON Projection}\label{json-projection}}

A JSON source is parsed deterministically.

JSON locations use JSON Pointer references.

Duplicate object keys are rejected because they would make the source meaning and traceability ambiguous.

Semantic interpretation of JSON values remains outside the deterministic projection step.

\hypertarget{csv-and-tsv-projection}{%
\subsubsection{CSV and TSV Projection}\label{csv-and-tsv-projection}}

CSV uses:

\begin{itemize}
\tightlist
\item
  RFC 4180-compatible parsing
\item
  comma as the fixed delimiter
\end{itemize}

TSV uses:

\begin{itemize}
\tightlist
\item
  tab as the fixed delimiter
\end{itemize}

Dialect sniffing is prohibited.

Row and column locations remain available for traceability.

\hypertarget{pdf-projection}{%
\subsubsection{PDF Projection}\label{pdf-projection}}

PDF support uses an explicit and pinned \texttt{pypdf} dependency.

Only machine-readable text is extracted.

Page boundaries are retained as source locators.

P4 does not perform:

\begin{itemize}
\tightlist
\item
  OCR
\item
  image extraction for semantic processing
\item
  drawing interpretation
\item
  diagram interpretation
\end{itemize}

An empty or unextractable page produces a visible projection issue.

A PDF containing both extractable text and unsupported visual information may produce a partial Source Projection.

\hypertarget{source-projection-result}{%
\subsubsection{Source Projection Result}\label{source-projection-result}}

A Source Projection has exactly one result:

\begin{verbatim}
complete
partial
unavailable
\end{verbatim}

\texttt{complete} means that the supported textual content was projected without a known projection loss.

\texttt{partial} means that usable textual content exists, but unsupported, unextractable or failed content was detected.

\texttt{unavailable} means that no usable deterministic textual projection could be created.

A partial result remains visible and requires its issues to be retained.

An unavailable result shall not enter semantic extraction.

\hypertarget{source-projection-fingerprint}{%
\subsubsection{Source Projection Fingerprint}\label{source-projection-fingerprint}}

A Source Projection fingerprint is derived from at least:

\begin{itemize}
\tightlist
\item
  original source SHA-256
\item
  adapter identifier
\item
  adapter version
\item
  deterministic adapter configuration
\end{itemize}

The fingerprint supports reproducibility and idempotence.

It does not replace:

\begin{itemize}
\tightlist
\item
  \texttt{project\_id}
\item
  \texttt{source\_id}
\item
  \texttt{source\_projection\_id}
\end{itemize}

\hypertarget{information-classification}{%
\subsubsection{Information Classification}\label{information-classification}}

P4 separates three classification dimensions:

\begin{itemize}
\tightlist
\item
  Information Type
\item
  Statement Modality
\item
  Epistemic Class
\end{itemize}

These dimensions shall not be collapsed into one field.

\hypertarget{information-types}{%
\subsubsection{Information Types}\label{information-types}}

The accepted Information Types are:

\begin{verbatim}
stakeholder
actor
user_need
requirement
use_case
function
logical_element
physical_element
interface
constraint
information_item
definition
rationale
decision
risk
ambiguity
gap
open_question
unclassified
\end{verbatim}

\texttt{subsystem} is not an Information Type.

Subsystem is a framework level and shall not be confused with the semantic kind of an Information Unit.

\texttt{unclassified} is an explicit valid result. The system shall not force an unsupported classification.

\hypertarget{statement-modalities}{%
\subsubsection{Statement Modalities}\label{statement-modalities}}

The accepted Statement Modalities are:

\begin{verbatim}
descriptive
normative
definitional
interrogative
\end{verbatim}

Statement Modality describes how the statement is expressed.

It does not determine its framework mapping or Engineering Approval.

\hypertarget{epistemic-classes}{%
\subsubsection{Epistemic Classes}\label{epistemic-classes}}

The accepted Epistemic Classes are:

\begin{verbatim}
explicit
interpretation
derivation
assumption
\end{verbatim}

\texttt{explicit} means that the professional statement is directly expressed in the source.

\texttt{interpretation} means that semantic interpretation is necessary, but the statement remains grounded in the source.

\texttt{derivation} means that the statement is derived from other Information Units.

\texttt{assumption} means that information not established by the source is introduced to make an interpretation possible or visible.

A derivation requires:

\begin{itemize}
\tightlist
\item
  referenced supporting Information Unit identifiers
\item
  a derivation rationale
\end{itemize}

During P4, supporting Information Units for a derivation shall belong to the same source.

An assumption requires:

\begin{itemize}
\tightlist
\item
  an explicit missing-evidence explanation
\item
  a visible assumption classification
\end{itemize}

The source anchor of an assumption records what triggered the assumption. It does not turn the assumption into source evidence.

\hypertarget{preliminary-coverage-exclusions}{%
\subsubsection{Preliminary Coverage Exclusions}\label{preliminary-coverage-exclusions}}

The following Information Types shall never create positive preliminary coverage:

\begin{verbatim}
risk
ambiguity
gap
open_question
\end{verbatim}

Information Units classified as \texttt{assumption} shall never create positive preliminary coverage.

They may still be associated with a relevant framework node for diagnostic traceability.

The complete deterministic coverage policy belongs to P6.

\hypertarget{human-review-separation}{%
\subsubsection{Human Review Separation}\label{human-review-separation}}

An immutable Information Unit shall not contain a mutable Human Review status.

Human Review decisions are separate records.

Terminology acceptance and Engineering Approval are separate processes.

Engineering Approval remains assigned to Phase G.

\hypertarget{semantic-confidence}{%
\subsubsection{Semantic Confidence}\label{semantic-confidence}}

An individual LLM agent shall not determine the final confidence of an Information Unit, terminology mapping or Framework Assignment.

Final confidence is derived deterministically from the variance of independent persona-agent results.

The accepted confidence values are:

\begin{verbatim}
high
medium
low
\end{verbatim}

Confidence is ordinal.

It is not:

\begin{itemize}
\tightlist
\item
  a statistical probability
\item
  a truth value
\item
  a Human Review decision
\item
  an Engineering Approval
\item
  a readiness decision
\end{itemize}

A confidence value requires an auditable rationale.

\hypertarget{multiagent-execution}{%
\subsubsection{Multiagent Execution}\label{multiagent-execution}}

P4 reuses the existing Phase F multiagent architecture.

All required members of one semantic team receive:

\begin{itemize}
\tightlist
\item
  the same task
\item
  the same structured input
\item
  the same role
\item
  different explicit personas or professional perspectives
\end{itemize}

Team members execute independently.

They shall not coordinate before consensus analysis.

P4 may define dedicated teams for:

\begin{itemize}
\tightlist
\item
  semantic extraction
\item
  terminology and ontology mapping
\item
  framework mapping
\end{itemize}

Existing Phase F outputs are not automatically P4 Information Units.

P4 requires P4-specific structured output schemas and validators.

\hypertarget{consensus-and-variance}{%
\subsubsection{Consensus and Variance}\label{consensus-and-variance}}

A deterministic consensus analyzer compares structured persona results.

The analyzer shall preserve:

\begin{verbatim}
consensus_level
variance_level
confidence
total_personas
supporting_personas
dissenting_personas
value_distribution
review_required
confidence_rationale
consensus_report_id
\end{verbatim}

The confidence is derived according to the following policy.

\hypertarget{high-confidence}{%
\paragraph{High Confidence}\label{high-confidence}}

\texttt{high} requires agreement of all required personas on the comparable result, without a contradictory result or unexplained omission.

This corresponds to low inter-persona variance.

\hypertarget{medium-confidence}{%
\paragraph{Medium Confidence}\label{medium-confidence}}

\texttt{medium} may be assigned when:

\begin{itemize}
\tightlist
\item
  a strict majority supports the same result and other personas omit it
\item
  a strict majority supports the same result but a minority provides a differing result
\end{itemize}

A differing result requires Human Review even when a majority exists.

\hypertarget{low-confidence}{%
\paragraph{Low Confidence}\label{low-confidence}}

\texttt{low} is assigned when:

\begin{itemize}
\tightlist
\item
  only one persona proposes the result
\item
  no clear majority exists
\item
  conflicting results exist
\item
  the team run is technically incomplete
\item
  the result cannot be compared reliably
\end{itemize}

Low confidence requires Human Review.

\hypertarget{inter-persona-and-intra-persona-variance}{%
\subsubsection{Inter-Persona and Intra-Persona Variance}\label{inter-persona-and-intra-persona-variance}}

Inter-persona variance measures differences between distinct professional perspectives.

Intra-persona variance measures differences between repeated runs of the same persona.

Repeated runs of the same persona shall not be counted as additional independent votes.

Each persona receives at most one vote in the team consensus.

Repeated runs measure stability and may reduce confidence when the persona produces inconsistent results.

They shall not create an artificial majority.

The persona team is intentionally heterogeneous and is not a random, independent statistical sample.

The resulting variance is an epistemic uncertainty signal, not a statistically calibrated probability.

\hypertarget{field-level-confidence}{%
\subsubsection{Field-Level Confidence}\label{field-level-confidence}}

Consensus and variance may be calculated separately for:

\begin{itemize}
\tightlist
\item
  existence of an Information Unit
\item
  semantic interpretation
\item
  Information Type
\item
  Statement Modality
\item
  Epistemic Class
\item
  terminology mapping
\item
  ontology mapping
\item
  framework target
\end{itemize}

A Framework Assignment receives its confidence from the agreement about its specific target and mapping rationale.

The lowest confidence of a critical required field conservatively limits the overall confidence of the resulting candidate.

\hypertarget{semantic-reference-stack}{%
\subsubsection{Semantic Reference Stack}\label{semantic-reference-stack}}

The accepted semantic reference stack is:

\begin{enumerate}
\def\labelenumi{\arabic{enumi}.}
\tightlist
\item
  SysML v2 and KerML
\item
  BFO 2020
\item
  IOF Core
\item
  Turing Core Vocabulary
\item
  Project Glossary
\item
  LLM candidate generation
\end{enumerate}

These layers have distinct responsibilities.

\hypertarget{sysml-v2-and-kerml}{%
\subsubsection{SysML v2 and KerML}\label{sysml-v2-and-kerml}}

SysML v2 and KerML define the normative target-model and target-notation semantics.

The curated repository references remain:

\begin{verbatim}
context/sysml/sysml_v2_spec_reference.json
context/sysml/sysml_v2_target_notation.json
\end{verbatim}

These files constrain later model generation and valid target notation.

They are not treated as a complete industrial-domain ontology.

\hypertarget{bfo-2020}{%
\subsubsection{BFO 2020}\label{bfo-2020}}

BFO 2020 is the selected top-level ontology reference.

It provides general distinctions for entities, processes, qualities, roles and related foundational categories.

The selected reference corresponds to BFO 2020 and ISO/IEC 21838-2.

BFO does not directly provide the complete MBSE or project vocabulary.

\hypertarget{iof-core}{%
\subsubsection{IOF Core}\label{iof-core}}

IOF Core is the selected industrial mid-level ontology reference.

The pinned reference version uses:

\begin{verbatim}
versionIRI: 202602
maturity: Released
\end{verbatim}

IOF Core provides industrial semantic structure above BFO and below project-specific terminology.

IOF Core does not replace:

\begin{itemize}
\tightlist
\item
  SysML v2
\item
  the accepted framework
\item
  Turing Core
\item
  the Project Glossary
\item
  human engineering decisions
\end{itemize}

\hypertarget{local-ontology-snapshots}{%
\subsubsection{Local Ontology Snapshots}\label{local-ontology-snapshots}}

BFO and IOF are used through pinned local snapshots.

The intended structure is:

\begin{verbatim}
external/ontologies/bfo/2020/bfo-core.owl
external/ontologies/iof/202602/AboutIOFProd.rdf
external/ontologies/iof/202602/Core.rdf
external/ontologies/iof/202602/AnnotationVocabulary.rdf
\end{verbatim}

The corresponding license material shall be stored with the snapshots.

P4 shall not perform:

\begin{itemize}
\tightlist
\item
  live ontology queries
\item
  automatic ontology downloads
\item
  automatic ontology updates
\item
  remote runtime dependency resolution
\end{itemize}

Updating an ontology snapshot requires an explicit reviewed change.

\hypertarget{ontology-registry}{%
\subsubsection{Ontology Registry}\label{ontology-registry}}

The selected ontology snapshots are registered in:

\begin{verbatim}
context/semantics/ontology_registry.json
\end{verbatim}

The registry records at least:

\begin{itemize}
\tightlist
\item
  reference-system identifier
\item
  version
\item
  version IRI where applicable
\item
  local file path
\item
  source authority
\item
  maturity
\item
  license reference
\item
  enabled runtime role
\item
  checksum
\end{itemize}

The registry is curated and versioned.

\hypertarget{reference-concept-index}{%
\subsubsection{Reference Concept Index}\label{reference-concept-index}}

P4 does not require an OWL reasoner or triple store.

A deterministic generated read-only concept index is stored in:

\begin{verbatim}
context/semantics/reference_concept_index.json
\end{verbatim}

The index supports runtime retrieval of selected:

\begin{itemize}
\tightlist
\item
  concept identifiers
\item
  IRIs
\item
  preferred labels
\item
  alternative labels
\item
  definitions
\item
  parent relationships
\item
  source-system references
\item
  version references
\end{itemize}

The generated index is not the ontology authority.

The pinned RDF and OWL snapshots remain the auditable external-reference artifacts.

The index can be regenerated deterministically.

\hypertarget{runtime-ontology-boundary}{%
\subsubsection{Runtime Ontology Boundary}\label{runtime-ontology-boundary}}

P4 shall not introduce:

\begin{itemize}
\tightlist
\item
  a triple store
\item
  an OWL reasoner
\item
  live SPARQL endpoints
\item
  unrestricted ontology graph traversal
\item
  automatic inference from the complete external ontology graph
\item
  automatic ontology updates
\end{itemize}

Only relevant retrieved concepts are supplied to an LLM.

Complete ontology files shall not be loaded into every prompt.

\hypertarget{turing-core-vocabulary}{%
\subsubsection{Turing Core Vocabulary}\label{turing-core-vocabulary}}

A curated Turing Core Vocabulary is stored in:

\begin{verbatim}
context/semantics/turing_core_vocabulary.json
\end{verbatim}

Turing Core provides the controlled bridge between:

\begin{itemize}
\tightlist
\item
  MBSE terminology
\item
  the accepted RFLP-oriented framework
\item
  SysML v2 target concepts
\item
  selected external ontology concepts
\item
  project-specific concepts
\end{itemize}

Turing Core concepts receive stable identifiers matching:

\begin{verbatim}
^TC-[0-9]{6}$
\end{verbatim}

Valid examples include:

\begin{itemize}
\tightlist
\item
  \texttt{TC-000001}
\item
  \texttt{TC-004281}
\item
  \texttt{TC-999999}
\end{itemize}

Turing Core is global to the Turing Generator.

It shall not be modified automatically from one project.

\hypertarget{sysdo-boundary}{%
\subsubsection{SysDO Boundary}\label{sysdo-boundary}}

SysDO is retained as a non-normative research and structuring reference.

It is not the primary semantic authority because its demonstrated scope and maturity do not satisfy the complete requirements of the Turing Generator.

In particular, it does not replace the selected combination of:

\begin{itemize}
\tightlist
\item
  SysML v2 and KerML
\item
  BFO
\item
  IOF Core
\item
  Turing Core
\item
  Project Glossary
\end{itemize}

SysDO remains disabled as a runtime ontology source during the P4 MVP.

Nothing is copied from SysDO without explicit review.

\hypertarget{external-ontology-mappings}{%
\subsubsection{External Ontology Mappings}\label{external-ontology-mappings}}

A Turing or Project Concept may map to an external concept using one of the following relations:

\begin{verbatim}
exact_match
narrower_than
broader_than
related_to
no_equivalent
\end{verbatim}

Every mapping requires:

\begin{itemize}
\tightlist
\item
  the referenced external concept or IRI
\item
  the referenced ontology and version
\item
  a mapping relation
\item
  a rationale
\item
  provenance
\end{itemize}

An unmapped concept is valid.

The system shall not force an \texttt{exact\_match}.

External ontology mappings cannot:

\begin{itemize}
\tightlist
\item
  create engineering requirements
\item
  create framework coverage
\item
  approve engineering information
\item
  overwrite a project meaning
\item
  overwrite Turing Core
\item
  alter the authoritative engineering model
\end{itemize}

\hypertarget{project-glossary}{%
\subsubsection{Project Glossary}\label{project-glossary}}

Every project owns a versioned Project Glossary under:

\begin{verbatim}
data/projects/<project_id>/semantics/project_glossary.json
\end{verbatim}

The Project Glossary captures reviewed project-specific vocabulary.

It supplements the external reference stack.

It does not replace or overwrite BFO, IOF Core or Turing Core.

\hypertarget{project-concept-identity}{%
\subsubsection{Project Concept Identity}\label{project-concept-identity}}

Every Project Concept receives a stable project-local identifier named \texttt{project\_concept\_id}.

The identifier:

\begin{itemize}
\tightlist
\item
  matches \texttt{\^{}PC-{[}0-9{]}\{6\}\$}
\item
  is unique within one project
\item
  is allocated sequentially
\item
  is not reused
\end{itemize}

The globally unambiguous reference is the pair:

\begin{verbatim}
<project_id>/<project_concept_id>
\end{verbatim}

Example:

\begin{verbatim}
318604/PC-000001
\end{verbatim}

Project Concept identity remains stable across revisions.

\hypertarget{project-concept-content}{%
\subsubsection{Project Concept Content}\label{project-concept-content}}

A Project Concept may contain:

\begin{itemize}
\tightlist
\item
  multilingual preferred labels
\item
  multilingual alternative labels
\item
  multilingual definitions
\item
  broader Project Concept references
\item
  related Project Concept references
\item
  Turing Core mappings
\item
  BFO or IOF mappings
\item
  provenance
\item
  lifecycle status
\item
  revision
\item
  rationale
\end{itemize}

Each project defines a default language.

\hypertarget{preferred-label-uniqueness}{%
\subsubsection{Preferred-Label Uniqueness}\label{preferred-label-uniqueness}}

Accepted preferred labels shall be unique within one project and language.

Uniqueness comparison applies:

\begin{itemize}
\tightlist
\item
  Unicode NFKC normalization
\item
  whitespace trimming
\item
  Unicode case folding
\end{itemize}

Labels that differ only through these comparison operations conflict.

For example, case-only or compatible full-width variants shall not silently create separate accepted concepts.

\hypertarget{ambiguous-alternative-labels}{%
\subsubsection{Ambiguous Alternative Labels}\label{ambiguous-alternative-labels}}

A true homonym may be retained as an alternative label only through an explicit Ambiguity Group.

An Ambiguity Group records at least:

\begin{itemize}
\tightlist
\item
  label
\item
  language
\item
  candidate Project Concept identifiers
\item
  \texttt{resolution\_rule:\ context\_required}
\end{itemize}

Ambiguous labels shall never be resolved automatically.

The user interface may display disambiguated names such as:

\begin{verbatim}
Port (Physical Interface)
Port (Network Endpoint)
\end{verbatim}

The display text does not replace the stable Project Concept identifiers.

\hypertarget{project-concept-provenance}{%
\subsubsection{Project Concept Provenance}\label{project-concept-provenance}}

Every Project Concept or revision requires provenance from at least one of:

\begin{itemize}
\tightlist
\item
  an engineering source
\item
  a context-only source
\item
  an explicit human terminology decision
\item
  a selected external reference system
\item
  Turing Core
\end{itemize}

Context-only provenance is permitted for terminology.

It does not create engineering evidence or framework coverage.

\hypertarget{project-concept-lifecycle}{%
\subsubsection{Project Concept Lifecycle}\label{project-concept-lifecycle}}

The accepted Project Concept lifecycle states are:

\begin{verbatim}
candidate
accepted
rejected
deprecated
\end{verbatim}

Concepts shall not be deleted or silently overwritten.

A concept cannot become accepted solely through an LLM result.

\hypertarget{terminology-decisions}{%
\subsubsection{Terminology Decisions}\label{terminology-decisions}}

Only a human may create the authoritative decision that accepts, rejects or later replaces a Project Concept.

Terminology Decisions receive immutable project-local identifiers matching:

\begin{verbatim}
^TD-[0-9]{6}$
\end{verbatim}

A Terminology Decision records at least:

\begin{itemize}
\tightlist
\item
  project identifier
\item
  Terminology Decision identifier
\item
  Project Concept identifier
\item
  Project Concept revision
\item
  decision
\item
  reviewer identity
\item
  decision timestamp
\item
  rationale
\end{itemize}

Reviewer identity is documented for auditability.

It is not treated as cryptographic identity proof.

\hypertarget{terminology-and-engineering-approval}{%
\subsubsection{Terminology and Engineering Approval}\label{terminology-and-engineering-approval}}

Terminology acceptance confirms the intended meaning of a Project Concept.

It does not approve an engineering statement.

Engineering Review and Approved Input Promotion remain assigned to Phase G.

A terminology-approved concept may be used consistently in later processing, but it does not create an approved Information Unit.

\hypertarget{llm-terminology-permissions}{%
\subsubsection{LLM Terminology Permissions}\label{llm-terminology-permissions}}

An LLM may propose:

\begin{itemize}
\tightlist
\item
  Project Concepts
\item
  preferred labels
\item
  alternative labels
\item
  definitions
\item
  ambiguity candidates
\item
  Turing Core mappings
\item
  external ontology mappings
\end{itemize}

An LLM shall not:

\begin{itemize}
\tightlist
\item
  accept a Project Concept
\item
  reject a Project Concept authoritatively
\item
  merge concepts authoritatively
\item
  resolve a homonym silently
\item
  declare an \texttt{exact\_match} authoritatively
\item
  overwrite an accepted concept
\item
  promote a Project Concept into Turing Core
\item
  share a Project Concept with another project
\end{itemize}

\hypertarget{terminology-conflicts}{%
\subsubsection{Terminology Conflicts}\label{terminology-conflicts}}

When new source usage conflicts with an accepted Project Glossary meaning, the system creates a visible Terminology Conflict.

The conflict shall not be resolved automatically.

A human determines whether the result requires:

\begin{itemize}
\tightlist
\item
  a new concept
\item
  a new concept revision
\item
  an additional alternative label
\item
  an Ambiguity Group
\item
  deprecation
\item
  rejection of the candidate interpretation
\item
  no glossary change
\end{itemize}

\hypertarget{project-concept-revisions}{%
\subsubsection{Project Concept Revisions}\label{project-concept-revisions}}

A Project Concept identifier remains stable.

A material edit creates a new revision.

Revision history remains preserved.

Only accepted revisions are authoritative in later semantic processing.

Candidate revisions may be shown to an LLM and reviewer but are not authoritative.

Rejected revisions shall not be used as positive semantic evidence.

Deprecated revisions remain available only for traceability.

P5 defines the detailed lifecycle-event and invalidation storage.

\hypertarget{project-isolation}{%
\subsubsection{Project Isolation}\label{project-isolation}}

Project Concepts remain project-specific.

They shall not be automatically:

\begin{itemize}
\tightlist
\item
  copied to another project
\item
  shared across projects
\item
  promoted into Turing Core
\item
  treated as globally authoritative
\end{itemize}

A future cross-project curation workflow requires a separate architecture decision.

\hypertarget{framework-assignment-identity}{%
\subsubsection{Framework Assignment Identity}\label{framework-assignment-identity}}

Every Framework Assignment receives an immutable project-local identifier named \texttt{framework\_assignment\_id}.

The identifier:

\begin{itemize}
\tightlist
\item
  matches \texttt{\^{}FA-{[}0-9{]}\{6\}\$}
\item
  is unique within one project
\item
  is allocated sequentially
\item
  is not reused
\end{itemize}

The globally unambiguous reference is the pair:

\begin{verbatim}
<project_id>/<framework_assignment_id>
\end{verbatim}

Example:

\begin{verbatim}
318604/FA-000001
\end{verbatim}

\hypertarget{framework-assignment-semantics}{%
\subsubsection{Framework Assignment Semantics}\label{framework-assignment-semantics}}

A Framework Assignment connects:

\begin{itemize}
\tightlist
\item
  exactly one Information Unit
\item
  to exactly one valid framework mapping target
\item
  in exactly one framework-template version
\end{itemize}

An Information Unit may have:

\begin{itemize}
\tightlist
\item
  no Framework Assignment
\item
  one Framework Assignment
\item
  multiple Framework Assignments
\end{itemize}

Every individual assignment requires its own:

\begin{itemize}
\tightlist
\item
  identifier
\item
  mapping rationale
\item
  mapping basis
\item
  confidence evidence
\end{itemize}

Multiple assignments are not represented by one ambiguous target field.

\hypertarget{framework-assignment-content}{%
\subsubsection{Framework Assignment Content}\label{framework-assignment-content}}

A Framework Assignment records at least:

\begin{verbatim}
schema_version
project_id
framework_assignment_id
information_unit_id
framework_template_id
framework_template_version
framework_node_id
mapping_bases
mapping_rationale
confidence
confidence_rationale
consensus_report_id
created_at
supersedes_assignment_id
\end{verbatim}

\texttt{supersedes\_assignment\_id} is absent when the assignment does not supersede an earlier assignment.

\hypertarget{framework-target-validation}{%
\subsubsection{Framework Target Validation}\label{framework-target-validation}}

\texttt{framework\_node\_id} must identify a valid mapping target in the referenced framework-template version.

Free-form target names are prohibited.

Unknown framework targets are rejected.

The initial accepted framework remains:

\begin{verbatim}
TURING_RFLP_FRAMEWORK
version 1.0.0
\end{verbatim}

\hypertarget{multiple-framework-assignments}{%
\subsubsection{Multiple Framework Assignments}\label{multiple-framework-assignments}}

Multiple Framework Assignments are permitted only when each target is professionally justified.

A mapping to a child node shall not automatically create a mapping to its parent.

A mapping shall not be propagated automatically to:

\begin{itemize}
\tightlist
\item
  another framework level
\item
  a sibling node
\item
  another subsystem
\item
  another Information Unit
\end{itemize}

Framework hierarchy and semantic relevance remain separate concerns.

\hypertarget{framework-mapping-bases}{%
\subsubsection{Framework Mapping Bases}\label{framework-mapping-bases}}

An assignment may reference one or multiple mapping bases:

\begin{verbatim}
direct_semantic
structural_context
project_glossary
turing_core
external_ontology
\end{verbatim}

Referenced concepts, Project Concept revisions and ontology IRIs remain explicitly traceable.

Label similarity alone is not a sufficient mapping rationale.

An accepted Project Glossary concept may support a Framework Assignment.

It shall not force one.

\hypertarget{framework-mapping-evaluation}{%
\subsubsection{Framework Mapping Evaluation}\label{framework-mapping-evaluation}}

The absence of a Framework Assignment is ambiguous because it could mean either:

\begin{itemize}
\tightlist
\item
  mapping has not been performed
\item
  mapping was performed and no valid target was found
\end{itemize}

Every completed mapping attempt therefore creates an immutable Framework Mapping Evaluation.

Framework Mapping Evaluation identifiers match:

\begin{verbatim}
^FME-[0-9]{6}$
\end{verbatim}

The accepted evaluation outcomes are:

\begin{verbatim}
mapped
unmapped
ambiguous
failed
\end{verbatim}

An evaluation references zero or multiple Framework Assignments.

\texttt{unmapped}, \texttt{ambiguous} and \texttt{failed} require an explicit rationale or error description.

\hypertarget{framework-mapping-governance}{%
\subsubsection{Framework-Mapping Governance}\label{framework-mapping-governance}}

The LLM may:

\begin{itemize}
\tightlist
\item
  propose Framework Assignments
\item
  propose multiple possible targets
\item
  report an ambiguous mapping
\item
  report that no valid mapping was found
\item
  provide concise mapping rationales
\end{itemize}

The LLM shall not:

\begin{itemize}
\tightlist
\item
  approve a Framework Assignment
\item
  silently choose one target from an ambiguous result
\item
  create a framework node
\item
  change the framework hierarchy
\item
  treat confidence as approval
\item
  perform Engineering Approval
\end{itemize}

P4 produces mapping candidates.

Engineering Approval remains assigned to Phase G.

\hypertarget{framework-version-binding}{%
\subsubsection{Framework Version Binding}\label{framework-version-binding}}

Every Framework Assignment remains permanently bound to the framework-template version used during its creation.

A new framework-template version requires a new mapping evaluation.

Existing assignments shall not be silently migrated.

\hypertarget{mapping-correction}{%
\subsubsection{Mapping Correction}\label{mapping-correction}}

A corrected mapping produces:

\begin{itemize}
\tightlist
\item
  a new Framework Mapping Evaluation
\item
  new Framework Assignments where required
\item
  explicit later supersession references
\end{itemize}

Existing assignments and evaluations remain auditable.

\hypertarget{semantic-persistence-root}{%
\subsubsection{Semantic Persistence Root}\label{semantic-persistence-root}}

Project-specific semantic records are stored logically below:

\begin{verbatim}
data/projects/<project_id>/semantics/
\end{verbatim}

The accepted semantic areas are:

\begin{verbatim}
data/projects/<project_id>/semantics/
├── source_projections/
├── information_units/
├── project_glossary.json
├── terminology_decisions/
├── terminology_conflicts/
└── framework_mappings/
\end{verbatim}

The exact organization of processing runs, temporary artifacts and derived indexes belongs to P5.

\hypertarget{source-projection-persistence}{%
\subsubsection{Source Projection Persistence}\label{source-projection-persistence}}

A persisted Source Projection contains:

\begin{verbatim}
source_projections/<source_projection_id>/
├── projection.json
└── content.txt
\end{verbatim}

\texttt{projection.json} records at least:

\begin{itemize}
\tightlist
\item
  schema version
\item
  project identifier
\item
  source identifier
\item
  Source Projection identifier
\item
  source SHA-256
\item
  adapter identifier
\item
  adapter version
\item
  adapter configuration
\item
  projection fingerprint
\item
  projection result
\item
  segments
\item
  source locators
\item
  issues
\item
  projected-content SHA-256
\item
  creation timestamp
\end{itemize}

\texttt{content.txt} contains the deterministic UTF-8 textual projection.

It does not replace the immutable original source in the P3 Source Registry.

\hypertarget{persistence-rules}{%
\subsubsection{Persistence Rules}\label{persistence-rules}}

P4 records are:

\begin{itemize}
\tightlist
\item
  schema-validated
\item
  project-isolated
\item
  written atomically
\item
  immutable after publication unless explicitly defined as versioned
\item
  never silently overwritten
\item
  never silently deleted
\end{itemize}

A new record shall not reference a different project.

Derived indexes may be regenerated.

They are not authoritative.

The validated P4 records remain authoritative for the P4 semantic repository.

\hypertarget{reproducibility-metadata}{%
\subsubsection{Reproducibility Metadata}\label{reproducibility-metadata}}

Semantic processing remains traceable to:

\begin{itemize}
\tightlist
\item
  project identifier
\item
  source identifier
\item
  Source Projection and fingerprint
\item
  adapter identifier and version
\item
  team configuration
\item
  persona configuration
\item
  LLM provider
\item
  LLM model
\item
  prompt or schema version
\item
  Consensus Report
\item
  Ontology Registry version
\item
  Reference Concept Index version
\item
  Turing Core version
\item
  Project Glossary concept revisions
\item
  framework-template identifier and version
\end{itemize}

BFO and IOF files are not copied into every project.

Project records reference the registered pinned versions and concept IRIs.

\hypertarget{reprocessing}{%
\subsubsection{Reprocessing}\label{reprocessing}}

Reprocessing the same Source Projection creates a new Processing Run.

An identical validated result may be recognized through fingerprints.

A materially different semantic result creates a new semantic record with a new identifier or revision.

Fingerprints support comparison and idempotence.

They do not replace semantic identities.

\hypertarget{p4-responsibility}{%
\subsubsection{P4 Responsibility}\label{p4-responsibility}}

P4 owns:

\begin{itemize}
\tightlist
\item
  deterministic Source Projection adapters
\item
  Source Projection validation
\item
  Information Unit schema and validation
\item
  semantic multiagent extraction
\item
  terminology and ontology candidates
\item
  Project Glossary domain rules
\item
  Framework Assignment and Evaluation rules
\item
  field-level consensus and variance
\item
  confidence derivation
\item
  persistence of validated semantic records
\end{itemize}

\hypertarget{p5-boundary}{%
\subsubsection{P5 Boundary}\label{p5-boundary}}

P5 owns:

\begin{itemize}
\tightlist
\item
  Processing Run identity
\item
  processing states
\item
  run resumption
\item
  organization of raw agent outputs
\item
  Consensus Report organization
\item
  temporary processing artifacts
\item
  lifecycle events
\item
  supersession
\item
  invalidation
\item
  derived operational indexes
\item
  detection of incomplete publication
\end{itemize}

P5 shall not rewrite the professional content of an Information Unit, Framework Assignment or terminology decision.

A Processing Run is complete only after its required P4 records have been validated and published successfully.

Partial processing results shall not appear as a completed run.

\hypertarget{p6-boundary}{%
\subsubsection{P6 Boundary}\label{p6-boundary}}

P6 reads the P4 semantic records that P5 identifies as usable.

P6 calculates:

\begin{itemize}
\tightlist
\item
  Preliminary Coverage
\item
  coverage gaps
\item
  mapping conflicts
\item
  preliminary framework support
\end{itemize}

P6 shall not modify P4 records.

P6 does not treat an Information Unit as approved engineering input.

\hypertarget{phase-g-boundary}{%
\subsubsection{Phase G Boundary}\label{phase-g-boundary}}

P4 records are candidates.

Phase G creates separate Human Review and Engineering Approval records.

An approval references the exact immutable record or concept revision that was reviewed.

Approval of one version does not automatically approve:

\begin{itemize}
\tightlist
\item
  a successor Information Unit
\item
  a new Framework Assignment
\item
  a new Project Concept revision
\item
  a remapped framework-template version
\end{itemize}

Only Phase-G-approved engineering information may enter Phases H through J.

\hypertarget{phase-f-integration}{%
\subsubsection{Phase F Integration}\label{phase-f-integration}}

P4 reuses the existing Phase F infrastructure for:

\begin{itemize}
\tightlist
\item
  team execution
\item
  persona instructions
\item
  LLM-provider access
\item
  raw agent artifacts
\item
  deterministic consensus analysis
\end{itemize}

P4 extends this infrastructure with:

\begin{itemize}
\tightlist
\item
  P4-specific structured result schemas
\item
  P4-specific teams and personas where required
\item
  field-level consensus comparison
\item
  inter-persona variance
\item
  intra-persona stability
\item
  semantic validators
\item
  controlled creation of P4 domain records
\end{itemize}

Existing Phase F reports and agent outputs are not automatically converted into P4 Information Units.

\hypertarget{model-generation-boundary}{%
\subsubsection{Model-Generation Boundary}\label{model-generation-boundary}}

P4 shall not:

\begin{itemize}
\tightlist
\item
  create approved engineering input
\item
  create model candidates
\item
  generate an internal engineering model
\item
  generate SysML v2 code
\item
  update CATIA
\item
  override the authoritative engineering model
\end{itemize}

Approved Input Promotion belongs to Phase G.

Model Candidate creation belongs to Phase H.

Internal model generation belongs to Phase I.

SysML v2 code generation belongs to Phase J.

Consequences

\hypertarget{positive-consequences}{%
\subsubsection{Positive Consequences}\label{positive-consequences}}

The architecture provides:

\begin{itemize}
\tightlist
\item
  deterministic traceability from source bytes to semantic interpretation
\item
  explicit separation of syntactic projection and semantic processing
\item
  stable project-local semantic identifiers
\item
  immutable and auditable Information Units
\item
  visible assumptions, ambiguity and disagreement
\item
  controlled use of external ontology references
\item
  project-specific terminology without global contamination
\item
  reproducible framework assignments
\item
  multiagent variance as an auditable uncertainty signal
\item
  separation of confidence and Engineering Approval
\item
  clear boundaries between P4, P5, P6 and Phase G
\item
  continued reuse of the implemented Phase F agent infrastructure
\item
  a bounded MVP without an ontology server or multimodal interpretation
\end{itemize}

\hypertarget{costs-and-limitations}{%
\subsubsection{Costs and Limitations}\label{costs-and-limitations}}

The architecture introduces:

\begin{itemize}
\tightlist
\item
  multiple semantic artifact types
\item
  additional validators and repositories
\item
  explicit ontology and vocabulary curation
\item
  more storage for immutable records and revisions
\item
  multiple LLM calls per semantic task
\item
  increased processing cost for multiagent execution
\item
  a Human Review requirement for terminology and ambiguous mappings
\item
  no support for visual engineering information
\item
  no automatic cross-source synthesis
\item
  no statistical calibration of confidence
\item
  no automatic ontology reasoning
\item
  no automatic cross-project terminology reuse
\end{itemize}

These costs are accepted in exchange for traceability, auditability and controlled semantic processing.

Alternatives Considered

\hypertarget{single-universal-mbse-ontology}{%
\subsubsection{Single Universal MBSE Ontology}\label{single-universal-mbse-ontology}}

Rejected.

No single reference system provides the complete combination of:

\begin{itemize}
\tightlist
\item
  top-level semantics
\item
  industrial concepts
\item
  MBSE and RFLP terminology
\item
  SysML v2 target semantics
\item
  project-specific vocabulary
\end{itemize}

The layered reference stack is more explicit and adaptable.

\hypertarget{bfo-or-iof-as-the-complete-project-vocabulary}{%
\subsubsection{BFO or IOF as the Complete Project Vocabulary}\label{bfo-or-iof-as-the-complete-project-vocabulary}}

Rejected.

BFO is too general and IOF Core is not a complete project or SysML vocabulary.

Both remain useful reference layers.

\hypertarget{sysdo-as-the-primary-runtime-ontology}{%
\subsubsection{SysDO as the Primary Runtime Ontology}\label{sysdo-as-the-primary-runtime-ontology}}

Rejected.

SysDO is useful as a non-normative research and structuring reference, but its demonstrated maturity and scope are insufficient for primary runtime authority.

\hypertarget{project-glossary-without-external-references}{%
\subsubsection{Project Glossary Without External References}\label{project-glossary-without-external-references}}

Rejected.

A project-only vocabulary would improve local consistency but reduce interoperability and make concept mappings harder to justify.

\hypertarget{live-ontology-service}{%
\subsubsection{Live Ontology Service}\label{live-ontology-service}}

Rejected for the MVP.

Live services would reduce reproducibility and introduce availability, versioning and security dependencies.

\hypertarget{full-owl-reasoner-or-triple-store}{%
\subsubsection{Full OWL Reasoner or Triple Store}\label{full-owl-reasoner-or-triple-store}}

Rejected for the MVP.

The accepted P4 use cases require controlled concept retrieval and explicit mapping, not unrestricted automated inference.

\hypertarget{loading-complete-ontologies-into-every-prompt}{%
\subsubsection{Loading Complete Ontologies Into Every Prompt}\label{loading-complete-ontologies-into-every-prompt}}

Rejected.

It would increase token use and reduce prompt focus without providing a controlled semantic guarantee.

\hypertarget{direct-original-file-processing-by-the-llm}{%
\subsubsection{Direct Original-File Processing by the LLM}\label{direct-original-file-processing-by-the-llm}}

Rejected.

It would make text extraction provider-dependent and weaken deterministic traceability.

\hypertarget{semantic-normalization-during-source-projection}{%
\subsubsection{Semantic Normalization During Source Projection}\label{semantic-normalization-during-source-projection}}

Rejected.

Spelling correction, synonym replacement, unit conversion and rewriting are semantic operations and could change source meaning.

\hypertarget{mutable-information-units}{%
\subsubsection{Mutable Information Units}\label{mutable-information-units}}

Rejected.

Overwriting extracted information would destroy auditability and make existing reviews and mappings ambiguous.

\hypertarget{automatic-cross-source-claim-merging}{%
\subsubsection{Automatic Cross-Source Claim Merging}\label{automatic-cross-source-claim-merging}}

Rejected.

Apparently equivalent statements may have different authority, wording, constraints or context.

They remain separate and traceable to their individual sources.

\hypertarget{framework-assignment-embedded-directly-in-the-information-unit}{%
\subsubsection{Framework Assignment Embedded Directly in the Information Unit}\label{framework-assignment-embedded-directly-in-the-information-unit}}

Rejected.

Framework mappings have their own rationale, confidence, lifecycle and framework-version binding.

They remain separate records.

\hypertarget{automatic-parent-or-level-propagation}{%
\subsubsection{Automatic Parent or Level Propagation}\label{automatic-parent-or-level-propagation}}

Rejected.

A semantic mapping to one node does not automatically prove relevance to its parent, siblings or another framework level.

\hypertarget{single-agent-confidence}{%
\subsubsection{Single-Agent Confidence}\label{single-agent-confidence}}

Rejected.

An individual LLM self-assessment is not a sufficiently auditable confidence basis.

Confidence is derived from independent persona results.

\hypertarget{treating-repeated-runs-as-independent-votes}{%
\subsubsection{Treating Repeated Runs as Independent Votes}\label{treating-repeated-runs-as-independent-votes}}

Rejected.

Repeated runs of the same persona measure stability but do not create additional independent professional perspectives.

\hypertarget{numeric-confidence-as-probability}{%
\subsubsection{Numeric Confidence as Probability}\label{numeric-confidence-as-probability}}

Rejected.

The persona team is not a statistically independent random sample.

The accepted confidence is an ordinal epistemic signal.

\hypertarget{automatic-terminology-acceptance}{%
\subsubsection{Automatic Terminology Acceptance}\label{automatic-terminology-acceptance}}

Rejected.

Terminology meaning requires explicit Human-in-the-Loop governance.

\hypertarget{engineering-approval-during-p4}{%
\subsubsection{Engineering Approval During P4}\label{engineering-approval-during-p4}}

Rejected.

P4 produces semantic and framework-mapping candidates.

Engineering Approval remains a separate Phase G responsibility.

Implementation Constraints

P4 implementation shall:

\begin{itemize}
\tightlist
\item
  preserve all P1 framework validation rules
\item
  preserve P2 project isolation
\item
  preserve P3 source immutability
\item
  enforce ADR-009 textual-modality restrictions
\item
  reject unsupported adapter behavior
\item
  reject cross-project semantic references
\item
  reject unknown framework targets
\item
  reject malformed semantic identifiers
\item
  reject direct context-only engineering contribution
\item
  reject silent overwrites
\item
  retain projection and semantic issues
\item
  retain agent disagreement
\item
  keep confidence separate from approval
\item
  keep terminology approval separate from Engineering Approval
\item
  remain compatible with the existing Phase F pipeline
\end{itemize}

P4 implementation shall not begin to depend on a different architecture without a new explicitly accepted Architecture Decision Record.

Verification Criteria

P4 is not complete until automated tests demonstrate at least:

\begin{itemize}
\tightlist
\item
  deterministic projection for every supported adapter
\item
  rejection of unsupported encodings
\item
  rejection of JSON duplicate keys
\item
  fixed CSV and TSV parsing behavior
\item
  PDF text-layer extraction with page traceability
\item
  visible partial and unavailable projection results
\item
  stable Source Projection and segment identifiers
\item
  stable Information Unit identifiers
\item
  Information Unit atomicity validation
\item
  same-source anchor enforcement
\item
  rejection of cross-source Information Unit synthesis
\item
  rejection of context-only Information Units
\item
  validation of all Information Types
\item
  validation of all Statement Modalities
\item
  validation of all Epistemic Classes
\item
  derivation-support reference validation
\item
  assumption missing-evidence validation
\item
  Project Concept identifier and revision validation
\item
  preferred-label uniqueness by language
\item
  Ambiguity Group validation
\item
  Terminology Decision validation
\item
  external ontology mapping validation
\item
  framework-template version validation
\item
  rejection of unknown framework targets
\item
  zero-to-many Framework Assignments
\item
  explicit unmapped, ambiguous and failed Mapping Evaluations
\item
  no automatic framework hierarchy propagation
\item
  deterministic consensus classification
\item
  separation of inter-persona and intra-persona variance
\item
  one vote per persona
\item
  deterministic ordinal confidence derivation
\item
  mandatory review for conflicting or low-confidence results
\item
  immutable semantic persistence
\item
  atomic record publication
\item
  rejection of cross-project semantic references
\item
  preservation of Phase F behavior
\item
  complete project test-suite compatibility
\end{itemize}

Related Decisions

\begin{itemize}
\tightlist
\item
  ADR-005 --- Project Workspace Architecture
\item
  ADR-009 --- Textual Source Processing Boundary
\item
  ADR-010 --- Project Source Registry Architecture
\end{itemize}

Implementation Status

Architecture accepted.

P4 implementation not started.

\textbackslash newpage

\hypertarget{adr-012--processing-state-and-artifact-organization}{%
\subsection{ADR-012 --- Processing State and Artifact Organization}\label{adr-012--processing-state-and-artifact-organization}}

Processing State and Artifact Organization

Status

Accepted

Date

2026-07-25

Context

Phase P introduces project-oriented processing around the completed Phase F agentic ingestion pipeline.

P1 implemented the versioned Turing RFLP Framework.

P2 implemented the persistent Project Workspace.

P3 implemented mandatory project assignment, immutable registered sources and explicit source roles.

P4 implemented deterministic Source Projections, source-traceable Information Units, semantic consensus, terminology mappings, Framework Assignments and Human Review Decisions.

The P1 through P4 repositories persist their individual artifacts independently. They do not yet provide:

\begin{itemize}
\tightlist
\item
  project-local Processing Run identity
\item
  a canonical operational processing state
\item
  an auditable state-transition history
\item
  project-local organization of agent outputs and Consensus Reports
\item
  retry, resumption and recovery behavior
\item
  supersession and invalidation behavior
\item
  source-level and project-level processing aggregation
\end{itemize}

P5 shall introduce these capabilities without replacing or duplicating the authority of existing Project, Source or Semantic Manifests.

P5 shall not perform Approved Input Promotion, model-candidate generation, internal model generation or SysML v2 generation.

Decision

\hypertarget{source-bound-processing-runs}{%
\subsubsection{Source-bound Processing Runs}\label{source-bound-processing-runs}}

A Source Processing Run processes exactly one primary registered source within exactly one project.

A Processing Run is not equivalent to one LLM invocation.

One Processing Run may contain:

\begin{itemize}
\tightlist
\item
  multiple processing stages
\item
  multiple agent executions
\item
  multiple persona runs
\item
  multiple LLM invocations
\item
  deterministic consensus calculations
\item
  Human Review waits
\item
  technical retries
\end{itemize}

Every Information Unit remains traceable to exactly one registered engineering source.

P5 shall not synthesize one Information Unit from multiple sources.

Equivalent, overlapping or contradictory statements from different sources remain separate, source-traceable Information Units.

Project-level comparison may later reference the results of multiple Source Processing Runs while preserving every individual source reference.

P5 provides the operational and traceability foundation for that comparison.

Project-wide coverage, overlap and conflict analysis belongs to P6.

The human decision concerning which reviewed information becomes Approved Input belongs to Phase G.

\hypertarget{source-processing-disposition}{%
\subsubsection{Source Processing Disposition}\label{source-processing-disposition}}

P5 introduces explicit project-local Processing Decisions for operational source treatment.

Supported Source Processing Dispositions are:

\begin{verbatim}
in_scope
context_only
out_of_scope
\end{verbatim}

\texttt{in\_scope} identifies a source that remains eligible for processing according to its registered source role.

\texttt{context\_only} restricts a source to contextual and terminology use.

An effectively context-only source shall not create:

\begin{itemize}
\tightlist
\item
  engineering Information Units
\item
  Framework Assignments
\item
  preliminary engineering coverage
\item
  approved readiness
\item
  model-generation input
\end{itemize}

\texttt{out\_of\_scope} identifies a registered source that does not belong to the currently processed system or project scope.

Out-of-scope sources remain registered and auditable. They are not deleted or silently ignored.

A Processing Decision is an operational Human-in-the-Loop decision.

It is not:

\begin{itemize}
\tightlist
\item
  Engineering Approval
\item
  Approved Input Promotion
\item
  a framework assignment
\item
  a terminology acceptance decision
\item
  model-generation authority
\end{itemize}

The existing P3 Source Manifest remains authoritative for the registered source and its stored source role.

A Processing Decision does not silently rewrite the Source Manifest.

\hypertarget{processing-run-identity}{%
\subsubsection{Processing Run Identity}\label{processing-run-identity}}

Every Processing Run receives an immutable project-local identifier named \texttt{processing\_run\_id}.

The identifier:

\begin{itemize}
\tightlist
\item
  matches \texttt{\^{}RUN-{[}0-9{]}\{6\}\$}
\item
  is unique within one project
\item
  is allocated sequentially
\item
  is not reused
\item
  remains unchanged for the lifetime of the run
\end{itemize}

The globally unambiguous reference is:

\begin{verbatim}
<project_id>/<processing_run_id>
\end{verbatim}

Example:

\begin{verbatim}
318604/RUN-000001
\end{verbatim}

\hypertarget{run-manifest}{%
\subsubsection{Run Manifest}\label{run-manifest}}

Every Processing Run contains an immutable \texttt{run\_manifest.json}.

The Run Manifest binds the run to at least:

\begin{verbatim}
schema_version
project_id
processing_run_id
source_id
source_sha256
source_role_snapshot
workflow_profile
configuration_fingerprint
framework_template_id
framework_template_version
semantic_reference_versions
created_at
supersedes_run_id
\end{verbatim}

\texttt{supersedes\_run\_id} is absent when the run has no predecessor.

The Run Manifest contains identity, scope and reproducibility bindings.

It does not contain a mutable current-state field.

\hypertarget{workflow-profiles}{%
\subsubsection{Workflow Profiles}\label{workflow-profiles}}

The required stages of one Processing Run are determined by an explicit, versioned workflow profile.

An engineering-source workflow may include:

\begin{verbatim}
source_projection
semantic_extraction
semantic_consensus
terminology_mapping
framework_assignment
human_review
publication
\end{verbatim}

A context-only workflow may omit engineering Information Unit extraction and Framework Assignment.

A run shall not be considered incomplete merely because a stage that is ineligible for its workflow profile was not executed.

\hypertarget{event-history}{%
\subsubsection{Event History}\label{event-history}}

The authoritative operational state of a Processing Run is represented by an immutable Event History.

Every state transition creates a new event.

Existing events shall not be overwritten, reordered or deleted.

Event identifiers:

\begin{itemize}
\tightlist
\item
  match \texttt{\^{}EVT-{[}0-9{]}\{6\}\$}
\item
  are unique within one Processing Run
\item
  are allocated sequentially
\item
  are not reused
\end{itemize}

An event records at least:

\begin{verbatim}
schema_version
project_id
processing_run_id
event_id
event_sequence
previous_state
next_state
processing_stage
event_type
attempt_id
reason_code
artifact_references
occurred_at
previous_event_fingerprint
event_fingerprint
\end{verbatim}

The fingerprint chain makes missing, reordered or modified events detectable.

An artifact reference contains only the information required to identify and validate the referenced artifact, including:

\begin{verbatim}
artifact_type
artifact_id
content_fingerprint
repository_relative_path
\end{verbatim}

The event does not duplicate the professional content of the referenced artifact.

\hypertarget{current-state-projection}{%
\subsubsection{Current-state Projection}\label{current-state-projection}}

The current Processing Run state is derived deterministically from the valid Event History.

The derived current-state view is not an independent authority.

It may be calculated during project loading or exposed through a regenerable derived index.

A derived index may be deleted and regenerated without loss of authoritative processing information.

\hypertarget{run-states}{%
\subsubsection{Run States}\label{run-states}}

The canonical Processing Run states are:

\begin{verbatim}
created
running
awaiting_review
blocked
failed
completed
superseded
\end{verbatim}

\texttt{created} means that the run has been persistently created but processing has not started.

\texttt{running} means that an allowed processing stage is being executed.

\texttt{awaiting\_review} means that an exact Human Review decision is required before processing may continue.

\texttt{blocked} means that a known deterministic prerequisite is missing or invalid.

Examples include:

\begin{itemize}
\tightlist
\item
  unavailable Source Projection
\item
  required context exceeding the token budget
\item
  invalid mandatory references
\item
  unresolved processing disposition
\item
  incomplete publication recovery
\item
  inconsistent Event History
\end{itemize}

\texttt{failed} means that a technical execution or artifact-validation failure occurred.

\texttt{completed} means that the complete workflow required by the selected workflow profile has been resolved.

\texttt{completed} does not mean:

\begin{itemize}
\tightlist
\item
  Engineering Approval
\item
  Approved Input
\item
  generation readiness
\item
  model acceptance
\end{itemize}

A completely reviewed rejection may result in a completed run without publication.

\texttt{superseded} means that an explicitly linked successor run replaces the run for current operational use.

A superseded run remains fully available for traceability.

\hypertarget{processing-stages}{%
\subsubsection{Processing Stages}\label{processing-stages}}

Run state and active Processing Stage are separate dimensions.

Supported P5 stages include:

\begin{verbatim}
source_projection
semantic_extraction
semantic_consensus
terminology_mapping
framework_assignment
human_review
publication
\end{verbatim}

This separation prevents stage-specific combinations from becoming independent state values.

\hypertarget{allowed-state-transitions}{%
\subsubsection{Allowed State Transitions}\label{allowed-state-transitions}}

The allowed core transitions are:

\begin{verbatim}
created
→ running
→ blocked
→ failed
→ superseded

running
→ awaiting_review
→ blocked
→ failed
→ completed
→ superseded

awaiting_review
→ running
→ completed
→ blocked
→ superseded

blocked
→ running
→ failed
→ superseded

failed
→ running
→ superseded

completed
→ superseded

superseded
→ no further transition
\end{verbatim}

Every transition requires an explicit reason and valid transition evidence.

\hypertarget{human-review-transitions}{%
\subsubsection{Human Review Transitions}\label{human-review-transitions}}

Human Review controls operational continuation but does not replace Phase G Engineering Approval.

The supported behavior is:

\begin{verbatim}
confirm
→ continue processing or publication

reject
→ do not publish the rejected target
→ complete the run when all required targets are resolved

request_changes
→ return to the required processing stage
→ create a new attempt when the run bindings remain unchanged
→ create a successor run when the bindings changed
\end{verbatim}

Only an exact decision bound to the current target and validation fingerprints may control the corresponding transition.

Consensus, confidence and variance remain review evidence.

They shall not create an automatic state transition that bypasses Human Review.

\hypertarget{project-local-artifact-organization}{%
\subsubsection{Project-local Artifact Organization}\label{project-local-artifact-organization}}

P5 introduces:

\begin{verbatim}
data/projects/<project_id>/
├── runs/
│   └── RUN-000001/
│       ├── run_manifest.json
│       ├── events/
│       │   ├── EVT-000001.json
│       │   └── EVT-000002.json
│       ├── artifacts/
│       │   ├── agent_outputs/
│       │   └── consensus_reports/
│       └── work/
└── processing_decisions/
    └── PD-000001.json
\end{verbatim}

The directories are created only when required.

\hypertarget{existing-artifact-authority}{%
\subsubsection{Existing Artifact Authority}\label{existing-artifact-authority}}

Existing P2 through P4 repositories remain authoritative for their own domain records.

These include:

\begin{verbatim}
project_manifest.json
sources/<source_id>/source_manifest.json
sources/<source_id>/content.<suffix>
semantics/source_projections/
semantics/information_units/
semantics/project_glossary.json
semantics/terminology_decisions/
semantics/terminology_mappings/
semantics/framework_assignments/
semantics/human_reviews/
\end{verbatim}

P5 shall not:

\begin{itemize}
\tightlist
\item
  move these artifacts into Processing Run directories
\item
  copy them into Processing Run directories as a second maintained record
\item
  rewrite their professional content
\item
  replace their identifiers
\item
  replace their fingerprints
\item
  weaken their individual validation contracts
\end{itemize}

P5 references these artifacts through stable identifiers and fingerprints.

\hypertarget{run-owned-artifacts}{%
\subsubsection{Run-owned Artifacts}\label{run-owned-artifacts}}

\texttt{artifacts/agent\_outputs/} stores immutable technical outputs of individual agent and persona executions.

\texttt{artifacts/consensus\_reports/} stores deterministic consensus and variance evidence.

Run-owned artifacts remain traceable to:

\begin{itemize}
\tightlist
\item
  project
\item
  Processing Run
\item
  primary source
\item
  processing stage
\item
  attempt
\item
  team
\item
  agent
\item
  persona
\item
  provider and model
\item
  prompt and schema version
\end{itemize}

These technical artifacts are execution evidence.

They are not authoritative engineering information.

\hypertarget{temporary-work-artifacts}{%
\subsubsection{Temporary Work Artifacts}\label{temporary-work-artifacts}}

\texttt{work/} contains incomplete and non-authoritative processing artifacts.

Temporary work artifacts shall not:

\begin{itemize}
\tightlist
\item
  satisfy run completion
\item
  create preliminary coverage
\item
  pass a publication gate
\item
  be treated as published semantic records
\end{itemize}

Temporary artifacts may be retained for failure diagnosis.

They shall not be silently promoted.

\hypertarget{phase-f-run-boundary}{%
\subsubsection{Phase F Run Boundary}\label{phase-f-run-boundary}}

Existing Phase F artifacts under:

\begin{verbatim}
data/team_runs/
\end{verbatim}

remain available for existing demonstrations, regression evidence and Phase F inspection.

They do not become authoritative Project Workspace state.

New project-oriented P5 Processing Runs are stored only below their containing project.

P5 shall not require destructive migration or deletion of existing Phase F artifacts.

\hypertarget{attempt-identity}{%
\subsubsection{Attempt Identity}\label{attempt-identity}}

Every retry attempt receives an immutable attempt identifier matching:

\begin{verbatim}
^ATT-[0-9]{6}$
\end{verbatim}

Attempts are ordered within their Processing Run and stage.

Attempt artifacts are stored separately.

Example:

\begin{verbatim}
artifacts/agent_outputs/
└── semantic_extraction/
    ├── ATT-000001/
    └── ATT-000002/
\end{verbatim}

No retry overwrites output from an earlier attempt.

\hypertarget{retry-behavior}{%
\subsubsection{Retry Behavior}\label{retry-behavior}}

A retry may remain within the same Processing Run only when all material run bindings remain unchanged.

These bindings include:

\begin{verbatim}
source_id
source_sha256
source_role_snapshot
workflow_profile
adapter configuration
prompt and schema versions
framework-template version
semantic reference versions
\end{verbatim}

A retry creates:

\begin{itemize}
\tightlist
\item
  a new Attempt
\item
  new technical artifacts
\item
  new state-transition events
\item
  explicit retry rationale
\end{itemize}

A retry shall not replace or edit earlier attempts.

\hypertarget{successor-runs}{%
\subsubsection{Successor Runs}\label{successor-runs}}

A material change to a run binding requires a new Processing Run.

Examples include:

\begin{itemize}
\tightlist
\item
  changed source content
\item
  changed source identity
\item
  changed effective source role
\item
  changed workflow profile
\item
  changed adapter configuration
\item
  changed prompt or output schema
\item
  changed framework-template version
\item
  changed semantic reference version
\end{itemize}

The new Run Manifest references its predecessor through \texttt{supersedes\_run\_id}.

A completed run is never reopened for a material change.

\hypertarget{supersession}{%
\subsubsection{Supersession}\label{supersession}}

Supersession follows this order:

\begin{enumerate}
\def\labelenumi{\arabic{enumi}.}
\tightlist
\item
  validate the predecessor
\item
  create and validate the successor Run Manifest
\item
  persist the successor run
\item
  append the supersession event to the predecessor
\item
  validate the complete relationship
\end{enumerate}

An interrupted supersession is detected during project reopening.

It is not silently completed or ignored.

A run marked as superseded remains immutable and available.

\hypertarget{artifact-lifecycle}{%
\subsubsection{Artifact Lifecycle}\label{artifact-lifecycle}}

P5 may derive the following operational lifecycle states for existing immutable artifacts:

\begin{verbatim}
active
superseded
invalidated
\end{verbatim}

Lifecycle state remains separate from artifact content.

Supersession or invalidation:

\begin{itemize}
\tightlist
\item
  does not delete the artifact
\item
  does not alter its content
\item
  does not alter its original Human Review Decision
\item
  does not turn a rejected artifact into an accepted artifact
\item
  does not create Approved Input
\end{itemize}

A lifecycle event references the exact artifact identity and fingerprint.

\hypertarget{source-disposition-changes}{%
\subsubsection{Source-disposition Changes}\label{source-disposition-changes}}

A Processing Decision that changes the effective treatment of a source requires dependency analysis.

When a source becomes effectively \texttt{context\_only} or \texttt{out\_of\_scope}, dependent engineering runs and their use in later P6 processing are invalidated.

The affected records remain available for traceability.

A source-disposition change shall not silently reinterpret existing artifacts.

\hypertarget{project-reopening}{%
\subsubsection{Project Reopening}\label{project-reopening}}

Project reopening validates at least:

\begin{itemize}
\tightlist
\item
  Project Manifest
\item
  Source Manifests
\item
  Run Manifests
\item
  Event identities and sequences
\item
  event fingerprint chains
\item
  allowed state transitions
\item
  Attempt identities
\item
  artifact references and fingerprints
\item
  Processing Decisions
\item
  supersession relationships
\item
  incomplete publication conditions
\end{itemize}

A valid history is reconstructed deterministically.

No project state is reconstructed from filenames alone.

\hypertarget{recovery-behavior}{%
\subsubsection{Recovery Behavior}\label{recovery-behavior}}

Recovery is explicit and fail-closed.

A consistent history produces the derived current state.

A recoverable interrupted operation produces:

\begin{verbatim}
blocked
\end{verbatim}

together with an explicit recovery diagnostic.

A technical execution failure produces:

\begin{verbatim}
failed
\end{verbatim}

An inconsistent, incomplete or manipulated Event History produces:

\begin{verbatim}
blocked
\end{verbatim}

and shall not be silently repaired.

A published artifact that exists without the expected completion event is reported as an incomplete publication.

An explicit recovery operation may validate the exact existing artifact and complete the Event History.

Recovery shall not create a duplicate artifact.

\hypertarget{source-level-aggregation}{%
\subsubsection{Source-level Aggregation}\label{source-level-aggregation}}

For each source, P5 identifies the current non-superseded Processing Run.

The source-processing view includes at least:

\begin{verbatim}
processing disposition
current processing run
current run state
current processing stage
latest attempt
blocking issues
failure issues
pending Human Review
superseded runs
invalidated artifacts
\end{verbatim}

\hypertarget{project-level-aggregation}{%
\subsubsection{Project-level Aggregation}\label{project-level-aggregation}}

Project Processing State is derived from validated project, source, run, event and Processing Decision records.

The supported project states are:

\begin{verbatim}
empty
not_started
in_progress
awaiting_review
attention_required
partially_processed
processed
\end{verbatim}

\texttt{empty} means that the project contains no registered sources.

\texttt{not\_started} means that relevant sources exist but no active processing has started.

\texttt{in\_progress} means that one or more active Processing Runs are being executed.

\texttt{awaiting\_review} means that an active run requires Human Review.

\texttt{attention\_required} means that one or more relevant sources are blocked, failed or inconsistent.

\texttt{partially\_processed} means that some relevant sources are completed while others remain unprocessed.

\texttt{processed} means that all relevant source-processing workflows have reached resolved terminal outcomes.

Out-of-scope sources remain visible but do not prevent a project from becoming processed.

\hypertarget{aggregation-counts}{%
\subsubsection{Aggregation Counts}\label{aggregation-counts}}

The project aggregation exposes separate counts for at least:

\begin{verbatim}
total sources
in-scope sources
context-only sources
out-of-scope sources
not-started sources
running sources
sources awaiting review
blocked sources
failed sources
completed sources
superseded runs
invalidated artifacts
\end{verbatim}

A headline Project Processing State shall never hide its underlying counts and issues.

\hypertarget{aggregation-authority-boundary}{%
\subsubsection{Aggregation Authority Boundary}\label{aggregation-authority-boundary}}

Project Processing State is an operational and dashboard-oriented view.

It is not:

\begin{itemize}
\tightlist
\item
  preliminary coverage
\item
  Approved Input
\item
  approved readiness
\item
  Engineering Approval
\item
  generation readiness
\item
  model-generation authority
\end{itemize}

P6 calculates Preliminary Coverage from eligible P4 records that P5 identifies as operationally usable.

Phase G performs Approved Input Promotion using exact Human Review evidence.

\hypertarget{p5-scope-boundary}{%
\subsubsection{P5 Scope Boundary}\label{p5-scope-boundary}}

P5 includes:

\begin{itemize}
\tightlist
\item
  Source Processing Run identity
\item
  immutable Run Manifests
\item
  immutable Event History
\item
  deterministic current-state derivation
\item
  project-local agent-output organization
\item
  project-local Consensus Report organization
\item
  Processing Decisions
\item
  retry and Attempt organization
\item
  supersession and invalidation
\item
  reopening and explicit recovery behavior
\item
  source-level processing aggregation
\item
  project-level processing aggregation
\end{itemize}

P5 does not include:

\begin{itemize}
\tightlist
\item
  automatic multi-source Information Unit synthesis
\item
  Approved Input Promotion
\item
  Engineering Approval
\item
  coverage calculation
\item
  model-candidate creation
\item
  internal model generation
\item
  SysML v2 generation
\item
  CATIA updates
\end{itemize}

Consequences

Positive consequences:

\begin{itemize}
\tightlist
\item
  Every processing transition remains auditable.
\item
  Current state can be reconstructed deterministically.
\item
  Multiple LLM and persona executions remain grouped within one source-bound workflow.
\item
  Multiple sources remain independently traceable.
\item
  Overlap and contradictions can later be compared without merging source authority.
\item
  Existing P2 through P4 repositories remain stable.
\item
  Retry does not destroy earlier execution evidence.
\item
  Material changes produce explicit successor runs.
\item
  Failed and interrupted processing can be diagnosed safely.
\item
  Out-of-scope sources remain visible without blocking project progress.
\item
  Project processing state can support the future dashboard.
\item
  Processing completion remains separate from Engineering Approval.
\end{itemize}

Trade-offs:

\begin{itemize}
\tightlist
\item
  Event persistence introduces additional artifacts.
\item
  Current-state reconstruction requires deterministic event validation.
\item
  Retry and supersession require strict fingerprint comparison.
\item
  Run-owned technical artifacts increase project storage.
\item
  Recovery requires explicit operator action.
\item
  Project aggregation must preserve detailed counts and issues.
\item
  Source-disposition changes require dependency analysis.
\end{itemize}

Alternatives Considered

Treating one Processing Run as one LLM invocation was rejected because one source workflow requires multiple stages, personas and deterministic analyses.

Creating multi-source Information Units was rejected because it would weaken source traceability and make contradictory evidence difficult to review.

Storing only a mutable current-state file was rejected because it would destroy transition, retry and recovery history.

Using only an Event History without a derived current-state view was rejected because project reopening and the future dashboard require efficient state access.

Copying existing semantic artifacts into Run directories was rejected because it would duplicate authority and create synchronization risk.

Moving existing P4 artifacts into Run directories was rejected because it would break established repositories and public contracts.

Using \texttt{data/team\_runs/} as Project Workspace authority was rejected because those artifacts are not project-isolated.

Overwriting earlier Agent Outputs during retry was rejected because it would destroy reproducibility evidence.

Reopening a completed run for material changes was rejected because it would make previous review and publication evidence ambiguous.

Automatically repairing invalid Event Histories was rejected because recovery must remain explicit and auditable.

Using Project Processing State as approval or generation authority was rejected because operational completion and Engineering Approval are separate concerns.

Affected Components

\begin{itemize}
\tightlist
\item
  future P5 processing-state modules
\item
  \texttt{data/projects/\textless{}project\_id\textgreater{}/runs/}
\item
  \texttt{data/projects/\textless{}project\_id\textgreater{}/processing\_decisions/}
\item
  existing P2 through P4 repositories as referenced dependencies
\item
  future P6 coverage processing
\item
  future P7 Project Dashboard
\item
  P5 automated tests and phase review
\end{itemize}

Implementation Constraints

P5 implementation shall:

\begin{itemize}
\tightlist
\item
  preserve P2 project isolation
\item
  preserve P3 source identity and source immutability
\item
  preserve P4 artifact immutability
\item
  reject cross-project run and artifact references
\item
  reject invalid identifiers
\item
  reject invalid state transitions
\item
  reject broken event sequences
\item
  reject invalid fingerprint chains
\item
  reject silent artifact overwrite
\item
  preserve earlier attempts
\item
  preserve superseded and invalidated artifacts
\item
  separate processing completion from Engineering Approval
\item
  remain compatible with existing Phase F behavior
\item
  avoid destructive migration of existing Phase F artifacts
\end{itemize}

Verification Criteria

P5 is not complete until automated tests demonstrate at least:

\begin{itemize}
\tightlist
\item
  project-local Processing Run identifier allocation
\item
  strict Run Manifest validation
\item
  exactly one primary source per Source Processing Run
\item
  multiple Agent and LLM executions within one run
\item
  workflow-profile validation
\item
  immutable Event persistence
\item
  deterministic event ordering
\item
  event fingerprint-chain validation
\item
  rejection of invalid state transitions
\item
  deterministic current-state reconstruction
\item
  Processing Decision validation
\item
  source-disposition behavior
\item
  separate Attempt persistence
\item
  retry with unchanged bindings
\item
  rejection of same-run retry with changed bindings
\item
  successor-run creation for changed bindings
\item
  supersession validation
\item
  artifact invalidation without deletion
\item
  recovery from incomplete publication
\item
  blocking of inconsistent Event Histories
\item
  source-level aggregation
\item
  project-level aggregation
\item
  out-of-scope exclusion from project-progress blocking
\item
  separation of Project Processing State and approval
\item
  preservation of P1 through P4 tests
\item
  complete project test-suite compatibility
\end{itemize}

Supersedes

None

Related Roadmap Phase

P5 --- Processing State and Artifact Organization

Related Decisions

\begin{itemize}
\tightlist
\item
  ADR-005 --- Project Workspace Architecture
\item
  ADR-009 --- Textual Source Processing Boundary
\item
  ADR-010 --- Project Source Registry Architecture
\item
  ADR-011 --- Semantic Information Unit and Ontology Boundary
\end{itemize}

Related Implementation

Not yet implemented.

\textbackslash newpage

\hypertarget{adr-013--preliminary-coverage-and-potential-model-support-assessment}{%
\subsection{ADR-013 --- Preliminary Coverage and Potential Model Support Assessment}\label{adr-013--preliminary-coverage-and-potential-model-support-assessment}}

Preliminary Coverage and Potential Model Support Assessment

Status

Accepted

Date

2026-07-27

Context

Phase P introduces project-oriented processing around the completed Phase F agentic ingestion pipeline.

P1 implemented the versioned Turing RFLP Framework with three levels and twelve stable mapping targets.

P2 implemented the persistent Project Workspace.

P3 implemented mandatory project assignment, immutable registered sources and the explicit source roles \texttt{engineering\_source} and \texttt{context\_only}.

P4 implemented source-traceable Information Units, Framework Assignment Candidates and exact Human Review Decisions.

P5 implemented project-local Processing Runs, immutable Event Histories, Processing Decisions, artifact lifecycle, source disposition, recovery behavior and source-level and project-level processing aggregation.

The existing artifacts establish traceable candidate evidence. They do not yet provide a deterministic answer to the following project-level questions:

\begin{itemize}
\tightlist
\item
  which framework nodes have preliminary evidence
\item
  which framework nodes remain uncovered
\item
  which nodes contain ambiguous or conflicting evidence
\item
  which framework levels are partially or completely covered
\item
  which model scopes may be potentially supported by the available evidence
\item
  whether a result is candidate evidence or reviewed candidate evidence
\item
  whether Approved Generation Readiness is currently available
\end{itemize}

P6 introduces deterministic Preliminary Coverage and potential model-support assessment.

P6 shall not perform:

\begin{itemize}
\tightlist
\item
  Approved Input Promotion
\item
  Approved Generation Readiness assessment
\item
  model-candidate generation
\item
  internal model generation
\item
  SysML v2 generation
\item
  engineering approval
\item
  automatic CATIA model mutation
\end{itemize}

Approved Input Promotion and Approved Generation Readiness remain assigned to Phase G and later responsible phases.

Decision

\hypertarget{assessment-authority}{%
\subsubsection{Assessment Authority}\label{assessment-authority}}

Preliminary Coverage is a derived, non-authoritative project assessment.

It is calculated from validated project records and may be deleted and regenerated without loss of authoritative engineering or processing information.

The assessment shall not replace or modify the authority of:

\begin{itemize}
\tightlist
\item
  the CATIA SysML v2 engineering model
\item
  the Project Manifest
\item
  registered Source Manifests
\item
  Processing Run Manifests and Event Histories
\item
  Information Unit Manifests
\item
  Framework Assignment Candidates
\item
  Human Review Decisions
\item
  accepted architecture decisions
\end{itemize}

Implementation reality and derived coverage evidence do not automatically become normative engineering requirements.

\hypertarget{preliminary-coverage-and-approved-readiness}{%
\subsubsection{Preliminary Coverage and Approved Readiness}\label{preliminary-coverage-and-approved-readiness}}

P6 distinguishes strictly between:

\begin{verbatim}
Preliminary Coverage
Approved Generation Readiness
\end{verbatim}

Preliminary Coverage:

\begin{itemize}
\tightlist
\item
  is available in Phase P
\item
  is based on eligible candidate evidence
\item
  does not require Human Review confirmation
\item
  may distinguish unreviewed and reviewed candidate evidence
\item
  does not authorize Approved Input Promotion
\item
  does not authorize model generation
\end{itemize}

Approved Generation Readiness:

\begin{itemize}
\tightlist
\item
  is not available in Phase P
\item
  is not calculated by P6
\item
  becomes available from Phase G or its responsible successor phase
\item
  requires approved engineering input
\item
  requires the responsible Human-in-the-Loop authority
\end{itemize}

P6 shall expose:

\begin{verbatim}
approved_readiness_status = not_available
approved_readiness_available_from_phase = G
\end{verbatim}

P6 shall not expose \texttt{approved\_readiness\ =\ false}, because \texttt{false} could imply that an available readiness assessment was executed and failed.

\hypertarget{assessment-inputs}{%
\subsubsection{Assessment Inputs}\label{assessment-inputs}}

A project assessment may reference only validated project-local records.

Required input categories are:

\begin{verbatim}
Framework Template
Preliminary Support Profile
registered Sources
effective Source Processing Dispositions
Processing Run and artifact lifecycle state
Framework Assignment Candidates
Human Review Decisions
\end{verbatim}

Every referenced artifact shall remain bound through its stable identifier and content fingerprint where the source contract provides one.

Cross-project references are invalid.

Unknown sources, framework nodes, candidates, decisions, runs or fingerprints are invalid.

Blocking persistence or reference issues shall remain visible in the derived assessment.

\hypertarget{eligible-sources}{%
\subsubsection{Eligible Sources}\label{eligible-sources}}

A source may contribute Preliminary Coverage only when all of the following are true:

\begin{enumerate}
\def\labelenumi{\arabic{enumi}.}
\tightlist
\item
  the source belongs to the assessed project
\item
  the registered source role is \texttt{engineering\_source}
\item
  the effective P5 Source Processing Disposition is \texttt{in\_scope}
\item
  the source fingerprint matches the referenced processing and semantic records
\item
  the relevant Processing Run and artifact references are valid
\item
  the referenced evidence is not invalidated
\end{enumerate}

A source shall not contribute Preliminary Coverage when its effective disposition is:

\begin{verbatim}
context_only
out_of_scope
\end{verbatim}

A Processing Decision shall not elevate a source registered as \texttt{context\_only} into an engineering-coverage source.

Context-only and out-of-scope sources remain registered and auditable.

\hypertarget{eligible-framework-assignment-evidence}{%
\subsubsection{Eligible Framework Assignment Evidence}\label{eligible-framework-assignment-evidence}}

One Framework Assignment Candidate may contribute Preliminary Coverage only when all of the following are true:

\begin{enumerate}
\def\labelenumi{\arabic{enumi}.}
\tightlist
\item
  \texttt{project\_id} matches the assessed project
\item
  \texttt{source\_id} references an eligible source
\item
  \texttt{information\_unit\_id} references the exact source-traceable Information Unit
\item
  the Framework Template identifier and version match the assessment template
\item
  \texttt{assignment\_status} is \texttt{assigned}
\item
  every counted proposal references a valid mapping-target node
\item
  the candidate and its referenced artifacts are not invalidated
\item
  the latest exact Human Review Decision does not reject the candidate or request changes
\end{enumerate}

The candidate confidence, consensus and variance values remain review evidence.

They shall not independently authorize:

\begin{itemize}
\tightlist
\item
  coverage
\item
  publication
\item
  Approved Input Promotion
\item
  generation readiness
\item
  model generation
\end{itemize}

Confidence, consensus and variance may be displayed as supporting evidence but shall not be converted into a numeric maturity or readiness score.

\hypertarget{non-covering-assignment-states}{%
\subsubsection{Non-covering Assignment States}\label{non-covering-assignment-states}}

Framework Assignment Candidates with these states shall not create coverage:

\begin{verbatim}
unassigned
ambiguous
conflict
\end{verbatim}

Their interpretation is:

\begin{verbatim}
unassigned
→ explicit uncovered evidence

ambiguous
→ no coverage and attention evidence

conflict
→ no coverage and attention evidence
\end{verbatim}

A node may have valid covering evidence and separate ambiguous or conflicting evidence at the same time.

Coverage and attention are therefore separate dimensions.

\hypertarget{human-review-resolution}{%
\subsubsection{Human Review Resolution}\label{human-review-resolution}}

Preliminary Coverage does not require Human Review confirmation.

The latest exact Human Review Decision bound to the candidate content and reference-validation fingerprints determines the review state of that candidate.

The effects are:

\begin{verbatim}
no exact decision
→ unreviewed candidate evidence
→ may create candidate coverage

confirm
→ reviewed candidate evidence
→ may create reviewed candidate coverage

reject
→ excluded from coverage

request_changes
→ excluded from coverage
→ attention required
\end{verbatim}

A confirmation in P6 confirms only the exact Framework Assignment Candidate.

It does not mean:

\begin{itemize}
\tightlist
\item
  Approved Input
\item
  Engineering Approval
\item
  Approved Generation Readiness
\item
  model acceptance
\item
  generation authorization
\end{itemize}

Stale decisions bound to an older candidate or validation fingerprint shall not control current coverage.

\hypertarget{framework-node-coverage}{%
\subsubsection{Framework Node Coverage}\label{framework-node-coverage}}

P6 assesses every mapping-target node defined by the active Framework Template.

The canonical node coverage states are:

\begin{verbatim}
uncovered
candidate_covered
reviewed_candidate_covered
\end{verbatim}

\texttt{uncovered} means that no eligible covering candidate evidence exists.

\texttt{candidate\_covered} means that at least one eligible unreviewed Framework Assignment Candidate maps an Information Unit to the node.

\texttt{reviewed\_candidate\_covered} means that at least one eligible Framework Assignment Candidate with a latest exact \texttt{confirm} decision maps an Information Unit to the node.

When both reviewed and unreviewed eligible evidence exist, \texttt{reviewed\_candidate\_covered} is the displayed coverage state while all evidence counts remain available.

Attention is represented separately:

\begin{verbatim}
attention_required = true | false
\end{verbatim}

Attention may result from:

\begin{itemize}
\tightlist
\item
  ambiguous assignment evidence
\item
  conflicting assignment evidence
\item
  \texttt{request\_changes}
\item
  invalid or stale references
\item
  invalidated evidence
\item
  blocking source, processing, assignment or review issues
\item
  multiple incompatible current evidence chains
\end{itemize}

One node may therefore be covered and require attention simultaneously.

\hypertarget{framework-node-coverage-record}{%
\subsubsection{Framework Node Coverage Record}\label{framework-node-coverage-record}}

The derived Framework Node Coverage record contains at least:

\begin{verbatim}
framework_node_id
mapping_key
node_name
level_node_id
coverage_state
attention_required
eligible_source_count
information_unit_count
assignment_candidate_count
confirmed_candidate_count
unreviewed_candidate_count
rejected_candidate_count
ambiguous_candidate_count
conflicting_candidate_count
source_ids
information_unit_ids
framework_assignment_candidate_ids
human_review_decision_ids
issue_codes
\end{verbatim}

Identifiers shall be unique and deterministically sorted.

Counts shall equal the corresponding unique reference sets.

Rejected or request-changes candidates remain auditable but shall not be counted as covering candidates.

\hypertarget{framework-level-coverage}{%
\subsubsection{Framework Level Coverage}\label{framework-level-coverage}}

The Framework Template defines three assessment levels:

\begin{verbatim}
Stakeholder Level
System Level
Subsystem Level
\end{verbatim}

Each level contains four mapping-target nodes.

The canonical level coverage states are:

\begin{verbatim}
uncovered
partially_covered
covered
\end{verbatim}

\texttt{uncovered} means that none of the level's mapping-target nodes has candidate or reviewed candidate coverage.

\texttt{partially\_covered} means that at least one but not all mapping-target nodes have candidate or reviewed candidate coverage.

\texttt{covered} means that every mapping-target node in the level has candidate or reviewed candidate coverage.

The level assessment contains at least:

\begin{verbatim}
level_node_id
level_name
coverage_state
covered_node_count
total_node_count
candidate_covered_node_count
reviewed_candidate_covered_node_count
attention_node_count
covered_node_ids
uncovered_node_ids
attention_node_ids
\end{verbatim}

P6 shall not calculate a weighted maturity or readiness percentage.

A UI may display exact ratios such as:

\begin{verbatim}
3 of 4 framework nodes preliminarily covered
\end{verbatim}

\hypertarget{project-coverage}{%
\subsubsection{Project Coverage}\label{project-coverage}}

The project assessment aggregates all Framework Node and Framework Level Coverage records.

The canonical project coverage states are:

\begin{verbatim}
uncovered
partially_covered
covered
attention_required
\end{verbatim}

The displayed project state follows this precedence:

\begin{enumerate}
\def\labelenumi{\arabic{enumi}.}
\tightlist
\item
  \texttt{attention\_required} when a blocking assessment issue exists
\item
  \texttt{covered} when every mapping-target node is covered and no blocking issue exists
\item
  \texttt{partially\_covered} when at least one but not every node is covered
\item
  \texttt{uncovered} when no node is covered
\end{enumerate}

Non-blocking node attention remains visible even when the overall project state is \texttt{partially\_covered} or \texttt{covered}.

The project assessment shall preserve exact counts and identifiers rather than derive a synthetic maturity score.

\hypertarget{preliminary-support-profile}{%
\subsubsection{Preliminary Support Profile}\label{preliminary-support-profile}}

Potential model support shall be defined by a separate versioned profile.

The initial profile is:

\begin{verbatim}
context/frameworks/turing_preliminary_support_profile.json
\end{verbatim}

The profile is not engineering-model authority.

It defines deterministic dependency rules for the derived P6 support assessment.

The profile shall bind at least:

\begin{verbatim}
schema_version
profile_id
profile_version
framework_template_id
framework_template_version
support_targets
\end{verbatim}

Each support target contains at least:

\begin{verbatim}
support_target_id
name
support_target_type
required_framework_node_ids
required_support_target_ids
\end{verbatim}

Unknown or duplicate references are invalid.

A profile bound to another Framework Template identifier or version is invalid.

\hypertarget{conservative-support-chain}{%
\subsubsection{Conservative Support Chain}\label{conservative-support-chain}}

The initial potential-support chain is:

\begin{verbatim}
Stakeholder Model
→ System Model
→ Subsystem Model
\end{verbatim}

\hypertarget{stakeholder-model}{%
\paragraph{Stakeholder Model}\label{stakeholder-model}}

Potential support requires coverage for:

\begin{verbatim}
Stakeholders
User Needs
Stakeholder Requirements
Use Cases
\end{verbatim}

\hypertarget{system-model}{%
\paragraph{System Model}\label{system-model}}

Potential support requires:

\begin{verbatim}
potential Stakeholder Model support
System Requirements
System Functional
System Logical
System Physical
\end{verbatim}

\hypertarget{subsystem-model}{%
\paragraph{Subsystem Model}\label{subsystem-model}}

Potential support requires:

\begin{verbatim}
potential System Model support
Subsystem Requirements
Subsystem Functional
Subsystem Logical
Subsystem Physical
\end{verbatim}

The dependency chain is conservative and monotonic.

A downstream support target shall not become potentially supported while an upstream required support target is not potentially supported.

\hypertarget{potential-support-states}{%
\subsubsection{Potential Support States}\label{potential-support-states}}

The canonical potential-support states are:

\begin{verbatim}
not_supported
partially_supported
potentially_supported
attention_required
\end{verbatim}

\texttt{not\_supported} means that none of the directly required framework-node coverage exists and no required upstream support target is satisfied.

\texttt{partially\_supported} means that some but not all direct and upstream requirements are satisfied.

\texttt{potentially\_supported} means that every direct framework-node requirement and every required upstream support target is preliminarily satisfied.

\texttt{attention\_required} means that the support dependencies would otherwise be satisfied but blocking or conflicting evidence prevents an unqualified support indication.

\texttt{potentially\_supported} shall always be interpreted as:

\begin{quote}
The versioned Preliminary Support Profile requirements possess eligible Preliminary Coverage evidence.
\end{quote}

It shall not be interpreted as:

\begin{itemize}
\tightlist
\item
  complete
\item
  correct
\item
  validated
\item
  approved
\item
  generation-ready
\item
  accepted by the engineering authority
\end{itemize}

\hypertarget{support-assessment-record}{%
\subsubsection{Support Assessment Record}\label{support-assessment-record}}

A derived potential-support record contains at least:

\begin{verbatim}
support_target_id
name
support_target_type
support_state
required_framework_node_ids
covered_framework_node_ids
missing_framework_node_ids
required_support_target_ids
satisfied_support_target_ids
unsatisfied_support_target_ids
attention_required
issue_codes
\end{verbatim}

All identifiers shall be unique and deterministically sorted.

\hypertarget{issue-handling}{%
\subsubsection{Issue Handling}\label{issue-handling}}

P6 issue levels are:

\begin{verbatim}
warning
blocking
\end{verbatim}

Warnings remain visible but do not automatically suppress valid coverage.

Blocking issues prevent the affected evidence from contributing to coverage.

A project-wide blocking integrity or reference issue changes the displayed project coverage state to \texttt{attention\_required}.

Representative blocking issue categories include:

\begin{itemize}
\tightlist
\item
  mixed project references
\item
  unknown source
\item
  unknown Information Unit
\item
  unknown Framework Assignment Candidate
\item
  unknown Framework node
\item
  Framework Template version mismatch
\item
  Preliminary Support Profile mismatch
\item
  source fingerprint mismatch
\item
  invalidated current evidence
\item
  stale exact-review binding
\item
  inconsistent duplicate identities
\item
  invalid support dependency graph
\item
  support dependency cycle
\end{itemize}

\hypertarget{determinism-and-ordering}{%
\subsubsection{Determinism and Ordering}\label{determinism-and-ordering}}

All derived records shall use stable ordering.

Unless a stricter domain order is defined:

\begin{itemize}
\tightlist
\item
  framework levels follow Framework Template order
\item
  framework nodes follow parent level and node order
\item
  sources follow \texttt{source\_id}
\item
  Information Units follow \texttt{information\_unit\_id}
\item
  Framework Assignment Candidates follow their identifiers
\item
  Human Review Decisions follow their identifiers
\item
  support targets follow profile order
\item
  issue codes and identifier sets are sorted deterministically
\end{itemize}

Equivalent validated inputs shall produce exactly equivalent derived assessments.

\hypertarget{assessment-fingerprint}{%
\subsubsection{Assessment Fingerprint}\label{assessment-fingerprint}}

Every project assessment shall contain an \texttt{assessment\_input\_fingerprint}.

The fingerprint binds the canonical representation of at least:

\begin{verbatim}
Framework Template identifier and version
Preliminary Support Profile identifier and version
eligible registered Source identities and fingerprints
effective Source Processing Dispositions
relevant Processing Run and artifact lifecycle identities
Framework Assignment Candidate identities and content fingerprints
latest exact Human Review Decision identities and fingerprints
assessment algorithm identifier and version
\end{verbatim}

The fingerprint is evidence of reproducibility.

It is not a mutable project status and does not become an independent authority.

\hypertarget{persistence}{%
\subsubsection{Persistence}\label{persistence}}

P6 shall not persist a mutable current Coverage state.

Coverage and potential-support assessments are calculated on demand by:

\begin{verbatim}
ProjectCoverageService.assess_project(project_id)
\end{verbatim}

A caller may serialize or cache a complete derived assessment for UI or report purposes only when:

\begin{itemize}
\tightlist
\item
  the derived status is explicitly marked non-authoritative
\item
  the assessment input fingerprint is preserved
\item
  the cache may be deleted and regenerated
\item
  stale caches are never treated as current authority
\end{itemize}

P7 may use the P6 assessment as a read-only dashboard projection.

\hypertarget{module-structure}{%
\subsubsection{Module Structure}\label{module-structure}}

The initial P6 module structure is:

\begin{verbatim}
modules/project_coverage/
├── __init__.py
├── errors.py
├── types.py
├── profile.py
├── evidence.py
├── coverage.py
├── support.py
└── service.py
\end{verbatim}

Responsibilities are:

\begin{verbatim}
errors.py
→ P6 validation, reference, integrity and assessment errors

types.py
→ immutable coverage, support, issue and assessment types

profile.py
→ versioned Preliminary Support Profile parsing and validation

evidence.py
→ candidate eligibility and latest exact Human Review resolution

coverage.py
→ node, level and project Preliminary Coverage

support.py
→ potential Model and SubModel support

service.py
→ project-local repository scans and complete assessment assembly
\end{verbatim}

The public package API is added only after the internal contracts are stable.

\hypertarget{implementation-sequence}{%
\subsubsection{Implementation Sequence}\label{implementation-sequence}}

P6 implementation follows:

\begin{verbatim}
P6-A1  Architecture and ADR-013
P6-I1  Errors, types and Preliminary Support Profile
P6-I2  Evidence eligibility and Human Review resolution
P6-I3  Node, level and project Preliminary Coverage
P6-I4  Preliminary Model and SubModel support
P6-I5  Coverage Service and public API
P6-I6  Integration, regression and P6 acceptance
\end{verbatim}

Implementation staging and commit may be deferred until the complete P6 implementation is ready, except that this accepted architecture decision shall be recorded before implementation depends on it.

Consequences

Positive consequences:

\begin{itemize}
\tightlist
\item
  Preliminary Coverage becomes deterministic and reproducible.
\item
  Coverage remains distinct from Approved Generation Readiness.
\item
  Context-only and out-of-scope sources cannot silently create engineering evidence.
\item
  Ambiguity and conflict remain visible without being misrepresented as coverage.
\item
  Human Review improves evidence classification without becoming model generation authority.
\item
  Model and SubModel support indications are based on an explicit versioned dependency profile.
\item
  P7 receives a stable read-only data contract for dashboard visualization.
\item
  Existing P1 through P5 authority boundaries remain unchanged.
\end{itemize}

Trade-offs:

\begin{itemize}
\tightlist
\item
  The conservative support chain may understate potential support.
\item
  A single uncovered required node prevents \texttt{potentially\_supported}.
\item
  Coverage does not measure semantic completeness inside one framework node.
\item
  The absence of weighted scoring provides less visual simplicity but avoids false precision.
\item
  Recalculation requires scanning several validated project repositories.
\end{itemize}

Risks and mitigations:

\begin{itemize}
\item
  Risk: candidate evidence may be misunderstood as approved readiness. Mitigation: explicit state names and \texttt{approved\_readiness\_status\ =\ not\_available}.
\item
  Risk: stale Human Review decisions could affect current coverage. Mitigation: only the latest exact content- and validation-fingerprint binding is relevant.
\item
  Risk: invalidated artifacts remain referenced by old candidates. Mitigation: P5 lifecycle validation excludes invalidated current evidence.
\item
  Risk: support rules become hidden implementation logic. Mitigation: versioned Preliminary Support Profile.
\item
  Risk: generated assessment caches become a second authority. Mitigation: derived, regenerable, fingerprint-bound and non-authoritative cache rules.
\end{itemize}

Rejected Alternatives

\hypertarget{count-only-confirmed-framework-assignment-candidates}{%
\subsubsection{Count only confirmed Framework Assignment Candidates}\label{count-only-confirmed-framework-assignment-candidates}}

Rejected because Preliminary Coverage is explicitly available without Human Approval.

This alternative would collapse Preliminary Coverage into a later approval concept.

\hypertarget{count-every-framework-assignment-candidate}{%
\subsubsection{Count every Framework Assignment Candidate}\label{count-every-framework-assignment-candidate}}

Rejected because \texttt{unassigned}, \texttt{ambiguous}, \texttt{conflict}, rejected and request-changes candidates do not provide valid covering evidence.

\hypertarget{confidence-weighted-coverage-score}{%
\subsubsection{Confidence-weighted coverage score}\label{confidence-weighted-coverage-score}}

Rejected because confidence is review evidence and not an authority boundary.

A weighted percentage would imply a precision and engineering maturity model that has not been defined.

\hypertarget{persist-mutable-coverage-status}{%
\subsubsection{Persist mutable Coverage status}\label{persist-mutable-coverage-status}}

Rejected because derived coverage could become stale and duplicate the authority of existing project records.

\hypertarget{calculate-approved-generation-readiness-in-p6}{%
\subsubsection{Calculate Approved Generation Readiness in P6}\label{calculate-approved-generation-readiness-in-p6}}

Rejected because Approved Generation Readiness is unavailable in Phase P and belongs to Phase G or its responsible successor phase.

\hypertarget{infer-potential-model-support-directly-in-the-dashboard}{%
\subsubsection{Infer potential model support directly in the dashboard}\label{infer-potential-model-support-directly-in-the-dashboard}}

Rejected because support dependencies would become hidden UI logic.

The rules belong in a versioned support profile and deterministic P6 service.

\hypertarget{allow-system-or-subsystem-support-without-upstream-support}{%
\subsubsection{Allow System or Subsystem support without upstream support}\label{allow-system-or-subsystem-support-without-upstream-support}}

Rejected for the initial profile.

The accepted initial profile uses the conservative dependency chain:

\begin{verbatim}
Stakeholder Model
→ System Model
→ Subsystem Model
\end{verbatim}

Future profiles may define different support targets only through an explicit versioned architecture and validation change.

\textbackslash newpage

\hypertarget{adr-014--project-dashboard-architecture-and-evidence-navigation}{%
\subsection{ADR-014 --- Project Dashboard Architecture and Evidence Navigation}\label{adr-014--project-dashboard-architecture-and-evidence-navigation}}

Project Dashboard Architecture and Evidence Navigation

Status

Accepted

Date

2026-07-27

Amendments

2026-07-27 -- First-Project Workspace Bootstrap

2026-07-27 -- Project Creation Availability with Existing Workspaces

2026-07-27 -- Phase-P P9 Project-bound Ingestion Integration Boundary

Context

Phase P introduces a project-oriented engineering workspace around the completed Phase F agentic ingestion pipeline.

P1 implemented the versioned Turing RFLP Framework.

P2 implemented persistent Project Workspaces and six-digit project identities.

P3 implemented immutable registered Sources with explicit source roles.

P4 implemented source-traceable Information Units, semantic candidates, Framework Assignment Candidates and exact Human Review Decisions.

P5 implemented immutable Processing Runs, Event Histories, Processing Decisions, artifact lifecycle, source disposition and project processing aggregation.

P6 implemented deterministic Preliminary Coverage and potential model-support assessment.

P7 introduces a read-only Project Dashboard for navigating the current project state and its supporting evidence.

Phase P was subsequently extended with P9, Project-bound Agentic Ingestion Integration. P9 follows the P8 Tests and Integration Readiness Review and connects the existing Phase F ingestion capability to the project-oriented contracts established by P2-P7. This amendment defines only the P7 integration boundary. P9 requires its own architecture decision before implementation.

The existing Team Agentic Ingestion UI remains an execution-oriented interface. It selects legacy inputs, configures ingestion runs, starts the pipeline and browses produced artifacts. It is not the project-level dashboard introduced by P7.

The dashboard must answer the following questions without creating a new source of truth:

\begin{itemize}
\tightlist
\item
  which project is currently selected
\item
  which Sources are registered and how they are classified
\item
  what the current project and source processing states are
\item
  which framework nodes and levels have Preliminary Coverage
\item
  which model and submodel scopes are potentially supported
\item
  which issues require attention
\item
  which Human Review Decisions affect the displayed state
\item
  which exact artifacts and documents support every displayed result
\item
  how the user can open those supporting artifacts directly from the dashboard
\end{itemize}

The dashboard shall be visually concise and shall use color only to communicate status.

Decision

\hypertarget{dashboard-authority}{%
\subsubsection{Dashboard Authority}\label{dashboard-authority}}

The Project Dashboard is a derived, read-only presentation layer.

It shall not become an authority for:

\begin{itemize}
\tightlist
\item
  engineering requirements
\item
  CATIA SysML v2 model content
\item
  project identity
\item
  Source registration
\item
  Processing State
\item
  Information Units
\item
  semantic candidates
\item
  Framework Assignment Candidates
\item
  Human Review Decisions
\item
  Preliminary Coverage
\item
  potential model support
\item
  Approved Generation Readiness
\end{itemize}

The dashboard may be deleted and regenerated without loss of authoritative engineering, semantic or processing information.

The dashboard shall obtain its information exclusively through the existing P2-P6 public services, repositories and immutable data types.

The UI shall not reimplement business rules already owned by P2-P6.

\hypertarget{read-only-scope}{%
\subsubsection{Read-only Scope}\label{read-only-scope}}

The five Project Dashboard views shall provide navigation and presentation only.

The application shell may expose one constrained P2 Project Workspace creation action. It is presented prominently when no valid Project Workspace exists and remains available as a secondary action when existing Project Workspaces can be selected.

P7 shall not:

\begin{itemize}
\tightlist
\item
  modify or delete an existing Project Workspace
\item
  accept a manually chosen Project ID
\item
  register or modify Sources
\item
  start, retry, supersede or mutate Processing Runs
\item
  create or modify Information Units
\item
  create or modify semantic candidates
\item
  create or modify Framework Assignment Candidates
\item
  create or modify Human Review Decisions
\item
  promote Approved Inputs
\item
  calculate Approved Generation Readiness
\item
  generate model candidates
\item
  generate SysML v2
\item
  mutate CATIA
\item
  persist mutable dashboard state as project authority
\end{itemize}

The existing Team Agentic Ingestion UI remains separate and unchanged during P7. Any adaptation that binds ingestion to a Project Workspace belongs to P9 and shall be governed by a separate architecture decision.

The only P7 write exception is the Project Workspace creation action. It shall:

\begin{itemize}
\tightlist
\item
  call the existing P2 \texttt{ProjectWorkspace.create\_project} contract
\item
  require a human-readable display name
\item
  allow an optional description
\item
  generate the six-digit Project ID through P2
\item
  pin the accepted Framework Template through the P2 Project Manifest
\item
  remain available after one or more valid Project Workspaces exist
\item
  create a new Project Workspace without modifying an existing one
\item
  create no Sources, Processing Runs, semantic artifacts or review decisions
\end{itemize}

After successful Project Workspace creation, every Project Dashboard view remains read-only.

Beginning with P9, the application shell may expose project-bound navigation from the dashboard to a separately governed ingestion execution view and back. The P7 dashboard views themselves remain read-only. Navigation does not authorize the dashboard presenter or renderer to register Sources, start Processing Runs or publish ingestion artifacts.

\hypertarget{dashboard-composition}{%
\subsubsection{Dashboard Composition}\label{dashboard-composition}}

The initial Project Dashboard contains five primary views:

\begin{verbatim}
Project Overview
Preliminary Coverage
Sources and Processing
Attention and Review
Traceability and Documents
\end{verbatim}

The selected project remains visible across all views.

\hypertarget{project-overview}{%
\paragraph{Project Overview}\label{project-overview}}

The overview presents a compact project summary:

\begin{itemize}
\tightlist
\item
  Project display name
\item
  six-digit Project ID
\item
  Framework Template identifier and version
\item
  Preliminary Support Profile identifier and version
\item
  registered Source count
\item
  project processing state
\item
  project Preliminary Coverage state
\item
  attention summary
\item
  potential Stakeholder Model support
\item
  potential System Model support
\item
  potential Subsystem Model support
\item
  explicit Approved Generation Readiness unavailability
\end{itemize}

The overview shall not display a synthetic maturity score.

Exact ratios may be displayed, for example:

\begin{verbatim}
3 of 4 System framework nodes preliminarily covered
\end{verbatim}

\hypertarget{preliminary-coverage}{%
\paragraph{Preliminary Coverage}\label{preliminary-coverage}}

The coverage view presents:

\begin{itemize}
\tightlist
\item
  Stakeholder, System and Subsystem framework levels
\item
  all twelve stable framework mapping targets
\item
  node coverage state
\item
  level coverage state
\item
  reviewed and unreviewed candidate counts
\item
  eligible Source and Information Unit counts
\item
  independent attention state
\item
  related issue codes
\item
  direct evidence navigation
\end{itemize}

Canonical node states remain owned by P6:

\begin{verbatim}
uncovered
candidate_covered
reviewed_candidate_covered
\end{verbatim}

Canonical level states remain owned by P6:

\begin{verbatim}
uncovered
partially_covered
covered
\end{verbatim}

Coverage and attention shall remain separate dimensions.

\hypertarget{sources-and-processing}{%
\paragraph{Sources and Processing}\label{sources-and-processing}}

The processing view presents each registered Source with:

\begin{itemize}
\tightlist
\item
  Source ID
\item
  original filename
\item
  source role
\item
  content fingerprint
\item
  effective processing disposition
\item
  current Processing Run ID
\item
  run state
\item
  processing stage
\item
  latest Attempt ID
\item
  pending review state
\item
  blocking issue codes
\item
  failure issue codes
\item
  superseded Run IDs
\item
  invalidated artifact count
\item
  direct Source Manifest and processing-history navigation
\end{itemize}

The dashboard shall not infer a processing state that is not supplied by P5.

\hypertarget{attention-and-review}{%
\paragraph{Attention and Review}\label{attention-and-review}}

The attention view presents:

\begin{itemize}
\tightlist
\item
  blocking and warning issues
\item
  affected project, Source, Information Unit, candidate, node or support target
\item
  Human Review target type
\item
  Human Review decision
\item
  decision identity
\item
  exact target content fingerprint
\item
  exact reference-validation fingerprint
\item
  reviewer identity where available
\item
  direct navigation to affected artifacts and decisions
\end{itemize}

A Human Review confirmation shown in P7 confirms only the exact reviewed target.

It shall not be presented as:

\begin{itemize}
\tightlist
\item
  Approved Input
\item
  Engineering Approval
\item
  Approved Generation Readiness
\item
  model acceptance
\item
  generation authorization
\end{itemize}

\hypertarget{traceability-and-documents}{%
\paragraph{Traceability and Documents}\label{traceability-and-documents}}

The traceability view presents the evidence chain behind dashboard results.

A typical chain is:

\begin{verbatim}
Source
→ Processing Run and Events
→ Information Unit
→ Framework Assignment Candidate
→ Reference Validation
→ Human Review Decision
→ Preliminary Coverage
→ Potential Model Support
\end{verbatim}

The view shall allow navigation in both directions where exact references exist.

\hypertarget{smart-evidence-navigation}{%
\subsubsection{Smart Evidence Navigation}\label{smart-evidence-navigation}}

Every dashboard element that presents a traceable project fact shall be capable of exposing the exact supporting artifacts.

Examples include:

\begin{itemize}
\tightlist
\item
  a project processing state
\item
  a Source disposition
\item
  a framework-node coverage state
\item
  an attention indicator
\item
  a support assessment
\item
  a Human Review status
\item
  an issue
\end{itemize}

The presenter layer shall bind such values to immutable Evidence References.

An Evidence Reference contains at least:

\begin{verbatim}
reference_type
reference_id
display_label
repository_relative_path
content_fingerprint
media_type
source_role
relationship
\end{verbatim}

Optional navigation metadata may include:

\begin{verbatim}
section_anchor
line_start
line_end
json_pointer
table_row_key
\end{verbatim}

The dashboard shall not construct arbitrary filesystem paths in the UI.

Evidence paths shall be resolved by trusted repository-aware resolvers.

Each resolved path must:

\begin{itemize}
\tightlist
\item
  remain within the configured repository or Project Workspace root
\item
  reject symbolic links where the owning repository rejects them
\item
  correspond to an existing authoritative or derived artifact
\item
  preserve project isolation
\item
  preserve the artifact identity and fingerprint where available
\end{itemize}

\hypertarget{one-supporting-artifact}{%
\paragraph{One Supporting Artifact}\label{one-supporting-artifact}}

When exactly one Evidence Reference supports a displayed value, activating the navigation control shall open that artifact directly in the internal document viewer.

\hypertarget{multiple-supporting-artifacts}{%
\paragraph{Multiple Supporting Artifacts}\label{multiple-supporting-artifacts}}

When multiple Evidence References support a displayed value, activating the navigation control shall open a compact evidence chooser.

The chooser shall:

\begin{itemize}
\tightlist
\item
  identify the relationship of each artifact to the displayed value
\item
  show artifact type, stable ID and concise label
\item
  preserve deterministic ordering
\item
  distinguish direct evidence from contextual evidence
\item
  allow one artifact to be opened without leaving the dashboard context
\end{itemize}

The chooser shall not silently choose one artifact when several materially contribute to the result.

\hypertarget{no-supporting-artifact}{%
\paragraph{No Supporting Artifact}\label{no-supporting-artifact}}

When a displayed value has no resolvable supporting artifact, the dashboard shall show that evidence navigation is unavailable.

It shall not create or guess a path.

Missing expected evidence shall be visible as an issue where the responsible P2-P6 contract classifies it as invalid or blocking.

\hypertarget{internal-document-viewer}{%
\subsubsection{Internal Document Viewer}\label{internal-document-viewer}}

The dashboard shall use an internal document viewer rather than browser \texttt{file://} links.

This avoids platform-specific file access, browser security restrictions and uncontrolled navigation outside the application.

The viewer shall support at least:

\begin{verbatim}
JSON
Markdown
plain text
CSV
\end{verbatim}

Viewer behavior:

\begin{itemize}
\tightlist
\item
  JSON is displayed in formatted, readable form
\item
  Markdown is rendered and may also be shown as source text
\item
  plain text preserves line structure
\item
  CSV may be displayed as a table and as raw text
\item
  unsupported binary content shows metadata and an explicit file action
\item
  large files use bounded previews and explicit expansion
\item
  fingerprints and repository-relative paths remain visible
\item
  opening a document never changes its review or approval state
\end{itemize}

Where optional navigation metadata is available, the viewer may highlight:

\begin{itemize}
\tightlist
\item
  a Markdown section
\item
  a line range
\item
  a JSON field or object
\item
  a table row
\end{itemize}

Highlighting is a navigation aid only and does not create a new persisted reference authority.

\hypertarget{navigation-context}{%
\subsubsection{Navigation Context}\label{navigation-context}}

Opening evidence shall preserve the user's dashboard context.

The application shall retain at least:

\begin{verbatim}
selected project
current dashboard view
selected node, source, issue or support target
opened evidence reference
\end{verbatim}

The user shall be able to return to the previous dashboard location without reconstructing filters manually.

Deep links may use application query parameters for navigation state.

Deep links shall contain stable project and artifact identities, not unrestricted filesystem paths.

\hypertarget{dashboard-presentation-model}{%
\subsubsection{Dashboard Presentation Model}\label{dashboard-presentation-model}}

A new \texttt{modules.project\_dashboard} package shall separate data collection, presentation and Streamlit rendering.

The planned package is:

\begin{verbatim}
modules/project_dashboard/
├── __init__.py
├── errors.py
├── types.py
├── references.py
├── presenter.py
└── service.py
\end{verbatim}

The Streamlit layer is planned as:

\begin{verbatim}
app/project_dashboard_ui.py
app/project_dashboard_app.py
\end{verbatim}

\hypertarget{dashboard-service}{%
\paragraph{Dashboard Service}\label{dashboard-service}}

The dashboard service coordinates read-only calls to P2-P6.

It shall return a complete immutable dashboard snapshot.

It shall not expose repository objects directly to the UI.

\hypertarget{presenter}{%
\paragraph{Presenter}\label{presenter}}

The presenter transforms domain records into display-ready immutable view models.

The presenter may:

\begin{itemize}
\tightlist
\item
  define stable ordering
\item
  format concise labels
\item
  group related records
\item
  bind Evidence References
\item
  derive display-only counts from already validated records
\item
  select status labels and status semantics
\end{itemize}

The presenter shall not:

\begin{itemize}
\tightlist
\item
  recalculate Preliminary Coverage
\item
  recalculate potential model support
\item
  change issue severity
\item
  infer approval
\item
  infer missing evidence
\item
  change processing state
\end{itemize}

\hypertarget{streamlit-ui}{%
\paragraph{Streamlit UI}\label{streamlit-ui}}

The Streamlit UI renders the supplied view models.

The UI shall contain minimal domain logic.

UI event handlers may:

\begin{itemize}
\tightlist
\item
  change the selected project
\item
  change the active dashboard view
\item
  open an Evidence Reference
\item
  choose among multiple Evidence References
\item
  apply presentation filters
\item
  request project-bound navigation to or from a P9 execution view
\end{itemize}

Dashboard-view event handlers shall not mutate P2-P6 project artifacts. The constrained Project Workspace creation action remains the sole P7 write exception. Any P9 execution action is owned by the separately governed P9 integration layer, not by the dashboard view or presenter.

\hypertarget{project-selection-and-project-creation}{%
\subsubsection{Project Selection and Project Creation}\label{project-selection-and-project-creation}}

When no valid Project Workspace exists, the application shall present a prominent first-project form instead of a terminal empty state.

When one or more valid Project Workspaces exist, the same creation capability shall remain available as a secondary, collapsed action alongside the deterministic project selector.

The form shall contain:

\begin{verbatim}
Project name
Description (optional)
Create project
\end{verbatim}

The user shall not enter the internal Project ID. P2 generates the six-digit identifier and persists the Project Manifest atomically.

After successful creation, the application shall select the new project, open the Overview and rerun the presentation layer.

The dashboard shall list existing Project Workspaces.

A project option shall display:

\begin{verbatim}
<display name> · <six-digit Project ID>
\end{verbatim}

The six-digit ID remains the stable internal project identity.

Project selection shall be deterministic.

An invalid or unsafe Project Workspace shall not be silently presented as a valid project.

Project scan issues shall remain visible.

\hypertarget{p9-navigation-and-integration-boundary}{%
\subsubsection{P9 Navigation and Integration Boundary}\label{p9-navigation-and-integration-boundary}}

P9 may connect the application shell to a separately defined project-bound Agentic Ingestion execution view. The integration shall preserve the separation between inspection and execution:

\begin{verbatim}
Project Dashboard
→ read-only inspection and evidence navigation

Project-bound Agentic Ingestion
→ explicit execution workflow governed by P9
\end{verbatim}

The application shell may expose navigation controls such as:

\begin{verbatim}
Start ingestion for this project
Return to Project Dashboard
\end{verbatim}

A navigation request shall carry only stable application identities and state, at least:

\begin{verbatim}
project_id
return_view
optional selected entity identity
\end{verbatim}

It shall not carry unrestricted filesystem paths.

The selected six-digit Project ID is mandatory for project-bound ingestion. P9 shall not silently fall back to a global, unassigned or different project when a project binding is unavailable or invalid.

The separately governed P9 execution layer may coordinate existing authoritative contracts, including:

\begin{itemize}
\tightlist
\item
  P3 Project Source Registry for Source registration and Source role assignment
\item
  P5 Processing operations for Processing Run and event persistence
\item
  the existing Phase F Team Agentic Ingestion pipeline as an execution engine
\item
  project-local publication of traceable reports and agent outputs
\end{itemize}

ADR-014 does not authorize or define those write operations. Their exact transaction boundaries, failure behavior, artifact mapping and recovery rules shall be specified in the P9 architecture decision.

After an ingestion execution returns to the dashboard, the application shall discard any stale dashboard snapshot and regenerate the selected project's view from the P2-P6 authorities. The dashboard shall not accept execution results directly as presentation truth.

A completed ingestion run shall not create or imply Preliminary Coverage unless the required valid P4 Information Units, Framework Assignment Candidates, reference validations and Human Review Decisions actually exist. Processing visibility and Preliminary Coverage remain separate states.

The intended demonstrator flow is:

\begin{verbatim}
Create or select Project
→ start project-bound ingestion
→ register Source
→ execute ingestion
→ persist Processing evidence
→ return to Dashboard
→ regenerate and inspect project state
\end{verbatim}

\hypertarget{visual-design-principles}{%
\subsubsection{Visual Design Principles}\label{visual-design-principles}}

The dashboard shall be visually concise, calm and information-dense without being cramped.

The default visual language shall use:

\begin{itemize}
\tightlist
\item
  neutral backgrounds
\item
  neutral borders
\item
  whitespace
\item
  typography
\item
  hierarchy
\item
  alignment
\item
  concise labels
\item
  restrained icons
\end{itemize}

Color shall be used only to communicate status.

Color shall not be used merely for decoration, section identity, navigation or branding emphasis.

Large decorative gradients, multicolored cards and unrelated accent colors are out of scope.

\hypertarget{status-semantics}{%
\paragraph{Status Semantics}\label{status-semantics}}

The initial status families are:

\begin{verbatim}
neutral
informational
candidate
reviewed
attention
blocking
unavailable
\end{verbatim}

Status presentation shall combine:

\begin{verbatim}
text label
icon or shape
color
\end{verbatim}

Color shall never be the only carrier of meaning.

Suggested semantic intent:

\begin{verbatim}
neutral or unavailable
→ gray

informational or unreviewed candidate
→ blue

reviewed candidate or covered
→ green

partial coverage or attention
→ amber

blocking, rejected or invalid
→ red
\end{verbatim}

The exact color tokens belong to the UI implementation and tests.

Status colors shall be applied only to compact status-bearing elements such as:

\begin{itemize}
\tightlist
\item
  badges
\item
  small indicators
\item
  icons
\item
  narrow status borders
\item
  status text
\end{itemize}

Entire cards, pages or large table regions shall not be filled with status colors.

The dashboard shall remain understandable in monochrome and for users with color-vision deficiencies.

\hypertarget{progressive-disclosure}{%
\subsubsection{Progressive Disclosure}\label{progressive-disclosure}}

The dashboard shall present summary before detail.

The initial view shows concise project-level status.

Detailed records are revealed through:

\begin{itemize}
\tightlist
\item
  expandable sections
\item
  filtered tables
\item
  evidence navigation
\item
  the internal document viewer
\end{itemize}

The dashboard shall avoid displaying all identifiers and fingerprints in the primary summary.

Exact identifiers, fingerprints and paths remain available in detailed views.

\hypertarget{status-language}{%
\subsubsection{Status Language}\label{status-language}}

The dashboard shall preserve the exact meaning of P5 and P6 states.

The UI may provide concise explanatory text, but it shall not rename a state in a way that changes its meaning.

In particular:

\begin{verbatim}
potentially_supported
\end{verbatim}

shall never be displayed as:

\begin{verbatim}
ready
approved
complete
valid
generation-ready
\end{verbatim}

The dashboard shall visibly expose:

\begin{verbatim}
approved_readiness_status = not_available
approved_readiness_available_from_phase = G
\end{verbatim}

The preferred user-facing explanation is:

\begin{verbatim}
Approved Generation Readiness is not assessed in Phase P.
It becomes available from Phase G.
\end{verbatim}

\hypertarget{determinism}{%
\subsubsection{Determinism}\label{determinism}}

Equivalent project records shall produce equivalent dashboard snapshots.

Stable ordering shall be defined for:

\begin{itemize}
\tightlist
\item
  projects
\item
  framework levels
\item
  framework nodes
\item
  Sources
\item
  Processing Runs
\item
  issues
\item
  Human Review Decisions
\item
  support targets
\item
  Evidence References
\end{itemize}

The dashboard shall not depend on filesystem iteration order.

\hypertarget{failure-behavior}{%
\subsubsection{Failure Behavior}\label{failure-behavior}}

The dashboard shall fail closed.

A failure to resolve one section shall not cause another section to invent data.

Where possible, the dashboard may render a partial snapshot with explicit issues.

Examples:

\begin{verbatim}
Coverage assessment unavailable
Source scan available
Processing scan available
\end{verbatim}

A section failure shall show:

\begin{itemize}
\tightlist
\item
  affected section
\item
  concise error
\item
  relevant issue code
\item
  available evidence navigation
\item
  no fabricated status
\end{itemize}

Unexpected exceptions shall be converted into safe presentation errors before reaching raw Streamlit output.

\hypertarget{testing-strategy}{%
\subsubsection{Testing Strategy}\label{testing-strategy}}

P7 shall include tests for:

\begin{itemize}
\tightlist
\item
  immutable dashboard types
\item
  project selection ordering
\item
  presenter determinism
\item
  status-label mapping
\item
  status-color restriction
\item
  project isolation
\item
  Evidence Reference validation
\item
  safe repository-relative path resolution
\item
  single-evidence direct navigation
\item
  multiple-evidence chooser behavior
\item
  missing-evidence behavior
\item
  internal viewer selection
\item
  section and line navigation metadata
\item
  source and processing presentation
\item
  coverage and support presentation
\item
  attention and Human Review presentation
\item
  partial failure behavior
\item
  public API exports
\item
  Streamlit-independent rendering contracts
\end{itemize}

Core dashboard logic shall be testable without starting Streamlit.

Streamlit tests shall focus on thin integration behavior rather than duplicate domain tests.

\hypertarget{implementation-sequence-1}{%
\subsubsection{Implementation Sequence}\label{implementation-sequence-1}}

P7 is implemented in six steps:

\begin{verbatim}
P7 Step 1 of 6
Architecture and ADR-014

P7 Step 2 of 6
Dashboard types, Evidence References and presenter foundation

P7 Step 3 of 6
Project selection and Project Overview

P7 Step 4 of 6
Processing, Coverage and Potential Support views

P7 Step 5 of 6
Attention, Human Review, traceability and document viewer

P7 Step 6 of 6
UI integration, tests, full regression and P7 acceptance
\end{verbatim}

P8 performs the Tests and Integration Readiness Review for P1-P7. P9 then implements Project-bound Agentic Ingestion Integration under its own ADR. This extension of Phase P does not change the six-step P7 implementation sequence.

Consequences

Positive consequences:

\begin{itemize}
\tightlist
\item
  project status becomes accessible through one coherent interface
\item
  every relevant displayed result can lead to its supporting artifacts
\item
  multiple evidence chains remain explicit rather than hidden
\item
  P2-P6 remain authoritative
\item
  the existing ingestion UI remains stable during P7
\item
  P7 provides a defined navigation seam for the later P9 integration
\item
  P9 can update project state without moving execution logic into dashboard views
\item
  the visual design communicates status without decorative noise
\item
  the dashboard remains suitable for demonstration and technical inspection
\item
  traceability becomes directly explorable
\item
  core behavior remains testable without Streamlit
\end{itemize}

Negative consequences:

\begin{itemize}
\tightlist
\item
  evidence navigation requires additional immutable view-model types
\item
  repository-aware path resolvers add implementation effort
\item
  the internal viewer must safely handle multiple text formats
\item
  deep-link and navigation context require explicit UI state handling
\item
  the dashboard views cannot provide write actions during P7
\item
  P9 requires a separate ADR and explicit adaptation of the existing ingestion flow
\item
  the integrated demonstrator remains split between inspection and execution views
\item
  some records may initially support only artifact-level rather than exact line-level navigation
\end{itemize}

Rejected Alternatives

\hypertarget{merge-p7-directly-into-the-existing-ingestion-ui}{%
\paragraph{Merge P7 directly into the existing ingestion UI}\label{merge-p7-directly-into-the-existing-ingestion-ui}}

Rejected because execution configuration and project-state inspection have different responsibilities. Merging them would create a large UI with mixed read and write semantics.

\hypertarget{embed-p9-execution-controls-inside-dashboard-views}{%
\paragraph{Embed P9 execution controls inside dashboard views}\label{embed-p9-execution-controls-inside-dashboard-views}}

Rejected because navigation from a dashboard view to an execution workflow is not equivalent to making the view itself writable. Source registration, Processing Run creation and ingestion execution belong to the separately governed P9 execution layer.

\hypertarget{treat-navigation-alone-as-project-bound-ingestion-integration}{%
\paragraph{Treat navigation alone as project-bound ingestion integration}\label{treat-navigation-alone-as-project-bound-ingestion-integration}}

Rejected because a link between two screens would not bind Sources, Processing Runs or generated artifacts to the selected Project Workspace. P9 must implement the repository and lifecycle bridge, not only navigation.

\hypertarget{use-direct-file-links}{%
\paragraph{\texorpdfstring{Use direct \texttt{file://} links}{Use direct file:// links}}\label{use-direct-file-links}}

Rejected because browser restrictions, platform differences and uncontrolled filesystem access make them unreliable and unsafe.

\hypertarget{open-only-repository-relative-paths-as-plain-text-labels}{%
\paragraph{Open only repository-relative paths as plain text labels}\label{open-only-repository-relative-paths-as-plain-text-labels}}

Rejected because the user must be able to navigate directly to supporting documents rather than manually locate them.

\hypertarget{automatically-open-the-first-artifact-when-multiple-artifacts-contribute}{%
\paragraph{Automatically open the first artifact when multiple artifacts contribute}\label{automatically-open-the-first-artifact-when-multiple-artifacts-contribute}}

Rejected because this would hide material evidence and could imply that one artifact alone supports the displayed result.

\hypertarget{recalculate-coverage-in-the-dashboard}{%
\paragraph{Recalculate coverage in the dashboard}\label{recalculate-coverage-in-the-dashboard}}

Rejected because P6 is the sole owner of Preliminary Coverage and potential model-support assessment.

\hypertarget{use-decorative-colors-for-dashboard-sections}{%
\paragraph{Use decorative colors for dashboard sections}\label{use-decorative-colors-for-dashboard-sections}}

Rejected because color is reserved for status semantics.

\hypertarget{add-approval-and-review-actions-to-p7}{%
\paragraph{Add approval and review actions to P7}\label{add-approval-and-review-actions-to-p7}}

Rejected because P7 dashboard views are read-only. Write workflows require their own explicit architecture and authority boundaries.

\hypertarget{leave-the-initial-empty-state-without-a-project-bootstrap}{%
\paragraph{Leave the initial empty state without a project bootstrap}\label{leave-the-initial-empty-state-without-a-project-bootstrap}}

Rejected because the application would have no valid starting action in a fresh repository. The constrained P2 bootstrap creates only the Project Workspace boundary and does not weaken the read-only status of dashboard views.

\hypertarget{add-full-project-lifecycle-management-to-p7}{%
\paragraph{Add full project lifecycle management to P7}\label{add-full-project-lifecycle-management-to-p7}}

Rejected because editing, deleting or arbitrarily managing existing projects would exceed the minimal bootstrap exception and mix project administration with evidence presentation.

Acceptance Criteria

ADR-014 is satisfied when:

\begin{enumerate}
\def\labelenumi{\arabic{enumi}.}
\tightlist
\item
  the five dashboard views are read-only and the only write action is constrained P2 Project Workspace creation
\item
  P2-P6 remain the sole domain authorities
\item
  every traceable displayed value can expose Evidence References
\item
  one Evidence Reference opens directly
\item
  multiple Evidence References produce an explicit chooser
\item
  arbitrary filesystem paths cannot be opened
\item
  an internal viewer supports JSON, Markdown, text and CSV
\item
  project and navigation context are preserved
\item
  status color is used only for status-bearing elements
\item
  status meaning is also conveyed through text and icon or shape
\item
  the dashboard exposes Preliminary Coverage and potential support accurately
\item
  Approved Generation Readiness remains explicitly unavailable in Phase P
\item
  the existing Team Agentic Ingestion UI remains unchanged during P7, and any P9 adaptation is governed by a separate ADR
\item
  core dashboard behavior is testable without Streamlit
\item
  the complete repository regression remains green
\item
  a fresh repository can create its first Project Workspace without a manually entered Project ID
\item
  an additional Project Workspace can be created while existing Project Workspaces remain selectable
\item
  successful Project Workspace creation selects the new project and opens the Overview
\item
  beginning with P9, the application shell may navigate to a separately governed execution view while dashboard views remain read-only
\item
  P9 navigation carries a valid selected Project ID and no unrestricted filesystem path
\item
  returning from P9 regenerates the dashboard snapshot from P2-P6 authorities
\item
  ADR-014 does not authorize P9 Source registration, Processing Run creation or artifact publication
\end{enumerate}

\textbackslash newpage

\hypertarget{adr-015--project-bound-agentic-ingestion-integration-architecture}{%
\subsection{ADR-015 --- Project-bound Agentic Ingestion Integration Architecture}\label{adr-015--project-bound-agentic-ingestion-integration-architecture}}

Project-bound Agentic Ingestion Integration Architecture

Status

Accepted

Date

2026-07-27

\hypertarget{context}{%
\subsubsection{Context}\label{context}}

Phase F implemented a working Team Agentic Ingestion pipeline and a standalone execution-oriented Streamlit interface.

Phase P introduced the project-oriented engineering workspace:

\begin{itemize}
\tightlist
\item
  P1 implemented the versioned Turing RFLP Framework Template.
\item
  P2 implemented persistent Project Workspaces and six-digit Project IDs.
\item
  P3 implemented immutable registered Sources with mandatory Project assignment and explicit Source roles.
\item
  P4 implemented project-local source projections, Information Units, semantic candidates, terminology mappings, Framework Assignment Candidates and Human Review Decisions.
\item
  P5 implemented immutable Processing Runs, Processing Events, Processing Decisions, artifact lifecycle, retry, supersession, recovery diagnostics and project-level Processing aggregation.
\item
  P6 implemented deterministic Preliminary Coverage and potential model-support assessment.
\item
  P7 implemented the read-only Project Dashboard with evidence navigation and a constrained Project Workspace creation action.
\item
  P8 verified the P1-P7 implementation baseline and confirmed that the existing public contracts can support a project-bound ingestion integration without a parallel project or processing architecture.
\end{itemize}

The existing Phase F ingestion pipeline currently operates with repository-global execution paths and is not bound to a selected Project Workspace, registered Source or P5 Processing Run.

A navigation link between the Project Dashboard and the existing ingestion UI would not be sufficient. A real integration must bind the uploaded Source, Processing Run, Attempts, generated reports and agent outputs to the selected Project Workspace and make their current state visible through the existing P5 and P7 authorities.

The integration must preserve the following accepted boundaries:

\begin{itemize}
\tightlist
\item
  Project identity is owned by P2.
\item
  Source registration and Source integrity are owned by P3.
\item
  Processing identity, lifecycle, artifact references, retry, supersession and recovery are owned by P5.
\item
  Preliminary Coverage and potential model support are owned by P6.
\item
  Project presentation is owned by P7.
\item
  Phase F remains the execution engine for Team Agentic Ingestion.
\item
  CATIA remains the authoritative engineering model.
\item
  A completed ingestion execution does not imply approved engineering knowledge, Preliminary Coverage, Approved Generation Readiness or model generation.
\end{itemize}

\hypertarget{decision}{%
\subsubsection{Decision}\label{decision}}

\hypertarget{1-introduce-a-dedicated-project-bound-ingestion-integration-layer}{%
\paragraph{1. Introduce a dedicated project-bound ingestion integration layer}\label{1-introduce-a-dedicated-project-bound-ingestion-integration-layer}}

P9 shall introduce a narrow orchestration and publication layer:

\begin{verbatim}
modules/project_ingestion/
├── __init__.py
├── errors.py
├── types.py
├── configuration.py
├── publisher.py
└── service.py
\end{verbatim}

The central public contract shall be exposed by:

\begin{Shaded}
\begin{Highlighting}[]
\NormalTok{ProjectBoundIngestionService}
\end{Highlighting}
\end{Shaded}

Its primary execution operation shall conceptually provide:

\begin{Shaded}
\begin{Highlighting}[]
\NormalTok{ProjectBoundIngestionService.execute(}
\NormalTok{    project\_id,}
\NormalTok{    source\_path,}
\NormalTok{    source\_role,}
\NormalTok{    ingestion\_configuration,}
\NormalTok{) }\OperatorTok{{-}\textgreater{}}\NormalTok{ ProjectBoundIngestionResult}
\end{Highlighting}
\end{Shaded}

The exact Python signature may be refined during implementation, but it shall preserve the following responsibilities:

\begin{enumerate}
\def\labelenumi{\arabic{enumi}.}
\tightlist
\item
  require one valid selected Project ID;
\item
  register the uploaded Source through P3;
\item
  use the registered project-local Source content as execution input;
\item
  create one P5 Processing Run;
\item
  create one Processing Attempt for project-bound agentic ingestion;
\item
  execute the existing Phase F Team Agentic Ingestion pipeline;
\item
  validate and fingerprint the complete publishable output set;
\item
  publish artifacts into the P5 run-owned artifact structure;
\item
  append immutable P5 Processing Events;
\item
  return a project-bound result containing stable identities and evidence references.
\end{enumerate}

The integration layer shall not replace, fork or duplicate the P2, P3, P5, P6 or P7 authorities.

\hypertarget{2-add-a-common-turing-generator-application-shell}{%
\paragraph{2. Add a common Turing Generator application shell}\label{2-add-a-common-turing-generator-application-shell}}

P9 shall introduce:

\begin{verbatim}
app/turing_generator_app.py
\end{verbatim}

The application shell shall provide at least:

\begin{verbatim}
Project Dashboard
Agentic Ingestion
\end{verbatim}

The Project Dashboard remains the P7 read-only inspection interface.

The Agentic Ingestion view is the P9 execution interface.

Navigation shall carry stable application identities and state only:

\begin{verbatim}
project_id
return_view
optional selected entity identity
\end{verbatim}

Navigation shall not carry unrestricted filesystem paths.

The selected six-digit Project ID is mandatory for P9 execution. The application shall fail closed when no valid project binding exists. It shall not silently use a global, unassigned or different project.

\hypertarget{3-preserve-the-existing-phase-f-pipeline-as-the-execution-engine}{%
\paragraph{3. Preserve the existing Phase F pipeline as the execution engine}\label{3-preserve-the-existing-phase-f-pipeline-as-the-execution-engine}}

The existing Team Agentic Ingestion implementation shall remain the execution engine.

P9 may add a backward-compatible output-root or execution-root option to the pipeline.

The default behavior shall remain unchanged:

\begin{verbatim}
no project-bound execution root supplied
→ existing Phase F global demo behavior
\end{verbatim}

Project-bound behavior shall be explicit:

\begin{verbatim}
project-bound execution root supplied
→ execution occurs inside the selected P5 Processing Run work directory
\end{verbatim}

Existing Phase F tests and the standalone Phase F demo shall remain operational.

The project-bound integration shall invoke the pipeline with the validated project-local Source content path returned by P3. The original temporary upload path shall not remain the authoritative execution input after registration.

\hypertarget{4-source-registration-boundary}{%
\paragraph{4. Source registration boundary}\label{4-source-registration-boundary}}

The Source shall be registered through:

\begin{Shaded}
\begin{Highlighting}[]
\NormalTok{ProjectSourceRegistry.register\_source(...)}
\end{Highlighting}
\end{Shaded}

The registration result shall provide the authoritative:

\begin{verbatim}
project_id
source_id
source_role
source_sha256
stored filename
registered content path
\end{verbatim}

Source registration is durable and precedes Processing Run creation.

If Source registration succeeds but Processing Run creation does not, the Source remains a valid registered Source and is presented as not started.

Duplicate Source content within the same project shall continue to be rejected by the P3 duplicate-content contract.

P9 shall not bypass Source integrity, Source role validation or project-local storage.

\hypertarget{5-processing-run-contract}{%
\paragraph{5. Processing Run contract}\label{5-processing-run-contract}}

P9 shall create Processing Runs through the existing P5 public contracts.

A new run shall bind at least:

\begin{verbatim}
project_id
processing_run_id
source_id
source_sha256
source_role_snapshot
workflow_profile
configuration_fingerprint
framework_template_id
framework_template_version
semantic_reference_versions
created_at
optional supersedes_run_id
\end{verbatim}

The Processing Run Manifest shall remain immutable.

The configuration fingerprint shall cover all material execution settings that affect reproducibility, including at least:

\begin{verbatim}
recipe ID
provider
model
dry-run mode
team-member limit
runs per member
relevant pipeline configuration version
\end{verbatim}

Secrets such as API keys shall never be persisted or included in fingerprints, reports, events or result objects.

\hypertarget{6-add-the-agentic-ingestion-processing-stage}{%
\paragraph{6. Add the agentic ingestion Processing Stage}\label{6-add-the-agentic-ingestion-processing-stage}}

P5 currently models the downstream engineering-processing stages.

P9 shall add:

\begin{verbatim}
agentic_ingestion
\end{verbatim}

to the canonical Processing Stage vocabulary.

This stage represents the existing Phase F team-based analysis and review-report generation.

It shall not be mislabeled as:

\begin{verbatim}
source_projection
semantic_extraction
semantic_consensus
terminology_mapping
framework_assignment
human_review
publication
\end{verbatim}

A successful P9 execution shall initially end with:

\begin{verbatim}
run_state:
awaiting_review

processing_stage:
agentic_ingestion
\end{verbatim}

\hypertarget{7-extend-run-owned-artifact-kinds}{%
\paragraph{7. Extend run-owned artifact kinds}\label{7-extend-run-owned-artifact-kinds}}

P5 currently supports run-owned Agent Outputs and Consensus Reports.

P9 shall extend the canonical run-owned artifact kinds to include:

\begin{verbatim}
agent_outputs
consensus_reports
review_reports
run_summaries
\end{verbatim}

The canonical structure shall be:

\begin{verbatim}
data/projects/<project_id>/runs/<processing_run_id>/
├── run_manifest.json
├── events/
├── artifacts/
│   ├── agent_outputs/
│   │   └── agentic_ingestion/<attempt_id>/
│   ├── consensus_reports/
│   │   └── agentic_ingestion/<attempt_id>/
│   ├── review_reports/
│   │   └── agentic_ingestion/<attempt_id>/
│   └── run_summaries/
│       └── agentic_ingestion/<attempt_id>/
└── work/
\end{verbatim}

The work directory is temporary and non-authoritative.

Published artifacts are immutable and shall be referenced by exact P5 \texttt{ProcessingArtifactReference} values containing:

\begin{verbatim}
artifact_type
artifact_id
content_fingerprint
repository_relative_path
\end{verbatim}

The artifact type shall reflect the actual artifact semantics. Review reports and run summaries shall not be mislabeled as consensus reports or agent outputs.

\hypertarget{8-processing-event-sequence}{%
\paragraph{8. Processing Event sequence}\label{8-processing-event-sequence}}

A normal first execution shall produce the following conceptual event sequence:

\begin{verbatim}
EVT-000001
event_type: run_created
previous_state: null
next_state: created

EVT-000002
event_type: stage_started
previous_state: created
next_state: running
processing_stage: agentic_ingestion
attempt_id: ATT-000001

EVT-000003
event_type: artifact_published
previous_state: running
next_state: running
processing_stage: agentic_ingestion
attempt_id: ATT-000001
artifact_references:
  - Agent Outputs
  - Consensus Reports
  - Review Report
  - Run Summaries

EVT-000004
event_type: review_requested
previous_state: running
next_state: awaiting_review
processing_stage: agentic_ingestion
attempt_id: ATT-000001
\end{verbatim}

Event IDs, sequences, fingerprints and previous-event fingerprints remain owned by P5.

P9 shall not persist mutable current-state files. Current state shall continue to be derived from the immutable Event History.

\hypertarget{9-output-validation-and-publication-transaction}{%
\paragraph{9. Output validation and publication transaction}\label{9-output-validation-and-publication-transaction}}

Pipeline output generation and authoritative artifact publication are separate steps.

The Phase F pipeline shall initially write into the run-owned temporary work directory.

Before publication, P9 shall validate the complete required output set.

Validation shall include at least:

\begin{verbatim}
all required files exist
all required files are regular files
no symbolic-link output is accepted
all output paths remain within the run work directory
all output files are readable
all published contents receive SHA-256 fingerprints
artifact identifiers are deterministic and valid
repository-relative target paths are safe
\end{verbatim}

No \texttt{artifact\_published} event shall be appended before the complete required output set has been validated and copied into its final immutable artifact directories.

If output generation is incomplete or validation fails:

\begin{verbatim}
no generated output is treated as published evidence
\end{verbatim}

Temporary work content may remain available for explicit recovery diagnostics, but it shall not be interpreted by P5, P6 or P7 as active published evidence.

\hypertarget{10-failure-and-recovery-behavior}{%
\paragraph{10. Failure and recovery behavior}\label{10-failure-and-recovery-behavior}}

The workflow cannot be one filesystem-atomic transaction because Source registration, Processing Run persistence, external LLM execution and artifact publication are separate durable operations.

The following behavior is required.

\hypertarget{source-registered-run-not-created}{%
\subparagraph{Source registered, Run not created}\label{source-registered-run-not-created}}

\begin{verbatim}
Source remains registered
Source Processing State remains not_started
No Processing Run is inferred
\end{verbatim}

\hypertarget{run-created-ingestion-execution-fails}{%
\subparagraph{Run created, ingestion execution fails}\label{run-created-ingestion-execution-fails}}

P9 shall append:

\begin{verbatim}
event_type: run_failed
next_state: failed
reason_code: team_agentic_ingestion_failed
\end{verbatim}

No unvalidated output shall be published.

\hypertarget{output-validation-fails}{%
\subparagraph{Output validation fails}\label{output-validation-fails}}

P9 shall append a failed or blocked transition with a stable reason code.

No \texttt{artifact\_published} event shall be written.

\hypertarget{artifact-files-published-publication-event-fails}{%
\subparagraph{Artifact files published, publication event fails}\label{artifact-files-published-publication-event-fails}}

This is a recovery-requiring partial transaction.

P9 shall raise or persist a recovery-required condition. It shall not report success.

The next scan shall expose the inconsistency as a blocking issue until recovery is completed.

\hypertarget{publication-event-succeeds-review-request-event-fails}{%
\subparagraph{Publication event succeeds, review-request event fails}\label{publication-event-succeeds-review-request-event-fails}}

The artifacts remain published and traceable.

The run shall not be reported as successfully awaiting review.

A recovery path shall append the missing valid transition or explicitly fail the run according to the accepted P5 transition rules.

\hypertarget{unexpected-exception}{%
\subparagraph{Unexpected exception}\label{unexpected-exception}}

Internal exception details shall not be exposed as authoritative user-facing state.

The integration result and UI shall report a safe failure summary while retaining technical diagnostics in appropriate developer logs or recovery evidence.

\hypertarget{11-retry-and-successor-rules}{%
\paragraph{11. Retry and successor rules}\label{11-retry-and-successor-rules}}

An unchanged material run binding shall use a retry within the same Processing Run.

A retry shall create a new attempt:

\begin{verbatim}
ATT-000002
ATT-000003
...
\end{verbatim}

and shall use the existing P5 retry operation.

Material changes require a successor Processing Run.

Material bindings include:

\begin{verbatim}
source_id
source_sha256
source_role_snapshot
workflow_profile
configuration_fingerprint
framework_template_id
framework_template_version
semantic_reference_versions
\end{verbatim}

Examples:

\begin{verbatim}
same Source and same configuration
→ retry in existing Run

changed model, recipe, team scope or runs per member
→ successor Run

changed Source role
→ successor Run

changed Source content
→ newly registered Source and new Run

changed Framework Template or semantic-reference version
→ successor Run
\end{verbatim}

P9 shall not create multiple concurrent current Runs for the same Source without an explicit valid supersession relationship.

\hypertarget{12-human-review-boundary}{%
\paragraph{12. Human Review boundary}\label{12-human-review-boundary}}

A successful P9 execution requests review but does not itself create an approved engineering decision.

The P9 result remains:

\begin{verbatim}
unreviewed
awaiting_review
\end{verbatim}

P9 shall not automatically create:

\begin{verbatim}
Approved Input
approved Human Review Decisions
approved terminology mappings
approved Framework Assignments
Approved Generation Readiness
model candidates
SysML v2
CATIA changes
\end{verbatim}

Existing and later Human Review contracts remain authoritative.

\hypertarget{13-coverage-boundary}{%
\paragraph{13. Coverage boundary}\label{13-coverage-boundary}}

P9 Processing visibility and P6 Preliminary Coverage are separate.

A successful project-bound ingestion run does not create Preliminary Coverage unless valid P4 artifacts required by P6 actually exist, including the required Information Units, Framework Assignment Candidates, reference validations and Human Review Decisions.

The Project Dashboard may therefore legitimately show:

\begin{verbatim}
Processing State:
awaiting_review

Preliminary Coverage:
uncovered

Potential Model Support:
not_supported
\end{verbatim}

This is not an inconsistency.

\hypertarget{14-dashboard-return-and-refresh}{%
\paragraph{14. Dashboard return and refresh}\label{14-dashboard-return-and-refresh}}

After returning from Agentic Ingestion, the application shall discard stale P7 presentation state for the affected project.

The dashboard shall regenerate its view from the P2-P6 authorities.

The P9 result object shall not become presentation truth.

The application may preserve stable navigation state such as:

\begin{verbatim}
selected project
return view
selected Source or Processing Run identity
\end{verbatim}

It shall clear open Evidence References that belong to a different project.

\hypertarget{15-project-isolation-and-path-safety}{%
\paragraph{15. Project isolation and path safety}\label{15-project-isolation-and-path-safety}}

Every P9 operation shall be explicitly project-bound.

The integration shall reject:

\begin{verbatim}
invalid Project IDs
unavailable Project Workspaces
cross-project Source references
cross-project Processing Run references
absolute published paths
parent traversal
symbolic-link escapes
repository escapes
artifact paths outside the selected project and run
\end{verbatim}

No P9 operation may silently fall back to repository-global Phase F folders when project-bound execution was requested.

\hypertarget{16-result-contract}{%
\paragraph{16. Result contract}\label{16-result-contract}}

The project-bound result shall be immutable and contain stable identities and safe evidence references.

It shall include at least:

\begin{verbatim}
project_id
source_id
processing_run_id
attempt_id
run_state
processing_stage
dry_run
published artifact references
safe failure or recovery status
\end{verbatim}

It shall not contain:

\begin{verbatim}
API keys
unrestricted filesystem paths
mutable Streamlit objects
raw exception objects as persisted state
\end{verbatim}

\hypertarget{17-ui-behavior}{%
\paragraph{17. UI behavior}\label{17-ui-behavior}}

The P9 execution interface shall make the active Project visible.

The intended demonstrator flow is:

\begin{verbatim}
Create or select Project
→ open Agentic Ingestion
→ upload legacy Source
→ select Source role
→ configure execution
→ register Source
→ start Processing Run
→ execute Team Agentic Ingestion
→ inspect execution result
→ return to Project Dashboard
→ inspect Source, Run, state and published evidence
\end{verbatim}

Real LLM execution shall continue to require explicit human confirmation.

Dry-run mode shall remain available and shall be visibly distinguished from an engineering assessment.

\hypertarget{18-no-project-lifecycle-management-in-p9}{%
\paragraph{18. No Project Lifecycle Management in P9}\label{18-no-project-lifecycle-management-in-p9}}

P9 shall not add:

\begin{verbatim}
Project display-name editing
Project description editing
Project deletion
Project archival
project-wide destructive lifecycle operations
\end{verbatim}

These capabilities require a separate Project Lifecycle Management decision and are not required for the P9 demonstrator.

\hypertarget{implementation-sequence-2}{%
\subsubsection{Implementation Sequence}\label{implementation-sequence-2}}

P9 shall proceed in six steps:

\begin{verbatim}
P9 Step 1 of 6
Integration Architecture and ADR-015

P9 Step 2 of 6
Common Turing Generator Navigation

P9 Step 3 of 6
Project-bound Source Upload and Source Registration

P9 Step 4 of 6
Bridge between Phase F Ingestion and P5 Processing Runs

P9 Step 5 of 6
Project-bound Artifacts, Dashboard Return and Refresh

P9 Step 6 of 6
End-to-End Demonstration, Full Regression and Phase-P Completion Review
\end{verbatim}

Implementation shall not proceed by creating a second Project Manifest, Source Registry, Processing State model or dashboard authority.

\hypertarget{consequences}{%
\subsubsection{Consequences}\label{consequences}}

\hypertarget{positive-consequences-1}{%
\paragraph{Positive consequences}\label{positive-consequences-1}}

\begin{itemize}
\tightlist
\item
  The existing Phase F capability becomes demonstrable inside a real Project Workspace.
\item
  P2-P7 remain authoritative.
\item
  The integration is narrow and testable.
\item
  Source, Run, Attempt and artifact identities become fully traceable.
\item
  Existing Phase F behavior remains backward compatible.
\item
  Failures remain visible rather than being converted into false success.
\item
  The Project Dashboard can inspect project-bound ingestion evidence without becoming writable.
\item
  Processing visibility remains semantically separate from engineering approval and model readiness.
\item
  Retry and supersession use the existing P5 lifecycle.
\item
  The implementation can later be extended toward P4 artifact production without replacing the P9 project-binding layer.
\end{itemize}

\hypertarget{negative-consequences}{%
\paragraph{Negative consequences}\label{negative-consequences}}

\begin{itemize}
\tightlist
\item
  The integration spans several durable operations and cannot be completely filesystem-atomic.
\item
  Explicit recovery handling is required for partial publication failures.
\item
  P5 Processing Stage and artifact-kind vocabularies must be extended.
\item
  The Phase F pipeline requires a backward-compatible execution-root adaptation.
\item
  A common application shell adds navigation and session-state complexity.
\item
  The first P9 demonstrator may show successful ingestion while Preliminary Coverage remains uncovered.
\item
  Project metadata editing and deletion remain unavailable.
\end{itemize}

\hypertarget{rejected-alternatives}{%
\subsubsection{Rejected Alternatives}\label{rejected-alternatives}}

\hypertarget{add-only-a-link-between-the-two-existing-streamlit-apps}{%
\paragraph{Add only a link between the two existing Streamlit apps}\label{add-only-a-link-between-the-two-existing-streamlit-apps}}

Rejected because navigation alone would not bind Sources, Processing Runs, Attempts or generated artifacts to the selected Project Workspace.

\hypertarget{move-all-p5-processing-logic-into-the-streamlit-ui}{%
\paragraph{Move all P5 Processing logic into the Streamlit UI}\label{move-all-p5-processing-logic-into-the-streamlit-ui}}

Rejected because UI code must not become the owner of Processing lifecycle, persistence or recovery rules.

\hypertarget{write-project-bound-artifacts-directly-into-global-phase-f-folders}{%
\paragraph{Write project-bound artifacts directly into global Phase F folders}\label{write-project-bound-artifacts-directly-into-global-phase-f-folders}}

Rejected because global folders cannot provide project isolation or authoritative P5 artifact ownership.

\hypertarget{treat-phase-f-run-directories-as-p5-processing-runs}{%
\paragraph{Treat Phase F run directories as P5 Processing Runs}\label{treat-phase-f-run-directories-as-p5-processing-runs}}

Rejected because Phase F run identity and P5 Processing Run identity have different contracts and lifecycle semantics.

\hypertarget{copy-outputs-directly-into-final-artifact-directories-during-pipeline-execution}{%
\paragraph{Copy outputs directly into final artifact directories during pipeline execution}\label{copy-outputs-directly-into-final-artifact-directories-during-pipeline-execution}}

Rejected because partially generated output could appear as published evidence before the complete output set has been validated.

\hypertarget{mark-a-successful-ingestion-run-as-completed}{%
\paragraph{Mark a successful ingestion run as completed}\label{mark-a-successful-ingestion-run-as-completed}}

Rejected because the generated reports and agent outputs remain unreviewed. The correct initial terminal state is \texttt{awaiting\_review}.

\hypertarget{treat-successful-ingestion-as-preliminary-coverage}{%
\paragraph{Treat successful ingestion as Preliminary Coverage}\label{treat-successful-ingestion-as-preliminary-coverage}}

Rejected because P6 requires specific valid P4 evidence and Human Review contracts.

\hypertarget{start-a-new-run-for-every-technical-retry}{%
\paragraph{Start a new Run for every technical retry}\label{start-a-new-run-for-every-technical-retry}}

Rejected because unchanged material bindings are handled by P5 Attempts within the same Run.

\hypertarget{reuse-the-same-run-after-material-configuration-changes}{%
\paragraph{Reuse the same Run after material configuration changes}\label{reuse-the-same-run-after-material-configuration-changes}}

Rejected because the immutable Run Manifest must preserve reproducibility. Material changes require a successor Run.

\hypertarget{store-api-keys-in-the-run-manifest-or-configuration-fingerprint}{%
\paragraph{Store API keys in the Run Manifest or configuration fingerprint}\label{store-api-keys-in-the-run-manifest-or-configuration-fingerprint}}

Rejected because secrets must never be persisted as project evidence.

\hypertarget{implement-project-editing-and-deletion-in-p9}{%
\paragraph{Implement Project editing and deletion in P9}\label{implement-project-editing-and-deletion-in-p9}}

Rejected because Project Lifecycle Management is separate from project-bound ingestion integration and destructive deletion requires its own safety, transaction and audit decision.

\hypertarget{acceptance-criteria}{%
\subsubsection{Acceptance Criteria}\label{acceptance-criteria}}

ADR-015 is satisfied when:

\begin{enumerate}
\def\labelenumi{\arabic{enumi}.}
\tightlist
\item
  a common application shell exposes Project Dashboard and project-bound Agentic Ingestion views;
\item
  every P9 execution requires one valid selected six-digit Project ID;
\item
  uploaded content is registered through P3 before processing;
\item
  the registered project-local Source content is the authoritative pipeline input;
\item
  one P5 Processing Run is created with immutable material bindings;
\item
  \texttt{agentic\_ingestion} is a canonical Processing Stage;
\item
  Agent Outputs, Consensus Reports, Review Reports and Run Summaries have distinct run-owned artifact kinds;
\item
  the Phase F pipeline remains backward compatible when no project execution root is supplied;
\item
  project-bound execution uses the selected Run work directory;
\item
  a complete first execution produces created, running, artifact-published and review-requested lifecycle evidence;
\item
  a successful execution ends in \texttt{awaiting\_review};
\item
  failed execution produces a visible P5 failure state;
\item
  incomplete or invalid output is never published;
\item
  every published artifact has an exact fingerprinted \texttt{ProcessingArtifactReference};
\item
  partial publication failures produce explicit recovery behavior rather than false success;
\item
  unchanged material bindings use a retry Attempt in the same Run;
\item
  material binding changes use a valid successor Run;
\item
  P9 does not create Approved Input, Approved Generation Readiness, model candidates or SysML v2;
\item
  P9 execution does not imply Preliminary Coverage;
\item
  returning to the dashboard regenerates P7 views from P2-P6 authorities;
\item
  cross-project references and unsafe paths are rejected;
\item
  API keys and unrestricted filesystem paths are not persisted;
\item
  core integration behavior is testable without Streamlit;
\item
  the existing Phase F regression remains green;
\item
  the complete repository regression remains green;
\item
  an end-to-end dry-run demonstrator can execute the intended project-bound flow;
\item
  Project metadata editing and Project deletion remain outside P9.
\end{enumerate}

\textbackslash newpage

\hypertarget{adr-016--human-review-workspace-and-approved-input-promotion-architecture}{%
\subsection{ADR-016 --- Human Review Workspace and Approved Input Promotion Architecture}\label{adr-016--human-review-workspace-and-approved-input-promotion-architecture}}

Human Review Workspace and Approved Input Promotion Architecture

Status

Accepted

Date

2026-07-31

\hypertarget{context-1}{%
\subsubsection{Context}\label{context-1}}

Phase P established the project-oriented processing and evidence architecture of the Turing Generator.

The relevant existing capabilities are:

\begin{itemize}
\tightlist
\item
  project-isolated Source registration;
\item
  immutable Processing Runs and Processing Events;
\item
  immutable run-owned Agent Outputs, Consensus Reports, Review Reports and Run Summaries;
\item
  source projections and independently reviewable Information Units;
\item
  terminology and Framework Assignment Candidates;
\item
  immutable Human Review Decisions bound to exact content and validation fingerprints;
\item
  project-level Processing and evidence presentation;
\item
  project-bound Agentic Ingestion ending in \texttt{awaiting\_review}.
\end{itemize}

A completed Processing Run produces authoritative evidence of what the system processed and generated. It does not produce approved engineering information.

The authoritative CATIA model assigns the required behavior to:

\begin{verbatim}
SF_007 Support Human Review and Approval
\end{verbatim}

and to the Logical Component:

\begin{verbatim}
LC_05 Candidate and Review Governance
\end{verbatim}

The primary applicable System Requirements include:

\begin{verbatim}
SYSR_014 Provide Generated Interpretations for Expert Review
SYSR_015 Require Review Before Subsequent Engineering Use
SYSR_016 Support Human Resolution of Conflicting Interpretations
SYSR_017 Record Human Conflict Resolution
SYSR_038 Capture Human Review Decision
SYSR_095 Bind Review Decision to Reviewed Content
SYSR_096 Require New Review after Content Change
SYSR_097 Retain Review Decision Accountability
SYSR_098 Distinguish Engineering Information Authority State
SYSR_099 Promote Approved Engineering Information
SYSR_100 Prevent Unauthorized Engineering Use
SYSR_101 Prevent Automated Approval Substitution
\end{verbatim}

The applicable SysML v2 representation requirements include:

\begin{verbatim}
SYSR_024 Generate SysML v2 Textual Representation
SYSR_072 Validate SysML v2 Textual Conformance
\end{verbatim}

The current Phase F/P review artifacts may contain dozens of detected engineering items. Different agents may produce equivalent, competing, partially overlapping or differently classified proposals.

A useful Human Review workflow therefore cannot treat a complete report as one binary approval target.

The user must be able to:

\begin{itemize}
\tightlist
\item
  review all agent proposals for the same detected subject;
\item
  accept one proposal with one click;
\item
  automatically mark competing variants as not selected;
\item
  edit a proposal inline;
\item
  combine content from multiple proposals;
\item
  reject one proposal or the complete detected subject;
\item
  apply document-level classification decisions where a document is homogeneous;
\item
  override document-level decisions for individual items;
\item
  apply focused decisions to a visible filtered result set;
\item
  review elements and relationships separately;
\item
  resolve open questions without losing overview;
\item
  inspect rejected content;
\item
  compare the original machine-generated report with the human-reviewed version;
\item
  continue editing until the review is explicitly finalized;
\item
  reopen a finalized document only by creating a documented successor version.
\end{itemize}

The original Source, Agent Outputs, Consensus Reports and Review Reports must remain immutable.

A human-reviewed document has higher engineering authority than its original machine-generated report only after explicit finalization and an exact persisted Human Review Decision.

Phase G must create the authoritative bridge:

\begin{verbatim}
Processing Evidence
→ Human Review Workspace
→ Finalized Reviewed Document
→ Approved Input
\end{verbatim}

Phase G shall not generate model candidates, an internal engineering model or SysML v2 architecture artifacts.

\hypertarget{decision-1}{%
\subsubsection{Decision}\label{decision-1}}

\hypertarget{1-implement-phase-g-as-the-concrete-realization-of-sf_007}{%
\paragraph{1. Implement Phase G as the concrete realization of SF\_007}\label{1-implement-phase-g-as-the-concrete-realization-of-sf_007}}

Phase G shall implement and decompose:

\begin{verbatim}
SF_007 Support Human Review and Approval
\end{verbatim}

inside the accepted responsibility of:

\begin{verbatim}
LC_05 Candidate and Review Governance
\end{verbatim}

No new top-level System Function or Logical Component is introduced by this decision.

The Phase G implementation shall provide:

\begin{enumerate}
\def\labelenumi{\arabic{enumi}.}
\tightlist
\item
  a versioned Human Review Workspace;
\item
  exact preservation of machine-generated review evidence;
\item
  document-, filtered-set- and item-level review operations;
\item
  explicit finalization of one reviewed document version;
\item
  exact Human Review Decision binding;
\item
  Approved Input promotion for accepted review items;
\item
  Approved Input invalidation, revocation and supersession;
\item
  a stable Approved Input read contract for Phase H.
\end{enumerate}

\hypertarget{2-preserve-an-explicit-authority-hierarchy}{%
\paragraph{2. Preserve an explicit authority hierarchy}\label{2-preserve-an-explicit-authority-hierarchy}}

The authority hierarchy for Phase G shall be:

\begin{verbatim}
Original Source
    immutable source authority

Agent Outputs and Consensus Reports
    immutable processing evidence

Original Review Report
    immutable machine-generated review presentation

Draft Review Version
    human working state, not approved engineering information

Finalized Reviewed Document Version
    immutable human-reviewed result

Approved Input
    authoritative project-local input for subsequent engineering processing
\end{verbatim}

No higher level may overwrite a lower level.

Human review shall create additional versioned artifacts and references. It shall never modify the original Source, processing artifacts, prior finalized review versions, prior Human Review Decisions, Approved Input manifests or Approved Input lifecycle events.

\hypertarget{3-introduce-human-review-workspace-identities}{%
\paragraph{3. Introduce Human Review Workspace identities}\label{3-introduce-human-review-workspace-identities}}

Phase G shall introduce the following project-local identities:

\begin{verbatim}
RVD-000001  Review Document
RVV-000001  Review Document Version
RVR-000001  Review Revision
RIT-000001  Review Item
SRA-000001  Scoped Review Action
AIN-000001  Approved Input
AIE-000001  Approved Input Event
\end{verbatim}

Identifiers are project-local, sequential, immutable and never reused.

\hypertarget{review-document}{%
\subparagraph{Review Document}\label{review-document}}

A Review Document identifies one review workspace created from one exact eligible evidence set.

It binds at least:

\begin{verbatim}
project_id
review_document_id
source_id
source_sha256
processing_run_id
attempt_id
primary_review_artifact_reference
supporting_artifact_references
framework_template_reference
semantic_reference_versions
created_at
content_fingerprint
\end{verbatim}

\hypertarget{review-document-version}{%
\subparagraph{Review Document Version}\label{review-document-version}}

A Review Document Version is one human-review version of a Review Document.

It binds at least:

\begin{verbatim}
review_document_version_id
review_document_id
version_number
predecessor_version_id
reopen_reason
opened_by
opened_at
version_state
head_revision_id
finalized_revision_id
finalized_at
finalization_decision_id
content_fingerprint
\end{verbatim}

Allowed version states are:

\begin{verbatim}
draft
finalized
\end{verbatim}

A finalized version is immutable.

\hypertarget{review-revision}{%
\subparagraph{Review Revision}\label{review-revision}}

A Review Revision is one immutable saved snapshot of a draft Review Document Version.

The user experience may present one continuously editable working copy. The persistence model shall remain append-only:

\begin{verbatim}
save draft
→ create next immutable Review Revision
→ preserve all previous revisions
\end{verbatim}

The current draft state is the latest valid revision derived from repository history. A mutable authoritative \texttt{current.json} file shall not be introduced.

\hypertarget{review-item}{%
\subparagraph{Review Item}\label{review-item}}

A Review Item represents one independently reviewable subject.

Allowed Review Item kinds are:

\begin{verbatim}
element
relationship
open_question
\end{verbatim}

\texttt{rejected\_content} is a presentation view derived from item decisions. It is not a separate source item kind.

One Review Item may contain multiple Agent Proposal References.

A Review Item shall preserve at least:

\begin{verbatim}
review_item_id
review_item_kind
stable_subject_key
section
original_report_locator
proposal_references
source_evidence_references
consensus_evidence_references
current_human_content
current_classification
current_framework_assignment
current_terminology_assignment
current_source_assignments
current_relationship_representation
effective_review_outcome
item_content_fingerprint
\end{verbatim}

The \texttt{stable\_subject\_key} supports continuity across document versions.

The system shall support explicit human correction when automatic grouping is wrong:

\begin{verbatim}
split Review Item
merge selected Review Items
\end{verbatim}

Split and merge operations shall preserve all original proposal references and shall be represented in Review Revision history.

\hypertarget{scoped-review-action}{%
\subparagraph{Scoped Review Action}\label{scoped-review-action}}

A Scoped Review Action represents one draft decision applied to more than one Review Item or to one review dimension.

It shall bind:

\begin{verbatim}
scoped_review_action_id
review_document_version_id
action_scope
decision_dimension
selected_value
filter_definition
materialized_review_item_ids
materialized_item_fingerprints
created_by
created_at
rationale
action_fingerprint
\end{verbatim}

Allowed action scopes are:

\begin{verbatim}
document_default
filtered_set
explicit_selection
\end{verbatim}

A filtered action shall store both the filter definition shown to the user and the exact Review Item IDs and fingerprints matched at action time.

It shall never remain a dynamic query that silently changes its effect when later Review Items are added, removed or edited.

\hypertarget{4-create-review-documents-from-exact-eligible-evidence}{%
\paragraph{4. Create Review Documents from exact eligible evidence}\label{4-create-review-documents-from-exact-eligible-evidence}}

Phase G shall support Review Documents derived from eligible P4 and P9 evidence.

\hypertarget{p4-evidence}{%
\subparagraph{P4 evidence}\label{p4-evidence}}

Eligible P4 review sources include:

\begin{verbatim}
Information Units
Terminology Mapping Candidates
Framework Assignment Candidates
associated reference validations
associated Human Review evidence
\end{verbatim}

An Information Unit remains the smallest independently understandable and reviewable source-derived claim.

Terminology and Framework Assignment Candidates are supporting classifications. They do not automatically become standalone Approved Inputs.

\hypertarget{p9-evidence}{%
\subparagraph{P9 evidence}\label{p9-evidence}}

Eligible P9 review sources include one exact active project-bound evidence set:

\begin{verbatim}
primary deterministic Review Report
underlying Agent Outputs
Consensus Reports
Run Summary
Processing Artifact References
Processing Run and Attempt identity
\end{verbatim}

The system shall construct Review Items from structured evidence. It shall not depend on free-form Markdown parsing as the sole source of Review Item identity.

The original Markdown Review Report remains the immutable human-readable source view.

The Review Workspace shall render structured Review Items inside a report-like document view so that editing feels local to the relevant report section.

A complete report shall never become one Approved Input.

One accepted Review Item may become one Approved Input.

\hypertarget{non-promotable-evidence}{%
\subparagraph{Non-promotable evidence}\label{non-promotable-evidence}}

The following are evidence or review-control information and shall not directly become Approved Input:

\begin{verbatim}
Agent confidence
Consensus level
variance indication
model buildability assessment
general recommendation
gap
risk
ambiguity
unanswered review question
run summary
processing completion
artifact publication
\end{verbatim}

An answer to an open question may become promotable only when the reviewer explicitly converts it into a reviewed engineering statement.

\hypertarget{5-present-all-agent-proposals-without-losing-evidence}{%
\paragraph{5. Present all Agent proposals without losing evidence}\label{5-present-all-agent-proposals-without-losing-evidence}}

For each Review Item, the user interface shall present all associated Agent Proposal References.

The UI may visually consolidate materially equivalent variants. It shall retain exact access to every original proposal.

Each proposal view shall expose at least:

\begin{verbatim}
agent identity
persona identity
candidate identity
proposed content
proposed classification
proposed Framework Assignment
proposed source assignments
rationale
confidence
generation-readiness statement
supporting and missing evidence
artifact reference
content fingerprint
\end{verbatim}

The following quick actions are required for each proposal:

\begin{verbatim}
Accept proposal
Edit and accept
Reject proposal
\end{verbatim}

The following item-level actions are required:

\begin{verbatim}
Combine proposals
Reject all proposals for this Review Item
Defer Review Item
Mark Review Item out of scope
Split Review Item
Merge selected Review Items
\end{verbatim}

Selecting one proposal shall mark competing proposals for the same Review Item as:

\begin{verbatim}
not_selected_due_to_human_selection
\end{verbatim}

This state is distinct from a finding that a competing proposal is factually incorrect.

Proposal-selection actions remain draft operations until document finalization.

\hypertarget{6-support-inline-editing-through-human-revisions}{%
\paragraph{6. Support inline editing through human revisions}\label{6-support-inline-editing-through-human-revisions}}

The Review Workspace shall allow direct inline editing of the current reviewed content.

The UI shall make the edited block appear in the position of the corresponding Review Item in the report-like document view.

Technically, the system shall create human-authored content in a new Review Revision. It shall not modify the original Agent Proposal.

Editable review dimensions include:

\begin{verbatim}
name
engineering statement
description
classification
information type
modality
epistemic status
Framework Assignment
terminology assignment
source assignments
human rationale
human confidence assessment
relationship representation
\end{verbatim}

Agent-generated fields remain immutable evidence. Human-edited values are explicit overrides.

A human revision shall preserve:

\begin{verbatim}
derived_from_proposal_ids
derived_from_review_item_ids
original content fingerprints
revised content fingerprint
changed fields
reviewer identity
change time
optional change rationale
\end{verbatim}

A rationale is mandatory when:

\begin{itemize}
\tightlist
\item
  a finalized document is reopened;
\item
  a proposal or complete Review Item is rejected;
\item
  a previous Approved Input is revoked;
\item
  a selected SysML v2 relationship construct differs materially from all proposed constructs;
\item
  a blocking finding is overridden where an override is permitted.
\end{itemize}

\hypertarget{7-separate-review-dimensions-and-define-precedence}{%
\paragraph{7. Separate review dimensions and define precedence}\label{7-separate-review-dimensions-and-define-precedence}}

Review decisions shall not collapse all aspects of one Review Item into one undifferentiated state.

The system shall support independent review dimensions, including:

\begin{verbatim}
content
classification
Framework Assignment
terminology assignment
source assignment
relationship representation
review outcome
\end{verbatim}

A document-level decision may set a default for one dimension without approving all content.

Example:

\begin{verbatim}
Document default:
Framework Assignment = System Requirements
\end{verbatim}

This does not automatically approve all Requirement statements.

Effective decision precedence is:

\begin{verbatim}
item override
> explicit-selection or materialized filtered-set action
> document default
> selected Agent proposal
\end{verbatim}

The UI shall show the origin of every effective value.

Recommended compact indicators are:

\begin{verbatim}
E  item-level explicit decision
F  materialized filtered-set or explicit-selection decision
D  document default
A  Agent proposal
\end{verbatim}

A scoped action shall not overwrite an existing higher-precedence decision unless the user explicitly requests that overwrite and confirms a preview of the affected Review Items.

\hypertarget{8-use-filters-as-focused-review-tools}{%
\paragraph{8. Use filters as focused review tools}\label{8-use-filters-as-focused-review-tools}}

Required filter dimensions include at least:

\begin{verbatim}
review status
Review Item kind
proposed classification
effective classification
proposed Framework Assignment
effective Framework Assignment
Agent identity
confidence
Consensus state
Agent disagreement
human modification state
Source identity
evidence sufficiency
relationship validation status
\end{verbatim}

Every scoped action shall visibly state whether it applies to:

\begin{verbatim}
the complete document
the currently materialized filtered result
an explicit manual selection
one Review Item
\end{verbatim}

Before applying a scoped action, the UI shall show:

\begin{verbatim}
number of affected Review Items
number with existing item overrides
number excluded because of higher-precedence decisions
number that would be overwritten after explicit confirmation
\end{verbatim}

Dangerous bulk rejection shall not be implemented as an unqualified one-click action.

A rejection affecting more than one Review Item requires a materialized target list, an impact preview, explicit confirmation and a rationale.

\hypertarget{9-use-four-primary-review-sections}{%
\paragraph{9. Use four primary review sections}\label{9-use-four-primary-review-sections}}

The Human Review Workspace shall use:

\begin{verbatim}
Elements
Relationships
Open Questions
Rejected Content
\end{verbatim}

\hypertarget{elements}{%
\subparagraph{Elements}\label{elements}}

The Elements section presents detected engineering-element proposals.

It shall support grouping and filtering by applicable model area, information type and classification.

The section shall not imply that Phase G has generated final model elements. It presents reviewed engineering information and modeling suggestions.

\hypertarget{relationships}{%
\subparagraph{Relationships}\label{relationships}}

The Relationships section presents only explicit relationship proposals already contained in eligible evidence.

Phase G shall not invent missing relationships.

The section may be empty when the source evidence does not contain an explicit relationship proposal.

Relationship proposals remain separately reviewable from element proposals.

\hypertarget{open-questions}{%
\subparagraph{Open Questions}\label{open-questions}}

The Open Questions section presents:

\begin{verbatim}
gaps
ambiguities
risks requiring a decision
missing evidence
review questions
unresolved classifications
unresolved relationship constructs
\end{verbatim}

An open question may be answered, closed as not relevant, deferred, marked out of scope, linked to Review Items or converted into a human-authored engineering statement.

\hypertarget{rejected-content}{%
\subparagraph{Rejected Content}\label{rejected-content}}

Rejected Content is an auditable view, not a deletion area.

It shall show the rejected proposal or item, original Agent and artifact reference, rejection rationale, replacement where applicable, document version, reviewer and time.

\hypertarget{10-require-sysml-v2-target-notation-conformance-for-relationships}{%
\paragraph{10. Require SysML v2 target-notation conformance for relationships}\label{10-require-sysml-v2-target-notation-conformance-for-relationships}}

The Relationship review interface shall use only constructs defined by the applicable versioned SysML v2 target-notation profile of the Turing Generator.

The UI shall not expose generic MBSE relationship labels or SysML v1 relationship names as authoritative relationship types.

The allowed construct vocabulary shall be loaded from the selected versioned target-notation profile. It shall not be maintained as an independent hard-coded UI list.

A relationship proposal shall preserve at least:

\begin{verbatim}
source Review Item or referenced element
target Review Item or referenced element
relationship semantic intent
selected SysML v2 construct
construct-specific properties
target-notation profile ID
target-notation profile version
textual-notation preview
supporting evidence
alternative proposed constructs
profile-validation status
profile-validation fingerprint
human review status
\end{verbatim}

A textual preview shall use the applicable SysML v2 notation.

Example for an applicable dependency construct:

\begin{verbatim}
dependency from 'Source Element' to 'Target Element';
\end{verbatim}

The exact syntax shall be produced and validated by the applicable target notation profile, not by free-form UI concatenation.

When no supported construct can be selected, the item state shall be:

\begin{verbatim}
unresolved_relationship_candidate
\end{verbatim}

An unresolved relationship candidate remains visible, may be edited, deferred, marked out of scope or rejected, and shall not become Approved Input for model-generation use.

This decision does not move SysML v2 artifact generation into Phase G.

Phase G reviews an existing relationship proposal and its intended applicable construct. Phases H to J remain responsible for model candidate creation, internal model representation and final SysML v2 artifact generation.

\hypertarget{11-keep-draft-actions-reversible-until-explicit-finalization}{%
\paragraph{11. Keep draft actions reversible until explicit finalization}\label{11-keep-draft-actions-reversible-until-explicit-finalization}}

The following actions modify only the current draft version:

\begin{verbatim}
accept proposal
edit proposal
combine proposals
reject proposal
reject Review Item
change classification
change Framework Assignment
change terminology assignment
change source assignment
change relationship representation
answer question
defer
mark out of scope
split
merge
apply document default
apply materialized filtered-set action
\end{verbatim}

These actions shall be automatically saved through immutable Review Revisions. They are reversible while the version remains \texttt{draft}.

They do not create Approved Input or an authoritative final Human Review Decision.

\hypertarget{12-finalize-one-exact-reviewed-document-version}{%
\paragraph{12. Finalize one exact reviewed document version}\label{12-finalize-one-exact-reviewed-document-version}}

The user finalizes a draft through an explicit operation conceptually labeled:

\begin{verbatim}
Review abschließen und freigeben
\end{verbatim}

Before finalization, the system shall present a summary including at least:

\begin{verbatim}
total Review Items
accepted as generated
accepted with human modification
combined
rejected
deferred
out of scope
unresolved
relationship validation results
open questions
changed items compared with predecessor version
\end{verbatim}

Finalization shall fail closed when:

\begin{itemize}
\tightlist
\item
  the Project, Source, Run or Artifact binding is invalid;
\item
  an authoritative evidence fingerprint has changed;
\item
  the selected Review Revision is not the current head revision;
\item
  a Review Item has no explicit or effective review outcome;
\item
  an unresolved item is not explicitly \texttt{deferred} or \texttt{out\_of\_scope};
\item
  a relationship selected for approval is not valid under the applicable SysML v2 target-notation profile;
\item
  a blocking evidence or integrity issue exists;
\item
  a cross-project reference exists;
\item
  a required supporting artifact is unavailable, superseded or invalidated;
\item
  a predecessor-version relation is inconsistent.
\end{itemize}

Finalization shall create:

\begin{verbatim}
immutable Finalized Reviewed Document manifest
immutable rendered Reviewed Report
immutable effective decision set
exact content fingerprint
exact validation fingerprint
\end{verbatim}

The rendered Reviewed Report is the primary human-readable reviewed artifact. The effective decision set is the primary machine-readable record.

The finalization target shall be added to the Human Review target vocabulary:

\begin{verbatim}
target_type: review_document_finalization
target_id: RVV-000001
\end{verbatim}

Finalization requires an exact persisted decision:

\begin{verbatim}
decision: confirm
review_mode: detailed_review
\end{verbatim}

The Human Review Decision shall bind the exact Review Document Version content fingerprint, validation fingerprint, reviewer identity, decision time and outcome.

Consensus, confidence, processing completion and artifact publication cannot substitute this decision.

\hypertarget{13-reopen-finalized-documents-only-through-a-successor-version}{%
\paragraph{13. Reopen finalized documents only through a successor version}\label{13-reopen-finalized-documents-only-through-a-successor-version}}

A finalized Review Document Version shall never become editable again.

A user may reopen the Review Document only by creating a successor Review Document Version.

Reopening requires a rationale and records:

\begin{verbatim}
predecessor version
reopen reason
opened by
opened at
baseline item fingerprints
baseline effective decisions
\end{verbatim}

The successor draft is initialized from the finalized predecessor.

The UI shall provide:

\begin{verbatim}
Review Copy
Original Report
Changes to Original
Changes to Predecessor Version
Version History
\end{verbatim}

A later finalized version shall not rewrite or delete its predecessor.

\hypertarget{14-define-approved-input-as-one-approved-review-item}{%
\paragraph{14. Define Approved Input as one approved Review Item}\label{14-define-approved-input-as-one-approved-review-item}}

One Approved Input represents one independently reviewable and approved engineering-information item.

Allowed initial Approved Input kinds are:

\begin{verbatim}
element_statement
relationship_statement
human_clarification
\end{verbatim}

An Approved Input Manifest shall bind at least:

\begin{verbatim}
schema_version
project_id
approved_input_id
approved_input_kind
authority_state
canonical_content
selected_classification
selected_framework_assignment
selected_terminology_assignment
selected_source_assignments
selected_relationship_representation
review_document_id
review_document_version_id
review_revision_id
review_item_id
review_item_fingerprint
finalization_decision_id
finalization_decision_fingerprint
finalization_validation_fingerprint
source_id
source_sha256
processing_run_id
attempt_id
primary_artifact_reference
supporting_artifact_references
proposal_references
created_at
content_fingerprint
\end{verbatim}

One finalized Reviewed Document Version may create zero, one or many Approved Inputs.

No Approved Input is created for an item whose effective outcome is:

\begin{verbatim}
rejected
deferred
out_of_scope
unresolved
\end{verbatim}

An open question creates no Approved Input unless the reviewer explicitly converts its answer into an approved clarification or engineering statement.

\hypertarget{15-define-approved-input-authority-and-lifecycle}{%
\paragraph{15. Define Approved Input authority and lifecycle}\label{15-define-approved-input-authority-and-lifecycle}}

Approved Input manifests are immutable. Lifecycle changes are represented by immutable Approved Input Events.

Allowed derived Authority States are:

\begin{verbatim}
active
invalidated
revoked
superseded
\end{verbatim}

\hypertarget{active}{%
\subparagraph{Active}\label{active}}

The Approved Input is valid for subsequent engineering use.

\hypertarget{invalidated}{%
\subparagraph{Invalidated}\label{invalidated}}

The Approved Input has lost technical or upstream integrity.

Examples include Source or artifact fingerprint mismatch, invalid Run binding, invalid finalization decision binding, relationship-profile validation failure or project-isolation failure.

Invalidation is system- or integrity-driven.

\hypertarget{revoked}{%
\subparagraph{Revoked}\label{revoked}}

A human-reviewed successor version explicitly withdraws the authority of a previous Approved Input without approving a replacement.

Examples include a previously accepted item that is now rejected, out of scope or explicitly withdrawn.

Revocation requires a rationale and reviewer accountability.

\hypertarget{superseded}{%
\subparagraph{Superseded}\label{superseded}}

A newer active Approved Input explicitly replaces an older Approved Input for the same stable review subject.

The old Approved Input remains immutable and traceable.

\texttt{rejected} is not an Approved Input lifecycle state. Rejected Review Items do not create Approved Inputs.

\hypertarget{16-reconcile-successor-versions-item-by-item}{%
\paragraph{16. Reconcile successor versions item by item}\label{16-reconcile-successor-versions-item-by-item}}

When a successor Review Document Version is finalized, the system shall compare stable Review Item identities and fingerprints with the predecessor version.

Required behavior:

\begin{verbatim}
unchanged accepted item
→ retain existing active Approved Input

changed item accepted again
→ create new Approved Input
→ mark previous Approved Input superseded

previously accepted item now rejected
→ revoke previous Approved Input

previously accepted item now out of scope
→ revoke previous Approved Input

new accepted item
→ create new Approved Input

upstream integrity failure
→ invalidate affected Approved Input
\end{verbatim}

Unchanged items shall not receive new Approved Input IDs merely because the containing document received a new version.

\hypertarget{17-extend-the-existing-human-review-repository}{%
\paragraph{17. Extend the existing Human Review repository}\label{17-extend-the-existing-human-review-repository}}

The existing Human Review Decision repository remains authoritative for final Human Review Decisions.

Phase G shall extend it with:

\begin{verbatim}
review_document_finalization
\end{verbatim}

It shall not introduce a parallel incompatible decision repository.

Existing target types remain valid:

\begin{verbatim}
information_unit_publication
terminology_mapping_candidate
framework_assignment_candidate
\end{verbatim}

The existing exact-confirmation principle remains:

\begin{verbatim}
target ID matches
target content fingerprint matches
validation fingerprint matches
latest exact decision is confirm
target validation is not invalid
\end{verbatim}

A stale Human Review Decision cannot authorize finalization or promotion of changed content.

\hypertarget{18-define-promotion-eligibility}{%
\paragraph{18. Define promotion eligibility}\label{18-define-promotion-eligibility}}

A finalized Review Item is eligible for Approved Input promotion only when all of the following are true:

\begin{verbatim}
Project exists and matches all references
Source exists and exact SHA-256 matches
Processing Run exists and matches Source and Project
Processing Run is valid for promotion
primary artifact reference is active
supporting artifact references are valid
Review Document Version is finalized
exact finalization Human Review Decision exists
exact finalization fingerprints match
effective item outcome is accepted
Review Item fingerprint is included in finalized version
no blocking evidence issue applies
required Framework and terminology references are valid
relationship representation is profile-valid where applicable
no cross-project reference exists
no newer lifecycle event blocks the item
\end{verbatim}

Eligibility shall be recalculated immediately before promotion.

A previously calculated eligibility result is advisory unless it is bound into the finalization validation fingerprint and revalidated at promotion time.

\hypertarget{19-define-project-local-persistence}{%
\paragraph{19. Define project-local persistence}\label{19-define-project-local-persistence}}

The conceptual project-local structure shall be:

\begin{verbatim}
data/projects/<project_id>/
├── reviews/
│   └── RVD-000001/
│       ├── review_document_manifest.json
│       └── versions/
│           ├── RVV-000001/
│           │   ├── review_version_manifest.json
│           │   ├── revisions/
│           │   │   ├── RVR-000001.json
│           │   │   └── RVR-000002.json
│           │   ├── scoped_actions/
│           │   │   └── SRA-000001.json
│           │   └── finalized/
│           │       ├── reviewed_document.json
│           │       ├── effective_decisions.json
│           │       └── reviewed_report.md
│           └── RVV-000002/
│               └── ...
├── semantics/
│   └── human_reviews/
│       └── HRD-000001.json
└── approved_inputs/
    ├── manifests/
    │   └── AIN-000001.json
    └── events/
        └── AIN-000001/
            ├── AIE-000001.json
            └── AIE-000002.json
\end{verbatim}

The exact internal file split may be refined during implementation.

Mandatory boundaries are:

\begin{itemize}
\tightlist
\item
  original run-owned artifacts stay under the Processing Run;
\item
  review artifacts are project-local;
\item
  Human Review Decisions stay in the existing repository;
\item
  Approved Inputs have one dedicated project-local repository;
\item
  persisted paths are repository-relative;
\item
  symbolic links and path traversal are rejected;
\item
  no operation may silently cross a Project boundary.
\end{itemize}

\hypertarget{20-define-public-application-contracts}{%
\paragraph{20. Define public application contracts}\label{20-define-public-application-contracts}}

The Human Review Workspace shall expose conceptual public operations equivalent to:

\begin{Shaded}
\begin{Highlighting}[]
\NormalTok{ReviewWorkspaceService.open\_review\_document(...)}
\NormalTok{ReviewWorkspaceService.create\_review\_document(...)}
\NormalTok{ReviewWorkspaceService.save\_revision(...)}
\NormalTok{ReviewWorkspaceService.apply\_scoped\_action(...)}
\NormalTok{ReviewWorkspaceService.split\_review\_item(...)}
\NormalTok{ReviewWorkspaceService.merge\_review\_items(...)}
\NormalTok{ReviewWorkspaceService.finalize\_version(...)}
\NormalTok{ReviewWorkspaceService.reopen\_finalized\_version(...)}
\end{Highlighting}
\end{Shaded}

Repository operations shall conceptually include:

\begin{Shaded}
\begin{Highlighting}[]
\NormalTok{ReviewWorkspaceRepository.load\_document(...)}
\NormalTok{ReviewWorkspaceRepository.load\_version(...)}
\NormalTok{ReviewWorkspaceRepository.load\_revision(...)}
\NormalTok{ReviewWorkspaceRepository.list\_versions(...)}
\NormalTok{ReviewWorkspaceRepository.scan(...)}
\end{Highlighting}
\end{Shaded}

Approved Input operations shall conceptually include:

\begin{Shaded}
\begin{Highlighting}[]
\NormalTok{ApprovedInputPromotionService.assess\_eligibility(...)}
\NormalTok{ApprovedInputPromotionService.promote\_finalized\_version(...)}
\NormalTok{ApprovedInputLifecycleService.invalidate(...)}
\NormalTok{ApprovedInputLifecycleService.revoke(...)}
\NormalTok{ApprovedInputLifecycleService.supersede(...)}
\end{Highlighting}
\end{Shaded}

The stable Phase H read contract shall be:

\begin{Shaded}
\begin{Highlighting}[]
\NormalTok{ApprovedInputRepository.list\_active\_approved\_inputs(}
\NormalTok{    project\_id,}
\NormalTok{) }\OperatorTok{{-}\textgreater{}} \BuiltInTok{tuple}\NormalTok{[ApprovedInputManifest, ...]}
\end{Highlighting}
\end{Shaded}

Phase H shall not read mutable draft Review Revisions, original Review Reports as model-generation authority, Agent confidence or Consensus as approval authority, Human Review UI state, or inactive Approved Inputs.

\hypertarget{21-define-finalization-and-promotion-failure-behavior}{%
\paragraph{21. Define finalization and promotion failure behavior}\label{21-define-finalization-and-promotion-failure-behavior}}

Finalization and promotion cannot be one filesystem-atomic transaction across all repositories.

The required sequence is:

\begin{verbatim}
1. validate current Review Revision
2. persist immutable finalized Review Document artifacts
3. record exact Human Review Decision
4. revalidate promotion eligibility
5. create Approved Input manifests
6. append required lifecycle events
7. report complete or recovery-required status
\end{verbatim}

No Approved Input may be created before the exact finalization Human Review Decision exists.

If finalization artifacts exist but the Human Review Decision cannot be recorded, the document is not treated as approved, no Approved Input is created and recovery is required.

If the Human Review Decision exists but promotion is only partially completed, existing promoted items remain traceable, the operation is not reported as fully successful and deterministic recovery resumes idempotently.

Promotion idempotence shall be based on at least:

\begin{verbatim}
project_id
review_document_version_id
review_item_id
review_item_fingerprint
finalization_decision_fingerprint
\end{verbatim}

Equivalent duplicate promotion shall not create a second active Approved Input.

\hypertarget{22-define-the-required-human-review-user-experience}{%
\paragraph{22. Define the required Human Review user experience}\label{22-define-the-required-human-review-user-experience}}

The Review Workspace is a critical engineering interface.

Its default layout shall preserve overview through:

\begin{verbatim}
left:
document outline, item list, filters and status counts

center:
report-like reviewed document with inline item editing

right:
all Agent proposals, comparison, evidence and original content
\end{verbatim}

The interface shall provide:

\begin{verbatim}
Elements
Relationships
Open Questions
Rejected Content
\end{verbatim}

The active Project, Source, Processing Run, Review Document and version shall remain visible.

The UI shall provide:

\begin{verbatim}
Review Copy
Original Report
Diff to Original
Diff to Predecessor
Version History
\end{verbatim}

Required status counts include:

\begin{verbatim}
open
accepted as generated
accepted with modification
combined
rejected
deferred
out of scope
unresolved
\end{verbatim}

The UI shall offer one-click item operations where safe.

One-click draft actions shall not be confused with final document approval. The final approval operation shall be visually distinct and shall always show a complete impact summary.

\hypertarget{23-preserve-review-and-rework-loops}{%
\paragraph{23. Preserve review and rework loops}\label{23-preserve-review-and-rework-loops}}

Human Review may determine that information must return to an earlier processing capability.

The workspace shall support explicit outcomes conceptually equivalent to:

\begin{verbatim}
request source-information rework
request semantic-governance rework
request candidate-regeneration
\end{verbatim}

These outcomes remain traceable to the affected Review Items and do not create Approved Input.

The exact Processing Run successor or retry behavior remains governed by the accepted Processing architecture.

\hypertarget{24-define-implementation-sequence}{%
\paragraph{24. Define implementation sequence}\label{24-define-implementation-sequence}}

Phase G shall proceed in the following steps:

\begin{verbatim}
G1
ADR-016 and Phase G contract acceptance

G2
Human Review Workspace identifiers, immutable types, manifests and repositories

G3
P4/P9 evidence adapters and deterministic Review Item construction

G4
Draft revisions, scoped actions, version finalization and reopening

G5
Approved Input manifests, repository, eligibility, promotion and lifecycle

G6
Human Review Workspace UI with Elements, Relationships, Open Questions and
Rejected Content

G7
Project-bound integration, recovery behavior, manual acceptance audit and full
regression
\end{verbatim}

The first executable implementation step after ADR-016 shall define contracts and tests before UI implementation.

\hypertarget{consequences-1}{%
\subsubsection{Consequences}\label{consequences-1}}

\hypertarget{positive-consequences-2}{%
\paragraph{Positive consequences}\label{positive-consequences-2}}

\begin{itemize}
\tightlist
\item
  The original machine-generated evidence remains permanently available.
\item
  The user can edit a report naturally without weakening traceability.
\item
  One large report can produce many independently reviewed Approved Inputs.
\item
  All Agent proposals remain visible and attributable.
\item
  One-click selection remains possible without making draft clicks immediately authoritative.
\item
  Document-level decisions reduce repetitive review effort.
\item
  Item-level overrides support mixed documents.
\item
  Materialized filter actions provide focused review without dynamic-scope risk.
\item
  Elements and relationships remain visually separated.
\item
  Open questions remain visible instead of disappearing into report prose.
\item
  Rejected content remains auditable.
\item
  Finalized documents are immutable.
\item
  Reopening creates an explicit successor version.
\item
  Unchanged Approved Inputs remain stable across document versions.
\item
  Changed Approved Inputs are superseded rather than overwritten.
\item
  Human withdrawal is distinguishable from technical invalidation.
\item
  Relationship review uses the applicable SysML v2 target notation instead of generic or obsolete relationship labels.
\item
  Phase H receives one stable, project-local Approved Input contract.
\end{itemize}

\hypertarget{negative-consequences-1}{%
\paragraph{Negative consequences}\label{negative-consequences-1}}

\begin{itemize}
\tightlist
\item
  The Phase G persistence model is more complex than a mutable edited Markdown file.
\item
  The UI requires structured Review Items and cannot rely only on Markdown.
\item
  Proposal grouping, split and merge behavior require deterministic identities and careful testing.
\item
  Document defaults and item overrides require transparent precedence rules.
\item
  Version finalization and Approved Input promotion require recovery handling.
\item
  SysML v2 relationship review depends on an available versioned target-notation profile.
\item
  A complete review can produce many immutable manifests and events.
\end{itemize}

\hypertarget{accepted-complexity}{%
\paragraph{Accepted complexity}\label{accepted-complexity}}

The additional complexity is accepted because the Human Review Workspace is an engineering authority boundary.

A simpler mutable-report design would not provide sufficient decision binding, version accountability, item-level approval, Agent-evidence preservation, project isolation, stale-decision prevention, relationship-notation conformance, Approved Input lifecycle or downstream traceability.

\hypertarget{rejected-alternatives-1}{%
\subsubsection{Rejected alternatives}\label{rejected-alternatives-1}}

\hypertarget{edit-the-original-review-report-in-place}{%
\paragraph{Edit the original Review Report in place}\label{edit-the-original-review-report-in-place}}

Rejected because it would destroy the immutable record of machine-generated evidence and make exact decision reconstruction unreliable.

\hypertarget{copy-and-edit-one-unstructured-markdown-document-only}{%
\paragraph{Copy and edit one unstructured Markdown document only}\label{copy-and-edit-one-unstructured-markdown-document-only}}

Rejected because item identities, proposal variants, scoped decisions, fingerprints and downstream Approved Input promotion would be unreliable.

\hypertarget{approve-or-reject-the-complete-report-as-one-target}{%
\paragraph{Approve or reject the complete report as one target}\label{approve-or-reject-the-complete-report-as-one-target}}

Rejected because reports may contain dozens of independently correct, incorrect or incomplete items.

\hypertarget{require-one-manual-click-for-every-classification-in-every-document}{%
\paragraph{Require one manual click for every classification in every document}\label{require-one-manual-click-for-every-classification-in-every-document}}

Rejected because homogeneous documents require efficient document-level decisions.

\hypertarget{apply-document-level-decisions-without-item-overrides}{%
\paragraph{Apply document-level decisions without item overrides}\label{apply-document-level-decisions-without-item-overrides}}

Rejected because mixed documents require explicit exceptions.

\hypertarget{keep-filter-actions-as-dynamic-saved-queries}{%
\paragraph{Keep filter actions as dynamic saved queries}\label{keep-filter-actions-as-dynamic-saved-queries}}

Rejected because later content changes could silently change the set affected by an earlier decision.

\hypertarget{treat-unselected-agent-proposals-as-factually-rejected}{%
\paragraph{Treat unselected Agent proposals as factually rejected}\label{treat-unselected-agent-proposals-as-factually-rejected}}

Rejected because choosing one variant does not necessarily prove all alternatives incorrect.

\hypertarget{let-agent-confidence-or-consensus-authorize-approval}{%
\paragraph{Let Agent confidence or Consensus authorize approval}\label{let-agent-confidence-or-consensus-authorize-approval}}

Rejected because automated evidence cannot replace explicit human approval.

\hypertarget{store-one-mutable-current-review-state}{%
\paragraph{Store one mutable current review state}\label{store-one-mutable-current-review-state}}

Rejected because it would weaken history reconstruction and recovery.

\hypertarget{reopen-and-modify-a-finalized-version-directly}{%
\paragraph{Reopen and modify a finalized version directly}\label{reopen-and-modify-a-finalized-version-directly}}

Rejected because it would invalidate the meaning of the original finalization decision.

\hypertarget{use-generic-relationship-types-in-the-ui}{%
\paragraph{Use generic relationship types in the UI}\label{use-generic-relationship-types-in-the-ui}}

Rejected because relationship review must use the applicable SysML v2 target-notation profile and remain compatible with later validated textual generation.

\hypertarget{generate-model-candidates-in-phase-g}{%
\paragraph{Generate model candidates in Phase G}\label{generate-model-candidates-in-phase-g}}

Rejected because model candidate generation belongs to Phase H and must consume Approved Input through the stable read contract.

\hypertarget{acceptance}{%
\subsubsection{Acceptance}\label{acceptance}}

The accepted decision comprises:

\begin{itemize}
\tightlist
\item
  immutable original reports;
\item
  editable draft Review Copies;
\item
  immutable finalized reviewed versions;
\item
  documented reopening through successor versions;
\item
  inline element editing;
\item
  all Agent proposals in one item-centered review view;
\item
  one-click proposal selection;
\item
  automatic non-selection of competing variants;
\item
  document-level and item-level decisions;
\item
  focused materialized filter actions;
\item
  separate Elements, Relationships, Open Questions and Rejected Content sections;
\item
  dangerous bulk rejection only with preview, confirmation and rationale;
\item
  versioned Approved Input promotion;
\item
  invalidation, revocation and supersession behavior;
\item
  mandatory use of applicable SysML v2 target notation for relationship review;
\item
  the complete Phase G architecture documented in this ADR.
\end{itemize}

\textbackslash newpage

\hypertarget{adr-017--simple-by-default-interaction-and-progressive-disclosure}{%
\subsection{ADR-017 --- Simple-by-Default Interaction and Progressive Disclosure}\label{adr-017--simple-by-default-interaction-and-progressive-disclosure}}

\textbf{Status:} Accepted \textbf{Date:} 2026-08-12 \textbf{Decision scope:} System-wide user interaction architecture for the Turing Generator prototype

\hypertarget{context-2}{%
\subsubsection{Context}\label{context-2}}

The Turing Generator intentionally preserves a rich internal architecture for authority, evidence, traceability, immutable revision history, validation and recovery.

Phase-G manual acceptance demonstrated that exposing this complete internal detail simultaneously in the primary Streamlit workflow creates unnecessary interaction cost.

The direct-prompt benchmark is relevant: a user can upload legacy engineering content to a general LLM and quickly receive plausible SysML-like output.

The Turing Generator therefore creates value only when its additional governance, reproducibility, traceability, Human Review authority and validation are preserved \textbf{without requiring the user to operate the audit architecture as the default interaction model}.

This decision does not reduce backend evidence, validation or authority.

\hypertarget{decision-2}{%
\subsubsection{Decision}\label{decision-2}}

The system shall apply the following interaction principle:

\begin{verbatim}
Simple by default.
Explainable on demand.
Fully traceable underneath.
\end{verbatim}

The interaction architecture separates three presentation levels.

\hypertarget{level-1--primary-task-oriented-workflow}{%
\paragraph{Level 1 --- Primary Task-oriented Workflow}\label{level-1--primary-task-oriented-workflow}}

Default views shall emphasize:

\begin{itemize}
\tightlist
\item
  engineering result
\item
  material uncertainty or unresolved issue
\item
  required human decision
\item
  next action
\end{itemize}

The primary workflow shall not require the user to interpret internal IDs, fingerprints, persistence paths or complete provenance graphs unless they are material to the immediate task.

\hypertarget{level-2--explanation-on-demand}{%
\paragraph{Level 2 --- Explanation on Demand}\label{level-2--explanation-on-demand}}

The user shall be able to inspect relevant explanation including:

\begin{itemize}
\tightlist
\item
  rationale
\item
  alternative proposals
\item
  relevant source evidence
\item
  disagreement / consensus where useful
\item
  confidence where useful
\item
  validation reason
\item
  material impact of a decision
\end{itemize}

Explanation is not authority by itself.

\hypertarget{level-3--audit-and-traceability}{%
\paragraph{Level 3 --- Audit and Traceability}\label{level-3--audit-and-traceability}}

Complete technical evidence remains inspectable, including:

\begin{itemize}
\tightlist
\item
  Project / Source identity
\item
  Processing Run / Attempt identity
\item
  artifact references
\item
  immutable revision lineage
\item
  Human Review Decision binding
\item
  fingerprints
\item
  validation evidence
\item
  Approved Input authority state
\item
  provenance and lifecycle events
\end{itemize}

This information remains persisted according to the owning contracts even when it is not shown by default.

\hypertarget{system-wide-scope}{%
\subsubsection{System-wide Scope}\label{system-wide-scope}}

The principle applies to:

\begin{itemize}
\tightlist
\item
  Source registration and Agentic Ingestion
\item
  semantic processing
\item
  Human Review
\item
  Approved Input promotion
\item
  Architecture / Model Candidate interaction
\item
  Internal Engineering Model generation
\item
  validation
\item
  SysML v2 generation and result presentation
\end{itemize}

\hypertarget{interaction-boundary}{%
\subsubsection{Interaction Boundary}\label{interaction-boundary}}

Task-oriented interaction and audit-oriented inspection are separate concerns.

A user shall not need to navigate audit-oriented information merely to continue a normal engineering workflow.

Human interaction should be exception-driven where possible.

Examples:

\begin{verbatim}
normal:
result + next action

uncertain:
result + material uncertainty + required decision

audit:
explicitly opened evidence / provenance / fingerprint detail
\end{verbatim}

\hypertarget{human-review-ux-consequences}{%
\subsubsection{Human Review UX Consequences}\label{human-review-ux-consequences}}

The current Human Review implementation remains valid as a Phase-G authority implementation.

A later guided-workflow pass should:

\begin{itemize}
\tightlist
\item
  place competing proposals side by side where useful
\item
  make the selected engineering statement and required decision primary
\item
  collapse detailed evidence by default
\item
  expose actionable relationship-validation reasons
\item
  provide clear write-action lifecycle feedback
\item
  keep complete traceability available on demand
\end{itemize}

No Phase-G authority contract is weakened to achieve this.

\hypertarget{agentic-ingestion-ux-consequences}{%
\subsubsection{Agentic Ingestion UX Consequences}\label{agentic-ingestion-ux-consequences}}

The default ingestion result should prefer:

\begin{itemize}
\tightlist
\item
  what was processed
\item
  whether the run succeeded
\item
  material findings / uncertainty
\item
  whether Human Review is required
\item
  the next action
\end{itemize}

Technical Run / Attempt / artifact detail remains available but need not dominate the default view.

\hypertarget{architecture--model-proposal-ux-consequences}{%
\subsubsection{Architecture / Model Proposal UX Consequences}\label{architecture--model-proposal-ux-consequences}}

Phase H and later proposal workflows should default to:

\begin{itemize}
\tightlist
\item
  proposed engineering content
\item
  relationships and structural implications
\item
  material alternatives
\item
  why the proposal matters
\item
  required human decision
\end{itemize}

Complete provenance remains expandable.

\hypertarget{sysml-v2-result-ux-consequences}{%
\subsubsection{SysML v2 Result UX Consequences}\label{sysml-v2-result-ux-consequences}}

The final workflow should present:

\begin{itemize}
\tightlist
\item
  model result
\item
  validation status
\item
  publication status
\item
  concise generation summary
\end{itemize}

Textual notation, full validation detail, fingerprints and traceability remain available on demand.

\hypertarget{technology-decision-for-the-demo}{%
\subsubsection{Technology Decision for the Demo}\label{technology-decision-for-the-demo}}

Streamlit remains the prototype UI technology through the 2026-08-18 product demo.

The following are explicitly not part of the demo critical path:

\begin{itemize}
\tightlist
\item
  React rewrite
\item
  Vue rewrite
\item
  separate FastAPI backend migration solely for UI restructuring
\end{itemize}

The current application may continue to call Python application services directly.

A future frontend/backend separation remains possible without changing the domain/service architecture.

\hypertarget{consequences-2}{%
\subsubsection{Consequences}\label{consequences-2}}

Positive:

\begin{itemize}
\tightlist
\item
  lower interaction burden
\item
  clearer engineering decisions
\item
  stronger contrast to direct ungoverned prompt-based generation
\item
  auditability remains intact
\item
  backend architecture remains reusable
\item
  UX can evolve independently from authority contracts
\end{itemize}

Trade-offs:

\begin{itemize}
\tightlist
\item
  the prototype must maintain both concise default views and deeper inspection
\item
  progressive disclosure requires deliberate information hierarchy
\item
  some existing Streamlit pages remain information-dense until the later UX pass
\end{itemize}

\hypertarget{non-goals}{%
\subsubsection{Non-goals}\label{non-goals}}

This ADR does not:

\begin{itemize}
\tightlist
\item
  remove persisted evidence
\item
  remove fingerprints
\item
  remove immutable history
\item
  weaken Human Review
\item
  make confidence or consensus authoritative
\item
  replace CATIA authority
\item
  redefine Phase-H Model Candidate architecture
\item
  mandate a production frontend framework
\end{itemize}

\hypertarget{presentation-implication}{%
\subsubsection{Presentation Implication}\label{presentation-implication}}

The architecture should be communicated as:

\begin{verbatim}
simple interaction surface
over
explicit governed engineering services
over
complete persisted evidence and traceability
\end{verbatim}

This supports the literature-derived Data / Process / Knowledge framing while keeping Knowledge and governance cross-cutting rather than forcing them into a single UI or deployment component.

\textbackslash newpage

\hypertarget{adr-018--model-candidate-layer-and-structural-comparability}{%
\subsection{ADR-018 --- Model Candidate Layer and Structural Comparability}\label{adr-018--model-candidate-layer-and-structural-comparability}}

\hypertarget{status}{%
\subsubsection{Status}\label{status}}

Accepted

\hypertarget{date}{%
\subsubsection{Date}\label{date}}

2026-08-12

\hypertarget{context-3}{%
\subsubsection{Context}\label{context-3}}

Phase G establishes the authoritative promotion boundary from reviewed project information to immutable Approved Inputs.

Phase H consumes those Approved Inputs and creates a coherent proposal for the target engineering model.

The Phase-H layer must support:

\begin{itemize}
\tightlist
\item
  joint interpretation of all active Approved Inputs of a project,
\item
  candidate model-element derivation,
\item
  candidate relationship derivation,
\item
  explicit relationship semantics and alternatives,
\item
  advisory relationship prioritization,
\item
  structural comparability across related models,
\item
  versioned structural modeling guidance,
\item
  Human Review,
\item
  complete provenance,
\item
  deterministic downstream consumption by Phase I,
\item
  and a comprehensible user-facing model proposal.
\end{itemize}

Phase H shall not generate SysML v2 textual notation and shall not create a new engineering authority source.

The system-wide interaction principle defined by ADR-017 applies:

\begin{quote}
Simple by default. Explainable on demand. Fully traceable underneath.
\end{quote}

A user shall therefore primarily see a coherent model proposition rather than individual technical candidate artifacts.

\begin{center}\rule{0.5\linewidth}{0.5pt}\end{center}

\hypertarget{decision-3}{%
\subsubsection{Decision}\label{decision-3}}

\hypertarget{1-phase-h-authority-boundary}{%
\paragraph{1. Phase-H authority boundary}\label{1-phase-h-authority-boundary}}

Phase H consumes engineering authority exclusively through:

\texttt{ApprovedInputRepository.list\_active\_approved\_inputs(project\_id)}

Draft Review State, Agent Confidence, Consensus data, inactive Approved Inputs, UI state, and original Review Reports shall not act as direct Phase-H authority sources.

Approved Inputs remain unchanged.

Phase H creates derived Model Candidates based on the exact active Approved-Input snapshot.

Model Candidates are not equivalent to Approved Inputs and are not equivalent to the authoritative CATIA engineering model.

Human acceptance of a Model Candidate authorizes its use for Phase-I model assembly only.

\begin{center}\rule{0.5\linewidth}{0.5pt}\end{center}

\hypertarget{2-model-candidate-set}{%
\paragraph{2. Model Candidate Set}\label{2-model-candidate-set}}

Each Phase-H execution creates an immutable Model Candidate Set.

Proposed identifier form:

\texttt{MCS-000001}

A Candidate Set represents a reproducible snapshot of the complete Phase-H model proposal for one project.

It shall bind at least:

\begin{itemize}
\tightlist
\item
  project identity,
\item
  Candidate Set identity,
\item
  predecessor Candidate Set where applicable,
\item
  regeneration reason where applicable,
\item
  exact Approved-Input references and fingerprints,
\item
  Approved-Input snapshot fingerprint,
\item
  framework-template reference,
\item
  Model Structure and Comparability Profile reference,
\item
  derivation-rules reference,
\item
  generation provenance,
\item
  Element Candidate references,
\item
  Relationship Candidate references,
\item
  creation timestamp,
\item
  content fingerprint.
\end{itemize}

There shall be no implicit "latest Candidate Set wins" behavior.

Downstream consumers must explicitly select a Candidate Set.

\begin{center}\rule{0.5\linewidth}{0.5pt}\end{center}

\hypertarget{3-element-candidates}{%
\paragraph{3. Element Candidates}\label{3-element-candidates}}

Candidate model elements are represented as immutable Element Candidates.

Proposed identifier form:

\texttt{MCE-000001}

Element Candidate instance identity is distinct from semantic continuity.

Each Element Candidate shall support at least:

\begin{itemize}
\tightlist
\item
  immutable candidate identity,
\item
  candidate subject key,
\item
  comparison anchor where applicable,
\item
  proposed name and description,
\item
  model area,
\item
  profile-defined element type,
\item
  framework assignment,
\item
  terminology assignment,
\item
  attributes,
\item
  exact Approved-Input provenance,
\item
  derivation rationale,
\item
  support level,
\item
  assumptions,
\item
  missing information,
\item
  structural-profile conformance,
\item
  predecessor candidate references,
\item
  immutable content fingerprint.
\end{itemize}

Where possible, a one-to-one derivation should retain continuity with the Approved Input stable subject key.

Aggregation or decomposition may create new candidate subject identities while preserving all contributing provenance.

\begin{center}\rule{0.5\linewidth}{0.5pt}\end{center}

\hypertarget{4-relationship-candidates}{%
\paragraph{4. Relationship Candidates}\label{4-relationship-candidates}}

Relationships are independent immutable Candidate artifacts.

Proposed identifier form:

\texttt{MCR-000001}

Relationships shall not be embedded implicitly in Element Candidates.

Every Relationship Candidate shall support at least:

\begin{itemize}
\tightlist
\item
  immutable relationship candidate identity,
\item
  Candidate Set identity,
\item
  relationship choice key where alternatives belong to the same decision,
\item
  exact source Element Candidate reference,
\item
  exact target Element Candidate reference,
\item
  source and target semantic subject keys,
\item
  relationship family,
\item
  semantic intent,
\item
  directionality,
\item
  exact Approved-Input provenance,
\item
  derivation rationale,
\item
  supporting evidence,
\item
  assumptions,
\item
  missing information,
\item
  priority assessment,
\item
  structural-comparability assessment,
\item
  structural-profile conformance,
\item
  upstream Approved-Input relationship representation where applicable,
\item
  predecessor candidate references,
\item
  immutable content fingerprint.
\end{itemize}

Every relationship eligible for Phase I must resolve its source and target to exact Element Candidates within the same Candidate Set.

Zero matching endpoints are unresolved.

More than one matching endpoint is ambiguous.

Exactly one matching endpoint is resolved.

\begin{center}\rule{0.5\linewidth}{0.5pt}\end{center}

\hypertarget{5-relationship-semantics}{%
\paragraph{5. Relationship semantics}\label{5-relationship-semantics}}

Relationship family, engineering semantic intent, and eventual SysML v2 serialization are separate concepts.

Relationship families may include, among others:

\begin{itemize}
\tightlist
\item
  dependency,
\item
  allocation,
\item
  flow,
\item
  refinement-related,
\item
  derivation-related,
\item
  framework-specific relationships.
\end{itemize}

The actual controlled vocabulary shall be defined by versioned configuration or profile data and shall not be silently hardcoded into processing logic.

Distinct semantic meanings shall never be silently collapsed into a generic relationship.

Where multiple materially different interpretations remain plausible, separate Relationship Candidates shall be created and grouped through one relationship choice key.

Existing Approved-Input relationship representations may be carried forward as evidence.

Phase H shall not silently override an explicit Approved-Input relationship interpretation.

Any changed interpretation requires explicit rationale and traceability.

Final semantic-to-SysML-v2 serialization belongs to Phase J.

Formal SysML-v2 validation belongs to Phase K.

\begin{center}\rule{0.5\linewidth}{0.5pt}\end{center}

\hypertarget{6-relationship-prioritization}{%
\paragraph{6. Relationship prioritization}\label{6-relationship-prioritization}}

Automated relationship prioritization is advisory only.

No automated priority shall constitute engineering approval.

The default priority classes are:

\begin{itemize}
\tightlist
\item
  preferred,
\item
  supported\_alternative,
\item
  exception\_candidate.
\end{itemize}

Prioritization shall record machine-readable criterion results and a human-readable rationale.

The prioritization criteria are evaluated in the following conceptual order:

\begin{enumerate}
\def\labelenumi{\arabic{enumi}.}
\tightlist
\item
  evidence directness,
\item
  semantic fit,
\item
  endpoint certainty,
\item
  structural-profile preference,
\item
  structural-comparability impact,
\item
  assumption burden,
\item
  conformance.
\end{enumerate}

An explicit Approved-Input relationship normally has stronger evidence directness than a relationship inferred solely from separate Approved Inputs.

Human Review remains the authorization boundary.

\begin{center}\rule{0.5\linewidth}{0.5pt}\end{center}

\hypertarget{7-model-structure-and-comparability-profile}{%
\paragraph{7. Model Structure and Comparability Profile}\label{7-model-structure-and-comparability-profile}}

Phase H introduces a versioned Model Structure and Comparability Profile.

This profile is separate from the SysML v2 Target Notation Profile.

The Model Structure and Comparability Profile defines, at minimum:

\begin{itemize}
\tightlist
\item
  preferred model areas,
\item
  permitted element types,
\item
  comparison anchors,
\item
  canonical structural patterns,
\item
  canonical relationship semantics,
\item
  relationship preferences,
\item
  permitted structural variation,
\item
  prioritization criteria,
\item
  exception rules,
\item
  review requirements.
\end{itemize}

The purpose of the profile is to improve structural comparability between:

\begin{itemize}
\tightlist
\item
  related products,
\item
  product variants,
\item
  independently generated models,
\item
  repeated model-generation executions.
\end{itemize}

The guiding principle is:

Where engineering meaning and modeling context are equivalent, the same structural representation should be preferred.

Different engineering meaning may justify different structures.

Intentional deviations remain possible but require explicit rationale and Human Review.

\begin{center}\rule{0.5\linewidth}{0.5pt}\end{center}

\hypertarget{8-structural-comparability-assessment}{%
\paragraph{8. Structural-comparability assessment}\label{8-structural-comparability-assessment}}

Candidates may record structural-comparability effects including:

\begin{itemize}
\tightlist
\item
  improves,
\item
  neutral,
\item
  reduces,
\item
  unknown.
\end{itemize}

The assessment may include:

\begin{itemize}
\tightlist
\item
  affected comparison anchors,
\item
  canonical-pattern match,
\item
  structural deviations,
\item
  rationale.
\end{itemize}

Comparability findings are advisory evidence for Human Review and shall not replace engineering judgment.

\begin{center}\rule{0.5\linewidth}{0.5pt}\end{center}

\hypertarget{9-human-review}{%
\paragraph{9. Human Review}\label{9-human-review}}

Model Candidate content is immutable.

Human Review decisions are persisted separately and are bound to the exact Candidate content fingerprint.

Supported decision concepts include:

\begin{itemize}
\tightlist
\item
  accepted,
\item
  rejected,
\item
  deferred,
\item
  accepted\_exception.
\end{itemize}

An accepted exception requires explicit rationale.

Human Review shall bind at least:

\begin{itemize}
\tightlist
\item
  project identity,
\item
  Candidate Set identity,
\item
  Candidate identity or explicitly defined Candidate Set selection scope,
\item
  exact content fingerprint,
\item
  reviewer identity,
\item
  review timestamp,
\item
  decision,
\item
  rationale where required,
\item
  structural-profile reference,
\item
  relevant conformance fingerprint,
\item
  Approved-Input snapshot fingerprint,
\item
  immutable decision fingerprint.
\end{itemize}

Existing generic Human Review infrastructure shall be reused where compatible rather than duplicated.

The exact review granularity may support review-by-exception, provided that every authoritative Phase-I selection remains explicitly persisted and bound to exact candidate content.

\begin{center}\rule{0.5\linewidth}{0.5pt}\end{center}

\hypertarget{10-lifecycle-and-regeneration}{%
\paragraph{10. Lifecycle and regeneration}\label{10-lifecycle-and-regeneration}}

Candidate artifacts and review decisions remain immutable.

If any material Approved Input referenced by a Candidate Set later becomes inactive, the historical Candidate Set remains preserved but becomes ineligible for new Phase-I assembly.

A new analysis creates a new Candidate Set.

Regeneration shall:

\begin{itemize}
\tightlist
\item
  create a new Candidate Set identity,
\item
  create new Candidate instance identities,
\item
  preserve predecessor references,
\item
  preserve semantic subject continuity where applicable,
\item
  preserve complete provenance.
\end{itemize}

Historical Candidate Sets are never silently replaced.

\begin{center}\rule{0.5\linewidth}{0.5pt}\end{center}

\hypertarget{11-phase-i-read-contract}{%
\paragraph{11. Phase-I read contract}\label{11-phase-i-read-contract}}

Phase I shall consume Phase-H output only through an explicit validated read contract.

Conceptually:

\texttt{ModelCandidateReadService.load\_phase\_i\_input(project\_id,\ candidate\_set\_id)}

The read boundary shall verify at least:

\begin{itemize}
\tightlist
\item
  Candidate Set integrity,
\item
  project isolation,
\item
  immutable fingerprints,
\item
  current activity of material Approved Inputs,
\item
  exact Human Review decisions,
\item
  absence of stale decisions,
\item
  endpoint resolution,
\item
  acceptance of required Element Candidates,
\item
  acceptance of required Relationship Candidates,
\item
  absence of blocking unresolved relationships,
\item
  structural-profile conformance or an explicitly accepted exception.
\end{itemize}

The resulting Phase-I input shall contain only explicitly authorized Candidate content and relevant provenance.

Phase I may assemble accepted Candidates into the Internal Engineering Model.

Phase I shall not invent new engineering semantics from unaccepted Phase-H Candidates.

\begin{center}\rule{0.5\linewidth}{0.5pt}\end{center}

\hypertarget{12-user-facing-model-proposal}{%
\paragraph{12. User-facing Model Proposal}\label{12-user-facing-model-proposal}}

The primary user-facing Phase-H result is a coherent Model Proposal rather than a collection of technical manifests.

The Model Proposal is a deterministic projection of one explicit Candidate Set.

It is not an independent authority source.

Conceptually, the presentation contract is:

\texttt{ModelProposalView}

It shall support at least:

\begin{itemize}
\tightlist
\item
  Candidate Set identity,
\item
  proposal summary,
\item
  proposed model elements,
\item
  proposed relationships,
\item
  structural overview,
\item
  relationship choice groups,
\item
  structural-comparability summary,
\item
  profile deviations,
\item
  required Human Review decisions,
\item
  summarized rationale,
\item
  next action.
\end{itemize}

All displayed information must resolve back to immutable Candidate artifacts and review evidence.

\begin{center}\rule{0.5\linewidth}{0.5pt}\end{center}

\hypertarget{13-visual-model-projections}{%
\paragraph{13. Visual model projections}\label{13-visual-model-projections}}

The Model Proposal may and should use visual model projections where they improve comprehension.

Examples include:

\begin{itemize}
\tightlist
\item
  structural overview diagrams,
\item
  UML-like component views,
\item
  SysML-like block views,
\item
  requirement-to-function views,
\item
  function-to-logical-component allocation views,
\item
  relationship networks,
\item
  alternative relationship visualizations.
\end{itemize}

These visualizations are explanatory projections of Candidate data.

They are not themselves the machine-readable model authority.

A visual relationship may therefore replace verbose technical relationship descriptions in the primary UI while retaining the complete technical relationship representation underneath.

The default user interaction shall prefer:

\begin{enumerate}
\def\labelenumi{\arabic{enumi}.}
\tightlist
\item
  understandable graphical/model result,
\item
  material uncertainty,
\item
  required human decision,
\item
  next action.
\end{enumerate}

Detailed rationale and alternatives shall be available on demand.

Complete provenance and audit information shall remain accessible underneath.

The Phase-H visualization does not have to constitute formally valid SysML v2 notation.

Formal model assembly belongs to Phase I.

Formal SysML v2 textual generation belongs to Phase J.

\begin{center}\rule{0.5\linewidth}{0.5pt}\end{center}

\hypertarget{14-optional-model-proposal-report}{%
\paragraph{14. Optional Model Proposal report}\label{14-optional-model-proposal-report}}

A human-readable or exportable Model Proposal Summary may be generated from an explicit Candidate Set.

Such a report may be retained as presentation, thesis, or audit evidence.

It shall remain a reproducible projection and shall never replace the Candidate Set as the machine-readable source of truth.

\begin{center}\rule{0.5\linewidth}{0.5pt}\end{center}

\hypertarget{15-phase-boundaries}{%
\paragraph{15. Phase boundaries}\label{15-phase-boundaries}}

Phase H:

\begin{itemize}
\tightlist
\item
  jointly interprets active Approved Inputs,
\item
  proposes model elements,
\item
  proposes relationships,
\item
  represents alternatives,
\item
  prioritizes relationships advisorially,
\item
  assesses structural comparability,
\item
  performs Candidate Human Review,
\item
  exposes the Model Proposal View.
\end{itemize}

Phase I:

\begin{itemize}
\tightlist
\item
  assembles explicitly accepted Candidates into the Internal Engineering Model.
\end{itemize}

Phase J:

\begin{itemize}
\tightlist
\item
  generates SysML v2 textual notation from the Internal Engineering Model.
\end{itemize}

Phase K:

\begin{itemize}
\tightlist
\item
  validates syntax, semantics, structure, target-profile conformance, and applicable comparability constraints.
\end{itemize}

\begin{center}\rule{0.5\linewidth}{0.5pt}\end{center}

\hypertarget{consequences-3}{%
\subsubsection{Consequences}\label{consequences-3}}

\hypertarget{positive-consequences-3}{%
\paragraph{Positive consequences}\label{positive-consequences-3}}

\begin{itemize}
\tightlist
\item
  Project information is synthesized into one coherent model proposal.
\item
  Approved Inputs remain immutable and authoritative within their existing scope.
\item
  Model derivation remains completely traceable.
\item
  Alternative modeling interpretations remain explicit.
\item
  Relationship semantics are not silently flattened.
\item
  Structural comparability becomes a first-class modeling concern.
\item
  Human Review remains the engineering authorization boundary.
\item
  The default UX can remain simple despite detailed underlying traceability.
\item
  Visual diagrams can communicate architecture considerably more efficiently than raw relationship manifests.
\item
  Phase H remains independent from final SysML v2 textual serialization.
\item
  Phase I receives a deterministic and validated input contract.
\end{itemize}

\hypertarget{trade-offs}{%
\paragraph{Trade-offs}\label{trade-offs}}

\begin{itemize}
\tightlist
\item
  Candidate Set lifecycle introduces additional persistent identities.
\item
  Candidate review and regeneration require explicit lifecycle handling.
\item
  Structural comparability requires a maintained versioned profile.
\item
  Visual proposal rendering introduces a presentation layer that must remain strictly derived from machine-readable Candidate data.
\item
  Relationship alternatives increase artifact count but prevent semantic information loss.
\end{itemize}

\begin{center}\rule{0.5\linewidth}{0.5pt}\end{center}

\hypertarget{rejected-alternatives-2}{%
\subsubsection{Rejected alternatives}\label{rejected-alternatives-2}}

\hypertarget{generate-sysml-v2-directly-from-approved-inputs-in-phase-h}{%
\paragraph{Generate SysML v2 directly from Approved Inputs in Phase H}\label{generate-sysml-v2-directly-from-approved-inputs-in-phase-h}}

Rejected because this would collapse modeling synthesis, engineering-model assembly, serialization, and validation into one stage and would bypass the planned Phase-I architecture boundary.

\hypertarget{treat-the-visual-diagram-as-the-authoritative-model}{%
\paragraph{Treat the visual diagram as the authoritative model}\label{treat-the-visual-diagram-as-the-authoritative-model}}

Rejected because graphical layout is a presentation concern and must remain reproducible from structured Candidate data.

\hypertarget{require-users-to-review-every-technical-candidate-manifest-directly}{%
\paragraph{Require users to review every technical Candidate manifest directly}\label{require-users-to-review-every-technical-candidate-manifest-directly}}

Rejected because this conflicts with ADR-017 and would create unnecessary UX complexity.

\hypertarget{collapse-all-relationship-semantics-into-generic-links}{%
\paragraph{Collapse all relationship semantics into generic links}\label{collapse-all-relationship-semantics-into-generic-links}}

Rejected because semantic distinctions are relevant for engineering correctness, structural comparability, and later SysML v2 generation.

\hypertarget{automatically-accept-the-highest-ranked-relationship}{%
\paragraph{Automatically accept the highest-ranked relationship}\label{automatically-accept-the-highest-ranked-relationship}}

Rejected because prioritization is advisory and may not substitute explicit Human Review.

\begin{center}\rule{0.5\linewidth}{0.5pt}\end{center}

\hypertarget{related-decisions}{%
\subsubsection{Related decisions}\label{related-decisions}}

\begin{itemize}
\tightlist
\item
  ADR-016 --- Human Review Workspace and Approved Input Promotion Architecture
\item
  ADR-017 --- Simple-by-Default Interaction and Progressive Disclosure
\end{itemize}

\begin{center}\rule{0.5\linewidth}{0.5pt}\end{center}

\hypertarget{implementation-note}{%
\subsubsection{Implementation note}\label{implementation-note}}

No Phase-H production implementation shall precede acceptance and persistence of this ADR.

Implementation shall reuse existing repository, Human Review, validation, identifier, fingerprint, and project-isolation contracts where applicable rather than duplicating upstream logic.

\textbackslash newpage

\hypertarget{adr-019--internal-engineering-model-assembly-architecture}{%
\subsection{ADR-019 --- Internal Engineering Model Assembly Architecture}\label{adr-019--internal-engineering-model-assembly-architecture}}

\hypertarget{status-1}{%
\subsubsection{Status}\label{status-1}}

Accepted

\hypertarget{date-1}{%
\subsubsection{Date}\label{date-1}}

2026-08-13

\hypertarget{context-4}{%
\subsubsection{Context}\label{context-4}}

Phase H establishes the reviewed Model Candidate Layer and exposes the sole validated Phase-H → Phase-I authority transfer through:

\texttt{ModelCandidateReadService.load\_phase\_i\_input(project\_id,\ candidate\_set\_id)}

The resulting \texttt{ModelCandidateAssemblyInput} contains only explicitly authorized Element Candidates and Relationship Candidates together with Candidate-Set, Approved-Input, Human-Review, exception, profile and generation provenance.

Phase I is responsible for assembling this authorized content into one coherent Internal Engineering Model.

Phase I shall not repeat the semantic interpretation performed in Phase H.

Phase I shall not bypass Human Review.

Phase I shall not generate SysML v2 textual notation.

The Framework Template and Model Structure and Comparability Profile already define the structural areas into which accepted Candidate content has been classified.

The accepted Phase-H architecture requires Phase I to preserve Candidate semantics exactly and to consume only the validated H→I read contract.

The implementation therefore requires a deterministic assembly layer that:

\begin{itemize}
\tightlist
\item
  materializes the selected framework hierarchy,
\item
  maps accepted Element Candidates into that hierarchy,
\item
  maps accepted Relationship Candidates to exact internal model elements,
\item
  preserves Human Review and Approved-Input traceability,
\item
  preserves accepted exceptions,
\item
  validates internal assembly integrity,
\item
  persists an immutable Internal Engineering Model snapshot,
\item
  and provides an explicit read boundary for Phase J.
\end{itemize}

The CATIA engineering model describes architecture derivation/structuring and architecture validation as related but distinct responsibilities. The implementation roadmap separates these temporally:

\begin{itemize}
\tightlist
\item
  Phase I performs deterministic architecture assembly,
\item
  Phase K performs broader model validation,
\item
  Phase J performs SysML v2 textual serialization.
\end{itemize}

\begin{center}\rule{0.5\linewidth}{0.5pt}\end{center}

\hypertarget{decision-4}{%
\subsubsection{Decision}\label{decision-4}}

\hypertarget{i-01--sole-upstream-authority}{%
\paragraph{I-01 --- Sole upstream authority}\label{i-01--sole-upstream-authority}}

Phase I shall consume Phase-H engineering content exclusively through:

\texttt{ModelCandidateReadService.load\_phase\_i\_input(project\_id,\ candidate\_set\_id)}

Phase-I production assembly shall not independently read:

\begin{itemize}
\tightlist
\item
  raw Approved Inputs,
\item
  Candidate Review persistence,
\item
  unreviewed Candidate Sets,
\item
  rejected Candidates,
\item
  deferred Candidates,
\item
  Model Proposal presentation state,
\item
  or other upstream processing evidence
\end{itemize}

for the purpose of deciding which engineering content is assembled.

The H→I read contract remains the authorization boundary.

\begin{center}\rule{0.5\linewidth}{0.5pt}\end{center}

\hypertarget{i-02--handoff-enrichment}{%
\paragraph{I-02 --- Handoff enrichment}\label{i-02--handoff-enrichment}}

\texttt{ModelCandidateAssemblyInput} shall be extended to include the exact:

\begin{itemize}
\tightlist
\item
  \texttt{framework\_template\_reference}
\item
  \texttt{derivation\_rules\_reference}
\end{itemize}

already bound by the source Model Candidate Set.

Phase I shall not independently rediscover these references from unrelated repository state.

This extension does not change Phase-H engineering semantics.

It strengthens the existing H→I authority transfer so Phase I receives all required pinned assembly context through one validated contract.

\begin{center}\rule{0.5\linewidth}{0.5pt}\end{center}

\hypertarget{i-03--deterministic-assembly}{%
\paragraph{I-03 --- Deterministic assembly}\label{i-03--deterministic-assembly}}

The normative Phase-I path shall be deterministic.

Phase I shall not require an LLM to reinterpret reviewed Candidate content.

The roadmap term "Model Generation Agent" is implemented architecturally as a deterministic Internal Model Assembly capability.

Any optional future agent assistance shall remain non-authoritative and shall not bypass this deterministic assembly contract.

\begin{center}\rule{0.5\linewidth}{0.5pt}\end{center}

\hypertarget{i-04--immutable-internal-engineering-model-snapshot}{%
\paragraph{I-04 --- Immutable Internal Engineering Model snapshot}\label{i-04--immutable-internal-engineering-model-snapshot}}

Every materially distinct authorized assembly shall produce an immutable Internal Engineering Model snapshot.

Identifier form:

\texttt{IEM-000001}

The Internal Engineering Model snapshot is a reproducible model-assembly boundary.

Published IEM snapshots shall not be modified in place.

\begin{center}\rule{0.5\linewidth}{0.5pt}\end{center}

\hypertarget{i-05--exact-authorization-binding}{%
\paragraph{I-05 --- Exact authorization binding}\label{i-05--exact-authorization-binding}}

Every IEM snapshot shall bind the exact Phase-I assembly input.

A deterministic \texttt{assembly\_input\_fingerprint} shall cover the authority-bearing H→I input required to reproduce the assembly, including at least:

\begin{itemize}
\tightlist
\item
  project identity,
\item
  Candidate Set identity and fingerprint,
\item
  Approved-Input snapshot fingerprint,
\item
  Framework Template reference,
\item
  Model Structure Profile reference,
\item
  derivation-rules reference,
\item
  accepted Element Candidate identities and fingerprints,
\item
  accepted Relationship Candidate identities and fingerprints,
\item
  Candidate Review Decision references and fingerprints,
\item
  accepted-exception references.
\end{itemize}

A later change in Human Review authorization therefore produces a different assembly input even when the Candidate Set itself remains unchanged.

\begin{center}\rule{0.5\linewidth}{0.5pt}\end{center}

\hypertarget{i-06--idempotent-exact-reassembly}{%
\paragraph{I-06 --- Idempotent exact reassembly}\label{i-06--idempotent-exact-reassembly}}

The combination of:

\begin{itemize}
\tightlist
\item
  identical \texttt{assembly\_input\_fingerprint}, and
\item
  identical assembly-rule reference/version
\end{itemize}

shall be treated idempotently.

The system shall not silently create multiple semantically identical IEM snapshots for the exact same authorized assembly state.

No implicit "latest IEM wins" behavior is permitted.

\begin{center}\rule{0.5\linewidth}{0.5pt}\end{center}

\hypertarget{i-07--separate-internal-model-element-identity}{%
\paragraph{I-07 --- Separate Internal Model Element identity}\label{i-07--separate-internal-model-element-identity}}

Accepted Element Candidates shall be assembled as immutable Internal Model Elements.

Identifier form:

\texttt{IME-000001}

The following identities remain distinct:

\begin{itemize}
\tightlist
\item
  Element Candidate identity (\texttt{MCE-...})
\item
  Internal Model Element identity (\texttt{IME-...})
\item
  semantic subject identity
\end{itemize}

An IME shall preserve explicit traceability to its source MCE.

\begin{center}\rule{0.5\linewidth}{0.5pt}\end{center}

\hypertarget{i-08--separate-internal-model-relationship-identity}{%
\paragraph{I-08 --- Separate Internal Model Relationship identity}\label{i-08--separate-internal-model-relationship-identity}}

Accepted Relationship Candidates shall be assembled as immutable Internal Model Relationships.

Identifier form:

\texttt{IMR-000001}

Every IMR shall:

\begin{itemize}
\tightlist
\item
  preserve its source MCR identity,
\item
  reference exact source and target IMEs in the same IEM snapshot,
\item
  preserve source and target semantic subject identities,
\item
  preserve Approved-Input and Human-Review traceability.
\end{itemize}

\begin{center}\rule{0.5\linewidth}{0.5pt}\end{center}

\hypertarget{i-09--preserve-accepted-relationship-semantics}{%
\paragraph{I-09 --- Preserve accepted relationship semantics}\label{i-09--preserve-accepted-relationship-semantics}}

Phase I shall preserve the accepted Candidate relationship semantics without reinterpretation.

The following remain distinct and shall be retained:

\begin{itemize}
\tightlist
\item
  \texttt{relationship\_family}
\item
  \texttt{semantic\_intent}
\item
  \texttt{directionality}
\end{itemize}

Phase I shall not convert one accepted semantic intent into another.

The semantic-to-SysML-v2 serialization decision belongs to Phase J.

\begin{center}\rule{0.5\linewidth}{0.5pt}\end{center}

\hypertarget{i-10--template-derived-structural-skeleton}{%
\paragraph{I-10 --- Template-derived structural skeleton}\label{i-10--template-derived-structural-skeleton}}

Phase I shall materialize the pinned Framework Template hierarchy as an internal structural skeleton.

For the accepted Turing RFLP framework this includes:

\begin{verbatim}
Stakeholder Level
├── Stakeholders
├── User Needs
├── Stakeholder Requirements
└── Use Cases

System Level
├── Requirements
├── Functional
├── Logical
└── Physical

Subsystem Level
├── Requirements
├── Functional
├── Logical
└── Physical
\end{verbatim}

Internal structure nodes are organizational model structure.

They are not additional inferred engineering elements.

\begin{center}\rule{0.5\linewidth}{0.5pt}\end{center}

\hypertarget{i-11--deterministic-containment}{%
\paragraph{I-11 --- Deterministic containment}\label{i-11--deterministic-containment}}

Every accepted IME shall be assigned to exactly one applicable Framework Template mapping node.

The assignment shall be derived from the reviewed Candidate information already established in Phase H, including:

\begin{itemize}
\tightlist
\item
  \texttt{model\_area}
\item
  \texttt{element\_type}
\item
  \texttt{framework\_assignment}
\end{itemize}

Phase I shall verify this mapping against the pinned:

\begin{itemize}
\tightlist
\item
  Framework Template
\item
  Model Structure and Comparability Profile
\end{itemize}

Phase I shall not semantically reclassify the element.

\begin{center}\rule{0.5\linewidth}{0.5pt}\end{center}

\hypertarget{i-12--no-invented-engineering-hierarchy}{%
\paragraph{I-12 --- No invented engineering hierarchy}\label{i-12--no-invented-engineering-hierarchy}}

Phase I shall not infer additional engineering containment, decomposition or ownership relationships merely from:

\begin{itemize}
\tightlist
\item
  element names,
\item
  similarity,
\item
  model area proximity,
\item
  ordering,
\item
  or structural convenience.
\end{itemize}

Any engineering hierarchy beyond the Framework Template structure requires explicit reviewed Candidate semantics.

The Framework Template hierarchy itself may be materialized because it is pinned configuration, not inferred engineering content.

\begin{center}\rule{0.5\linewidth}{0.5pt}\end{center}

\hypertarget{i-13--empty-structural-nodes-allowed}{%
\paragraph{I-13 --- Empty structural nodes allowed}\label{i-13--empty-structural-nodes-allowed}}

The complete selected Framework Template structure may be materialized even when one or more structural areas contain no accepted engineering elements.

Phase I shall not decide whether empty structural areas become explicit SysML v2 packages or other textual constructs.

That representation decision belongs downstream.

\begin{center}\rule{0.5\linewidth}{0.5pt}\end{center}

\hypertarget{i-14--accepted-exceptions-remain-explicit}{%
\paragraph{I-14 --- Accepted exceptions remain explicit}\label{i-14--accepted-exceptions-remain-explicit}}

\texttt{accepted\_exception} authorization shall remain explicit in the Internal Engineering Model.

The affected IME or IMR shall preserve the exact associated Human Review Decision reference.

Phase I shall not silently normalize an accepted exception into ordinary conformance.

An accepted exception may authorize a reviewed structural deviation.

It shall not disable fundamental assembly-integrity requirements.

\begin{center}\rule{0.5\linewidth}{0.5pt}\end{center}

\hypertarget{i-15--fail-closed-on-non-deterministically-assemblable-content}{%
\paragraph{I-15 --- Fail closed on non-deterministically assemblable content}\label{i-15--fail-closed-on-non-deterministically-assemblable-content}}

If accepted Candidate content cannot be assembled without introducing new engineering semantics or making an unauthorized choice, Phase I shall fail closed.

The system shall produce an explicit assembly finding/diagnostic rather than inventing a solution.

The required resolution shall return to the applicable upstream authority, which may include:

\begin{itemize}
\tightlist
\item
  Candidate interpretation,
\item
  Human Review,
\item
  Structure Profile,
\item
  Framework Template,
\item
  or architecture configuration.
\end{itemize}

\begin{center}\rule{0.5\linewidth}{0.5pt}\end{center}

\hypertarget{i-16--phase-i-assembly-integrity-validation}{%
\paragraph{I-16 --- Phase-I assembly-integrity validation}\label{i-16--phase-i-assembly-integrity-validation}}

Phase I shall validate internal assembly integrity.

This validation includes at least:

\begin{itemize}
\tightlist
\item
  unique IEM / IME / IMR identities,
\item
  project isolation,
\item
  exact Candidate traceability,
\item
  exact Human Review traceability,
\item
  exact Approved-Input traceability,
\item
  every IME assigned to a valid structural node,
\item
  \texttt{element\_type} compatibility with the selected model area/profile,
\item
  every IMR endpoint resolving to an IME in the same IEM,
\item
  no dangling Internal Model Relationships,
\item
  no unauthorized Candidate content,
\item
  preservation of accepted exceptions,
\item
  no duplicate internal identities,
\item
  deterministic snapshot fingerprint integrity.
\end{itemize}

Phase-I assembly validation shall not silently modify engineering content in order to pass.

\begin{center}\rule{0.5\linewidth}{0.5pt}\end{center}

\hypertarget{i-17--phase-i-does-not-replace-phase-k}{%
\paragraph{I-17 --- Phase I does not replace Phase K}\label{i-17--phase-i-does-not-replace-phase-k}}

Phase-I validation is limited to assembly and internal representation integrity.

Broader model validation remains Phase K, including applicable:

\begin{itemize}
\tightlist
\item
  architecture-constraint validation,
\item
  structural-pattern validation,
\item
  larger-context compatibility,
\item
  interface consistency,
\item
  relationship-semantic validation,
\item
  target-model compatibility,
\item
  and publication-blocking validation.
\end{itemize}

Phase I may report that content is not assemblable.

It shall not absorb the complete Phase-K validation responsibility.

\begin{center}\rule{0.5\linewidth}{0.5pt}\end{center}

\hypertarget{i-18--no-sysml-v2-textual-constructs-in-phase-i}{%
\paragraph{I-18 --- No SysML v2 textual constructs in Phase I}\label{i-18--no-sysml-v2-textual-constructs-in-phase-i}}

The Internal Engineering Model remains representation-neutral with respect to concrete SysML v2 textual syntax.

Phase I may retain engineering concepts such as:

\begin{itemize}
\tightlist
\item
  \texttt{system\_requirement}
\item
  \texttt{function}
\item
  \texttt{logical\_component}
\item
  \texttt{allocated\_to}
\item
  \texttt{dependency}
\item
  \texttt{flows\_to}
\end{itemize}

Phase I shall not generate concrete textual constructs such as:

\begin{itemize}
\tightlist
\item
  \texttt{requirement}
\item
  \texttt{part\ def}
\item
  \texttt{action\ def}
\item
  \texttt{dependency\ from\ ...}
\item
  package serialization
\item
  target-specific SysML v2 syntax
\end{itemize}

SysML v2 textual generation belongs to Phase J.

\begin{center}\rule{0.5\linewidth}{0.5pt}\end{center}

\hypertarget{i-19--explicit-phase-j-read-boundary}{%
\paragraph{I-19 --- Explicit Phase-J read boundary}\label{i-19--explicit-phase-j-read-boundary}}

Phase J shall consume an explicitly selected Internal Engineering Model snapshot through a validated read contract.

Conceptually:

\begin{Shaded}
\begin{Highlighting}[]
\NormalTok{InternalModelReadService.load\_phase\_j\_input(}
\NormalTok{    project\_id,}
\NormalTok{    internal\_engineering\_model\_id,}
\NormalTok{) }\OperatorTok{{-}\textgreater{}}\NormalTok{ InternalEngineeringModelSnapshot}
\end{Highlighting}
\end{Shaded}

There shall be no implicit latest-IEM selection.

The Phase-J read contract shall verify snapshot integrity and project isolation before exposing the Internal Engineering Model downstream.

\begin{center}\rule{0.5\linewidth}{0.5pt}\end{center}

\hypertarget{i-20--catia-and-logical-architecture-alignment}{%
\paragraph{I-20 --- CATIA and logical-architecture alignment}\label{i-20--catia-and-logical-architecture-alignment}}

The Internal Engineering Model corresponds conceptually to the architecture result produced by the CATIA system behavior "Derive and Structure Architecture Information".

The implementation separation is:

\begin{verbatim}
LC_07 Architecture Synthesis and Validation
├── Phase I — synthesis / deterministic assembly
└── Phase K — broader architecture/model validation

LC_08 SysML v2 Artifact Generation
└── Phase J — SysML v2 textual generation
\end{verbatim}

Implementation phases therefore do not map one-to-one to Logical Components.

This separation is intentional and shall be documented rather than treated as a conflict between implementation and CATIA architecture.

\begin{center}\rule{0.5\linewidth}{0.5pt}\end{center}

\hypertarget{internal-engineering-model-artifacts}{%
\subsubsection{Internal Engineering Model artifacts}\label{internal-engineering-model-artifacts}}

\hypertarget{internal-engineering-model-manifest}{%
\paragraph{Internal Engineering Model Manifest}\label{internal-engineering-model-manifest}}

The IEM manifest shall bind at least:

\begin{itemize}
\tightlist
\item
  schema version
\item
  project identity
\item
  IEM identity
\item
  assembly-input fingerprint
\item
  Candidate Set identity
\item
  Candidate Set content fingerprint
\item
  Approved-Input snapshot fingerprint
\item
  Framework Template reference
\item
  Model Structure Profile reference
\item
  derivation-rules reference
\item
  assembly-rules reference
\item
  assembly provenance
\item
  Internal Model Structure reference
\item
  IME references
\item
  IMR references
\item
  Candidate Review Decision references
\item
  accepted-exception references
\item
  creation timestamp
\item
  immutable content fingerprint
\end{itemize}

\begin{center}\rule{0.5\linewidth}{0.5pt}\end{center}

\hypertarget{internal-model-element}{%
\paragraph{Internal Model Element}\label{internal-model-element}}

Each IME shall support at least:

\begin{itemize}
\tightlist
\item
  schema version
\item
  project identity
\item
  IEM identity
\item
  IME identity
\item
  semantic subject identity
\item
  source MCE identity
\item
  source MCE fingerprint
\item
  name
\item
  description
\item
  model area
\item
  element type
\item
  Framework assignment
\item
  terminology assignment
\item
  attributes
\item
  comparison anchor where applicable
\item
  Approved-Input references
\item
  Human Review Decision reference
\item
  accepted-exception reference where applicable
\item
  immutable content fingerprint
\end{itemize}

\begin{center}\rule{0.5\linewidth}{0.5pt}\end{center}

\hypertarget{internal-model-relationship}{%
\paragraph{Internal Model Relationship}\label{internal-model-relationship}}

Each IMR shall support at least:

\begin{itemize}
\tightlist
\item
  schema version
\item
  project identity
\item
  IEM identity
\item
  IMR identity
\item
  source IME identity
\item
  target IME identity
\item
  source semantic subject identity
\item
  target semantic subject identity
\item
  relationship family
\item
  semantic intent
\item
  directionality
\item
  source MCR identity
\item
  source MCR fingerprint
\item
  Approved-Input references
\item
  Human Review Decision reference
\item
  accepted-exception reference where applicable
\item
  immutable content fingerprint
\end{itemize}

\begin{center}\rule{0.5\linewidth}{0.5pt}\end{center}

\hypertarget{internal-model-structure}{%
\paragraph{Internal Model Structure}\label{internal-model-structure}}

The Internal Model Structure shall contain:

\begin{itemize}
\tightlist
\item
  Framework Template identity
\item
  deterministic materialized structure nodes
\item
  parent/child structure-node relationships
\item
  stable node ordering
\item
  exact IME membership per structural node
\item
  immutable structure fingerprint
\end{itemize}

The structure artifact shall not duplicate the complete engineering content of the IME artifacts.

It represents model organization and containment.

\begin{center}\rule{0.5\linewidth}{0.5pt}\end{center}

\hypertarget{persistence-1}{%
\subsubsection{Persistence}\label{persistence-1}}

The default project-local persistence layout is:

\begin{verbatim}
data/projects/<project_id>/internal_models/
└── IEM-000001/
    ├── manifest.json
    ├── structure.json
    ├── elements/
    │   ├── IME-000001.json
    │   └── ...
    └── relationships/
        ├── IMR-000001.json
        └── ...
\end{verbatim}

The complete IEM bundle shall be persisted atomically and immutably.

Partial publication shall not constitute a valid IEM snapshot.

Repository scanning shall fail closed on:

\begin{itemize}
\tightlist
\item
  corrupt manifests,
\item
  missing referenced artifacts,
\item
  unexpected artifact identities,
\item
  cross-project references,
\item
  fingerprint mismatch,
\item
  symlink/path-safety violations,
\item
  or interrupted temporary publication state.
\end{itemize}

\begin{center}\rule{0.5\linewidth}{0.5pt}\end{center}

\hypertarget{assembly-flow}{%
\subsubsection{Assembly flow}\label{assembly-flow}}

The normative Phase-I flow is:

\begin{verbatim}
ModelCandidateReadService
        │
        ▼
ModelCandidateAssemblyInput
        │
        ├── pinned Framework Template
        ├── pinned Model Structure Profile
        └── pinned assembly / derivation context
        │
        ▼
InternalModelAssemblyService
        │
        ├── materialize structural skeleton
        ├── MCE → IME
        ├── MCR → IMR
        ├── preserve exceptions
        ├── preserve traceability
        └── validate assembly integrity
        │
        ▼
Immutable InternalEngineeringModelSnapshot
        │
        ▼
InternalModelReadService
        │
        ▼
Phase J
\end{verbatim}

\begin{center}\rule{0.5\linewidth}{0.5pt}\end{center}

\hypertarget{human-review-boundary}{%
\subsubsection{Human Review boundary}\label{human-review-boundary}}

Phase I introduces no new normative Human Review layer for semantic reinterpretation.

Human Review of engineering meaning occurs upstream.

If deterministic assembly cannot continue without a new engineering decision, Phase I fails closed and surfaces the condition for upstream resolution.

Phase I shall not use an additional AI or Human Review step to silently modify already accepted Candidate semantics.

\begin{center}\rule{0.5\linewidth}{0.5pt}\end{center}

\hypertarget{phase-boundaries}{%
\subsubsection{Phase boundaries}\label{phase-boundaries}}

\hypertarget{phase-h}{%
\paragraph{Phase H}\label{phase-h}}

\begin{itemize}
\tightlist
\item
  interprets approved engineering information,
\item
  proposes model elements and relationships,
\item
  identifies alternatives,
\item
  assesses comparability,
\item
  performs Human Review,
\item
  authorizes Candidate content for assembly.
\end{itemize}

\hypertarget{phase-i}{%
\paragraph{Phase I}\label{phase-i}}

\begin{itemize}
\tightlist
\item
  consumes only authorized Candidate content,
\item
  materializes the framework structure,
\item
  assembles Internal Model Elements,
\item
  assembles Internal Model Relationships,
\item
  preserves exact semantics and traceability,
\item
  performs assembly-integrity validation,
\item
  persists the immutable Internal Engineering Model.
\end{itemize}

\hypertarget{phase-j}{%
\paragraph{Phase J}\label{phase-j}}

\begin{itemize}
\tightlist
\item
  maps the Internal Engineering Model to supported SysML v2 constructs,
\item
  applies the target-notation profile,
\item
  generates SysML v2 textual notation.
\end{itemize}

\hypertarget{phase-k}{%
\paragraph{Phase K}\label{phase-k}}

\begin{itemize}
\tightlist
\item
  performs broader syntax, semantic, structural, integration, constraint, traceability and comparability validation as applicable.
\end{itemize}

\begin{center}\rule{0.5\linewidth}{0.5pt}\end{center}

\hypertarget{consequences-4}{%
\subsubsection{Consequences}\label{consequences-4}}

\hypertarget{positive-consequences-4}{%
\paragraph{Positive consequences}\label{positive-consequences-4}}

\begin{itemize}
\tightlist
\item
  Phase-H Human Review remains authoritative.
\item
  Model assembly is deterministic and reproducible.
\item
  The Internal Engineering Model is independent of SysML v2 textual syntax.
\item
  Framework hierarchy and engineering semantics remain separate concepts.
\item
  Accepted exceptions remain visible.
\item
  Every internal element and relationship retains exact upstream traceability.
\item
  Phase J receives one coherent model rather than independent Candidate artifacts.
\item
  Phase K remains a distinct validation responsibility.
\item
  Repeated exact assembly is idempotent.
\item
  Historical IEM snapshots remain reproducible.
\end{itemize}

\hypertarget{trade-offs-1}{%
\paragraph{Trade-offs}\label{trade-offs-1}}

\begin{itemize}
\tightlist
\item
  Phase I introduces IEM / IME / IMR identities.
\item
  A separate internal structure artifact is required.
\item
  H→I transfer must be enriched with pinned Framework Template and derivation-rules references.
\item
  Repository and fingerprint validation add implementation work.
\item
  Some accepted Candidate combinations may fail closed and require upstream resolution rather than automatic correction.
\end{itemize}

\begin{center}\rule{0.5\linewidth}{0.5pt}\end{center}

\hypertarget{rejected-alternatives-3}{%
\subsubsection{Rejected alternatives}\label{rejected-alternatives-3}}

\hypertarget{reinterpret-accepted-candidates-in-phase-i}{%
\paragraph{Reinterpret accepted Candidates in Phase I}\label{reinterpret-accepted-candidates-in-phase-i}}

Rejected because this would bypass Phase-H Human Review and duplicate semantic model derivation.

\hypertarget{generate-sysml-v2-directly-in-phase-i}{%
\paragraph{Generate SysML v2 directly in Phase I}\label{generate-sysml-v2-directly-in-phase-i}}

Rejected because model assembly and target-language serialization are separate responsibilities and Phase J owns SysML v2 textual generation.

\hypertarget{use-candidate-ids-directly-as-final-internal-model-identities}{%
\paragraph{Use Candidate IDs directly as final internal model identities}\label{use-candidate-ids-directly-as-final-internal-model-identities}}

Rejected because Candidate identity, semantic continuity and assembled snapshot identity have different lifecycle semantics.

\hypertarget{treat-the-framework-template-as-engineering-content}{%
\paragraph{Treat the Framework Template as engineering content}\label{treat-the-framework-template-as-engineering-content}}

Rejected because template structure is organizational configuration and does not by itself create new engineering elements or engineering relationships.

\hypertarget{automatically-invent-missing-hierarchy}{%
\paragraph{Automatically invent missing hierarchy}\label{automatically-invent-missing-hierarchy}}

Rejected because additional engineering containment or decomposition requires reviewed semantic evidence.

\hypertarget{let-phase-i-perform-all-downstream-model-validation}{%
\paragraph{Let Phase I perform all downstream model validation}\label{let-phase-i-perform-all-downstream-model-validation}}

Rejected because this would collapse assembly and validation and would conflict with the planned Phase-K boundary.

\hypertarget{select-the-latest-internal-engineering-model-implicitly}{%
\paragraph{Select the latest Internal Engineering Model implicitly}\label{select-the-latest-internal-engineering-model-implicitly}}

Rejected because explicit snapshot identity is required for reproducibility and traceability.

\begin{center}\rule{0.5\linewidth}{0.5pt}\end{center}

\hypertarget{related-decisions-1}{%
\subsubsection{Related decisions}\label{related-decisions-1}}

\begin{itemize}
\tightlist
\item
  ADR-016 --- Human Review Workspace and Approved Input Promotion Architecture
\item
  ADR-017 --- Simple-by-Default Interaction and Progressive Disclosure
\item
  ADR-018 --- Model Candidate Layer and Structural Comparability
\end{itemize}

\begin{center}\rule{0.5\linewidth}{0.5pt}\end{center}

\hypertarget{initial-implementation-decomposition}{%
\subsubsection{Initial implementation decomposition}\label{initial-implementation-decomposition}}

After this ADR is committed, Phase I may proceed in the following internal implementation slices:

\begin{verbatim}
I1  identifiers + immutable Internal Model domain types
I2  manifests + fingerprints + H→I contract enrichment
I3  Framework/Profile resolution + structure materialization
I4  deterministic MCE/MCR → IME/IMR assembly
I5  repository + immutable persistence + assembly-integrity validation
I6  Phase-J read contract + regression + SSOT closeout
\end{verbatim}

This decomposition is an implementation planning aid.

The official roadmap phase remains Phase I --- Model Generation Agent / Internal Engineering Model.

\begin{center}\rule{0.5\linewidth}{0.5pt}\end{center}

\hypertarget{implementation-note-1}{%
\subsubsection{Implementation note}\label{implementation-note-1}}

No Phase-I production implementation shall precede acceptance and persistence of this ADR.

Implementation shall reuse the existing:

\begin{itemize}
\tightlist
\item
  project-isolation rules,
\item
  safe-path behavior,
\item
  atomic-persistence patterns,
\item
  identifier conventions,
\item
  deterministic manifest/fingerprint patterns,
\item
  Framework Template loader,
\item
  Model Structure Profile loader,
\item
  and H→I authorization contract
\end{itemize}

where applicable rather than duplicating upstream logic.

\textbackslash newpage

\hypertarget{adr-020--hybrid-target-projection-and-coverage-architecture}{%
\subsection{ADR-020 --- Hybrid Target Projection and Coverage Architecture}\label{adr-020--hybrid-target-projection-and-coverage-architecture}}

\hypertarget{status-2}{%
\subsubsection{Status}\label{status-2}}

Accepted

\hypertarget{date-2}{%
\subsubsection{Date}\label{date-2}}

2026-08-13

\hypertarget{context-5}{%
\subsubsection{Context}\label{context-5}}

Phase H derives target-model Candidates from reviewed Approved Inputs.

The existing \texttt{ProfileDrivenModelCandidateDeriver} is intentionally conservative and deterministic. It maps reviewed information only when the selected Model Structure Profile supports the mapping and fails closed when no supported mapping exists or a mapping is ambiguous.

This behavior is valuable as a fast, reproducible projection path, but it is not sufficient as the only target-projection strategy for heterogeneous engineering information. A selected modeling framework must not force valid engineering information into an unsuitable target-model shape, and information that cannot be projected must not disappear silently.

The Phase-H orchestration already depends on the \texttt{ModelCandidateDeriver} protocol rather than one concrete derivation implementation. This allows the existing deterministic strategy to remain available while a second, LLM-assisted strategy is introduced without changing the Phase-H → Phase-I authority boundary.

Phase I remains responsible only for assembling human-authorized Model Candidates into the Internal Engineering Model. Phase I shall not be reopened for target-projection reasoning.

The executable prototype schedule requires a bounded extension. The existing deterministic path shall therefore be preserved and reused as the first stage of the hybrid path.

CATIA requirement/function reconciliation is intentionally deferred to Phase N2, where capabilities and architecture decisions introduced during Phases G--L are reconciled as one coherent model update.

\begin{center}\rule{0.5\linewidth}{0.5pt}\end{center}

\hypertarget{decision-5}{%
\subsubsection{Decision}\label{decision-5}}

\hypertarget{h9-01--preserve-strict-deterministic-projection}{%
\paragraph{H9-01 --- Preserve strict deterministic projection}\label{h9-01--preserve-strict-deterministic-projection}}

\texttt{ProfileDrivenModelCandidateDeriver} remains a supported target-projection strategy.

It shall continue to:

\begin{itemize}
\tightlist
\item
  use only the pinned Framework Template, Model Structure Profile and derivation rules,
\item
  derive only profile-supported Candidate semantics,
\item
  fail closed on ambiguous or unsupported target mappings,
\item
  and require no LLM execution.
\end{itemize}

This path is the fast / reproducible projection option and may be used for quick checks, dry runs and inputs that are already sufficiently classified.

\begin{center}\rule{0.5\linewidth}{0.5pt}\end{center}

\hypertarget{h9-02--explicit-projection-dispositions}{%
\paragraph{H9-02 --- Explicit projection dispositions}\label{h9-02--explicit-projection-dispositions}}

Every active Approved Input considered for target projection shall receive exactly one projection disposition:

\begin{itemize}
\tightlist
\item
  \texttt{mapped}
\item
  \texttt{ambiguous}
\item
  \texttt{unmapped}
\item
  \texttt{intentionally\_not\_projected}
\end{itemize}

\texttt{human\_clarification} Approved Inputs remain reviewed context and are \texttt{intentionally\_not\_projected} unless a later accepted architecture explicitly changes that rule.

\begin{center}\rule{0.5\linewidth}{0.5pt}\end{center}

\hypertarget{h9-03--complete-projection-coverage}{%
\paragraph{H9-03 --- Complete projection coverage}\label{h9-03--complete-projection-coverage}}

Projection coverage shall account for the complete active Approved-Input snapshot.

The invariant is:

\begin{verbatim}
mapped
+ ambiguous
+ unmapped
+ intentionally_not_projected
= total Approved Inputs considered
\end{verbatim}

No Approved Input may disappear silently because it does not fit the selected Framework or Model Structure Profile.

\texttt{unmapped} describes a limitation of the selected target projection, not a rejection of the underlying approved engineering information.

\begin{center}\rule{0.5\linewidth}{0.5pt}\end{center}

\hypertarget{h9-04--shared-deterministic-profile-resolver}{%
\paragraph{H9-04 --- Shared deterministic profile resolver}\label{h9-04--shared-deterministic-profile-resolver}}

The profile-matching logic shall be factored into one reusable deterministic resolver.

Both strict deterministic projection and later LLM-assisted projection shall use this resolver.

The system shall not maintain separate deterministic mapping semantics for the two modes.

\begin{center}\rule{0.5\linewidth}{0.5pt}\end{center}

\hypertarget{h9-05--deterministic-first-hybrid-projection}{%
\paragraph{H9-05 --- Deterministic-first hybrid projection}\label{h9-05--deterministic-first-hybrid-projection}}

The LLM-assisted strategy shall execute deterministic profile resolution first.

Inputs resolved uniquely by the deterministic resolver shall not be sent to the LLM in the normal hybrid path.

Only inputs classified as \texttt{ambiguous} or \texttt{unmapped} are eligible for LLM-assisted target-projection reasoning.

This limits token use, request rate and unnecessary semantic reinterpretation.

\begin{center}\rule{0.5\linewidth}{0.5pt}\end{center}

\hypertarget{h9-06--llm-output-remains-candidate-level-modeling}{%
\paragraph{H9-06 --- LLM output remains Candidate-level modeling}\label{h9-06--llm-output-remains-candidate-level-modeling}}

The LLM-assisted mapper shall not generate normative SysML v2 text.

It may propose structured target-model semantics such as target model area, element type, framework assignment, relationship semantic intent, rationale, alternatives, or an explicit \texttt{unmapped} result.

Generated proposals remain Phase-H Candidate information.

\begin{center}\rule{0.5\linewidth}{0.5pt}\end{center}

\hypertarget{h9-07--no-forced-mapping}{%
\paragraph{H9-07 --- No forced mapping}\label{h9-07--no-forced-mapping}}

The LLM-assisted path is not required to map every Approved Input.

If the selected Framework/Profile has no defensible representation, the result shall remain explicitly \texttt{unmapped}.

The LLM shall not invent an engineering meaning merely to achieve target coverage.

\begin{center}\rule{0.5\linewidth}{0.5pt}\end{center}

\hypertarget{h9-08--existing-human-review-authority-remains-unchanged}{%
\paragraph{H9-08 --- Existing Human Review authority remains unchanged}\label{h9-08--existing-human-review-authority-remains-unchanged}}

LLM-assisted projections do not become approved engineering information merely because a model produced them.

They shall pass the same Candidate validation and Human Review authority boundary as deterministic Model Candidates.

The sole Phase-H → Phase-I boundary remains:

\texttt{ModelCandidateReadService.load\_phase\_i\_input(project\_id,\ candidate\_set\_id)}

No Phase-I contract change is required by H9.

\begin{center}\rule{0.5\linewidth}{0.5pt}\end{center}

\hypertarget{h9-09--projection-provenance}{%
\paragraph{H9-09 --- Projection provenance}\label{h9-09--projection-provenance}}

The derivation path shall remain traceable through \texttt{ModelCandidateGenerationProvenance}.

For LLM-assisted projection, provenance shall identify at least the applicable derivation method, model reference, recipe and/or agent reference where applicable, and context fingerprint.

Per-input LLM projection evidence shall remain attributable to the affected Approved Input.

\begin{center}\rule{0.5\linewidth}{0.5pt}\end{center}

\hypertarget{h9-10--token-and-request-rate-protection}{%
\paragraph{H9-10 --- Token and request-rate protection}\label{h9-10--token-and-request-rate-protection}}

The normal LLM-assisted path shall minimize model traffic by design.

It shall:

\begin{itemize}
\tightlist
\item
  perform deterministic resolution first,
\item
  send only unresolved inputs,
\item
  avoid resending complete source documents when bounded Approved-Input content is sufficient,
\item
  send only relevant target-profile context,
\item
  support bounded batching where appropriate,
\item
  avoid reinterpreting uniquely mapped deterministic results,
\item
  and avoid unbounded automatic retry loops.
\end{itemize}

Correctness and explicit uncertainty remain more important than forcing a mapping.

\begin{center}\rule{0.5\linewidth}{0.5pt}\end{center}

\hypertarget{h9-11--phase-j-remains-deterministic-serialization}{%
\paragraph{H9-11 --- Phase J remains deterministic serialization}\label{h9-11--phase-j-remains-deterministic-serialization}}

H9 changes target-model projection in Phase H.

It does not move semantic interpretation into Phase J.

Phase J continues to receive a validated Internal Engineering Model and serialize already accepted model semantics to the selected SysML v2 target notation.

\begin{center}\rule{0.5\linewidth}{0.5pt}\end{center}

\hypertarget{h9-12--catia-reconciliation-is-deferred-to-phase-n2}{%
\paragraph{H9-12 --- CATIA reconciliation is deferred to Phase N2}\label{h9-12--catia-reconciliation-is-deferred-to-phase-n2}}

H9 may introduce capabilities that require requirement/function reconciliation, including selectable deterministic and LLM-assisted target projection, projection coverage assessment, explicit ambiguous/unmapped preservation, deterministic-first LLM routing, and LLM projection provenance.

These are reconciliation candidates only.

No CATIA Requirement or Function shall be added or modified as part of the H9 implementation closeout. Final coverage assessment and model changes belong to Phase N2.

\begin{center}\rule{0.5\linewidth}{0.5pt}\end{center}

\hypertarget{initial-implementation-decomposition-1}{%
\subsubsection{Initial implementation decomposition}\label{initial-implementation-decomposition-1}}

\begin{verbatim}
H9.1  Projection disposition and coverage model
H9.2  Shared deterministic profile resolver
H9.3  Strict deterministic deriver migrated to shared resolver
H9.4  Structured LLM projection contract
H9.5  Token-conscious HybridModelCandidateDeriver
H9.6  Merge, validation and provenance
H9.7  Human Review / Phase-I compatibility regression
H9.8  SSOT closeout and Phase-N2 reconciliation candidate recording
\end{verbatim}

H9.1--H9.3 require no LLM request.

\begin{center}\rule{0.5\linewidth}{0.5pt}\end{center}

\hypertarget{consequences-5}{%
\subsubsection{Consequences}\label{consequences-5}}

\hypertarget{positive}{%
\paragraph{Positive}\label{positive}}

\begin{itemize}
\tightlist
\item
  deterministic projection remains available and fast,
\item
  framework limitations become visible instead of silently losing information,
\item
  LLM usage is concentrated on cases that actually require semantic reasoning,
\item
  Phase I remains stable,
\item
  Human Review authority remains stable,
\item
  the architecture supports future alternative Framework/Profile combinations,
\item
  and Phase J can remain reproducible.
\end{itemize}

\hypertarget{trade-offs-2}{%
\paragraph{Trade-offs}\label{trade-offs-2}}

\begin{itemize}
\tightlist
\item
  Phase H gains an additional projection-coverage concept,
\item
  hybrid execution will require explicit LLM response validation,
\item
  unresolved information can legitimately remain unmapped,
\item
  and final CATIA capability reconciliation is deferred until Phase N2.
\end{itemize}

\begin{center}\rule{0.5\linewidth}{0.5pt}\end{center}

\hypertarget{acceptance-1}{%
\subsubsection{Acceptance}\label{acceptance-1}}

\textbackslash newpage

\hypertarget{adr-021--syside-compatible-sysml-v2-generation-architecture}{%
\subsection{ADR-021 --- SYSIDE-Compatible SysML v2 Generation Architecture}\label{adr-021--syside-compatible-sysml-v2-generation-architecture}}

\hypertarget{status-3}{%
\subsubsection{Status}\label{status-3}}

Accepted

\hypertarget{date-3}{%
\subsubsection{Date}\label{date-3}}

2026-08-13

\hypertarget{context-6}{%
\subsubsection{Context}\label{context-6}}

Phase H and the controlled H9 extension establish reviewed Model Candidates from Approved Inputs.

Phase I assembles the exact human-authorized Candidate selection into one immutable, representation-neutral Internal Engineering Model (IEM).

The current verified implementation baseline is:

\texttt{1cee49350dd2f24d6a1c80fb0aa1c0d2b5fd27fc}

Phase J now has one responsibility:

\begin{quote}
deterministically serialize one explicitly selected, validated Internal Engineering Model snapshot into SysML v2 textual artifacts that are compatible with SYSIDE.
\end{quote}

Phase J shall not repeat engineering interpretation.

Phase J shall not use an LLM to decide model meaning.

Phase J shall not use the CATIA model of the Turing Generator as a syntax template for generated target models.

The Turing Generator's CATIA model describes the Turing Generator itself. It remains the engineering authority for the product architecture and is relevant to later Architecture-to-Requirements reconciliation, especially Phase N2. It does not define the textual syntax that Phase J shall generate for arbitrary target models.

The SysML v2 source/reference hierarchy for Phase J is instead:

\begin{verbatim}
SysML v2 specification / release repository
        │
        │ primary language and syntax reference
        ▼
SYSIDE
        │
        │ target execution / validation environment
        ▼
validated local syntax experiments
        │
        │ implementation evidence for the supported subset
        ▼
Apollo 11 SysML v2 example project
        │
        │ non-normative structure / style reference only
        ▼
Turing Generator target-notation and generation profiles
\end{verbatim}

The repository already records the local reference locations:

\begin{verbatim}
external/sysml-v2-release/
external/apollo11-sysml-v2/
\end{verbatim}

\texttt{context/sources/source\_manifest.json} defines the SysML v2 Release repository as the primary specification reference and Apollo 11 as a non-normative example reference.

\texttt{context/sysml/sysml\_v2\_spec\_reference.json} is a curated working reference. It is intentionally not treated as the complete grammar and explicitly requires consulting the local SysML v2 release repository and validating uncertain syntax in the selected environment.

\texttt{context/examples/apollo11\_structure\_reference.md} records useful Apollo organization and modeling patterns, but explicitly prevents Apollo from overriding the SysML v2 specification or the Turing target notation.

The current target-notation artifact is:

\texttt{context/sysml/sysml\_v2\_target\_notation.json}

Version:

\texttt{0.1.0}

It already constrains generation to a deliberately limited SysML v2 subset, but its generation-state language predates the completed Phase-H/Phase-I architecture. It still refers to \texttt{approved\_model\_data} as direct generation input and associates generation directly with \texttt{data/output/}.

The accepted architecture is now:

\begin{verbatim}
Approved Input
→ reviewed Model Candidates
→ Internal Engineering Model
→ Phase J generation
→ Phase K validation
→ Phase L publication
\end{verbatim}

Therefore Phase J requires a clean serialization architecture that preserves the IEM authority boundary, separates syntax policy from semantic mapping and artifact organization, fails closed when a mapping is unsupported, and produces an exact validation-ready artifact set for Phase K.

\begin{center}\rule{0.5\linewidth}{0.5pt}\end{center}

\hypertarget{decision-6}{%
\subsubsection{Decision}\label{decision-6}}

\hypertarget{j-01--sole-upstream-authority}{%
\paragraph{J-01 --- Sole upstream authority}\label{j-01--sole-upstream-authority}}

Phase J shall consume engineering-model content exclusively through:

\begin{Shaded}
\begin{Highlighting}[]
\NormalTok{InternalModelReadService.load\_phase\_j\_input(}
\NormalTok{    project\_id,}
\NormalTok{    internal\_engineering\_model\_id,}
\NormalTok{) }\OperatorTok{{-}\textgreater{}}\NormalTok{ InternalEngineeringModelSnapshot}
\end{Highlighting}
\end{Shaded}

The IEM shall be explicitly addressed.

No implicit "latest IEM" selection is allowed.

Phase-J generation shall not independently read Model Candidates, Approved Inputs, Review persistence, raw ingestion artifacts or presentation state in order to decide what model content is generated.

The Phase-I read boundary remains authoritative for the exact model snapshot supplied to Phase J.

\begin{center}\rule{0.5\linewidth}{0.5pt}\end{center}

\hypertarget{j-02--deterministic-serialization-only}{%
\paragraph{J-02 --- Deterministic serialization only}\label{j-02--deterministic-serialization-only}}

The normative Phase-J generation path shall be deterministic.

Phase J shall not use an LLM or agent to:

\begin{itemize}
\tightlist
\item
  reinterpret engineering meaning,
\item
  choose between semantic alternatives,
\item
  invent relationships,
\item
  infer additional hierarchy,
\item
  choose a target-model area,
\item
  repair unsupported semantics,
\item
  or generate free-form SysML v2 text.
\end{itemize}

Engineering interpretation occurs in Phase H/H9 and is authorized through Human Review.

Phase I assembles that authorization without semantic reinterpretation.

Phase J is therefore a compiler/serializer from accepted internal semantics to a controlled textual representation.

For identical:

\begin{itemize}
\tightlist
\item
  IEM content,
\item
  target-notation reference,
\item
  generation-profile reference,
\item
  artifact-structure reference,
\item
  and generator implementation/rule version,
\end{itemize}

the generated textual content shall be byte-identical.

\begin{center}\rule{0.5\linewidth}{0.5pt}\end{center}

\hypertarget{j-03--phase-j-syntax-and-reference-authority-hierarchy}{%
\paragraph{J-03 --- Phase-J syntax and reference authority hierarchy}\label{j-03--phase-j-syntax-and-reference-authority-hierarchy}}

Phase J shall apply the following hierarchy when defining or extending generated SysML v2 syntax:

\begin{enumerate}
\def\labelenumi{\arabic{enumi}.}
\tightlist
\item
  local SysML v2 specification/release repository
\item
  successful validation in SYSIDE
\item
  intentionally maintained local syntax fixtures/experiments
\item
  Apollo 11 as a non-normative example reference
\item
  project-specific style preferences
\end{enumerate}

The local SysML v2 release repository remains the primary language and syntax reference.

SYSIDE is the target environment against which the generated subset shall be validated.

Apollo 11 may inform:

\begin{itemize}
\tightlist
\item
  package organization,
\item
  separation of definitions and usages,
\item
  naming style,
\item
  large-model structuring,
\item
  and examples of valid modeling patterns.
\end{itemize}

Apollo 11 shall not override the specification reference or a SYSIDE validation result.

\begin{center}\rule{0.5\linewidth}{0.5pt}\end{center}

\hypertarget{j-04--catia-remains-separate-from-the-phase-j-syntax-authority}{%
\paragraph{J-04 --- CATIA remains separate from the Phase-J syntax authority}\label{j-04--catia-remains-separate-from-the-phase-j-syntax-authority}}

The CATIA model of the Turing Generator is the engineering model of the Turing Generator product itself.

It may define or constrain what the Turing Generator shall be capable of doing.

It shall not be used as the normative textual syntax source for Phase J.

In particular:

\begin{verbatim}
CATIA Turing Generator model
→ product requirements / functions / logical architecture
→ Phase-N2 reconciliation authority

SysML v2 specification + SYSIDE
→ generated language / syntax authority for Phase J
\end{verbatim}

A future independent CATIA example model may be used as an interoperability test case.

Such a model shall remain non-normative for the Phase-J syntax mapping unless a later accepted architecture decision explicitly changes that rule.

\begin{center}\rule{0.5\linewidth}{0.5pt}\end{center}

\hypertarget{j-05--three-separate-versioned-generation-policy-artifacts}{%
\paragraph{J-05 --- Three separate versioned generation-policy artifacts}\label{j-05--three-separate-versioned-generation-policy-artifacts}}

Phase J shall separate three concerns that are currently partially mixed.

\hypertarget{target-notation-profile}{%
\subparagraph{Target Notation Profile}\label{target-notation-profile}}

Path:

\texttt{context/sysml/sysml\_v2\_target\_notation.json}

Responsibility:

\begin{quote}
Which SysML v2 textual constructs are allowed to be generated?
\end{quote}

The target notation defines the permitted language subset and syntax patterns.

It shall not decide which IEM semantic maps to which construct.

It shall not decide file/package organization.

The current \texttt{0.1.0} artifact shall be intentionally revised during J1 before it is used as the production Phase-J target-notation contract.

The revision shall update stale pre-H/I generation-state language and preserve the principle:

\begin{quote}
Reduce MVP scope by limiting the generated SysML v2 subset, not by accepting invalid syntax.
\end{quote}

\hypertarget{generation-profile}{%
\subparagraph{Generation Profile}\label{generation-profile}}

Proposed path:

\texttt{context/sysml/turing\_sysml\_v2\_generation\_profile.json}

Responsibility:

\begin{quote}
How is each supported IEM semantic represented using allowed target-notation constructs?
\end{quote}

It shall define explicit mappings for:

\begin{itemize}
\tightlist
\item
  IEM element types,
\item
  model areas where required to disambiguate representation,
\item
  relationship family,
\item
  semantic intent,
\item
  directionality,
\item
  permitted attribute projections,
\item
  documentation projection,
\item
  and accepted-exception representation.
\end{itemize}

\hypertarget{artifact-structure-profile}{%
\subparagraph{Artifact Structure Profile}\label{artifact-structure-profile}}

Proposed path:

\texttt{context/sysml/turing\_sysml\_v2\_artifact\_structure.json}

Responsibility:

\begin{quote}
How is one generated model organized into output units, packages, names and deterministic ordering?
\end{quote}

It shall define at least:

\begin{itemize}
\tightlist
\item
  root-package policy,
\item
  Framework-node → package projection,
\item
  default output-unit strategy,
\item
  package naming,
\item
  deterministic ordering,
\item
  relative output path rules,
\item
  and future multi-file extensibility.
\end{itemize}

These three references shall be independently versioned and fingerprinted.

\begin{center}\rule{0.5\linewidth}{0.5pt}\end{center}

\hypertarget{j-06--exact-generation-context-is-pinned}{%
\paragraph{J-06 --- Exact generation context is pinned}\label{j-06--exact-generation-context-is-pinned}}

Every Phase-J generation attempt shall bind the exact:

\begin{itemize}
\tightlist
\item
  source IEM identity and content fingerprint,
\item
  Target Notation Profile identity/version/fingerprint,
\item
  Generation Profile identity/version/fingerprint,
\item
  Artifact Structure Profile identity/version/fingerprint,
\item
  generator-rule/implementation reference where applicable.
\end{itemize}

Generation shall not silently load a newer profile after an artifact set has been created.

The exact generation context shall be sufficient to reproduce the generated content.

\begin{center}\rule{0.5\linewidth}{0.5pt}\end{center}

\hypertarget{j-07--preflight-mapping-completeness-before-rendering}{%
\paragraph{J-07 --- Preflight mapping completeness before rendering}\label{j-07--preflight-mapping-completeness-before-rendering}}

Phase J shall perform deterministic generation preflight before producing a successful artifact set.

The preflight shall verify at least:

\begin{itemize}
\tightlist
\item
  IEM snapshot integrity has already passed the I→J boundary,
\item
  every materialized Framework node has a deterministic package projection,
\item
  every used IEM \texttt{element\_type} has one applicable explicit mapping rule,
\item
  every used relationship semantic has one applicable explicit mapping rule,
\item
  mapping rules reference only allowed target-notation constructs,
\item
  relationship directionality is supported by the selected mapping,
\item
  every IMR endpoint resolves to a generated IME symbol,
\item
  generated symbols are unique,
\item
  names and documentation can be safely serialized,
\item
  accepted exceptions remain representable and traceable,
\item
  and the requested artifact structure can be rendered without an implicit semantic choice.
\end{itemize}

Any missing or ambiguous required mapping is blocking.

Phase J shall not partially succeed by silently omitting unsupported IEM content.

\begin{center}\rule{0.5\linewidth}{0.5pt}\end{center}

\hypertarget{j-08--unsupported-semantics-fail-closed}{%
\paragraph{J-08 --- Unsupported semantics fail closed}\label{j-08--unsupported-semantics-fail-closed}}

An unsupported IEM semantic shall produce an explicit blocking generation finding.

Examples include:

\begin{itemize}
\tightlist
\item
  unsupported element type,
\item
  unsupported relationship semantic intent,
\item
  unsupported directionality,
\item
  target construct not allowed by the pinned Target Notation Profile,
\item
  or a semantic mapping that would require a new engineering decision.
\end{itemize}

Phase J shall not apply a generic fallback such as:

\begin{verbatim}
unknown relationship
→ dependency
\end{verbatim}

merely to produce syntactically valid output.

Phase J shall not replace a required formal relationship with a documentation note when doing so would change or discard the accepted engineering semantics.

Resolution belongs to the relevant upstream profile, modeling decision or target notation extension.

\begin{center}\rule{0.5\linewidth}{0.5pt}\end{center}

\hypertarget{j-09--target-notation-extension-requires-validated-syntax-evidence}{%
\paragraph{J-09 --- Target-notation extension requires validated syntax evidence}\label{j-09--target-notation-extension-requires-validated-syntax-evidence}}

A new SysML v2 construct shall not become generator output merely because a plausible syntax pattern is known.

Before a previously unsupported construct is activated in the Target Notation Profile, J1 shall:

\begin{enumerate}
\def\labelenumi{\arabic{enumi}.}
\tightlist
\item
  identify the required semantic and construct,
\item
  inspect the local SysML v2 release/specification material,
\item
  create the smallest useful local syntax fixture,
\item
  validate the fixture in SYSIDE,
\item
  record the accepted syntax pattern,
\item
  add the construct to the Target Notation Profile,
\item
  add its mapping to the Generation Profile,
\item
  add automated regression tests.
\end{enumerate}

This applies especially to currently unresolved Phase-J needs such as:

\begin{itemize}
\tightlist
\item
  Use Case representation,
\item
  allocation,
\item
  dependency variants,
\item
  derivation,
\item
  refinement,
\item
  satisfaction,
\item
  interaction,
\item
  flow,
\item
  and traceability relationships.
\end{itemize}

No construct shall be guessed from its English name.

\begin{center}\rule{0.5\linewidth}{0.5pt}\end{center}

\hypertarget{j-10--framework-structure-projects-deterministically-to-package-structure}{%
\paragraph{J-10 --- Framework structure projects deterministically to package structure}\label{j-10--framework-structure-projects-deterministically-to-package-structure}}

The materialized \texttt{InternalModelStructure} is organizational structure already authorized by Phase I.

Phase J may project that exact structure into SysML v2 packages without inventing engineering hierarchy.

The default structure shall retain:

\begin{verbatim}
Stakeholder Level
├── Stakeholders
├── User Needs
├── Stakeholder Requirements
└── Use Cases

System Level
├── Requirements
├── Functional
├── Logical
└── Physical

Subsystem Level
├── Requirements
├── Functional
├── Logical
└── Physical
\end{verbatim}

Package names are a textual representation of the pinned Framework structure.

They are not additional engineering elements.

Phase J shall not infer additional containment from:

\begin{itemize}
\tightlist
\item
  names,
\item
  IME ordering,
\item
  element similarity,
\item
  relationship proximity,
\item
  or convenience.
\end{itemize}

\begin{center}\rule{0.5\linewidth}{0.5pt}\end{center}

\hypertarget{j-11--empty-configured-structure-may-remain-explicit}{%
\paragraph{J-11 --- Empty configured structure may remain explicit}\label{j-11--empty-configured-structure-may-remain-explicit}}

The default artifact-structure profile may render the complete configured Framework skeleton, including empty nodes, as packages.

This supports structural comparability across generated model snapshots.

An empty package shall not be interpreted as proof that the corresponding engineering scope is complete.

If a later Artifact Structure Profile chooses to omit empty packages, that behavior shall be explicit, versioned and deterministic.

\begin{center}\rule{0.5\linewidth}{0.5pt}\end{center}

\hypertarget{j-12--default-mvp-output-uses-one-sysml-v2-unit}{%
\paragraph{J-12 --- Default MVP output uses one SysML v2 unit}\label{j-12--default-mvp-output-uses-one-sysml-v2-unit}}

The domain contract shall permit one or more generated units.

The default MVP Artifact Structure Profile shall initially generate one \texttt{.sysml} text unit containing the complete package hierarchy.

Reason:

\begin{itemize}
\tightlist
\item
  minimizes fragile cross-file references,
\item
  provides one simple validation target,
\item
  keeps the first closed vertical slice small,
\item
  and preserves future extensibility through the artifact-set abstraction.
\end{itemize}

The architecture therefore supports:

\begin{verbatim}
GeneratedSysMLArtifactSet
├── Unit 1
├── Unit 2
└── ...
\end{verbatim}

while the initial profile selects:

\begin{verbatim}
GeneratedSysMLArtifactSet
└── generated_model.sysml
\end{verbatim}

A later multi-file profile shall not require redesigning the IEM or generation domain model.

\begin{center}\rule{0.5\linewidth}{0.5pt}\end{center}

\hypertarget{j-13--element-mapping-is-explicit-and-profile-controlled}{%
\paragraph{J-13 --- Element mapping is explicit and profile-controlled}\label{j-13--element-mapping-is-explicit-and-profile-controlled}}

Every generated engineering element shall originate from exactly one IEM element and one applicable Generation Profile rule.

The profile may distinguish identical \texttt{element\_type} values by \texttt{model\_area}.

This is required because, for example, \texttt{function} is used for both:

\begin{itemize}
\tightlist
\item
  \texttt{system.functional}
\item
  \texttt{subsystem.functional}
\end{itemize}

The generator shall not infer representation from an element name.

Likely MVP families include:

\begin{verbatim}
requirements         → requirement constructs
functions            → action constructs
logical components   → part constructs
physical components  → part constructs
\end{verbatim}

However, exact mappings become normative only after their Target Notation constructs and SYSIDE syntax patterns are validated and recorded.

Stakeholders, User Needs and Use Cases shall likewise receive explicit validated rules before production generation.

\begin{center}\rule{0.5\linewidth}{0.5pt}\end{center}

\hypertarget{j-14--stable-generated-symbols-are-separate-from-engineering-names}{%
\paragraph{J-14 --- Stable generated symbols are separate from engineering names}\label{j-14--stable-generated-symbols-are-separate-from-engineering-names}}

Each generated IEM element shall receive a deterministic machine-safe generated symbol derived from immutable internal identity, not from a potentially mutable or colliding display name.

Conceptually:

\begin{verbatim}
IME-000042
→ IME_000042
\end{verbatim}

The exact textual symbol pattern shall be specified in the Artifact Structure or Generation Profile and validated against SYSIDE identifier rules.

The engineering \texttt{name} remains a separate human-readable property.

Phase J shall not make technical identity depend solely on:

\begin{itemize}
\tightlist
\item
  display-name uniqueness,
\item
  capitalization,
\item
  whitespace,
\item
  punctuation,
\item
  or automatic name normalization.
\end{itemize}

This preserves stable relationship endpoint references even when engineering names are similar.

\begin{center}\rule{0.5\linewidth}{0.5pt}\end{center}

\hypertarget{j-15--untrusted-text-is-escaped-and-never-inserted-raw}{%
\paragraph{J-15 --- Untrusted text is escaped and never inserted raw}\label{j-15--untrusted-text-is-escaped-and-never-inserted-raw}}

IEM names, descriptions and generic attributes may originate from heterogeneous legacy information.

Phase J shall treat them as data, not as trusted SysML v2 syntax.

The renderer shall have explicit deterministic escaping/sanitization rules for:

\begin{itemize}
\tightlist
\item
  generated identifiers,
\item
  quoted names where used,
\item
  documentation blocks,
\item
  line breaks,
\item
  comment delimiters,
\item
  and other syntax-sensitive characters.
\end{itemize}

No source-provided text may inject additional SysML v2 statements through raw string concatenation.

If text cannot be represented under the selected rule set without loss or ambiguity, generation shall fail closed with a finding.

\begin{center}\rule{0.5\linewidth}{0.5pt}\end{center}

\hypertarget{j-16--generic-iem-attributes-are-not-automatically-formal-sysml-attributes}{%
\paragraph{J-16 --- Generic IEM attributes are not automatically formal SysML attributes}\label{j-16--generic-iem-attributes-are-not-automatically-formal-sysml-attributes}}

\texttt{InternalModelElement.attributes} contains generic name/value data.

Phase J shall not blindly transform every generic pair into formal SysML v2 attribute syntax.

The Generation Profile shall distinguish:

\begin{verbatim}
explicitly supported formal attribute
→ formal target construct

generic metadata / trace information
→ documentation and/or machine-readable traceability

unsupported semantic attribute
→ blocking finding when omission would lose engineering meaning
\end{verbatim}

This prevents arbitrary legacy metadata from silently becoming formal model semantics.

\begin{center}\rule{0.5\linewidth}{0.5pt}\end{center}

\hypertarget{j-17--relationship-serialization-preserves-accepted-semantics-exactly}{%
\paragraph{J-17 --- Relationship serialization preserves accepted semantics exactly}\label{j-17--relationship-serialization-preserves-accepted-semantics-exactly}}

Every generated relationship shall originate from exactly one IEM relationship.

The renderer shall preserve the accepted distinction between:

\begin{itemize}
\tightlist
\item
  \texttt{relationship\_family}
\item
  \texttt{semantic\_intent}
\item
  \texttt{directionality}
\end{itemize}

The Generation Profile shall map an exact supported tuple to an exact Target-Notation construct.

Conceptually:

\begin{verbatim}
(
    relationship_family,
    semantic_intent,
    directionality,
)
→ one explicit rendering rule
\end{verbatim}

A different SysML construct may not be substituted merely because it looks similar.

Phase J shall not alter relationship direction to simplify rendering.

\begin{center}\rule{0.5\linewidth}{0.5pt}\end{center}

\hypertarget{j-18--accepted-exceptions-remain-explicit}{%
\paragraph{J-18 --- Accepted exceptions remain explicit}\label{j-18--accepted-exceptions-remain-explicit}}

An accepted exception in the IEM remains an accepted reviewed deviation.

Phase J shall not normalize it away.

The generated artifact set shall preserve exact traceability to the associated Human Review Decision.

Where the selected target notation supports a faithful formal representation, the content may be rendered normally while the exception remains explicit in traceability.

If the accepted exception cannot be represented without semantic loss, Phase J shall fail closed.

\begin{center}\rule{0.5\linewidth}{0.5pt}\end{center}

\hypertarget{j-19--deterministic-ordering-and-formatting}{%
\paragraph{J-19 --- Deterministic ordering and formatting}\label{j-19--deterministic-ordering-and-formatting}}

Generated content shall have one canonical deterministic ordering.

The default ordering shall derive only from stable configured or internal information such as:

\begin{enumerate}
\def\labelenumi{\arabic{enumi}.}
\tightlist
\item
  Artifact Structure Profile order
\item
  Framework node order
\item
  IEM element identity
\item
  IEM relationship identity
\end{enumerate}

The renderer shall not depend on:

\begin{itemize}
\tightlist
\item
  filesystem enumeration order,
\item
  dictionary insertion from uncontrolled input,
\item
  current timestamp,
\item
  random values,
\item
  LLM response order,
\item
  or locale-specific sorting.
\end{itemize}

Formatting shall be canonical so that exact regeneration can be compared byte-for-byte.

\begin{center}\rule{0.5\linewidth}{0.5pt}\end{center}

\hypertarget{j-20--machine-readable-traceability-accompanies-the-sysml-v2-text}{%
\paragraph{J-20 --- Machine-readable traceability accompanies the SysML v2 text}\label{j-20--machine-readable-traceability-accompanies-the-sysml-v2-text}}

Phase J shall produce structured traceability together with textual SysML v2.

Each generated element/relationship trace entry shall support the chain:

\begin{verbatim}
generated unit / symbol / location
→ IME or IMR
→ source MCE or MCR
→ Approved Input reference(s)
→ Candidate Human Review Decision
→ accepted exception where applicable
\end{verbatim}

Traceability shall not rely only on human-readable \texttt{doc} text.

The generated SysML v2 text may contain concise traceability documentation when allowed by the profile, but the machine-readable traceability artifact remains the exact Phase-J evidence boundary.

Where practical, trace entries may include deterministic line ranges after rendering.

\begin{center}\rule{0.5\linewidth}{0.5pt}\end{center}

\hypertarget{j-21--generated-artifact-set-is-a-distinct-immutable-domain-contract}{%
\paragraph{J-21 --- Generated artifact set is a distinct immutable domain contract}\label{j-21--generated-artifact-set-is-a-distinct-immutable-domain-contract}}

A successful Phase-J generation returns one immutable \texttt{GeneratedSysMLArtifactSet}.

Conceptually it shall contain at least:

\begin{verbatim}
schema_version
project_id
source_internal_engineering_model_id
source_iem_content_fingerprint

target_notation_reference
generation_profile_reference
artifact_structure_reference

generation_input_fingerprint
generation_provenance

units[]
traceability_entries[]
nonblocking_diagnostics[]

content_fingerprint
\end{verbatim}

Each generated unit shall contain at least:

\begin{verbatim}
unit_id
relative_path
content
content_fingerprint
generated_symbol_ids
source_ime_ids
source_imr_ids
\end{verbatim}

The artifact-set fingerprint shall bind the complete successful generated state.

Timestamps, temporary paths and other non-semantic runtime metadata shall not affect deterministic generation identity.

\begin{center}\rule{0.5\linewidth}{0.5pt}\end{center}

\hypertarget{j-22--exact-generation-identity-and-idempotence}{%
\paragraph{J-22 --- Exact generation identity and idempotence}\label{j-22--exact-generation-identity-and-idempotence}}

Phase J shall calculate a deterministic \texttt{generation\_input\_fingerprint} covering at least:

\begin{itemize}
\tightlist
\item
  exact source IEM identity and fingerprint,
\item
  exact Target Notation Profile reference/fingerprint,
\item
  exact Generation Profile reference/fingerprint,
\item
  exact Artifact Structure Profile reference/fingerprint,
\item
  exact generator rules/implementation reference required for reproducibility.
\end{itemize}

The same exact generation identity shall produce byte-identical artifact content and artifact fingerprints.

Phase J shall not silently select newer profiles under the same prior generation identity.

\begin{center}\rule{0.5\linewidth}{0.5pt}\end{center}

\hypertarget{j-23--generation-findings-are-explicit-and-typed}{%
\paragraph{J-23 --- Generation findings are explicit and typed}\label{j-23--generation-findings-are-explicit-and-typed}}

Phase J shall expose structured generation findings rather than embedding all errors in exception strings.

A finding shall support at least:

\begin{verbatim}
code
message
issue_level
target_type
target_id
profile_rule_id
blocking
\end{verbatim}

Expected finding families include:

\begin{verbatim}
UNSUPPORTED_ELEMENT_MAPPING
UNSUPPORTED_RELATIONSHIP_MAPPING
UNSUPPORTED_DIRECTIONALITY
TARGET_CONSTRUCT_NOT_ALLOWED
UNRENDERABLE_IDENTIFIER
UNSAFE_DOCUMENTATION_CONTENT
DUPLICATE_GENERATED_SYMBOL
UNRESOLVED_GENERATED_ENDPOINT
STRUCTURE_PROJECTION_ERROR
PROFILE_REFERENCE_MISMATCH
\end{verbatim}

A blocking finding prevents a successful \texttt{GeneratedSysMLArtifactSet}.

\begin{center}\rule{0.5\linewidth}{0.5pt}\end{center}

\hypertarget{j-24--phase-j-validates-generation-integrity-not-full-sysml-correctness}{%
\paragraph{J-24 --- Phase J validates generation integrity, not full SysML correctness}\label{j-24--phase-j-validates-generation-integrity-not-full-sysml-correctness}}

Phase J owns deterministic generation-integrity checks, including:

\begin{itemize}
\tightlist
\item
  complete mapping coverage,
\item
  allowed-construct use,
\item
  symbol uniqueness,
\item
  endpoint resolvability,
\item
  deterministic package membership,
\item
  canonical formatting,
\item
  and internal traceability completeness.
\end{itemize}

Phase J does not replace Phase K.

In particular, Phase J does not claim full validation merely because the renderer emitted text matching an expected template.

\begin{center}\rule{0.5\linewidth}{0.5pt}\end{center}

\hypertarget{j-25--explicit-phase-j--phase-k-boundary}{%
\paragraph{J-25 --- Explicit Phase-J → Phase-K boundary}\label{j-25--explicit-phase-j--phase-k-boundary}}

The successful output of Phase J is the immutable \texttt{GeneratedSysMLArtifactSet}.

Conceptually Phase K consumes it through:

\begin{Shaded}
\begin{Highlighting}[]
\NormalTok{SysMLValidationService.validate(}
\NormalTok{    artifact\_set: GeneratedSysMLArtifactSet,}
\NormalTok{) }\OperatorTok{{-}\textgreater{}}\NormalTok{ SysMLValidationResult}
\end{Highlighting}
\end{Shaded}

The exact Phase-K service contract may be refined in the Phase-K architecture decision, but Phase J shall expose all data required for:

\begin{itemize}
\tightlist
\item
  SYSIDE syntax/parser validation,
\item
  reference resolution,
\item
  Target Notation Profile conformance,
\item
  structural validation,
\item
  semantic validation where applicable,
\item
  traceability validation,
\item
  and publication-blocking findings.
\end{itemize}

Phase K shall not silently rewrite generated text to make validation pass.

\begin{center}\rule{0.5\linewidth}{0.5pt}\end{center}

\hypertarget{j-26--syside-validation-belongs-to-phase-k-while-j1-provides-syntax-evidence}{%
\paragraph{J-26 --- SYSIDE validation belongs to Phase K, while J1 provides syntax evidence}\label{j-26--syside-validation-belongs-to-phase-k-while-j1-provides-syntax-evidence}}

There is an intentional distinction between:

\hypertarget{j1-syntax-experiments}{%
\subparagraph{J1 syntax experiments}\label{j1-syntax-experiments}}

Small controlled fixtures used to establish supported generator patterns before a mapping is activated.

and:

\hypertarget{phase-k-product-validation}{%
\subparagraph{Phase K product validation}\label{phase-k-product-validation}}

Validation of an actual complete generated artifact set.

Therefore:

\begin{itemize}
\tightlist
\item
  J1 may use SYSIDE to validate candidate syntax patterns.
\item
  Phase J production does not claim an artifact is valid merely because its individual templates were previously validated.
\item
  Phase K validates the actual generated result.
\end{itemize}

\begin{center}\rule{0.5\linewidth}{0.5pt}\end{center}

\hypertarget{j-27--phase-l-alone-performs-final-publication}{%
\paragraph{J-27 --- Phase L alone performs final publication}\label{j-27--phase-l-alone-performs-final-publication}}

Phase J shall not publish final generated output directly to \texttt{data/output/}.

Phase J creates the validation-ready artifact-set contract.

Phase K validates it.

Phase L owns final versioned publication of accepted generated output.

The existing Source Manifest rule that generated SysML v2 output belongs in:

\texttt{data/output/}

remains valid for the final published product artifact.

It does not require Phase J to bypass validation and publish directly.

The closed path is:

\begin{verbatim}
IEM
→ Phase J GeneratedSysMLArtifactSet
→ Phase K SysMLValidationResult
→ Phase L versioned data/output/ publication
\end{verbatim}

\begin{center}\rule{0.5\linewidth}{0.5pt}\end{center}

\hypertarget{j-28--generated-output-is-not-the-turing-generator-architecture-model}{%
\paragraph{J-28 --- Generated output is not the Turing Generator architecture model}\label{j-28--generated-output-is-not-the-turing-generator-architecture-model}}

Generated customer/target-model output and the engineering model of the Turing Generator are separate artifact classes.

Phase J shall never overwrite or treat as target output:

\begin{itemize}
\tightlist
\item
  the CATIA model of the Turing Generator,
\item
  its textual export,
\item
  or the temporary SYSIDE shadow model of the Turing Generator.
\end{itemize}

Generated model output belongs to the project/output publication boundary defined downstream.

\begin{center}\rule{0.5\linewidth}{0.5pt}\end{center}

\hypertarget{j-29--no-hidden-semantic-fallback-through-documentation}{%
\paragraph{J-29 --- No hidden semantic fallback through documentation}\label{j-29--no-hidden-semantic-fallback-through-documentation}}

Documentation blocks are allowed for:

\begin{itemize}
\tightlist
\item
  descriptions,
\item
  rationale,
\item
  traceability,
\item
  assumptions,
\item
  review references,
\item
  and metadata explicitly designated as documentation.
\end{itemize}

Documentation shall not be used as a hidden fallback to claim that an unsupported formal semantic has been preserved.

Example:

\begin{verbatim}
accepted IMR = satisfies
\end{verbatim}

If the active Generation Profile requires a formal \texttt{satisfies} representation and no validated supported target construct exists, writing:

\begin{verbatim}
doc /* A satisfies B */
\end{verbatim}

does not count as successful semantic serialization.

The generation is blocked until the mapping is explicitly supported or the upstream modeling decision changes.

\begin{center}\rule{0.5\linewidth}{0.5pt}\end{center}

\hypertarget{j-30--quality-is-reduced-by-subset-not-by-semantic-loss}{%
\paragraph{J-30 --- Quality is reduced by subset, not by semantic loss}\label{j-30--quality-is-reduced-by-subset-not-by-semantic-loss}}

The Phase-J MVP may support only a controlled subset of SysML v2.

It may therefore block generation of an IEM that requires unsupported constructs.

It shall not claim success by:

\begin{itemize}
\tightlist
\item
  dropping accepted elements,
\item
  dropping accepted relationships,
\item
  weakening relationship types,
\item
  flattening engineering semantics,
\item
  inventing equivalent-looking constructs,
\item
  or embedding unsupported model meaning only in prose.
\end{itemize}

A smaller valid and faithful supported subset is preferred over broader apparently successful but semantically lossy output.

\begin{center}\rule{0.5\linewidth}{0.5pt}\end{center}

\hypertarget{initial-generation-domain-contracts}{%
\subsubsection{Initial generation-domain contracts}\label{initial-generation-domain-contracts}}

The following contracts are proposed for Phase J.

\hypertarget{targetnotationreference}{%
\paragraph{TargetNotationReference}\label{targetnotationreference}}

Pinned identity of the active Target Notation Profile.

At least:

\begin{verbatim}
context_id
version
content_fingerprint
\end{verbatim}

\hypertarget{sysmlgenerationprofilereference}{%
\paragraph{SysMLGenerationProfileReference}\label{sysmlgenerationprofilereference}}

Pinned identity of the IEM → SysML semantic mapping profile.

At least:

\begin{verbatim}
profile_id
profile_version
profile_fingerprint
\end{verbatim}

\hypertarget{sysmlartifactstructurereference}{%
\paragraph{SysMLArtifactStructureReference}\label{sysmlartifactstructurereference}}

Pinned identity of the artifact/package layout profile.

At least:

\begin{verbatim}
profile_id
profile_version
profile_fingerprint
\end{verbatim}

\hypertarget{sysmlgenerationcontext}{%
\paragraph{SysMLGenerationContext}\label{sysmlgenerationcontext}}

At least:

\begin{verbatim}
target_notation_reference
generation_profile_reference
artifact_structure_reference
generator_rules_reference
\end{verbatim}

\hypertarget{sysmlgenerationfinding}{%
\paragraph{SysMLGenerationFinding}\label{sysmlgenerationfinding}}

At least:

\begin{verbatim}
code
message
issue_level
target_type
target_id
profile_rule_id
blocking
\end{verbatim}

\hypertarget{generatedsysmlunit}{%
\paragraph{GeneratedSysMLUnit}\label{generatedsysmlunit}}

At least:

\begin{verbatim}
unit_id
relative_path
content
content_fingerprint
generated_symbol_ids
source_internal_model_element_ids
source_internal_model_relationship_ids
\end{verbatim}

\hypertarget{generatedsysmltraceabilityentry}{%
\paragraph{GeneratedSysMLTraceabilityEntry}\label{generatedsysmltraceabilityentry}}

At least:

\begin{verbatim}
generated_unit_id
generated_symbol_id
generated_location

source_internal_engineering_model_id
source_internal_model_element_id | null
source_internal_model_relationship_id | null

source_model_candidate_id
approved_input_references
review_decision_reference
accepted_exception_reference | null
\end{verbatim}

\hypertarget{generatedsysmlartifactset}{%
\paragraph{GeneratedSysMLArtifactSet}\label{generatedsysmlartifactset}}

At least:

\begin{verbatim}
schema_version
project_id

source_internal_engineering_model_id
source_iem_content_fingerprint

generation_context
generation_input_fingerprint
generation_provenance

units
traceability_entries
nonblocking_diagnostics

content_fingerprint
\end{verbatim}

\begin{center}\rule{0.5\linewidth}{0.5pt}\end{center}

\hypertarget{generation-flow}{%
\subsubsection{Generation flow}\label{generation-flow}}

The normative Phase-J production flow is:

\begin{verbatim}
InternalModelReadService
        │
        ▼
validated explicit InternalEngineeringModelSnapshot
        │
        ▼
Generation Context Resolver
        │
        ├── Target Notation Profile
        ├── Generation Profile
        └── Artifact Structure Profile
        │
        ▼
Generation Preflight
        │
        ├── mapping completeness
        ├── construct allow-list
        ├── symbol planning
        ├── escaping/renderability
        └── endpoint planning
        │
        ▼
Package Projection
        │
        ▼
Element Rendering
        │
        ▼
Relationship Rendering
        │
        ▼
Traceability Projection
        │
        ▼
Canonical Unit Formatting
        │
        ▼
Immutable GeneratedSysMLArtifactSet
        │
        ▼
Phase K validation
\end{verbatim}

No LLM exists in this production flow.

\begin{center}\rule{0.5\linewidth}{0.5pt}\end{center}

\hypertarget{reference-and-validation-flow-for-new-syntax}{%
\subsubsection{Reference and validation flow for new syntax}\label{reference-and-validation-flow-for-new-syntax}}

When Phase J requires a construct not yet supported by the active target notation:

\begin{verbatim}
required IEM semantic
        │
        ▼
inspect local SysML v2 release repository
        │
        ▼
smallest useful syntax fixture
        │
        ▼
validate fixture in SYSIDE
        │
   ┌────┴────┐
   │         │
 valid     invalid
   │         │
   ▼         ▼
record      do not activate
pattern     mapping
   │
   ▼
extend Target Notation Profile
   │
   ▼
extend Generation Profile
   │
   ▼
automated generator regression
\end{verbatim}

Apollo 11 may be consulted to find representative usage patterns, but it does not replace the specification/SYSIDE validation steps.

\begin{center}\rule{0.5\linewidth}{0.5pt}\end{center}

\hypertarget{default-artifact-structure}{%
\subsubsection{Default artifact structure}\label{default-artifact-structure}}

The initial MVP profile shall target conceptually:

\begin{verbatim}
generated_model.sysml

package GeneratedModel {
    package StakeholderLevel {
        package Stakeholders { ... }
        package UserNeeds { ... }
        package StakeholderRequirements { ... }
        package UseCases { ... }
    }

    package SystemLevel {
        package Requirements { ... }
        package Functional { ... }
        package Logical { ... }
        package Physical { ... }
    }

    package SubsystemLevel {
        package Requirements { ... }
        package Functional { ... }
        package Logical { ... }
        package Physical { ... }
    }
}
\end{verbatim}

The example is architectural and not yet an accepted literal syntax fixture.

Exact package identifier/quoting rules shall be established through the Target-Notation and Artifact Structure Profiles.

\begin{center}\rule{0.5\linewidth}{0.5pt}\end{center}

\hypertarget{relationship-to-existing-context-artifacts}{%
\subsubsection{Relationship to existing context artifacts}\label{relationship-to-existing-context-artifacts}}

\hypertarget{contextsourcessource_manifestjson}{%
\paragraph{\texorpdfstring{\texttt{context/sources/source\_manifest.json}}{context/sources/source\_manifest.json}}\label{contextsourcessource_manifestjson}}

Continues to define:

\begin{itemize}
\tightlist
\item
  SysML v2 release repository as the primary specification reference,
\item
  Apollo 11 as non-normative example reference,
\item
  generated output as a distinct artifact class.
\end{itemize}

\hypertarget{contextsysmlsysml_v2_spec_referencejson}{%
\paragraph{\texorpdfstring{\texttt{context/sysml/sysml\_v2\_spec\_reference.json}}{context/sysml/sysml\_v2\_spec\_reference.json}}\label{contextsysmlsysml_v2_spec_referencejson}}

Remains the compact working reference for routine implementation.

It is not the complete grammar.

Uncertain or newly required constructs trigger direct reference to the local SysML v2 release repository.

\hypertarget{contextsysmlsysml_v2_target_notationjson}{%
\paragraph{\texorpdfstring{\texttt{context/sysml/sysml\_v2\_target\_notation.json}}{context/sysml/sysml\_v2\_target\_notation.json}}\label{contextsysmlsysml_v2_target_notationjson}}

Becomes the Phase-J allowed-construct contract.

J1 shall revise the current \texttt{0.1.0} profile before production generation.

\hypertarget{contextexamplesapollo11_structure_referencemd}{%
\paragraph{\texorpdfstring{\texttt{context/examples/apollo11\_structure\_reference.md}}{context/examples/apollo11\_structure\_reference.md}}\label{contextexamplesapollo11_structure_referencemd}}

Remains a non-normative structure/style reference.

No Apollo engineering content, CoSMA framework or package hierarchy becomes Turing output policy merely because it appears in Apollo.

\hypertarget{catia-turing-generator-model}{%
\paragraph{CATIA Turing Generator model}\label{catia-turing-generator-model}}

Remains the engineering authority for the Turing Generator product.

It is not a Phase-J target-syntax reference.

Any capability gaps introduced or discovered during J--L are recorded for later Phase-N2 Architecture-to-Requirements reconciliation.

\begin{center}\rule{0.5\linewidth}{0.5pt}\end{center}

\hypertarget{persistence-and-publication-boundary}{%
\subsubsection{Persistence and publication boundary}\label{persistence-and-publication-boundary}}

Phase J creates a deterministic immutable artifact-set value.

A successful Phase-J result is not yet final published output.

Diagnostic implementations may serialize temporary artifacts for tests or run-owned evidence, but such files do not become final generated product artifacts merely because they exist on disk.

Final product publication remains Phase L after required Phase-K validation.

\texttt{data/output/} is therefore the final published generated-output location, not a reason to bypass the validation boundary.

\begin{center}\rule{0.5\linewidth}{0.5pt}\end{center}

\hypertarget{phase-j-implementation-slices}{%
\subsubsection{Phase-J implementation slices}\label{phase-j-implementation-slices}}

The accepted implementation decomposition after this ADR is intended to be:

\begin{verbatim}
J1  Generation foundation
    - errors and identifiers
    - immutable generation domain contracts
    - profile references / fingerprints
    - Target Notation 0.2 cleanup
    - controlled syntax fixtures for currently required unsupported constructs

J2  Generation profiles and preflight
    - Generation Profile
    - Artifact Structure Profile
    - deterministic profile loaders
    - mapping completeness
    - construct allow-list validation
    - blocking generation findings

J3  Package, symbol and canonical ordering projection
    - Framework structure → packages
    - stable generated symbols
    - deterministic naming / escaping
    - canonical ordering

J4  Element rendering
    - explicit element mapping rules
    - safe documentation projection
    - supported attribute projection
    - element traceability

J5  Relationship rendering
    - explicit relationship semantic mapping
    - directionality
    - exact generated endpoint binding
    - no generic fallback

J6  Artifact-set and traceability assembly
    - GeneratedSysMLUnit
    - GeneratedSysMLTraceabilityEntry
    - GeneratedSysMLArtifactSet
    - deterministic generation fingerprints
    - exact regeneration/idempotence tests

J7  Phase-K boundary and completion
    - explicit validation handoff
    - H9/I→J compatibility regression
    - full repository regression
    - SSOT update
    - Phase-N2 reconciliation candidate recording
\end{verbatim}

The slices may be internally subdivided without changing the accepted architecture.

\begin{center}\rule{0.5\linewidth}{0.5pt}\end{center}

\hypertarget{phase-j-acceptance-criteria}{%
\subsubsection{Phase-J acceptance criteria}\label{phase-j-acceptance-criteria}}

Phase J is complete only when:

\begin{itemize}
\tightlist
\item
  one explicit validated IEM can be generated without bypassing the I→J boundary,
\item
  generation is deterministic and requires no LLM,
\item
  every generated construct is allowed by the pinned Target Notation Profile,
\item
  every rendered IEM semantic is covered by one explicit Generation Profile rule,
\item
  unsupported required semantics fail closed,
\item
  no accepted IEM element or relationship disappears silently,
\item
  source-provided text cannot inject SysML syntax,
\item
  generated relationships bind exact generated endpoints,
\item
  machine-readable IEM → generated-artifact traceability is complete,
\item
  repeated exact generation is byte-identical,
\item
  the artifact set is suitable for Phase-K validation,
\item
  J does not publish final output directly,
\item
  representative supported syntax patterns have SYSIDE validation evidence,
\item
  focused Phase-J tests pass,
\item
  the complete repository regression passes,
\item
  \texttt{git\ diff\ -\/-check} passes,
\item
  and SSOT is synchronized.
\end{itemize}

\begin{center}\rule{0.5\linewidth}{0.5pt}\end{center}

\hypertarget{consequences-6}{%
\subsubsection{Consequences}\label{consequences-6}}

\hypertarget{positive-consequences-5}{%
\paragraph{Positive consequences}\label{positive-consequences-5}}

\begin{itemize}
\tightlist
\item
  Human-reviewed engineering semantics remain authoritative.
\item
  Phase J is reproducible and testable.
\item
  LLM variability cannot alter generated syntax or model meaning.
\item
  SYSIDE compatibility becomes an explicit target rather than an assumed by-product.
\item
  SysML v2 specification authority remains separate from example-project style.
\item
  Apollo 11 remains useful without becoming normative.
\item
  The CATIA Turing Generator model is no longer conflated with target syntax.
\item
  Unsupported semantics are visible instead of silently weakened.
\item
  Target notation, semantic mapping and artifact layout can evolve independently.
\item
  The same IEM may later be serialized using a different accepted artifact structure without re-running semantic interpretation.
\item
  The generated artifact set provides a precise input to Phase K.
\item
  Final publication remains protected by Phase K and Phase L.
\end{itemize}

\hypertarget{trade-offs-3}{%
\paragraph{Trade-offs}\label{trade-offs-3}}

\begin{itemize}
\tightlist
\item
  Phase J introduces additional versioned profiles.
\item
  Not every IEM may initially be generatable with the limited MVP subset.
\item
  New semantic constructs require explicit syntax experiments and SYSIDE validation.
\item
  One-file MVP output favors robustness over immediate large-project modularity.
\item
  Stable generated symbols may be less visually elegant than display-name-based identifiers but provide stronger identity and endpoint stability.
\item
  Maintaining machine-readable traceability adds artifact volume.
\item
  Strict failure on unsupported semantics may require target-notation/profile extensions before a complete model can be generated.
\end{itemize}

These trade-offs are accepted because the objective is valid, faithful, traceable SYSIDE-compatible SysML v2 generation rather than broad but lossy text emission.

\begin{center}\rule{0.5\linewidth}{0.5pt}\end{center}

\hypertarget{alternatives-considered}{%
\subsubsection{Alternatives considered}\label{alternatives-considered}}

\hypertarget{alternative-a--llm-generates-sysml-v2-directly-from-the-iem}{%
\paragraph{Alternative A --- LLM generates SysML v2 directly from the IEM}\label{alternative-a--llm-generates-sysml-v2-directly-from-the-iem}}

Rejected.

Reason:

\begin{itemize}
\tightlist
\item
  creates a second semantic interpretation step after Human Review,
\item
  reduces reproducibility,
\item
  complicates validation,
\item
  can silently alter accepted relationship meaning,
\item
  increases token/request use,
\item
  and makes exact regeneration difficult.
\end{itemize}

\hypertarget{alternative-b--generate-directly-from-model-candidates}{%
\paragraph{Alternative B --- Generate directly from Model Candidates}\label{alternative-b--generate-directly-from-model-candidates}}

Rejected.

Reason:

\begin{itemize}
\tightlist
\item
  bypasses the accepted Phase-I authority boundary,
\item
  duplicates internal-model assembly responsibilities,
\item
  and risks generation from content outside the exact IEM snapshot.
\end{itemize}

\hypertarget{alternative-c--use-apollo-11-as-the-syntax-template}{%
\paragraph{Alternative C --- Use Apollo 11 as the syntax template}\label{alternative-c--use-apollo-11-as-the-syntax-template}}

Rejected.

Reason:

\begin{itemize}
\tightlist
\item
  Apollo 11 is intentionally non-normative,
\item
  contains domain- and methodology-specific patterns,
\item
  and shall not override the SysML v2 specification or SYSIDE validation.
\end{itemize}

\hypertarget{alternative-d--use-the-turing-generator-catia-model-as-the-generator-syntax-template}{%
\paragraph{Alternative D --- Use the Turing Generator CATIA model as the generator syntax template}\label{alternative-d--use-the-turing-generator-catia-model-as-the-generator-syntax-template}}

Rejected.

Reason:

\begin{itemize}
\tightlist
\item
  it models the Turing Generator product,
\item
  it is not the normative source for target-model textual syntax,
\item
  and conflating product architecture with serialization policy would create a false authority relationship.
\end{itemize}

\hypertarget{alternative-e--emit-generic-dependencies-for-unsupported-relationships}{%
\paragraph{Alternative E --- Emit generic dependencies for unsupported relationships}\label{alternative-e--emit-generic-dependencies-for-unsupported-relationships}}

Rejected.

Reason:

\begin{itemize}
\tightlist
\item
  syntactic success would hide semantic loss,
\item
  accepted Human Review meaning would be altered downstream,
\item
  and generated output would no longer faithfully represent the IEM.
\end{itemize}

\hypertarget{alternative-f--write-directly-to-dataoutput-in-phase-j}{%
\paragraph{\texorpdfstring{Alternative F --- Write directly to \texttt{data/output/} in Phase J}{Alternative F --- Write directly to data/output/ in Phase J}}\label{alternative-f--write-directly-to-dataoutput-in-phase-j}}

Rejected.

Reason:

\begin{itemize}
\tightlist
\item
  bypasses the planned validation/publication separation,
\item
  conflates generated work product with accepted published artifact,
\item
  and weakens the Phase-K/Phase-L gates.
\end{itemize}

\hypertarget{alternative-g--start-immediately-with-many-separate-sysml-files}{%
\paragraph{\texorpdfstring{Alternative G --- Start immediately with many separate \texttt{.sysml} files}{Alternative G --- Start immediately with many separate .sysml files}}\label{alternative-g--start-immediately-with-many-separate-sysml-files}}

Deferred.

Reason:

\begin{itemize}
\tightlist
\item
  increases cross-file reference complexity before the base mapping is proven,
\item
  broadens the first SYSIDE validation surface,
\item
  and is unnecessary for the initial closed vertical slice.
\end{itemize}

The artifact-set contract deliberately leaves multi-file generation open for a later Artifact Structure Profile.

\begin{center}\rule{0.5\linewidth}{0.5pt}\end{center}

\hypertarget{catia--phase-n2-reconciliation}{%
\subsubsection{CATIA / Phase-N2 reconciliation}\label{catia--phase-n2-reconciliation}}

This ADR makes an implementation-architecture decision for Phase J.

It does not modify the CATIA model.

Capabilities or constraints that may require CATIA Requirement/Function reconciliation shall be recorded during implementation and reconciled in Phase N2 together with the other G--L architecture changes.

Potential N2 reconciliation topics introduced or made explicit by ADR-021 include:

\begin{itemize}
\tightlist
\item
  deterministic generation from the validated Internal Engineering Model,
\item
  explicit SYSIDE-compatible target-notation generation,
\item
  versioned semantic-to-SysML generation profiles,
\item
  fail-closed unsupported semantic behavior,
\item
  machine-readable generated-artifact traceability,
\item
  deterministic generation identity/idempotence,
\item
  and separation of generation, validation and publication.
\end{itemize}

No SYSR, SF, Logical Component or allocation is created or changed by this ADR.

\begin{center}\rule{0.5\linewidth}{0.5pt}\end{center}

\hypertarget{decision-state-after-acceptance}{%
\subsubsection{Decision state after acceptance}\label{decision-state-after-acceptance}}

After explicit acceptance of ADR-021:

\begin{verbatim}
Phase H/H9: COMPLETE
Phase I:    COMPLETE
ADR-021:    ACCEPTED
Phase J:    architecture accepted
Next:       J1 — Generation foundation and validated syntax-fixture preparation
\end{verbatim}

No Phase-J implementation begins before explicit acceptance of this ADR.

\textbackslash newpage

\hypertarget{adr-022--sysml-v2-validation-layer-architecture}{%
\subsection{ADR-022 --- SysML v2 Validation Layer Architecture}\label{adr-022--sysml-v2-validation-layer-architecture}}

\hypertarget{status-4}{%
\subsubsection{Status}\label{status-4}}

Accepted

\hypertarget{date-4}{%
\subsubsection{Date}\label{date-4}}

2026-08-14

\hypertarget{context-7}{%
\subsubsection{Context}\label{context-7}}

Phase J is completed and provides one deterministic, immutable validation-ready \texttt{GeneratedSysMLArtifactSet}.

The accepted end-to-end boundary is:

\begin{verbatim}
Internal Engineering Model
→ Phase J — deterministic SysML v2 generation
→ GeneratedSysMLArtifactSet
→ Phase K — validation
→ SysMLValidationResult
→ Phase L — versioned publication
\end{verbatim}

Phase J already guarantees generation-time integrity including deterministic serialization, exact generation-policy references, generated unit fingerprints, artifact-set fingerprints, complete IME/IMR coverage, generated locations and machine-readable traceability.

Phase K has a different responsibility.

It shall validate the complete generated artifact before publication without:

\begin{itemize}
\tightlist
\item
  regenerating SysML v2,
\item
  changing generated text,
\item
  reinterpreting engineering semantics,
\item
  bypassing the Phase-J artifact contract,
\item
  or publishing output itself.
\end{itemize}

The roadmap requires Phase K to cover:

\begin{itemize}
\tightlist
\item
  syntax validation,
\item
  Target Notation validation,
\item
  Artifact Structure validation,
\item
  relationship and endpoint consistency,
\item
  traceability,
\item
  comparability/profile consistency,
\item
  deterministic findings,
\item
  and a fail-closed Phase-L publication gate.
\end{itemize}

The current Phase-J syntax evidence was established through small controlled SYSIDE fixtures. Those checks authorize supported syntax patterns but do not constitute validation of a complete generated target model.

Phase K therefore requires both deterministic Turing-specific validation and an explicit external SysML v2 compatibility-validation boundary.

\begin{center}\rule{0.5\linewidth}{0.5pt}\end{center}

\hypertarget{decision-7}{%
\subsubsection{Decision}\label{decision-7}}

\hypertarget{k-01--sole-upstream-validation-subject}{%
\paragraph{K-01 --- Sole upstream validation subject}\label{k-01--sole-upstream-validation-subject}}

The sole normative Phase-K input is one explicit immutable:

\texttt{GeneratedSysMLArtifactSet}

The service boundary is:

\begin{Shaded}
\begin{Highlighting}[]
\NormalTok{SysMLValidationService.validate(}
\NormalTok{    artifact\_set: GeneratedSysMLArtifactSet,}
\NormalTok{) }\OperatorTok{{-}\textgreater{}}\NormalTok{ SysMLValidationResult}
\end{Highlighting}
\end{Shaded}

Phase K shall not independently reload:

\begin{itemize}
\tightlist
\item
  Approved Inputs,
\item
  Model Candidates,
\item
  Candidate Review persistence,
\item
  raw Processing evidence,
\item
  UI state,
\item
  or the source IEM for semantic reinterpretation
\end{itemize}

in order to determine whether the generated model is valid.

There is no implicit latest-artifact selection.

The exact validation subject is bound by \texttt{GeneratedSysMLArtifactSet.content\_fingerprint}.

\begin{center}\rule{0.5\linewidth}{0.5pt}\end{center}

\hypertarget{k-02--validation-is-observational-and-never-repairs}{%
\paragraph{K-02 --- Validation is observational and never repairs}\label{k-02--validation-is-observational-and-never-repairs}}

Phase K shall never silently modify generated content in order to make validation pass.

Invalid content shall produce explicit findings and block publication.

Phase K shall not:

\begin{itemize}
\tightlist
\item
  rename generated symbols,
\item
  reorder generated packages,
\item
  alter relationship endpoints,
\item
  substitute target constructs,
\item
  normalize engineering descriptions,
\item
  auto-format and replace generated content,
\item
  or invoke Phase J to silently regenerate corrected output.
\end{itemize}

Resolution belongs to the phase, policy artifact or implementation responsible for the finding.

\begin{center}\rule{0.5\linewidth}{0.5pt}\end{center}

\hypertarget{k-03--two-layer-validation-architecture}{%
\paragraph{K-03 --- Two-layer validation architecture}\label{k-03--two-layer-validation-architecture}}

Phase K shall contain two intentionally separate validation layers.

\hypertarget{deterministic-internal-validation}{%
\subparagraph{Deterministic internal validation}\label{deterministic-internal-validation}}

Implemented inside the Turing Generator.

It validates Turing-specific generated-artifact invariants including:

\begin{itemize}
\tightlist
\item
  artifact identity and fingerprint integrity,
\item
  generation-context integrity,
\item
  Target Notation reference integrity,
\item
  Generation Profile reference integrity,
\item
  Artifact Structure reference integrity,
\item
  Generator Rules reference integrity,
\item
  generated-unit structure,
\item
  target-subset conformance,
\item
  package and artifact-layout rules,
\item
  relationship and endpoint consistency,
\item
  traceability integrity,
\item
  accepted-exception preservation,
\item
  and publication-policy invariants.
\end{itemize}

These checks require no external modeling tool.

\hypertarget{external-sysml-v2-compatibility-validation}{%
\subparagraph{External SysML v2 compatibility validation}\label{external-sysml-v2-compatibility-validation}}

Implemented through an explicit external-validator adapter.

The MVP target validator is the SYSIDE Modeler CLI using its non-mutating validation/check capability.

External validation establishes parser/tool compatibility for the actual complete generated artifact set.

Internal deterministic validation shall not attempt to become a complete SysML v2 parser.

External validation shall not replace Turing-specific deterministic contract validation.

Both layers are required for a publication-ready PASS.

\begin{center}\rule{0.5\linewidth}{0.5pt}\end{center}

\hypertarget{k-04--separate-versioned-validation-profile}{%
\paragraph{K-04 --- Separate versioned Validation Profile}\label{k-04--separate-versioned-validation-profile}}

Phase K introduces:

\begin{verbatim}
context/sysml/turing_sysml_v2_validation_profile.json
\end{verbatim}

Initial identity:

\begin{verbatim}
profile_id:      TURING_SYSML_V2_VALIDATION
profile_version: 1.0.0
\end{verbatim}

The Validation Profile defines:

\begin{quote}
Which validation checks are required, how findings are classified, which external validator is required, and what constitutes a Phase-K publication PASS.
\end{quote}

It shall define at least:

\begin{itemize}
\tightlist
\item
  required internal validators,
\item
  required external validator,
\item
  external-validator availability policy,
\item
  diagnostic severity policy,
\item
  publication-blocking policy,
\item
  diagnostic normalization policy,
\item
  and validation-result fingerprint policy.
\end{itemize}

It shall not redefine:

\begin{itemize}
\tightlist
\item
  SysML v2 target syntax,
\item
  IEM → SysML semantic mappings,
\item
  package organization,
\item
  generator formatting,
\item
  or engineering semantics.
\end{itemize}

Those responsibilities remain with the existing Target Notation, Generation Profile, Artifact Structure Profile, Generator Rules and Model Structure / Comparability Profile.

\begin{center}\rule{0.5\linewidth}{0.5pt}\end{center}

\hypertarget{k-05--exact-generation-policy-resolution}{%
\paragraph{K-05 --- Exact generation-policy resolution}\label{k-05--exact-generation-policy-resolution}}

\texttt{GeneratedSysMLArtifactSet.generation\_context} contains exact references and fingerprints for:

\begin{itemize}
\tightlist
\item
  Target Notation,
\item
  Generation Profile,
\item
  Artifact Structure Profile,
\item
  Generator Rules.
\end{itemize}

Phase K shall resolve these exact referenced policy artifacts and require exact fingerprint agreement.

Conceptually:

\begin{verbatim}
artifact-set pinned reference
        │
        ▼
resolve exact policy artifact
        │
        ▼
validate policy artifact
        │
        ▼
recalculate fingerprint
        │
        ├── exact match → continue
        └── mismatch / unavailable → blocking finding
\end{verbatim}

Phase K shall not silently substitute a newer policy version.

The Phase-J generation baseline is:

\begin{verbatim}
CTX_SYSML_V2_TARGET_NOTATION        0.2.0
TURING_SYSML_V2_GENERATION         1.0.0
TURING_SYSML_V2_ARTIFACT_STRUCTURE 1.0.0
TURING_SYSML_V2_GENERATOR_RULES    1.0.0
\end{verbatim}

A generated artifact whose exact generation policy cannot be resolved shall not be publishable.

\begin{center}\rule{0.5\linewidth}{0.5pt}\end{center}

\hypertarget{k-06--standalone-artifact-set-boundary-validation}{%
\paragraph{K-06 --- Standalone artifact-set boundary validation}\label{k-06--standalone-artifact-set-boundary-validation}}

Phase K shall independently validate the received Phase-J artifact contract without requiring the original IEM snapshot.

Checks shall include at least:

\begin{itemize}
\tightlist
\item
  \texttt{GeneratedSysMLArtifactSet} content fingerprint,
\item
  generated-unit content fingerprints,
\item
  generation-input fingerprint,
\item
  unit identity uniqueness,
\item
  unit path uniqueness,
\item
  safe relative paths,
\item
  generated-symbol uniqueness,
\item
  traceability-entry uniqueness,
\item
  traceability unit references,
\item
  traceability generated-symbol references,
\item
  traceability line ranges,
\item
  source-IEM identity consistency,
\item
  element/relationship trace coverage consistency,
\item
  controlled line-ending invariants,
\item
  and textual-unit integrity.
\end{itemize}

This is a boundary check of received evidence.

It does not remove Phase J's responsibility to build correct artifacts.

\begin{center}\rule{0.5\linewidth}{0.5pt}\end{center}

\hypertarget{k-07--target-notation-validation}{%
\paragraph{K-07 --- Target Notation validation}\label{k-07--target-notation-validation}}

Phase K shall confirm that generated output remains inside the exact permitted Target Notation subset.

The internal validator shall validate the deliberately constrained Turing Generator output forms rather than implement the complete SysML v2 grammar.

The MVP production subset currently includes generated forms for:

\begin{itemize}
\tightlist
\item
  Package,
\item
  Documentation block,
\item
  Requirement Usage,
\item
  Use Case Definition,
\item
  Action Usage,
\item
  Part Usage,
\item
  Dependency,
\item
  Allocation,
\item
  Satisfaction,
\item
  and qualified references.
\end{itemize}

A generated construct outside the resolved Target Notation contract is blocking.

Full syntax/model interpretation remains the responsibility of external SYSIDE validation.

\begin{center}\rule{0.5\linewidth}{0.5pt}\end{center}

\hypertarget{k-08--artifact-structure-validation}{%
\paragraph{K-08 --- Artifact Structure validation}\label{k-08--artifact-structure-validation}}

The generated artifact set shall be checked against the exact resolved Artifact Structure Profile.

For the current MVP this includes at least:

\begin{verbatim}
one generated unit
unit id GSU-000001
relative path generated_model.sysml
root package GeneratedModel
complete configured Framework package hierarchy
configured empty packages retained
canonical Framework ordering
canonical element ordering
canonical relationship ordering
relationships placed at the configured root location
safe relative output paths
\end{verbatim}

Phase K validates generated structure.

It shall not infer additional engineering containment.

\begin{center}\rule{0.5\linewidth}{0.5pt}\end{center}

\hypertarget{k-09--relationship-and-endpoint-validation}{%
\paragraph{K-09 --- Relationship and endpoint validation}\label{k-09--relationship-and-endpoint-validation}}

Phase K shall validate generated relationship consistency independently of Phase-J rendering.

For every generated relationship it shall validate, where deterministically observable:

\begin{itemize}
\tightlist
\item
  the relationship construct is permitted,
\item
  both referenced endpoints resolve,
\item
  endpoint resolution is unambiguous,
\item
  endpoint target constructs are compatible,
\item
  endpoint element kinds are compatible,
\item
  qualified endpoint references are valid,
\item
  and required endpoint rendering/order is preserved.
\end{itemize}

The active production rules include:

\begin{verbatim}
allocation
  source: Feature
  target: Feature

dependency / depends_on
  source: Feature or Definition
  target: Feature or Definition

satisfaction
  IEM semantic: source satisfies target
  generated form: satisfy TARGET by SOURCE
  source: PartUsage or ActionUsage Feature
  target: RequirementUsage Feature
\end{verbatim}

Phase K validates generated-model consistency.

It shall not reinterpret the original IMR semantic intent.

The reviewed IEM semantic → target-construct mapping remains a Phase-J generation responsibility.

\begin{center}\rule{0.5\linewidth}{0.5pt}\end{center}

\hypertarget{k-10--scope-of-semantic-constraint-and-comparability-validation}{%
\paragraph{K-10 --- Scope of semantic, constraint and comparability validation}\label{k-10--scope-of-semantic-constraint-and-comparability-validation}}

The roadmap terms:

\begin{itemize}
\tightlist
\item
  relationship semantic validation,
\item
  constraint validation,
\item
  comparability-profile validation
\end{itemize}

shall be interpreted according to the accepted Phase-H/I/J authority boundaries.

Phase K validates publication-relevant generated-model consistency.

It shall not rerun Phase-H semantic interpretation or Human Review.

For the MVP:

\hypertarget{relationship-semantic-validation}{%
\subparagraph{Relationship semantic validation}\label{relationship-semantic-validation}}

Means confirming that the generated relationship is permitted by the pinned Generation Profile and that endpoint roles are compatible with the target construct.

\hypertarget{constraint-validation}{%
\subparagraph{Constraint validation}\label{constraint-validation}}

Means deterministic generated-artifact and target-policy constraints applicable to publication.

It does not mean executing arbitrary product-design constraints or simulating engineering behavior.

\hypertarget{comparability-profile-validation}{%
\subparagraph{Comparability-profile validation}\label{comparability-profile-validation}}

Means validating consistency of the pinned generation-policy chain with the expected Model Structure / Comparability Profile and preservation of reviewed exception traceability.

Candidate-level comparability decisions shall not be recomputed in Phase K.

\begin{center}\rule{0.5\linewidth}{0.5pt}\end{center}

\hypertarget{k-11--traceability-validation}{%
\paragraph{K-11 --- Traceability validation}\label{k-11--traceability-validation}}

Every generated engineering representation must remain traceable.

Phase K shall validate that the transferred chain remains structurally complete:

\begin{verbatim}
generated location
→ GeneratedSysMLTraceabilityEntry
→ source IEM
→ IME or IMR
→ Model Candidate
→ Approved Input
→ Human Review Decision
→ accepted exception where applicable
\end{verbatim}

Every generated element and relationship representation shall have exactly one applicable traceability entry.

Generated line locations remain one-based and inclusive.

A finding associated with a generated element or relationship shall reference the existing Phase-J traceability key:

\begin{verbatim}
generated_unit_id
+
generated_symbol_id
\end{verbatim}

Findings applying to complete units or package structure may reference the unit and generated location only.

\begin{center}\rule{0.5\linewidth}{0.5pt}\end{center}

\hypertarget{k-12--immutable-normalized-validation-findings}{%
\paragraph{K-12 --- Immutable normalized validation findings}\label{k-12--immutable-normalized-validation-findings}}

Phase K shall produce immutable normalized findings.

Conceptually:

\begin{Shaded}
\begin{Highlighting}[]
\NormalTok{SysMLValidationFinding(}
\NormalTok{    code}\OperatorTok{=}\NormalTok{...,}
\NormalTok{    category}\OperatorTok{=}\NormalTok{...,}
\NormalTok{    severity}\OperatorTok{=}\NormalTok{...,}
\NormalTok{    blocking}\OperatorTok{=}\NormalTok{...,}
\NormalTok{    message}\OperatorTok{=}\NormalTok{...,}
\NormalTok{    generated\_unit\_id}\OperatorTok{=}\NormalTok{...,}
\NormalTok{    generated\_symbol\_id}\OperatorTok{=}\NormalTok{...,}
\NormalTok{    generated\_location}\OperatorTok{=}\NormalTok{...,}
\NormalTok{    validator\_id}\OperatorTok{=}\NormalTok{...,}
\NormalTok{    validator\_rule\_id}\OperatorTok{=}\NormalTok{...,}
\NormalTok{)}
\end{Highlighting}
\end{Shaded}

Finding categories shall distinguish at least:

\begin{verbatim}
artifact_integrity
validation_context
target_notation
artifact_structure
relationship_consistency
traceability
external_syntax
external_semantics
external_warning
validator_infrastructure
\end{verbatim}

Severity shall use a controlled vocabulary:

\begin{verbatim}
info
warning
error
\end{verbatim}

\texttt{blocking} is explicit and controlled by the Validation Profile.

Findings shall have deterministic canonical ordering.

Recommended ordering:

\begin{verbatim}
blocking first
→ category
→ generated unit
→ line
→ column where available
→ code
→ message
\end{verbatim}

\begin{center}\rule{0.5\linewidth}{0.5pt}\end{center}

\hypertarget{k-13--syside-external-validator-adapter}{%
\paragraph{K-13 --- SYSIDE external-validator adapter}\label{k-13--syside-external-validator-adapter}}

SYSIDE shall be integrated through a dedicated adapter rather than directly into the Phase-K orchestration service.

Conceptually:

\begin{verbatim}
SysMLExternalValidator
        ▲
        │
SysideCliValidator
\end{verbatim}

The adapter shall receive the exact generated units and materialize an isolated ephemeral validation workspace.

Generated units shall be written there byte-for-byte using their configured relative paths.

The workspace shall contain no unrelated project models.

The adapter shall then execute a controlled, non-mutating SYSIDE validation operation conceptually equivalent to:

\begin{verbatim}
syside check
\end{verbatim}

The concrete invocation shall use deterministic/noninteractive settings where supported.

Temporary absolute filesystem paths shall not become part of deterministic validation identity.

External diagnostics shall be normalized to generated-unit-relative locations.

\begin{center}\rule{0.5\linewidth}{0.5pt}\end{center}

\hypertarget{k-14--external-validator-identity}{%
\paragraph{K-14 --- External validator identity}\label{k-14--external-validator-identity}}

Every external validation execution shall record the actual validator identity.

At least:

\begin{itemize}
\tightlist
\item
  validator ID,
\item
  tool name,
\item
  tool version,
\item
  validation command contract,
\item
  validator configuration fingerprint.
\end{itemize}

For SYSIDE, the installed version shall be discovered at execution time and recorded.

The external-validator identity participates in the validation execution fingerprint.

A different SYSIDE version therefore represents a materially different validation environment.

\begin{center}\rule{0.5\linewidth}{0.5pt}\end{center}

\hypertarget{k-15--unavailable-external-validator-is-incomplete-never-pass}{%
\paragraph{K-15 --- Unavailable external validator is incomplete, never PASS}\label{k-15--unavailable-external-validator-is-incomplete-never-pass}}

External validator infrastructure may be unavailable because, for example:

\begin{itemize}
\tightlist
\item
  SYSIDE CLI is not installed,
\item
  required licensing is unavailable,
\item
  the executable cannot start,
\item
  the validator version is unsupported,
\item
  execution crashes or times out,
\item
  required standard-library context is unavailable.
\end{itemize}

Such states shall not be reported as \texttt{invalid}.

They shall also never be reported as \texttt{valid}.

Instead:

\begin{verbatim}
validation_status = incomplete
publication_gate = blocked
\end{verbatim}

with an explicit blocking \texttt{validator\_infrastructure} finding.

This preserves the distinction between:

\begin{verbatim}
invalid model
\end{verbatim}

and:

\begin{verbatim}
required validation could not be completed
\end{verbatim}

while remaining fail-closed.

\begin{center}\rule{0.5\linewidth}{0.5pt}\end{center}

\hypertarget{k-16--three-state-validation-result}{%
\paragraph{K-16 --- Three-state validation result}\label{k-16--three-state-validation-result}}

\texttt{SysMLValidationResult} shall support exactly three normative overall statuses:

\begin{verbatim}
valid
invalid
incomplete
\end{verbatim}

\hypertarget{valid}{%
\subparagraph{valid}\label{valid}}

All required internal checks completed.

All required external validators completed.

No blocking findings exist.

\begin{verbatim}
publication_gate = passed
\end{verbatim}

\hypertarget{invalid}{%
\subparagraph{invalid}\label{invalid}}

Validation completed sufficiently to establish one or more actual blocking model/artifact defects.

\begin{verbatim}
publication_gate = blocked
\end{verbatim}

\hypertarget{incomplete}{%
\subparagraph{incomplete}\label{incomplete}}

A required validator or validation context could not be resolved or executed sufficiently to establish validity.

\begin{verbatim}
publication_gate = blocked
\end{verbatim}

Only:

\begin{verbatim}
validation_status == valid
AND
publication_gate == passed
\end{verbatim}

is eligible for Phase L.

\begin{center}\rule{0.5\linewidth}{0.5pt}\end{center}

\hypertarget{k-17--warning-policy}{%
\paragraph{K-17 --- Warning policy}\label{k-17--warning-policy}}

Internal and external warnings shall remain visible.

Warnings are not automatically equivalent to errors.

The Validation Profile determines whether a warning category or specific validator rule blocks publication.

The MVP default policy is:

\begin{verbatim}
SYSIDE warning
→ severity = warning
→ blocking = false
\end{verbatim}

unless an explicit Turing validation rule elevates the finding.

Errors remain blocking.

The MVP shall therefore not globally use a \texttt{warnings-as-errors} publication policy.

\begin{center}\rule{0.5\linewidth}{0.5pt}\end{center}

\hypertarget{k-18--deterministic-validation-identity-and-fingerprints}{%
\paragraph{K-18 --- Deterministic validation identity and fingerprints}\label{k-18--deterministic-validation-identity-and-fingerprints}}

Phase K shall calculate an exact validation-input fingerprint from at least:

\begin{itemize}
\tightlist
\item
  \texttt{GeneratedSysMLArtifactSet.content\_fingerprint},
\item
  Validation Profile reference and fingerprint,
\item
  resolved external-validator identity/version,
\item
  external-validator configuration fingerprint,
\item
  validation command-contract identifier.
\end{itemize}

Conceptually:

\begin{verbatim}
validation_input_fingerprint =
SHA256(
    artifact-set fingerprint
  + validation policy
  + resolved validator environment
)
\end{verbatim}

If a required validator is unavailable, the resolved environment shall contain an explicit unavailable state rather than omit validator identity.

The final \texttt{SysMLValidationResult.content\_fingerprint} shall cover deterministic result content including:

\begin{itemize}
\tightlist
\item
  validation-input fingerprint,
\item
  validation status,
\item
  publication gate,
\item
  normalized findings,
\item
  external-validator evidence,
\item
  artifact-set reference,
\item
  Validation Profile reference.
\end{itemize}

Deterministic validation identity shall not include:

\begin{itemize}
\tightlist
\item
  wall-clock timestamps,
\item
  temporary workspace paths,
\item
  execution duration,
\item
  ANSI formatting,
\item
  machine-specific absolute paths.
\end{itemize}

For identical artifact, validation policy, validator version/configuration and normalized validation result, the same validation-result fingerprint shall be produced.

\begin{center}\rule{0.5\linewidth}{0.5pt}\end{center}

\hypertarget{k-19--no-separate-phase-k-persistence-authority-required-for-the-mvp}{%
\paragraph{K-19 --- No separate Phase-K persistence authority required for the MVP}\label{k-19--no-separate-phase-k-persistence-authority-required-for-the-mvp}}

Phase K does not require a separate mutable validation repository merely to perform validation.

The normative Phase-K result is the immutable:

\begin{verbatim}
SysMLValidationResult
\end{verbatim}

Its cryptographic fingerprint is the validation-result identity transferred to Phase L.

Phase L may persist the validation result/report together with the published versioned output package.

A future validation-history repository may be added without changing the Phase-K service boundary.

\begin{center}\rule{0.5\linewidth}{0.5pt}\end{center}

\hypertarget{k-20--deterministic-human-readable-validation-report}{%
\paragraph{K-20 --- Deterministic human-readable validation report}\label{k-20--deterministic-human-readable-validation-report}}

The machine-readable \texttt{SysMLValidationResult} is the authoritative Phase-K output.

A human-readable validation report may be projected deterministically from that result.

It shall summarize at least:

\begin{itemize}
\tightlist
\item
  project identity,
\item
  source IEM identity,
\item
  source GeneratedSysMLArtifactSet fingerprint,
\item
  Validation Profile identity,
\item
  SYSIDE validator identity/version,
\item
  internal validation status,
\item
  external validation status,
\item
  blocking findings,
\item
  warnings,
\item
  publication-gate result,
\item
  traceability references.
\end{itemize}

The human-readable report is a projection.

It is not separate validation authority.

\begin{center}\rule{0.5\linewidth}{0.5pt}\end{center}

\hypertarget{k-21--exact-fingerprint-bound-phase-l-publication-gate}{%
\paragraph{K-21 --- Exact fingerprint-bound Phase-L publication gate}\label{k-21--exact-fingerprint-bound-phase-l-publication-gate}}

Phase L shall accept only the exact artifact set covered by the successful validation result.

Conceptually:

\begin{Shaded}
\begin{Highlighting}[]
\NormalTok{OutputWriter.publish(}
\NormalTok{    artifact\_set: GeneratedSysMLArtifactSet,}
\NormalTok{    validation\_result: SysMLValidationResult,}
\NormalTok{)}
\end{Highlighting}
\end{Shaded}

Phase L must verify:

\begin{verbatim}
validation_result.source_artifact_set_fingerprint
==
artifact_set.content_fingerprint
\end{verbatim}

and:

\begin{verbatim}
validation_result.validation_status == "valid"
validation_result.publication_gate == "passed"
\end{verbatim}

A validation result for another artifact set shall never authorize publication.

Any modification of generated content invalidates the publication gate through the artifact-set fingerprint.

\begin{center}\rule{0.5\linewidth}{0.5pt}\end{center}

\hypertarget{k-22--findings-report-likely-resolution-ownership}{%
\paragraph{K-22 --- Findings report likely resolution ownership}\label{k-22--findings-report-likely-resolution-ownership}}

Validation findings may identify their likely resolution boundary without automatically changing that boundary.

Examples:

\begin{verbatim}
artifact fingerprint mismatch
→ corrupted/invalid Phase-J artifact transfer

Target Notation reference mismatch
→ generation-policy/reference issue

unsupported generated target form
→ Target Notation / Generation Profile / generator issue

package-layout mismatch
→ Artifact Structure / generator issue

unresolved relationship endpoint
→ generator / generated-artifact consistency issue

SYSIDE syntax error
→ generator or Target Notation syntax issue

SYSIDE semantic/type error
→ Generation Profile / endpoint mapping / generator issue

traceability mismatch
→ Phase-J artifact assembly issue

SYSIDE unavailable
→ validation infrastructure issue
\end{verbatim}

Phase K reports the issue.

It does not repair engineering content.

\begin{center}\rule{0.5\linewidth}{0.5pt}\end{center}

\hypertarget{k-23--controlled-larger-context-compatibility}{%
\paragraph{K-23 --- Controlled larger-context compatibility}\label{k-23--controlled-larger-context-compatibility}}

The MVP external validation workspace shall validate exactly the generated artifact set plus the validator's controlled SysML standard-library context.

No unrelated repository model shall be loaded implicitly.

Future compatibility validation against additional target-model libraries or larger engineering contexts may be added through versioned Validation Profile configuration.

Additional context shall be explicitly pinned and fingerprinted.

Ambient filesystem discovery is not permitted as hidden validation authority.

\begin{center}\rule{0.5\linewidth}{0.5pt}\end{center}

\hypertarget{k-24--no-llm-in-the-normative-validation-path}{%
\paragraph{K-24 --- No LLM in the normative validation path}\label{k-24--no-llm-in-the-normative-validation-path}}

The normative Phase-K path is deterministic/tool-based.

No LLM is required to:

\begin{itemize}
\tightlist
\item
  decide whether syntax is valid,
\item
  determine endpoint compatibility,
\item
  decide publication eligibility,
\item
  repair invalid models,
\item
  or reinterpret validation findings.
\end{itemize}

An LLM may later explain findings in a UI.

Such explanation remains non-authoritative and cannot alter the validation result or publication gate.

\begin{center}\rule{0.5\linewidth}{0.5pt}\end{center}

\hypertarget{phase-k-output-contract}{%
\subsubsection{Phase-K Output Contract}\label{phase-k-output-contract}}

Conceptually:

\begin{Shaded}
\begin{Highlighting}[]
\AttributeTok{@dataclass}\NormalTok{(frozen}\OperatorTok{=}\VariableTok{True}\NormalTok{, slots}\OperatorTok{=}\VariableTok{True}\NormalTok{)}
\KeywordTok{class}\NormalTok{ SysMLValidationResult:}
\NormalTok{    schema\_version: }\BuiltInTok{str}

\NormalTok{    project\_id: }\BuiltInTok{str}
\NormalTok{    source\_internal\_engineering\_model\_id: }\BuiltInTok{str}
\NormalTok{    source\_artifact\_set\_fingerprint: }\BuiltInTok{str}

\NormalTok{    validation\_profile\_reference: SysMLValidationProfileReference}
\NormalTok{    validation\_input\_fingerprint: }\BuiltInTok{str}

\NormalTok{    external\_validator\_evidence: (}
        \BuiltInTok{tuple}\NormalTok{[SysMLExternalValidationEvidence, ...]}
\NormalTok{    )}

\NormalTok{    findings: }\BuiltInTok{tuple}\NormalTok{[SysMLValidationFinding, ...]}

\NormalTok{    validation\_status: }\BuiltInTok{str}
\NormalTok{    publication\_gate: }\BuiltInTok{str}

\NormalTok{    content\_fingerprint: }\BuiltInTok{str}
\end{Highlighting}
\end{Shaded}

A wall-clock timestamp is not required in the deterministic domain contract.

Operational publication/run metadata may record time independently without changing validation identity.

\begin{center}\rule{0.5\linewidth}{0.5pt}\end{center}

\hypertarget{normative-phase-k-flow}{%
\subsubsection{Normative Phase-K Flow}\label{normative-phase-k-flow}}

\begin{verbatim}
GeneratedSysMLArtifactSet
        │
        ▼
Artifact-set boundary integrity
        │
        ▼
Resolve exact generation policy references
        │
        ▼
Resolve Validation Profile
        │
        ▼
Deterministic internal validation
        │
        ├── target notation
        ├── artifact structure
        ├── relationship/endpoints
        └── traceability
        │
        ▼
Materialize isolated temporary validation workspace
        │
        ▼
SYSIDE CLI adapter
        │
        ▼
Normalize external diagnostics
        │
        ▼
Merge + canonical-order findings
        │
        ▼
validation status
        │
        ├── valid
        ├── invalid
        └── incomplete
        │
        ▼
publication gate
        │
        ├── passed
        └── blocked
        │
        ▼
Immutable SysMLValidationResult
        │
        ▼
Phase L
\end{verbatim}

\begin{center}\rule{0.5\linewidth}{0.5pt}\end{center}

\hypertarget{initial-implementation-decomposition-2}{%
\subsubsection{Initial Implementation Decomposition}\label{initial-implementation-decomposition-2}}

After acceptance, the initial implementation decomposition is:

\begin{verbatim}
K1  Validation domain foundation + Validation Profile
K2  Artifact/context/Target-Notation/Structure/Traceability validators
K3  Relationship + endpoint consistency validator
K4  SYSIDE CLI adapter + deterministic diagnostic normalization
K5  SysMLValidationService + status/gate/fingerprint assembly
K6  J→K→L boundary regression + Phase-K closeout
\end{verbatim}

The SYSIDE integration path shall be tested separately from the fast deterministic unit-test suite.

\begin{center}\rule{0.5\linewidth}{0.5pt}\end{center}

\hypertarget{phase-boundaries-1}{%
\subsubsection{Phase Boundaries}\label{phase-boundaries-1}}

\hypertarget{phase-j-owns}{%
\paragraph{Phase J owns}\label{phase-j-owns}}

\begin{itemize}
\tightlist
\item
  IEM semantic → SysML construct mapping,
\item
  rendering,
\item
  escaping,
\item
  package projection,
\item
  canonical textual generation,
\item
  generated traceability,
\item
  \texttt{GeneratedSysMLArtifactSet} assembly.
\end{itemize}

\hypertarget{phase-k-owns}{%
\paragraph{Phase K owns}\label{phase-k-owns}}

\begin{itemize}
\tightlist
\item
  received artifact integrity,
\item
  target-policy conformance,
\item
  generated-model consistency,
\item
  external SYSIDE compatibility,
\item
  validation findings,
\item
  validation identity,
\item
  publication gate.
\end{itemize}

\hypertarget{phase-l-owns}{%
\paragraph{Phase L owns}\label{phase-l-owns}}

\begin{itemize}
\tightlist
\item
  version allocation,
\item
  \texttt{data/output/} persistence,
\item
  published package manifest,
\item
  persisted validation evidence,
\item
  immutable output references,
\item
  downloadable/inspectable packaging.
\end{itemize}

\begin{center}\rule{0.5\linewidth}{0.5pt}\end{center}

\hypertarget{explicit-non-goals}{%
\subsubsection{Explicit Non-Goals}\label{explicit-non-goals}}

Phase K shall not:

\begin{itemize}
\tightlist
\item
  reload Candidates to reinterpret them,
\item
  change Human Review authority,
\item
  mutate the IEM,
\item
  regenerate SysML v2,
\item
  repair generated text,
\item
  run formatting that replaces generated content,
\item
  invent missing model semantics,
\item
  publish files,
\item
  create implicit latest-artifact selection,
\item
  silently skip external validation,
\item
  treat validator unavailability as success,
\item
  load ambient unrelated model context,
\item
  replace SYSIDE with an ad-hoc regex implementation of the full SysML grammar.
\end{itemize}

\begin{center}\rule{0.5\linewidth}{0.5pt}\end{center}

\hypertarget{consequences-7}{%
\subsubsection{Consequences}\label{consequences-7}}

\hypertarget{positive-1}{%
\paragraph{Positive}\label{positive-1}}

\begin{itemize}
\tightlist
\item
  Phase J, K and L remain cleanly separated.
\item
  Validation is reproducible and fail-closed.
\item
  External validator unavailability cannot accidentally authorize publication.
\item
  Turing-specific invariants remain deterministic.
\item
  Complete generated models receive actual SYSIDE compatibility validation.
\item
  Validation findings remain exactly traceable to generated engineering evidence.
\item
  Warning policy remains explicit rather than tool-dependent.
\item
  The K→L gate is cryptographically bound to the exact validated artifact set.
\item
  No semantic Human Review authority is reopened during validation.
\item
  A future alternative external validator can be introduced behind the same adapter contract.
\end{itemize}

\hypertarget{trade-offs-4}{%
\paragraph{Trade-offs}\label{trade-offs-4}}

\begin{itemize}
\tightlist
\item
  Phase K gains its own versioned Validation Profile.
\item
  External validation depends on available SYSIDE infrastructure.
\item
  A valid internal result is insufficient for publication when required external validation cannot run.
\item
  Validator versions become part of reproducibility evidence.
\item
  Diagnostic normalization requires an explicit adapter layer.
\end{itemize}

\begin{center}\rule{0.5\linewidth}{0.5pt}\end{center}

\hypertarget{acceptance-2}{%
\subsubsection{Acceptance}\label{acceptance-2}}

The accepted key decisions are:

\begin{enumerate}
\def\labelenumi{\arabic{enumi}.}
\tightlist
\item
  \texttt{GeneratedSysMLArtifactSet} is the sole normative Phase-K input.
\item
  Phase K combines deterministic internal checks with external SYSIDE validation.
\item
  Phase K introduces a separate versioned Validation Profile.
\item
  SYSIDE CLI is the MVP external-validation boundary.
\item
  Validation has \texttt{valid}, \texttt{invalid} and \texttt{incomplete} overall states.
\item
  Required-validator unavailability means \texttt{incomplete} and blocks publication.
\item
  Warnings remain visible but are not automatically publication-blocking.
\item
  Phase K does not reinterpret source-IEM or Candidate semantics.
\item
  Phase-L publication authority is fingerprint-bound to the exact successfully validated artifact set.
\item
  No separate Phase-K persistence layer is required for the MVP.
\end{enumerate}

\textbackslash newpage

\hypertarget{adr-023--final-model-review-and-output-publication-architecture}{%
\subsection{ADR-023 --- Final Model Review and Output Publication Architecture}\label{adr-023--final-model-review-and-output-publication-architecture}}

\hypertarget{status-5}{%
\subsubsection{Status}\label{status-5}}

Accepted

\hypertarget{date-5}{%
\subsubsection{Date}\label{date-5}}

2026-08-14

\hypertarget{context-8}{%
\subsubsection{Context}\label{context-8}}

Phase J and Phase K are completed and provide an explicit generated-model and validation boundary:

\begin{verbatim}
Internal Engineering Model
→ Phase J — deterministic SysML v2 generation
→ GeneratedSysMLArtifactSet
→ Phase K — deterministic + external validation
→ SysMLValidationResult
\end{verbatim}

Phase J produces concrete \texttt{.sysml} units together with deterministic generation identity, fingerprints and machine-readable traceability.

Phase K validates the exact generated artifact without modifying it. It produces an immutable \texttt{SysMLValidationResult} with the normative states:

\begin{verbatim}
valid
invalid
incomplete
\end{verbatim}

and the publication gate:

\begin{verbatim}
passed
blocked
\end{verbatim}

ADR-022 already defines the exact fingerprint-bound K→L publication condition and Phase K implements:

\begin{Shaded}
\begin{Highlighting}[]
\NormalTok{validate\_phase\_l\_handoff(}
\NormalTok{    artifact\_set: GeneratedSysMLArtifactSet,}
\NormalTok{    validation\_result: SysMLValidationResult,}
\NormalTok{) }\OperatorTok{{-}\textgreater{}} \VariableTok{None}
\end{Highlighting}
\end{Shaded}

That gate correctly prevents publication unless the validation result:

\begin{itemize}
\tightlist
\item
  belongs to the same Project,
\item
  belongs to the same source IEM,
\item
  covers the exact generated artifact fingerprint,
\item
  has \texttt{validation\_status\ ==\ "valid"},
\item
  and has \texttt{publication\_gate\ ==\ "passed"}.
\end{itemize}

However, automated validation alone does not constitute final engineering release authority.

The generated SysML v2 is the first point in the workflow where a human can inspect the complete resulting model representation as:

\begin{itemize}
\tightlist
\item
  generated SysML v2 code,
\item
  structural/model diagrams,
\item
  generated relationships,
\item
  validation findings,
\item
  traceability,
\item
  upstream model-candidate evidence,
\item
  generation rationales,
\item
  and available agent/personality proposals.
\end{itemize}

A final Human-in-the-Loop review is therefore required before generated output becomes an authoritative published project output.

The user shall be able to:

\begin{itemize}
\tightlist
\item
  inspect the generated model directly in the UI,
\item
  inspect the exact generated \texttt{.sysml} code,
\item
  inspect a diagrammatic model projection,
\item
  navigate between model elements and generated code,
\item
  inspect validation findings and their locations,
\item
  inspect upstream agent/personality proposals and rationales where available,
\item
  accept or reject the generated result,
\item
  request changes,
\item
  propose changes through code or diagram interaction,
\item
  provide structured feedback,
\item
  optionally request another bounded LLM/agent proposal loop,
\item
  and explicitly approve one exact reviewed revision for final publication.
\end{itemize}

Generated-but-not-approved model results must therefore be persisted as project-local review evidence before final publication.

Only after explicit Human approval may an output become a final versioned project result under \texttt{data/output/}.

This ADR refines the original Phase-L concept from a pure "Output Writer" into a controlled:

\begin{verbatim}
Final Model Review
+
Human Release Gate
+
Output Publication
\end{verbatim}

architecture.

\begin{center}\rule{0.5\linewidth}{0.5pt}\end{center}

\hypertarget{decision-8}{%
\subsubsection{Decision}\label{decision-8}}

\hypertarget{l-01--phase-l-responsibility}{%
\paragraph{L-01 --- Phase-L responsibility}\label{l-01--phase-l-responsibility}}

Phase L owns the final model review and controlled publication boundary.

The accepted end-to-end flow is:

\begin{verbatim}
Approved Input
→ Model Candidates
→ Candidate Human Review
→ Internal Engineering Model
→ deterministic SysML v2 generation
→ automated validation
→ Final Model Human Review
→ explicit Human release approval
→ immutable versioned publication
\end{verbatim}

Phase L shall not replace or weaken the earlier Human Review gate in Phase H.

The two review gates have different purposes:

\begin{verbatim}
Phase H
candidate-level engineering approval
before authoritative Internal Engineering Model assembly

Phase L
complete generated-model review and release approval
before final project publication
\end{verbatim}

Both are required.

\begin{center}\rule{0.5\linewidth}{0.5pt}\end{center}

\hypertarget{l-02--review-entry-is-distinct-from-publication-eligibility}{%
\paragraph{L-02 --- Review entry is distinct from publication eligibility}\label{l-02--review-entry-is-distinct-from-publication-eligibility}}

A Phase-L review workspace may be created for an exact:

\begin{verbatim}
GeneratedSysMLArtifactSet
+
SysMLValidationResult
\end{verbatim}

even when the K result is:

\begin{verbatim}
invalid / blocked
\end{verbatim}

or:

\begin{verbatim}
incomplete / blocked
\end{verbatim}

This is intentional because validation findings are useful Human Review evidence and may drive a correction loop.

Review entry shall still require exact integrity and binding:

\begin{itemize}
\tightlist
\item
  artifact-set integrity,
\item
  validation-result integrity,
\item
  same Project,
\item
  same source IEM,
\item
  exact artifact-set fingerprint binding.
\end{itemize}

The existing:

\begin{Shaded}
\begin{Highlighting}[]
\NormalTok{validate\_phase\_l\_handoff(...)}
\end{Highlighting}
\end{Shaded}

remains the normative \textbf{final publication gate}.

It shall not be weakened merely to allow review of invalid or incomplete results.

A separate Phase-L review-subject boundary may validate exact subject binding without requiring \texttt{valid\ +\ passed}.

\begin{center}\rule{0.5\linewidth}{0.5pt}\end{center}

\hypertarget{l-03--generated-but-not-approved-output-is-review-evidence-not-published-output}{%
\paragraph{L-03 --- Generated-but-not-approved output is review evidence, not published output}\label{l-03--generated-but-not-approved-output-is-review-evidence-not-published-output}}

Generated \texttt{.sysml} files that have not received final Human release approval shall not be treated as final published model output.

They shall be persisted in a project-local Final Model Review workspace.

Conceptually:

\begin{verbatim}
data/projects/<project_id>/
└── final_model_reviews/
    └── FMR-000001/
        ├── manifest.json
        ├── revisions/
        │   ├── FRV-000001/
        │   └── FRV-000002/
        └── decisions/
\end{verbatim}

These artifacts are:

\begin{itemize}
\tightlist
\item
  project-bound,
\item
  inspectable,
\item
  immutable once a review revision is created,
\item
  traceable,
\item
  and non-authoritative for final publication.
\end{itemize}

The legacy global principle:

\begin{verbatim}
generated SysML v2 output → data/output/
\end{verbatim}

is refined as follows:

\begin{verbatim}
generated SysML under review
→ project-local Final Model Review workspace

Human-approved published SysML
→ data/output/<project_id>/OUT-xxxxxx/
\end{verbatim}

This refinement preserves \texttt{data/output/} as the location of final generated output while preventing unreleased drafts from being mistaken for published project results.

\begin{center}\rule{0.5\linewidth}{0.5pt}\end{center}

\hypertarget{l-04--stable-project-local-final-model-review-identities}{%
\paragraph{L-04 --- Stable project-local Final Model Review identities}\label{l-04--stable-project-local-final-model-review-identities}}

Phase L shall introduce explicit project-local identities using the established six-digit identifier policy.

Initial identity families are:

\begin{verbatim}
FMR-000001  Final Model Review
FRV-000001  Final Model Review Revision
FRI-000001  Final Model Review Item
FRD-000001  Final Model Review Decision
OUT-000001  Published Output Package
\end{verbatim}

Identifiers are:

\begin{itemize}
\tightlist
\item
  project-local,
\item
  monotonically increasing,
\item
  immutable,
\item
  never gap-reused after publication/persistence,
\item
  and independent from display names.
\end{itemize}

Exact identifier modules and allocation mechanics belong to implementation but shall follow the existing repository-wide deterministic identifier pattern.

\begin{center}\rule{0.5\linewidth}{0.5pt}\end{center}

\hypertarget{l-05--immutable-review-revisions}{%
\paragraph{L-05 --- Immutable review revisions}\label{l-05--immutable-review-revisions}}

A Final Model Review is a long-lived review container.

Each concrete review subject is an immutable Final Model Review Revision.

One revision binds exactly:

\begin{itemize}
\tightlist
\item
  Project ID,
\item
  source IEM ID,
\item
  \texttt{GeneratedSysMLArtifactSet.content\_fingerprint},
\item
  \texttt{SysMLValidationResult.content\_fingerprint},
\item
  generated units and their fingerprints,
\item
  generation-policy references,
\item
  Validation Profile reference,
\item
  validation status and publication gate,
\item
  traceability references,
\item
  applicable upstream Candidate/Review evidence references,
\item
  applicable agent/personality proposal evidence,
\item
  and the deterministic review-view fingerprint.
\end{itemize}

A revision shall never be edited in place.

If generated content or validation evidence changes, a new revision is created.

Conceptually:

\begin{verbatim}
FMR-000001
├── FRV-000001  generated + reviewed, changes requested
├── FRV-000002  regenerated + revalidated, changes requested
└── FRV-000003  regenerated + revalidated, approved
\end{verbatim}

Earlier revisions remain immutable evidence.

\begin{center}\rule{0.5\linewidth}{0.5pt}\end{center}

\hypertarget{l-06--final-model-review-ui-projection}{%
\paragraph{L-06 --- Final Model Review UI projection}\label{l-06--final-model-review-ui-projection}}

Phase L shall provide a deterministic non-authoritative review projection for the focused UI.

Conceptually:

\begin{Shaded}
\begin{Highlighting}[]
\NormalTok{FinalModelReviewView}
\end{Highlighting}
\end{Shaded}

It shall present at least:

\begin{itemize}
\tightlist
\item
  review identity and revision,
\item
  current review state,
\item
  model summary,
\item
  diagrammatic structural view,
\item
  exact generated SysML v2 code,
\item
  generated-unit selection where multiple units exist,
\item
  generated symbols,
\item
  relationship views,
\item
  validation status,
\item
  blocking findings,
\item
  warnings,
\item
  traceability,
\item
  upstream Candidate decisions,
\item
  generation rationale,
\item
  available agent/personality proposals,
\item
  unresolved review items,
\item
  required Human decisions,
\item
  and the next action.
\end{itemize}

The view is a projection over immutable evidence.

It shall not itself become engineering or publication authority.

\begin{center}\rule{0.5\linewidth}{0.5pt}\end{center}

\hypertarget{l-07--diagram-and-sysml-code-are-linked-review-surfaces}{%
\paragraph{L-07 --- Diagram and SysML code are linked review surfaces}\label{l-07--diagram-and-sysml-code-are-linked-review-surfaces}}

The UI shall support both:

\begin{verbatim}
diagram / model view
\end{verbatim}

and:

\begin{verbatim}
exact generated `.sysml` code
\end{verbatim}

as first-class review surfaces.

Where possible, selection shall be cross-linked through existing generated identity and traceability:

\begin{verbatim}
diagram element / relationship
→ generated_symbol_id
→ GeneratedSysMLTraceabilityEntry
→ generated unit + line location
→ exact SysML code
\end{verbatim}

and conversely:

\begin{verbatim}
generated code location
→ generated symbol
→ model element / relationship
→ upstream traceability
\end{verbatim}

The UI shall not reconstruct engineering authority by reverse-parsing the generated SysML.

Diagrammatic views shall preferentially use the authoritative upstream model structure and existing traceability while the code view displays the exact Phase-J output.

\begin{center}\rule{0.5\linewidth}{0.5pt}\end{center}

\hypertarget{l-08--direct-ui-edits-create-change-proposals-not-silent-authoritative-mutations}{%
\paragraph{L-08 --- Direct UI edits create change proposals, not silent authoritative mutations}\label{l-08--direct-ui-edits-create-change-proposals-not-silent-authoritative-mutations}}

The user may interactively modify or annotate:

\begin{itemize}
\tightlist
\item
  generated SysML code,
\item
  diagram relationships,
\item
  model structure,
\item
  attributes represented in the review UI,
\item
  and review comments.
\end{itemize}

However, an edit shall not silently mutate:

\begin{itemize}
\tightlist
\item
  the source IEM,
\item
  the existing \texttt{GeneratedSysMLArtifactSet},
\item
  the existing \texttt{SysMLValidationResult},
\item
  or a published output.
\end{itemize}

A UI edit shall create explicit review evidence such as a change proposal or structured change request bound to the current review revision.

For code editing, the UI may provide an editor experience conceptually similar to:

\begin{verbatim}
edit generated_model.sysml
→ inspect diff
→ save as change proposal
\end{verbatim}

The saved proposal shall preserve at least:

\begin{itemize}
\tightlist
\item
  exact base artifact fingerprint,
\item
  affected unit,
\item
  affected generated symbols where resolvable,
\item
  original content or location reference,
\item
  proposed content/diff,
\item
  reviewer identity,
\item
  reviewer rationale,
\item
  and change classification.
\end{itemize}

A proposed edit is not publishable generated output.

\begin{center}\rule{0.5\linewidth}{0.5pt}\end{center}

\hypertarget{l-09--change-requests-are-routed-to-the-owning-authority-boundary}{%
\paragraph{L-09 --- Change requests are routed to the owning authority boundary}\label{l-09--change-requests-are-routed-to-the-owning-authority-boundary}}

Phase L shall not use manual final-code editing to bypass the Internal Engineering Model.

A requested change shall be classified by its likely resolution boundary.

At minimum:

\begin{verbatim}
engineering_semantics
→ Phase H / Candidate Review / IEM revision path

generated_representation
→ Phase J generation policy / generator correction

validation_policy_or_tool
→ Phase K validation policy / validator correction

review_presentation_only
→ Phase-L UI/read-model correction without changing model authority
\end{verbatim}

Examples:

\begin{verbatim}
"Function A belongs to another Logical Component."
→ engineering_semantics

"The correct accepted relationship should be satisfaction, not allocation."
→ engineering_semantics

"The model is correct but the generator emitted the wrong target syntax."
→ generated_representation

"SYSIDE rejects a construct that the Turing validator considered valid."
→ generation / Target Notation / validation investigation

"The diagram label is misleading but the underlying model is correct."
→ review_presentation_only
\end{verbatim}

The correction shall occur at the owning layer and then re-enter the downstream pipeline.

Phase L shall not invent a hidden third engineering authority in edited SysML text.

\begin{center}\rule{0.5\linewidth}{0.5pt}\end{center}

\hypertarget{l-10--optional-bounded-llmagent-revision-loop}{%
\paragraph{L-10 --- Optional bounded LLM/agent revision loop}\label{l-10--optional-bounded-llmagent-revision-loop}}

A Human reviewer may request another agent/LLM-assisted model proposal cycle from Final Model Review.

The request shall be explicit and shall preserve the Human feedback that triggered it.

Conceptually:

\begin{verbatim}
Final Model Review
→ Human change request / feedback
→ optional bounded agent/LLM re-proposal
→ proposed Candidate changes / alternatives
→ existing Candidate Human Review
→ new authoritative IEM
→ deterministic J generation
→ K validation
→ new Final Model Review Revision
\end{verbatim}

LLM or agent execution may:

\begin{itemize}
\tightlist
\item
  propose alternative model elements,
\item
  propose alternative relationships,
\item
  explain tradeoffs,
\item
  identify gaps,
\item
  compare candidate structures,
\item
  and respond to reviewer feedback.
\end{itemize}

It shall not:

\begin{itemize}
\tightlist
\item
  directly rewrite final published SysML,
\item
  approve its own proposals,
\item
  bypass Candidate Human Review,
\item
  bypass IEM assembly,
\item
  bypass deterministic Phase J,
\item
  bypass Phase K,
\item
  or create final release authority.
\end{itemize}

The normative chain remains:

\begin{verbatim}
agent/LLM proposal
→ Human decision
→ authoritative model state
→ deterministic generation
→ validation
→ final Human release approval
\end{verbatim}

\begin{center}\rule{0.5\linewidth}{0.5pt}\end{center}

\hypertarget{l-11--agentpersonality-evidence-is-visible-but-non-authoritative}{%
\paragraph{L-11 --- Agent/personality evidence is visible but non-authoritative}\label{l-11--agentpersonality-evidence-is-visible-but-non-authoritative}}

Final Model Review shall surface available proposal evidence from different agent roles or modeling personalities where such evidence exists.

The review projection should preserve, where available:

\begin{itemize}
\tightlist
\item
  agent/personality identity,
\item
  proposal,
\item
  rationale,
\item
  confidence/support information,
\item
  alternatives,
\item
  semantic uncertainty,
\item
  source evidence,
\item
  generation provenance,
\item
  and prior Human decision.
\end{itemize}

The UI may present multiple perspectives side by side.

Example:

\begin{verbatim}
Current accepted model relation:
allocated_to

Systems Engineer perspective:
allocated_to
rationale: ...

Architecture Reviewer perspective:
dependency
rationale: ...

Human options:
keep current
request change
request another proposal cycle
inspect upstream evidence
\end{verbatim}

Selecting or favoring an alternative does not directly modify the final model.

It creates or contributes to a Human-reviewed upstream change request.

Agent agreement shall remain review evidence, not approval authority.

\begin{center}\rule{0.5\linewidth}{0.5pt}\end{center}

\hypertarget{l-12--every-material-revision-must-be-regenerated-and-revalidated}{%
\paragraph{L-12 --- Every material revision must be regenerated and revalidated}\label{l-12--every-material-revision-must-be-regenerated-and-revalidated}}

Any material change affecting engineering content or generated SysML requires a new downstream artifact chain.

The required loop is:

\begin{verbatim}
change accepted at owning authority boundary
→ new/updated authoritative model state
→ deterministic Phase-J generation
→ new GeneratedSysMLArtifactSet fingerprint
→ Phase-K validation
→ new SysMLValidationResult fingerprint
→ new Final Model Review Revision
\end{verbatim}

A prior K result never authorizes changed content.

A prior Final Model Review approval never authorizes changed content.

No "validation carry-over" is permitted across a changed artifact fingerprint.

\begin{center}\rule{0.5\linewidth}{0.5pt}\end{center}

\hypertarget{l-13--final-human-release-approval-is-mandatory}{%
\paragraph{L-13 --- Final Human release approval is mandatory}\label{l-13--final-human-release-approval-is-mandatory}}

Final publication requires an explicit immutable Human decision bound to one exact Final Model Review Revision.

Conceptually:

\begin{Shaded}
\begin{Highlighting}[]
\NormalTok{FinalModelReviewDecision(}
\NormalTok{    final\_model\_review\_id}\OperatorTok{=}\NormalTok{...,}
\NormalTok{    final\_model\_review\_revision\_id}\OperatorTok{=}\NormalTok{...,}
\NormalTok{    decision}\OperatorTok{=}\StringTok{"approved\_for\_publication"}\NormalTok{,}
\NormalTok{    generated\_artifact\_set\_fingerprint}\OperatorTok{=}\NormalTok{...,}
\NormalTok{    validation\_result\_fingerprint}\OperatorTok{=}\NormalTok{...,}
\NormalTok{    reviewer\_identity}\OperatorTok{=}\NormalTok{...,}
\NormalTok{    rationale}\OperatorTok{=}\NormalTok{...,}
\NormalTok{    content\_fingerprint}\OperatorTok{=}\NormalTok{...,}
\NormalTok{)}
\end{Highlighting}
\end{Shaded}

Final approval shall be possible only if:

\begin{itemize}
\tightlist
\item
  the review revision is internally complete,
\item
  the exact artifact-set integrity is valid,
\item
  the exact validation-result integrity is valid,
\item
  \texttt{validation\_status\ ==\ "valid"},
\item
  \texttt{publication\_gate\ ==\ "passed"},
\item
  the validation result covers the exact artifact-set fingerprint,
\item
  no mandatory review item remains unresolved,
\item
  no accepted change request is waiting for regeneration,
\item
  and the reviewer explicitly approves the revision for publication.
\end{itemize}

Final release approval is a Human authority action.

It shall not be inferred from:

\begin{itemize}
\tightlist
\item
  absence of findings,
\item
  agent consensus,
\item
  LLM confidence,
\item
  SYSIDE success alone,
\item
  UI navigation state,
\item
  or elapsed time.
\end{itemize}

\begin{center}\rule{0.5\linewidth}{0.5pt}\end{center}

\hypertarget{l-14--final-approval-becomes-stale-on-any-subject-change}{%
\paragraph{L-14 --- Final approval becomes stale on any subject change}\label{l-14--final-approval-becomes-stale-on-any-subject-change}}

A Final Model Review Decision is valid only for the exact fingerprints it authorizes.

If any of the following changes:

\begin{itemize}
\tightlist
\item
  generated unit content,
\item
  artifact-set fingerprint,
\item
  validation-result fingerprint,
\item
  required validation evidence,
\item
  review revision subject,
\item
  or an approval-relevant policy reference,
\end{itemize}

the previous approval shall not authorize publication of the changed result.

No post-approval content mutation is permitted.

\begin{center}\rule{0.5\linewidth}{0.5pt}\end{center}

\hypertarget{l-15--final-publication-service-boundary}{%
\paragraph{L-15 --- Final publication service boundary}\label{l-15--final-publication-service-boundary}}

Only after successful automated validation and explicit Human release approval may publication occur.

The Phase-L publication service boundary becomes conceptually:

\begin{Shaded}
\begin{Highlighting}[]
\NormalTok{OutputWriter.publish(}
\NormalTok{    artifact\_set: GeneratedSysMLArtifactSet,}
\NormalTok{    validation\_result: SysMLValidationResult,}
\NormalTok{    final\_review\_decision: FinalModelReviewDecision,}
\NormalTok{) }\OperatorTok{{-}\textgreater{}}\NormalTok{ PublishedOutputPackage}
\end{Highlighting}
\end{Shaded}

Before writing any final output, the service shall:

\begin{enumerate}
\def\labelenumi{\arabic{enumi}.}
\tightlist
\item
  validate artifact-set integrity,
\item
  validate validation-result integrity,
\item
  execute the existing \texttt{validate\_phase\_l\_handoff(...)},
\item
  validate Final Model Review Decision integrity,
\item
  confirm exact Project binding,
\item
  confirm exact review-revision binding,
\item
  confirm exact artifact-set fingerprint binding,
\item
  confirm exact validation-result fingerprint binding,
\item
  require \texttt{approved\_for\_publication},
\item
  and reject stale or mismatched Human approval.
\end{enumerate}

Phase L shall never implicitly select:

\begin{itemize}
\tightlist
\item
  latest artifact,
\item
  latest validation result,
\item
  latest review revision,
\item
  or latest Human decision.
\end{itemize}

All publication inputs are explicit.

\begin{center}\rule{0.5\linewidth}{0.5pt}\end{center}

\hypertarget{l-16--final-output-location}{%
\paragraph{L-16 --- Final output location}\label{l-16--final-output-location}}

Final Human-approved output shall be published under:

\begin{verbatim}
data/output/<project_id>/<output_package_id>/
\end{verbatim}

Example:

\begin{verbatim}
data/output/
└── 000001/
    └── OUT-000001/
        ├── manifest.json
        ├── generated_model.sysml
        ├── generation_summary.json
        ├── validation_result.json
        ├── validation_report.md
        └── traceability.json
\end{verbatim}

If Phase J later produces multiple generated units, their configured relative paths shall be preserved within the output package.

Only final approved packages belong in this publication namespace.

\begin{center}\rule{0.5\linewidth}{0.5pt}\end{center}

\hypertarget{l-17--published-output-identity-and-version-model}{%
\paragraph{L-17 --- Published output identity and version model}\label{l-17--published-output-identity-and-version-model}}

Final outputs use project-local immutable sequential IDs:

\begin{verbatim}
OUT-000001
OUT-000002
OUT-000003
\end{verbatim}

The OUT sequence represents publication identity and ordering.

It is not semantic versioning.

A later OUT ID does not inherently mean a semantic major/minor/patch change.

Engineering and validation identity remains fingerprint-based.

OUT IDs shall be:

\begin{itemize}
\tightlist
\item
  project-local,
\item
  sequential,
\item
  immutable,
\item
  never reused,
\item
  and stable after publication.
\end{itemize}

\begin{center}\rule{0.5\linewidth}{0.5pt}\end{center}

\hypertarget{l-18--versioned-output-publication-profile}{%
\paragraph{L-18 --- Versioned Output Publication Profile}\label{l-18--versioned-output-publication-profile}}

Phase L shall introduce a small versioned publication policy artifact:

\begin{verbatim}
context/sysml/turing_sysml_v2_output_profile.json
\end{verbatim}

Initial identity:

\begin{verbatim}
profile_id:      TURING_SYSML_V2_OUTPUT
profile_version: 1.0.0
\end{verbatim}

The Output Profile shall define publication concerns only, including:

\begin{itemize}
\tightlist
\item
  output root,
\item
  required package file roles,
\item
  unit placement rules,
\item
  manifest schema expectations,
\item
  fingerprint policy,
\item
  idempotence policy,
\item
  and derived archive/download policy.
\end{itemize}

It shall not redefine:

\begin{itemize}
\tightlist
\item
  engineering semantics,
\item
  Candidate Review rules,
\item
  IEM assembly,
\item
  Target Notation,
\item
  Generation Profile mappings,
\item
  validation semantics,
\item
  or Human release authority.
\end{itemize}

Human approval remains an architectural gate, not a configurable bypass.

\begin{center}\rule{0.5\linewidth}{0.5pt}\end{center}

\hypertarget{l-19--authoritative-published-package-contents}{%
\paragraph{L-19 --- Authoritative published package contents}\label{l-19--authoritative-published-package-contents}}

The MVP authoritative output package shall contain at least:

\begin{verbatim}
manifest.json
<all GeneratedSysMLUnit relative paths>
generation_summary.json
validation_result.json
validation_report.md
traceability.json
\end{verbatim}

\hypertarget{generated-sysml-units}{%
\subparagraph{Generated SysML units}\label{generated-sysml-units}}

Published \texttt{.sysml} units shall be byte-for-byte identical to the exact Human-approved \texttt{GeneratedSysMLArtifactSet}.

Phase L shall not:

\begin{itemize}
\tightlist
\item
  reformat,
\item
  rewrite,
\item
  normalize,
\item
  reorder,
\item
  rename,
\item
  repair,
\item
  or regenerate
\end{itemize}

the SysML content during publication.

\hypertarget{validation_resultjson}{%
\subparagraph{validation\_result.json}\label{validation_resultjson}}

This is the machine-readable authoritative Phase-K validation evidence.

\hypertarget{validation_reportmd}{%
\subparagraph{validation\_report.md}\label{validation_reportmd}}

This is a deterministic Human-readable projection of the validation result.

It is not independent validation authority.

\hypertarget{traceabilityjson}{%
\subparagraph{traceability.json}\label{traceabilityjson}}

This preserves the generated-output traceability needed to navigate from final output back through:

\begin{verbatim}
generated symbol / location
→ IEM element or relationship
→ Model Candidate
→ Approved Input
→ Human Review Decision
→ accepted exception where applicable
\end{verbatim}

\hypertarget{generation_summaryjson}{%
\subparagraph{generation\_summary.json}\label{generation_summaryjson}}

This records deterministic generation context and provenance required to understand the published artifact, without creating new engineering semantics.

\hypertarget{manifestjson}{%
\subparagraph{manifest.json}\label{manifestjson}}

This is the authoritative package index and integrity boundary.

\begin{center}\rule{0.5\linewidth}{0.5pt}\end{center}

\hypertarget{l-20--published-output-manifest}{%
\paragraph{L-20 --- Published output manifest}\label{l-20--published-output-manifest}}

Conceptually:

\begin{Shaded}
\begin{Highlighting}[]
\NormalTok{PublishedOutputManifest(}
\NormalTok{    schema\_version}\OperatorTok{=}\NormalTok{...,}
\NormalTok{    project\_id}\OperatorTok{=}\NormalTok{...,}
\NormalTok{    output\_package\_id}\OperatorTok{=}\NormalTok{...,}
\NormalTok{    source\_internal\_engineering\_model\_id}\OperatorTok{=}\NormalTok{...,}
\NormalTok{    source\_artifact\_set\_fingerprint}\OperatorTok{=}\NormalTok{...,}
\NormalTok{    validation\_result\_fingerprint}\OperatorTok{=}\NormalTok{...,}
\NormalTok{    final\_model\_review\_id}\OperatorTok{=}\NormalTok{...,}
\NormalTok{    final\_model\_review\_revision\_id}\OperatorTok{=}\NormalTok{...,}
\NormalTok{    final\_review\_decision\_id}\OperatorTok{=}\NormalTok{...,}
\NormalTok{    final\_review\_decision\_fingerprint}\OperatorTok{=}\NormalTok{...,}
\NormalTok{    output\_profile\_reference}\OperatorTok{=}\NormalTok{...,}
\NormalTok{    publication\_input\_fingerprint}\OperatorTok{=}\NormalTok{...,}
\NormalTok{    files}\OperatorTok{=}\NormalTok{(...),}
\NormalTok{    content\_fingerprint}\OperatorTok{=}\NormalTok{...,}
\NormalTok{)}
\end{Highlighting}
\end{Shaded}

Every persisted package file shall have:

\begin{itemize}
\tightlist
\item
  relative path,
\item
  controlled file role,
\item
  content fingerprint,
\item
  and source/generated reference where applicable.
\end{itemize}

The complete output package shall have one deterministic content fingerprint.

\begin{center}\rule{0.5\linewidth}{0.5pt}\end{center}

\hypertarget{l-21--publication-input-fingerprint-and-idempotence}{%
\paragraph{L-21 --- Publication input fingerprint and idempotence}\label{l-21--publication-input-fingerprint-and-idempotence}}

Phase L shall calculate a deterministic publication-input fingerprint covering at least:

\begin{itemize}
\tightlist
\item
  exact \texttt{GeneratedSysMLArtifactSet.content\_fingerprint},
\item
  exact \texttt{SysMLValidationResult.content\_fingerprint},
\item
  exact Final Model Review Decision fingerprint,
\item
  exact Final Model Review Revision fingerprint,
\item
  exact Output Profile reference and fingerprint.
\end{itemize}

Conceptually:

\begin{verbatim}
publication_input_fingerprint =
SHA256(
    artifact-set identity
  + validation identity
  + Human release identity
  + publication policy identity
)
\end{verbatim}

Publishing the exact same authorized input repeatedly shall return the same existing published output package rather than allocate new OUT IDs.

Example:

\begin{verbatim}
same exact authorized publication input
→ OUT-000001

same exact authorized publication input again
→ OUT-000001
\end{verbatim}

A materially different approved publication input receives a new OUT ID.

\begin{center}\rule{0.5\linewidth}{0.5pt}\end{center}

\hypertarget{l-22--atomic-immutable-publication}{%
\paragraph{L-22 --- Atomic immutable publication}\label{l-22--atomic-immutable-publication}}

Final publication shall be fail-closed and atomic.

Conceptually:

\begin{verbatim}
allocate OUT-000001
→ create hidden temporary package
→ write all files
→ verify file fingerprints
→ verify manifest
→ verify complete package integrity
→ atomic rename
→ OUT-000001 becomes visible
\end{verbatim}

A visible final output directory shall always represent a complete validated publication.

Interrupted or incomplete temporary publication state shall never be treated as a valid published output.

\begin{center}\rule{0.5\linewidth}{0.5pt}\end{center}

\hypertarget{l-23--recovery-and-integrity-scanning}{%
\paragraph{L-23 --- Recovery and integrity scanning}\label{l-23--recovery-and-integrity-scanning}}

The Output Repository shall detect at least:

\begin{itemize}
\tightlist
\item
  incomplete temporary publication directories,
\item
  malformed OUT identities,
\item
  missing manifest,
\item
  missing required files,
\item
  unexpected files where prohibited by profile,
\item
  file fingerprint mismatch,
\item
  manifest fingerprint mismatch,
\item
  unsafe paths,
\item
  symlink/path escape,
\item
  duplicate publication-input fingerprints,
\item
  and project/OUT identity mismatch.
\end{itemize}

Recovery diagnostics shall be explicit and fail-closed.

Interrupted publication evidence shall not be silently deleted during normal reads.

\begin{center}\rule{0.5\linewidth}{0.5pt}\end{center}

\hypertarget{l-24--directory-package-is-publication-authority-archive-is-derived}{%
\paragraph{L-24 --- Directory package is publication authority; archive is derived}\label{l-24--directory-package-is-publication-authority-archive-is-derived}}

The authoritative final output is the immutable directory package plus its manifest.

A ZIP or other archive may be produced for download convenience.

Such an archive is a derived transport representation.

It shall not replace the directory package as publication authority.

Archive metadata such as timestamps or compression differences shall not alter the identity of the authoritative publication.

\begin{center}\rule{0.5\linewidth}{0.5pt}\end{center}

\hypertarget{l-25--final-project-output-association-occurs-only-after-approval-and-publication}{%
\paragraph{L-25 --- Final project-output association occurs only after approval and publication}\label{l-25--final-project-output-association-occurs-only-after-approval-and-publication}}

Final Model Review artifacts are project-local review evidence but shall not be listed as released project outputs.

Only successfully published:

\begin{verbatim}
OUT-xxxxxx
\end{verbatim}

packages become final project outputs.

The Project Dashboard / Guided Workflow UI shall distinguish at least:

\begin{verbatim}
Generated / under review
Changes requested
Validation blocked
Ready for Human approval
Approved for publication
Published
\end{verbatim}

A generated model shall not appear as a final project model merely because J or K completed successfully.

\begin{center}\rule{0.5\linewidth}{0.5pt}\end{center}

\hypertarget{l-26--read-boundary-for-ui-and-downstream-presentation}{%
\paragraph{L-26 --- Read boundary for UI and downstream presentation}\label{l-26--read-boundary-for-ui-and-downstream-presentation}}

Publication persistence shall expose explicit read services conceptually like:

\begin{Shaded}
\begin{Highlighting}[]
\NormalTok{OutputRepository.list\_outputs(}
\NormalTok{    project\_id,}
\NormalTok{) }\OperatorTok{{-}\textgreater{}} \BuiltInTok{tuple}\NormalTok{[PublishedOutputManifest, ...]}
\end{Highlighting}
\end{Shaded}

and:

\begin{Shaded}
\begin{Highlighting}[]
\NormalTok{OutputRepository.load\_output(}
\NormalTok{    project\_id,}
\NormalTok{    output\_package\_id,}
\NormalTok{) }\OperatorTok{{-}\textgreater{}}\NormalTok{ PublishedOutputPackage}
\end{Highlighting}
\end{Shaded}

Final Model Review shall expose separate project-bound review read services.

The later Guided Workflow UI shall consume these read boundaries.

The UI shall not infer publication authority directly from filesystem presence or session state.

\begin{center}\rule{0.5\linewidth}{0.5pt}\end{center}

\hypertarget{l-27--syside-unavailability-blocks-approvalpublication-not-review}{%
\paragraph{L-27 --- SYSIDE unavailability blocks approval/publication, not review}\label{l-27--syside-unavailability-blocks-approvalpublication-not-review}}

The current verification workstation may not have the SYSIDE CLI available.

This does not prevent:

\begin{itemize}
\tightlist
\item
  Phase-L architecture,
\item
  review-workspace implementation,
\item
  review UI implementation,
\item
  review of generated code,
\item
  review of diagrams,
\item
  review of agent proposals,
\item
  change requests,
\item
  or deterministic unit testing.
\end{itemize}

It does prevent final Human release approval and final publication of a real model because the required K state cannot become:

\begin{verbatim}
valid / passed
\end{verbatim}

without required external validation.

An \texttt{incomplete\ /\ blocked} result may be reviewed, but shall not be approved for publication.

\begin{center}\rule{0.5\linewidth}{0.5pt}\end{center}

\hypertarget{l-28--human-authority-is-not-delegated-to-llms-or-agents}{%
\paragraph{L-28 --- Human authority is not delegated to LLMs or agents}\label{l-28--human-authority-is-not-delegated-to-llms-or-agents}}

No LLM or agent is normative final release authority.

Agents may support:

\begin{itemize}
\tightlist
\item
  proposal generation,
\item
  alternative generation,
\item
  completeness review,
\item
  relationship analysis,
\item
  rationale generation,
\item
  and response to Human feedback.
\end{itemize}

Only an explicit Human decision may authorize final publication.

The final gate is therefore:

\begin{verbatim}
automated validation PASS
+
explicit Human release approval
=
eligible for publication
\end{verbatim}

not:

\begin{verbatim}
agent consensus
=
publication
\end{verbatim}

\begin{center}\rule{0.5\linewidth}{0.5pt}\end{center}

\hypertarget{l-29--complete-revision-loop}{%
\paragraph{L-29 --- Complete revision loop}\label{l-29--complete-revision-loop}}

The normative Phase-L correction loop is:

\begin{verbatim}
J GeneratedSysMLArtifactSet
        ↓
K SysMLValidationResult
        ↓
Final Model Review Revision
        ↓
Human inspection:
  diagram
  code
  validation
  traceability
  agent/personality evidence
        ↓
        ├── approve
        │      ↓
        │   Human release decision
        │      ↓
        │   OutputWriter.publish(...)
        │      ↓
        │   OUT-xxxxxx
        │
        └── request changes
               ↓
            explicit change proposal
               ↓
            owning authority boundary
               ↓
        optional bounded LLM/agent proposal
               ↓
            Human review/acceptance
               ↓
          new authoritative model state
               ↓
        deterministic regeneration
               ↓
             validation
               ↓
       new Final Model Review Revision
\end{verbatim}

Every loop iteration remains traceable.

\begin{center}\rule{0.5\linewidth}{0.5pt}\end{center}

\hypertarget{l-30--non-goals-and-prohibited-shortcuts}{%
\paragraph{L-30 --- Non-goals and prohibited shortcuts}\label{l-30--non-goals-and-prohibited-shortcuts}}

Phase L shall not:

\begin{itemize}
\tightlist
\item
  treat generated SysML as final before Human release approval,
\item
  mutate a prior immutable review revision,
\item
  silently change validated generated content,
\item
  publish \texttt{invalid} or \texttt{incomplete} results,
\item
  publish without exact K fingerprint binding,
\item
  publish without exact Human approval binding,
\item
  use manually edited final SysML as a second engineering authority,
\item
  let an LLM directly rewrite and publish final output,
\item
  bypass Candidate Human Review for semantic changes,
\item
  bypass IEM assembly,
\item
  bypass deterministic Phase-J generation,
\item
  carry old validation across changed content,
\item
  carry old Human approval across changed content,
\item
  infer publication from "latest" artifacts,
\item
  or collapse review evidence and final published output into one storage state.
\end{itemize}

\begin{center}\rule{0.5\linewidth}{0.5pt}\end{center}

\hypertarget{final-model-review-state-model}{%
\subsubsection{Final Model Review state model}\label{final-model-review-state-model}}

The initial controlled review lifecycle shall support at least:

\begin{verbatim}
generated
validation_blocked
review_pending
changes_requested
regeneration_required
ready_for_approval
approved_for_publication
published
\end{verbatim}

These are workflow/read-model states.

Normative authority remains in immutable artifacts and Human decisions rather than in a mutable status string alone.

A review is \texttt{ready\_for\_approval} only when the exact review revision has:

\begin{itemize}
\tightlist
\item
  complete subject integrity,
\item
  \texttt{valid\ /\ passed} K evidence,
\item
  no unresolved mandatory review items,
\item
  and no outstanding accepted change request.
\end{itemize}

\begin{center}\rule{0.5\linewidth}{0.5pt}\end{center}

\hypertarget{publication-package-authority-chain}{%
\subsubsection{Publication package authority chain}\label{publication-package-authority-chain}}

A final published package shall be able to prove:

\begin{verbatim}
OUT-xxxxxx
→ PublishedOutputManifest
→ exact Final Model Review Decision
→ exact Final Model Review Revision
→ exact SysMLValidationResult
→ exact GeneratedSysMLArtifactSet
→ exact IEM
→ reviewed Model Candidates
→ Approved Input
→ Human Review evidence
\end{verbatim}

This chain is required for auditability and for later Project Dashboard / Guided Workflow presentation.

\begin{center}\rule{0.5\linewidth}{0.5pt}\end{center}

\hypertarget{ui-interaction-principle}{%
\subsubsection{UI interaction principle}\label{ui-interaction-principle}}

The established interaction principle remains:

\begin{verbatim}
Simple by default.
Explainable on demand.
Fully traceable underneath.
\end{verbatim}

The default Final Model Review UI should emphasize:

\begin{enumerate}
\def\labelenumi{\arabic{enumi}.}
\tightlist
\item
  what model was generated,
\item
  whether automated validation passed,
\item
  what requires Human attention,
\item
  relevant alternative proposals,
\item
  the current diagram/model structure,
\item
  the exact SysML code on demand,
\item
  and the next Human action.
\end{enumerate}

Detailed IDs, fingerprints, policy references, evidence chains and diagnostics remain available through progressive disclosure.

The UI shall support direct review interaction without weakening the immutable authority model underneath.

\begin{center}\rule{0.5\linewidth}{0.5pt}\end{center}

\hypertarget{implementation-decomposition}{%
\subsubsection{Implementation decomposition}\label{implementation-decomposition}}

Phase L shall be implemented in the following controlled slices.

\hypertarget{l1--final-model-review-domain-foundation}{%
\paragraph{L1 --- Final Model Review domain foundation}\label{l1--final-model-review-domain-foundation}}

\begin{itemize}
\tightlist
\item
  identifiers,
\item
  immutable review/revision/item/decision types,
\item
  validation,
\item
  fingerprints,
\item
  lifecycle vocabulary.
\end{itemize}

\hypertarget{l2--final-model-review-repository}{%
\paragraph{L2 --- Final Model Review repository}\label{l2--final-model-review-repository}}

\begin{itemize}
\tightlist
\item
  project-local review workspace,
\item
  immutable revisions,
\item
  decision persistence,
\item
  safe paths,
\item
  integrity,
\item
  scanning,
\item
  recovery.
\end{itemize}

\hypertarget{l3--final-model-review-read-model-and-ui-projection}{%
\paragraph{L3 --- Final Model Review read model and UI projection}\label{l3--final-model-review-read-model-and-ui-projection}}

\begin{itemize}
\tightlist
\item
  deterministic \texttt{FinalModelReviewView},
\item
  diagram/model projection,
\item
  exact SysML code projection,
\item
  validation findings,
\item
  traceability,
\item
  upstream Candidate evidence,
\item
  agent/personality proposal evidence,
\item
  required decisions,
\item
  next action.
\end{itemize}

\hypertarget{l4--change-proposal-and-revision-loop}{%
\paragraph{L4 --- Change proposal and revision loop}\label{l4--change-proposal-and-revision-loop}}

\begin{itemize}
\tightlist
\item
  code-edit proposals,
\item
  diagram/model change proposals,
\item
  structured Human feedback,
\item
  change classification,
\item
  owning-boundary routing,
\item
  optional bounded LLM/agent re-proposal,
\item
  regeneration/revalidation handoff,
\item
  successor review revision creation.
\end{itemize}

\hypertarget{l5--final-human-release-gate}{%
\paragraph{L5 --- Final Human release gate}\label{l5--final-human-release-gate}}

\begin{itemize}
\tightlist
\item
  exact revision binding,
\item
  exact artifact fingerprint binding,
\item
  exact validation fingerprint binding,
\item
  unresolved-item checks,
\item
  \texttt{approved\_for\_publication},
\item
  stale approval rejection.
\end{itemize}

\hypertarget{l6--output-publication-repository-and-outputwriter}{%
\paragraph{L6 --- Output publication repository and OutputWriter}\label{l6--output-publication-repository-and-outputwriter}}

\begin{itemize}
\tightlist
\item
  Output Profile,
\item
  OUT identity,
\item
  deterministic package projections,
\item
  publication-input fingerprint,
\item
  idempotence,
\item
  atomic publication,
\item
  bundle integrity,
\item
  read service,
\item
  derived download/archive support.
\end{itemize}

\hypertarget{l7--end-to-end-integration-and-acceptance}{%
\paragraph{L7 --- End-to-end integration and acceptance}\label{l7--end-to-end-integration-and-acceptance}}

\begin{itemize}
\tightlist
\item
  J→K→Final Review→Publication integration,
\item
  real SYSIDE validation environment,
\item
  real \texttt{valid\ /\ passed} result,
\item
  Human Review acceptance,
\item
  real \texttt{OUT-000001} publication,
\item
  review/change loop evidence,
\item
  focused regression,
\item
  complete repository regression,
\item
  \texttt{git\ diff\ -\/-check},
\item
  manual acceptance,
\item
  SSOT closeout.
\end{itemize}

No implementation slice may weaken the Human Review, fingerprint, project isolation, deterministic generation, validation or publication boundaries defined by this ADR.

\begin{center}\rule{0.5\linewidth}{0.5pt}\end{center}

\hypertarget{consequences-8}{%
\subsubsection{Consequences}\label{consequences-8}}

\hypertarget{positive-2}{%
\paragraph{Positive}\label{positive-2}}

\begin{itemize}
\tightlist
\item
  the final generated model receives an explicit Human release gate,
\item
  Human Review can evaluate the actual generated \texttt{.sysml} result rather than only upstream abstract Candidates,
\item
  validation results become visible review evidence,
\item
  generated code and model diagrams can be reviewed together,
\item
  agent/personality alternatives remain visible and comparable,
\item
  Human feedback can trigger controlled revision loops,
\item
  LLM assistance remains useful without becoming authority,
\item
  direct UI editing remains possible without creating a hidden engineering source of truth,
\item
  every revision remains auditable,
\item
  final output is clearly separated from generated-but-unreleased work,
\item
  only fully validated and explicitly approved results enter \texttt{data/output/},
\item
  and the final Project output can be traced through the complete engineering workflow.
\end{itemize}

\hypertarget{tradeoffs}{%
\paragraph{Tradeoffs}\label{tradeoffs}}

\begin{itemize}
\tightlist
\item
  Phase L becomes larger than a simple filesystem Output Writer,
\item
  corrections may require returning to earlier phases rather than patching final code directly,
\item
  every material revision requires regeneration and validation,
\item
  the UI needs clear distinction between proposed edits and authoritative changes,
\item
  and a live external validator remains necessary before real publication.
\end{itemize}

These costs are accepted because they preserve engineering authority, traceability and Human control.

\begin{center}\rule{0.5\linewidth}{0.5pt}\end{center}

\hypertarget{relationship-to-prior-adrs}{%
\subsubsection{Relationship to prior ADRs}\label{relationship-to-prior-adrs}}

This ADR complements and does not replace:

\begin{itemize}
\tightlist
\item
  ADR-016 --- Human Review Workspace and Approved Input Promotion Architecture,
\item
  ADR-017 --- Simple-by-Default Interaction and Progressive Disclosure,
\item
  ADR-018 --- Model Candidate Layer and Structural Comparability,
\item
  ADR-019 --- Internal Engineering Model Assembly Architecture,
\item
  ADR-020 --- Hybrid Target Projection and Coverage Architecture,
\item
  ADR-021 --- SYSIDE-Compatible SysML v2 Generation Architecture,
\item
  ADR-022 --- SysML v2 Validation Layer Architecture.
\end{itemize}

ADR-022 K-21 remains normative for final publication.

This ADR clarifies that the K→L \texttt{valid\ +\ passed} gate is a \textbf{publication gate}, not a prohibition against reviewing invalid or incomplete generated models.

This ADR also refines the previous workspace rule for generated SysML:

\begin{verbatim}
review-stage generated SysML
→ project-local Final Model Review workspace

final Human-approved generated SysML
→ data/output/
\end{verbatim}

CATIA remains authoritative for Turing Generator engineering knowledge and is not replaced by generated project output.

\begin{center}\rule{0.5\linewidth}{0.5pt}\end{center}

\hypertarget{decision-summary}{%
\subsubsection{Decision summary}\label{decision-summary}}

The accepted final prototype path is:

\begin{verbatim}
Source
→ Processing Run
→ Human Review
→ Approved Input
→ Model Candidates
→ Candidate Human Review
→ Internal Engineering Model
→ deterministic SysML v2 generation
→ automated validation
→ Final Model Human Review
→ optional controlled revision / LLM-agent loop
→ explicit Human release approval
→ immutable versioned Published Output Package
\end{verbatim}

The core release rule is:

\begin{verbatim}
Generated
≠ Final

Validated
≠ Final

Human-approved exact validated revision
= eligible for final publication
\end{verbatim}

Phase L shall implement this rule without weakening the established deterministic generation, validation, traceability or Human Review architecture.

\textbackslash newpage

\hypertarget{adr-024--guided-engineering-workflow-and-ux-projection-architecture}{%
\subsection{ADR-024 --- Guided Engineering Workflow and UX Projection Architecture}\label{adr-024--guided-engineering-workflow-and-ux-projection-architecture}}

\hypertarget{status-6}{%
\subsubsection{Status}\label{status-6}}

Accepted

\hypertarget{date-6}{%
\subsubsection{Date}\label{date-6}}

2026-08-16

\hypertarget{context-9}{%
\subsubsection{Context}\label{context-9}}

The Turing Generator has reached a functionally complete prototype baseline through Phase L.

Baseline implementation reference:

\begin{verbatim}
0fce9928a04e047e4b39484b421632fb64ed905f
Complete Phase L final model review and output publication
\end{verbatim}

At this baseline, the implemented workflow covers:

\begin{verbatim}
Source
→ Processing
→ Human Review
→ Approved Input
→ Model Candidates
→ Candidate Human Review
→ Internal Engineering Model
→ SysML v2 Generation
→ Validation
→ Final Model Human Review
→ Human Release Approval
→ Published Output
\end{verbatim}

The technical architecture provides immutable evidence, explicit Human authority, deterministic read models, validation gates and end-to-end traceability.

However, formative usability evaluation of the functional prototype showed that the existing user interface remains too strongly oriented around technical artifacts, workflow metadata and implementation structure.

The functionality is available, but the engineer has to spend unnecessary effort finding:

\begin{itemize}
\tightlist
\item
  what engineering information was ingested,
\item
  which processing result is relevant,
\item
  where Human action is required,
\item
  where multiple Agent / Persona results agree,
\item
  where they differ,
\item
  which alternatives are available,
\item
  why an Agent proposed a result,
\item
  and what the next engineering decision should be.
\end{itemize}

The issue is therefore not missing processing capability.

The issue is the mapping between the implemented processing architecture and the daily engineering task.

This ADR defines the UX architecture for WP-09 through WP-12.

\begin{center}\rule{0.5\linewidth}{0.5pt}\end{center}

\hypertarget{decision-9}{%
\subsubsection{Decision}\label{decision-9}}

\hypertarget{ux-01--engineer-centered-interaction}{%
\paragraph{UX-01 --- Engineer-centered interaction}\label{ux-01--engineer-centered-interaction}}

The primary UI shall be organized around the engineer's work rather than the internal implementation structure.

The default questions answered by the UI are:

\begin{verbatim}
What was provided?
What did the system derive?
Do the independent perspectives agree?
Where do I need to decide?
What happens next?
\end{verbatim}

Technical metadata shall remain available but shall not dominate the primary working surface.

\begin{center}\rule{0.5\linewidth}{0.5pt}\end{center}

\hypertarget{ux-02--guided-workflow-is-a-projection-not-authority}{%
\paragraph{UX-02 --- Guided Workflow is a projection, not authority}\label{ux-02--guided-workflow-is-a-projection-not-authority}}

The Guided Engineering Workflow shall not create another authoritative workflow state machine.

It is a deterministic projection of existing persisted authoritative state.

Conceptually:

\begin{verbatim}
authoritative repositories / read services
                ↓
      GuidedWorkflowReadService
                ↓
        GuidedWorkflowView
                ↓
             Streamlit
\end{verbatim}

The UI may maintain transient navigation state.

It shall not duplicate engineering, review, validation, release or publication authority inside Streamlit session state.

\begin{center}\rule{0.5\linewidth}{0.5pt}\end{center}

\hypertarget{ux-03--engineering-content-before-metadata}{%
\paragraph{UX-03 --- Engineering Content before Metadata}\label{ux-03--engineering-content-before-metadata}}

Primary views shall present engineering content first.

Examples include:

\begin{itemize}
\tightlist
\item
  source content,
\item
  extracted engineering statements,
\item
  proposed elements,
\item
  proposed relationships,
\item
  engineering classifications,
\item
  model structure,
\item
  validation findings,
\item
  generated SysML.
\end{itemize}

Implementation metadata such as:

\begin{itemize}
\tightlist
\item
  internal IDs,
\item
  hashes,
\item
  fingerprints,
\item
  artifact paths,
\item
  processing-run identifiers,
\item
  manifest versions,
\end{itemize}

shall remain available through progressive disclosure for traceability and audit.

The intended hierarchy is:

\begin{verbatim}
Primary layer
→ What is this?

Decision layer
→ What do I need to decide?

Explanation layer
→ Why was this proposed?

Traceability layer
→ Where exactly did it come from?
\end{verbatim}

\begin{center}\rule{0.5\linewidth}{0.5pt}\end{center}

\hypertarget{ux-04--decision-centered-interaction}{%
\paragraph{UX-04 --- Decision-Centered Interaction}\label{ux-04--decision-centered-interaction}}

Open Human decisions are first-class UI objects.

The Guided Workflow shall prioritize:

\begin{verbatim}
action required
→ relevant engineering content
→ alternatives
→ supporting rationale
→ Human decision
\end{verbatim}

The engineer shall not have to inspect technical processing state to discover that a decision is required.

The project entry view should therefore prioritize work such as:

\begin{verbatim}
Your work

4 decisions required
2 results contain relevant variance
11 results are confirmed / ready

Next action
→ Review extracted engineering information
\end{verbatim}

The exact counts are derived from authoritative persisted state.

\begin{center}\rule{0.5\linewidth}{0.5pt}\end{center}

\hypertarget{ux-05--variance-is-first-class-engineering-information}{%
\paragraph{UX-05 --- Variance is first-class engineering information}\label{ux-05--variance-is-first-class-engineering-information}}

Whenever a processing step contains redundant LLM / Agent / Persona execution, the UI shall expose the resulting variance explicitly.

Existing consensus concepts shall be reused.

Examples include:

\begin{verbatim}
unanimous
majority
single
none
incomparable
incomplete
\end{verbatim}

and:

\begin{verbatim}
low
medium
high
\end{verbatim}

variance.

Persona stability and incomplete Agent evidence shall remain distinguishable from inter-Persona disagreement.

Repeated executions of one Persona shall not be presented as additional independent votes.

\begin{center}\rule{0.5\linewidth}{0.5pt}\end{center}

\hypertarget{ux-06--side-by-side-before-aggregation}{%
\paragraph{UX-06 --- Side-by-Side before Aggregation}\label{ux-06--side-by-side-before-aggregation}}

When multiple Agent / Persona results exist for the same engineering subject, the preferred presentation is side-by-side comparison.

Conceptually:

\begin{verbatim}
┌────────────────┬────────────────┬────────────────┐
│ Persona A      │ Persona B      │ Persona C      │
│                │                │                │
│ Proposal A     │ Proposal A     │ Proposal B     │
│ rationale ▾    │ rationale ▾    │ rationale ▾    │
└────────────────┴────────────────┴────────────────┘

2 / 3 → Proposal A
1 / 3 → Proposal B

Human decision required
\end{verbatim}

Consensus summaries complement this comparison.

They shall not hide the individual Agent / Persona results.

\begin{center}\rule{0.5\linewidth}{0.5pt}\end{center}

\hypertarget{ux-07--visual-consensus-language}{%
\paragraph{UX-07 --- Visual consensus language}\label{ux-07--visual-consensus-language}}

Agreement and variance shall be visually recognizable at a glance.

The semantic presentation shall follow:

\begin{verbatim}
GREEN
unanimous / low variance

AMBER
majority / medium variance

RED
high variance / no consensus / blocking Human decision

NEUTRAL
incomplete / unavailable / not yet processed
\end{verbatim}

Color alone shall never carry the meaning.

Every visual state shall also contain an explicit textual label, for example:

\begin{verbatim}
Unanimous · 3 / 3 Personas agree
Majority · 2 / 3 Personas agree
High variance · Human decision required
\end{verbatim}

Green means agreement between independent perspectives.

Green does not mean Human approval.

\begin{center}\rule{0.5\linewidth}{0.5pt}\end{center}

\hypertarget{ux-08--recurring-interaction-pattern}{%
\paragraph{UX-08 --- Recurring interaction pattern}\label{ux-08--recurring-interaction-pattern}}

The same interaction pattern shall be reused across processing stages:

\begin{verbatim}
INPUT
↓
PROPOSED RESULT(S)
↓
VARIANCE / CONFIDENCE
↓
HUMAN DECISION — where required
↓
ACCEPTED RESULT
\end{verbatim}

This applies to:

\begin{itemize}
\tightlist
\item
  ingestion interpretation,
\item
  semantic extraction,
\item
  classification,
\item
  Human Review,
\item
  Model Candidate review,
\item
  architecture / Model Proposal review,
\item
  Final Model Review.
\end{itemize}

This consistency is intended to reduce cognitive switching between phases.

\begin{center}\rule{0.5\linewidth}{0.5pt}\end{center}

\hypertarget{ux-09--guided-workflow-stages}{%
\paragraph{UX-09 --- Guided workflow stages}\label{ux-09--guided-workflow-stages}}

The engineer-facing workflow shall expose the following conceptual stages:

\begin{verbatim}
1  Project & Sources
2  Processing
3  Human Review & Approved Input
4  Model Proposal & Candidate Review
5  Final Model Review
6  Published Output
\end{verbatim}

These stages are presentation concepts.

They do not replace existing Phase/domain authority.

\begin{center}\rule{0.5\linewidth}{0.5pt}\end{center}

\hypertarget{ux-10--simple-by-default-explainable-on-demand}{%
\paragraph{UX-10 --- Simple by default, explainable on demand}\label{ux-10--simple-by-default-explainable-on-demand}}

ADR-017 remains the primary presentation principle:

\begin{verbatim}
Simple by default.
Explainable on demand.
Fully traceable underneath.
\end{verbatim}

Default views shall show only information necessary for the current engineering task.

Rationale, evidence, technical metadata and complete traceability remain reachable without changing authority or losing auditability.

\begin{center}\rule{0.5\linewidth}{0.5pt}\end{center}

\hypertarget{ux-11--no-implicit-latest-authority}{%
\paragraph{UX-11 --- No implicit latest authority}\label{ux-11--no-implicit-latest-authority}}

The UI may identify useful display defaults.

It shall not silently select a latest artifact as authority for a write action.

Candidate decisions, review decisions, change proposals, release decisions and publication actions must remain bound to explicit immutable artifact identities.

Display convenience shall never weaken authority binding.

\begin{center}\rule{0.5\linewidth}{0.5pt}\end{center}

\hypertarget{ux-12--existing-services-remain-normative}{%
\paragraph{UX-12 --- Existing services remain normative}\label{ux-12--existing-services-remain-normative}}

The Guided Workflow shall delegate write actions to the already established domain services.

Examples include:

\begin{verbatim}
ReviewApprovalWorkflowService
ModelCandidateReviewRepository
FinalModelReviewChangeService
FinalModelReviewReleaseService
OutputWriter
\end{verbatim}

The UI shall not reproduce their validation or authority rules.

\begin{center}\rule{0.5\linewidth}{0.5pt}\end{center}

\hypertarget{ux-13--dual-layer-presentation}{%
\paragraph{UX-13 --- Dual-layer presentation}\label{ux-13--dual-layer-presentation}}

Existing functional views shall not be discarded merely because their current presentation is technically dense.

Where useful, engineer-facing workflow surfaces shall support two presentation depths over the same authoritative state:

\begin{verbatim}
Focused View
→ engineering content
→ Agent / Persona alternatives
→ consensus and variance
→ open Human decisions
→ accepted / approved engineering results
→ next engineering action

Technical View
→ processing identities
→ run / attempt state
→ artifact references
→ fingerprints
→ provenance
→ traceability
→ diagnostic and lifecycle information
\end{verbatim}

The default presentation shall be the Focused View.

The engineer may explicitly enable technical details when deeper analysis, diagnosis or auditability is required.

Conceptually:

\begin{verbatim}
same authoritative state
          │
          ├── Focused presentation
          │     └── task-oriented engineering information
          │
          └── Technical presentation
                └── deeper processing and traceability information
\end{verbatim}

The presentation-depth setting is UI state only.

It may be stored transiently in Streamlit session state because switching the presentation depth does not modify:

\begin{itemize}
\tightlist
\item
  engineering content,
\item
  Agent evidence,
\item
  Human decisions,
\item
  Candidate authority,
\item
  validation state,
\item
  Final Model Review state,
\item
  publication eligibility,
\item
  or published output.
\end{itemize}

The following distinction is normative:

\begin{verbatim}
Focused ≠ reduced evidence
Focused = reduced visible complexity

Technical ≠ separate workflow
Technical = deeper presentation of the same workflow
\end{verbatim}

Views that already provide useful functionality, including Project Dashboard, Processing and Human Review, shall therefore be adapted progressively to this dual-layer interaction model rather than replaced unnecessarily.

The Project Dashboard retains the distinct engineering purpose:

\begin{verbatim}
Which engineering / model information is already available?
\end{verbatim}

The Guided Engineering Workspace retains the distinct purpose:

\begin{verbatim}
What requires my attention or decision now?
\end{verbatim}

Individual LLM-backed processing steps may remain separate working surfaces when this improves clarity. They shall reuse the same content-first, variance-aware and Human-decision-centered interaction language.

\begin{center}\rule{0.5\linewidth}{0.5pt}\end{center}

\hypertarget{ux-14--global-application-context}{%
\paragraph{UX-14 --- Global application context}\label{ux-14--global-application-context}}

Project selection and presentation depth are application-level context.

They shall therefore remain directly accessible independently of the currently selected working surface.

The common application shell shall expose:

\begin{verbatim}
Project
→ explicit currently selected Project

Technical details
→ Focused / Technical presentation depth
\end{verbatim}

Conceptually:

\begin{verbatim}
APPLICATION CONTEXT
Project: [ selected Project ]
Technical details: [ on / off ]

        ↓ shared by all views

Engineering Workspace
Project Dashboard
Processing
Human Review
Model Proposal
Final Model Review
Published Output
\end{verbatim}

The engineer shall not have to navigate to the Project Dashboard merely to change the current Project.

Likewise, the engineer shall not have to return to the Engineering Workspace to change presentation depth.

Project selection changes UI working context.

It shall not modify persisted engineering state.

Changing the selected Project shall clear stale entity-level UI navigation from the previously selected Project.

No Project shall be silently selected merely because it is the newest or only recently used Project when no valid explicit selection exists.

Project creation remains a Project Workspace management function but shall be directly accessible from the global application context.

The Project selector and the Project creation action shall therefore be colocated in the common application shell:

\begin{verbatim}
Project
[ Selected Project ▼ ]
+ Create new project
\end{verbatim}

Project creation is an explicit action and shall not be represented as a pseudo-Project inside the Project selection list.

After successful creation, the newly created Project becomes the explicitly selected Project context.

The Project Dashboard shall not duplicate the creation control when it is rendered inside the common application shell.

The presentation-depth control shall use a persistent application-level toggle and shall affect all engineer-facing views progressively as WP-09 through WP-11 adopt the dual-layer presentation architecture.

\begin{center}\rule{0.5\linewidth}{0.5pt}\end{center}

\hypertarget{ux-15--authority-preserving-write-interaction}{%
\paragraph{UX-15 --- Authority-preserving write interaction}\label{ux-15--authority-preserving-write-interaction}}

The Guided Engineering Workflow may initiate authoritative actions.

It shall never become engineering or workflow authority itself.

Every write initiated from an engineer-facing UI surface shall follow:

\begin{verbatim}
explicit immutable target
+
explicit Human action
+
existing domain authority service
+
immutable authoritative persistence
+
read-side reconstruction after the write
\end{verbatim}

Conceptually:

\begin{verbatim}
Streamlit interaction
        ↓
GuidedWorkflowWriteService
        ↓
existing normative domain service
        ↓
authoritative immutable persistence
        ↓
existing read service / repository
        ↓
reconstructed UI state
\end{verbatim}

\texttt{GuidedWorkflowWriteService} is an application-level delegation boundary only.

It shall not:

\begin{itemize}
\tightlist
\item
  reproduce Candidate Review rules,
\item
  reproduce Final Model Review routing rules,
\item
  reproduce release-gate logic,
\item
  infer engineering approval,
\item
  infer a latest authoritative target,
\item
  mutate generated SysML directly,
\item
  or store authoritative decisions in Streamlit session state.
\end{itemize}

The following domain services remain normative:

\begin{verbatim}
ModelCandidateReviewRepository
FinalModelReviewChangeService
FinalModelReviewReleaseService
OutputWriter
\end{verbatim}

Candidate Review writes shall require an explicit:

\begin{verbatim}
Project
Candidate Set
Candidate type
Candidate identity
Human decision
Reviewer identity
\end{verbatim}

Final Model Review change requests shall require an explicit:

\begin{verbatim}
Project
Final Model Review
Final Model Review revision
affected review surface
change classification
Human feedback
reviewer identity
\end{verbatim}

Final release approval shall require an explicit:

\begin{verbatim}
Project
Final Model Review
Final Model Review revision
reviewer identity
\end{verbatim}

The UI may simplify labels and interaction controls.

It shall not simplify or weaken the underlying authority contract.

After a successful write, the UI shall discard transient assumptions and reconstruct the visible state from persisted authoritative evidence.

\begin{center}\rule{0.5\linewidth}{0.5pt}\end{center}

\hypertarget{ux-16--processing-and-human-review-projection}{%
\paragraph{UX-16 --- Processing and Human Review projection}\label{ux-16--processing-and-human-review-projection}}

WP-10 shall adapt Processing and Human Review to the same deterministic dual-layer presentation architecture established by this ADR.

The default Processing projection shall prioritize:

\begin{verbatim}
Source filename / readable Source identity
→ current Processing status
→ relevant result / attention state
→ next engineering action
\end{verbatim}

The default Processing surface shall not require the engineer to interpret Source IDs, Processing Run IDs, Attempt IDs, configuration fingerprints, artifact fingerprints or SHA-256 values to perform normal work. Those facts remain available in Technical View.

The default Human Review projection shall prioritize:

\begin{verbatim}
engineering statement
→ independent Persona proposals
→ consensus / material variance
→ required Human decision
→ accepted / approved result
\end{verbatim}

All Persona proposals remain individually inspectable. Repeated executions associated with one Persona remain grouped beneath that Persona and shall not be displayed as additional independent votes.

Consensus and variance shall be projected only from existing persisted consensus evidence. The presentation layer shall not infer semantic agreement when authoritative consensus evidence is absent. Relationship Review Items for which the existing Consensus Analyzer provides no explicit relationship consensus shall therefore present consensus as unavailable rather than inventing an agreement state.

The following distinction is normative:

\begin{verbatim}
Focused Processing
→ filename / Source role
→ Processing status
→ result or required attention
→ next action

Focused Human Review
→ engineering content
→ Persona alternatives
→ consensus / variance
→ Human decision
→ review / approval lifecycle state

Technical View
→ Source / Run / Attempt identities
→ Review Document / Version / Revision identities
→ artifact references
→ fingerprints
→ evidence locators
→ provenance and lifecycle diagnostics
\end{verbatim}

Processing and Human Review projection adapters are read-side presentation components only. They shall not own Source Registry authority, Processing Run state, Review Revision state, Human Review Decision authority, Approved Input authority, consensus authority or write validation.

Existing services remain normative, including:

\begin{verbatim}
ProjectBoundIngestionService
ReviewApprovalWorkflowService
\end{verbatim}

Human write actions continue to operate on exact immutable targets and shall be followed by read-side reconstruction.

\begin{center}\rule{0.5\linewidth}{0.5pt}\end{center}

\hypertarget{ux-17--architecture--model-proposal-projection}{%
\paragraph{UX-17 --- Architecture / Model Proposal projection}\label{ux-17--architecture--model-proposal-projection}}

WP-11 shall adapt Model Proposal and Candidate Review to the same deterministic, dual-layer presentation architecture established by this ADR.

The default Model Proposal projection shall prioritize:

\begin{verbatim}
proposed architecture / model structure
→ material alternatives and structural variance
→ required Human Candidate decisions
→ Candidate Review state
→ Phase-I readiness
→ next engineering action
\end{verbatim}

The engineer shall not have to reconstruct the proposed architecture mentally from Candidate IDs, repository state or independent technical tables.

The Focused Model Proposal shall therefore expose an engineer-readable structural projection of the exact selected Candidate Set, including:

\begin{verbatim}
model elements
→ readable name
→ model area / type
→ support state
→ Human review state

model relationships
→ source
→ semantic intent
→ target
→ resolution / priority state
→ Human review state

relationship choice groups
→ all alternatives
→ preferred alternatives where persisted
→ accepted alternatives where persisted
→ whether Human review remains required

structural / profile deviations
→ affected Candidate
→ conformance / comparability state
→ whether Human review remains required
\end{verbatim}

Relationship Choice Groups are alternatives, not independent Agent votes. The presentation layer shall not translate Candidate alternatives into Persona consensus or invent agreement evidence.

Phase-I readiness shall be projected only from the existing Model Proposal / Phase-I gate state. The UI shall not infer that a Candidate Set is ready merely because all currently visible controls appear resolved.

The following distinction is normative:

\begin{verbatim}
Focused Model Proposal
→ architecture / proposed model content
→ alternatives and material deviations
→ Human decisions
→ review progress
→ Phase-I readiness
→ next action

Technical View
→ Candidate Set identity / fingerprint
→ Candidate identities
→ Approved Input references
→ structure-profile conformance
→ comparison anchors / deviation identities
→ decision identities and exact traceability
\end{verbatim}

Model Proposal presentation adapters are read-side components only. They shall not own Candidate authority, Approved Input authority, Candidate Review authority or Internal Engineering Model assembly authority.

Existing services remain normative, including:

\begin{verbatim}
ModelProposalReadService
ModelCandidateReviewRepository
ModelCandidateReadService / Phase-I gate
\end{verbatim}

Every Candidate Review write remains bound to the explicit immutable Candidate Set and Candidate identity currently shown. After a write, the visible proposal shall be reconstructed from persisted authoritative state.

No graph or diagram rendering technology is made normative by this decision. WP-11 may use the existing structural overview to improve architecture readability without introducing a second model representation or model authority.

\begin{center}\rule{0.5\linewidth}{0.5pt}\end{center}

\hypertarget{formative-usability-evaluation}{%
\subsubsection{Formative usability evaluation}\label{formative-usability-evaluation}}

The UX redesign follows a formative evaluation of the functionally complete prototype.

The evaluation is treated as formative engineering feedback, not as a controlled quantitative usability study unless separate evidence supports such a claim.

Observed findings and resulting design responses are maintained in:

\begin{verbatim}
collaboration/ux/usability_findings.md
\end{verbatim}

The implemented baseline and redesigned UI are intentionally retained as separate development stages for later thesis documentation.

\begin{center}\rule{0.5\linewidth}{0.5pt}\end{center}

\hypertarget{thesis-evidence}{%
\subsubsection{Thesis evidence}\label{thesis-evidence}}

The UX iteration shall be documented so that the development process can later be reconstructed as:

\begin{verbatim}
technical feasibility
→ functionally complete prototype
→ formative usability evaluation
→ identified usability findings
→ UX requirements
→ wireframes
→ redesigned Guided Engineering Workflow
→ final demonstrator
\end{verbatim}

Stable wireframe identifiers shall be maintained under:

\begin{verbatim}
collaboration/ux/wireframes/
\end{verbatim}

The intended initial set is:

\begin{verbatim}
WF-01  Engineer Home / Your Work
WF-02  Ingested Engineering Content
WF-03  Unanimous Persona Results
WF-04  Variant Persona Results
WF-05  Human Decision
WF-06  Model Proposal Review
WF-07  Final Model Review
\end{verbatim}

Final thesis figures may be derived from these documented wireframes and from before/after screenshots of the actual prototype.

\begin{center}\rule{0.5\linewidth}{0.5pt}\end{center}

\hypertarget{work-package-mapping}{%
\subsubsection{Work-package mapping}\label{work-package-mapping}}

\hypertarget{wp-09--guided-workflow-ui}{%
\paragraph{WP-09 --- Guided Workflow UI}\label{wp-09--guided-workflow-ui}}

Establish:

\begin{itemize}
\tightlist
\item
  Guided Engineering Workflow projection,
\item
  shared navigation,
\item
  shared Decision presentation,
\item
  shared Variance presentation,
\item
  progressive disclosure,
\item
  reusable side-by-side result comparison.
\end{itemize}

\hypertarget{wp-10--ingestion--human-review-ux-simplification}{%
\paragraph{WP-10 --- Ingestion + Human Review UX Simplification}\label{wp-10--ingestion--human-review-ux-simplification}}

Apply the UX architecture to:

\begin{itemize}
\tightlist
\item
  Source presentation,
\item
  processing results,
\item
  redundant Agent / Persona outputs,
\item
  Human Review decisions.
\end{itemize}

\hypertarget{wp-11--architecture--model-proposal-ux}{%
\paragraph{WP-11 --- Architecture / Model Proposal UX}\label{wp-11--architecture--model-proposal-ux}}

Apply the same interaction model to:

\begin{itemize}
\tightlist
\item
  Model Candidate proposals,
\item
  structural comparison,
\item
  architecture/model representation,
\item
  candidate Human decisions.
\end{itemize}

\hypertarget{wp-12--end-to-end-demo-hardening}{%
\paragraph{WP-12 --- End-to-End Demo Hardening}\label{wp-12--end-to-end-demo-hardening}}

Integrate and harden the complete engineer-facing workflow without weakening any existing authority, validation or traceability contract.

\begin{center}\rule{0.5\linewidth}{0.5pt}\end{center}

\hypertarget{consequences-9}{%
\subsubsection{Consequences}\label{consequences-9}}

Positive:

\begin{itemize}
\tightlist
\item
  the UI is organized around daily engineering work,
\item
  Human decisions become immediately visible,
\item
  Agent variance becomes useful information instead of hidden processing detail,
\item
  redundant Agent results can be compared quickly,
\item
  existing functional views can be retained and progressively improved,
\item
  focused and technical work can use the same authoritative backend,
\item
  technical traceability remains fully available,
\item
  the same interaction language is reused across phases,
\item
  UX evolution is reproducibly documented for the thesis.
\end{itemize}

Trade-offs:

\begin{itemize}
\tightlist
\item
  presentation read models become more important,
\item
  some technical metadata moves behind additional interaction,
\item
  visual consensus representations require careful accessibility semantics,
\item
  responsive side-by-side layouts require deliberate UI design.
\end{itemize}

These trade-offs are accepted.

\begin{center}\rule{0.5\linewidth}{0.5pt}\end{center}

\hypertarget{prohibited-shortcuts}{%
\subsubsection{Prohibited shortcuts}\label{prohibited-shortcuts}}

The following are explicitly prohibited:

\begin{itemize}
\tightlist
\item
  replacing engineering authority with UI state,
\item
  treating consensus as Human approval,
\item
  hiding dissenting Agent / Persona results behind an aggregate score,
\item
  counting repeated runs of one Persona as independent votes,
\item
  presenting color without textual meaning,
\item
  direct mutation of generated SysML from the UI,
\item
  implicit latest-artifact authority,
\item
  bypassing Human Review,
\item
  bypassing Validation,
\item
  bypassing Final Human release approval,
\item
  removing traceability merely to simplify the visible UI,
\item
  maintaining separate engineering authority for Focused and Technical views,
\item
  changing engineering state merely by switching presentation depth.
\end{itemize}

\begin{center}\rule{0.5\linewidth}{0.5pt}\end{center}

\hypertarget{final-principle}{%
\subsubsection{Final principle}\label{final-principle}}

The UI shall minimize the engineer's effort to understand and decide.

It shall not minimize the engineering evidence available to justify that decision.

\textbackslash newpage

\hypertarget{adr-025--semantic-proposal-consolidation-and-persona-aware-consensus}{%
\subsection{ADR-025 --- Semantic Proposal Consolidation and Persona-Aware Consensus}\label{adr-025--semantic-proposal-consolidation-and-persona-aware-consensus}}

\hypertarget{status-7}{%
\subsubsection{Status}\label{status-7}}

Accepted

\hypertarget{date-7}{%
\subsubsection{Date}\label{date-7}}

2026-08-17

\hypertarget{context-10}{%
\subsubsection{Context}\label{context-10}}

The WP-12 controlled multi-document dry run exposed a structural mismatch between the ingestion proposal layer and the Human Review experience.

The current review projection can materialize several independent Review Items for Agent proposals that refer to the same engineering concept but use different wording or a different proposed element classification. This creates excessive Human decisions and can hide the most useful form of variance: different Persona interpretations of one semantic subject.

The same problem exists for relationships. Syntactically different predicates can express the same engineering relation, while current identity and consensus paths do not provide a stable semantic relationship subject.

This is not a reason to make Agent output contract-perfect before Human Review. Reviewable semantic uncertainty remains legitimate processing output. The Human Review boundary exists to resolve modeling ambiguity, classification variance, relationship ambiguity, and other engineering uncertainty. Hard authority and integrity failures remain fail-closed.

The architecture therefore needs a derived, immutable layer that answers one limited question before Human Review:

\begin{quote}
Which exact existing Agent proposals appear to refer to the same engineering subject?
\end{quote}

That layer must never invent a new engineering statement, silently discard source evidence, convert agreement into Human approval, or make uncertain equivalence authoritative.

\hypertarget{decision-10}{%
\subsubsection{Decision}\label{decision-10}}

Introduce a \textbf{Semantic Proposal Consolidation} layer between structured derivation proposals and the Persona-aware consensus / Human Review projection.

The target flow is:

\begin{verbatim}
Exact immutable Agent Proposals
+ exact Source Evidence
        ↓
Semantic Consolidation
        │
        ├─ Elements
        │   group by engineering meaning,
        │   not by exact name or element_type
        │
        ├─ Relationships
        │   group by consolidated endpoint subjects
        │   and semantic predicate meaning
        │
        └─ comparison outcomes
            equivalent | distinct | uncertain
        ↓
Immutable Semantic Consolidation Artifact
        ↓
Persona-aware Consensus Projection
        ↓
one Engineering Review subject per semantic subject
        ↓
Human Review
\end{verbatim}

Semantic consolidation executes during Processing, after structured derivation proposals exist and before the Processing Run reaches \texttt{awaiting\_review}. Opening Human Review reconstructs persisted results and does not initiate a new semantic-model call.

\hypertarget{authority-boundary}{%
\subsubsection{Authority boundary}\label{authority-boundary}}

The Semantic Consolidation Artifact is \textbf{derived processing evidence}. It is not an Approved Input, Model Candidate, Internal Engineering Model, or Human approval.

The consolidator may only classify relations among already-existing exact proposals. It must not:

\begin{itemize}
\tightlist
\item
  create a new engineering claim;
\item
  overwrite an Agent proposal;
\item
  mutate source evidence;
\item
  silently select an element classification;
\item
  silently resolve an uncertain relationship endpoint;
\item
  turn Persona agreement into Human approval;
\item
  infer that two proposals are equivalent when the comparison result is \texttt{uncertain} or unavailable.
\end{itemize}

Every semantic subject remains traceable to the exact immutable proposal, Agent/Persona/run, upstream artifact fingerprint, and source-evidence references from which it was derived.

\hypertarget{merge-authority}{%
\subsubsection{Merge authority}\label{merge-authority}}

A semantic merge is permitted only when explicit consolidation evidence classifies the relevant proposals as \texttt{equivalent}.

The contract is fail-closed \textbf{against merge authority}:

\begin{verbatim}
equivalent
→ may support one semantic subject

distinct
→ must remain separate

uncertain
→ must remain separate and remain visible for Human reasoning

missing / malformed comparison
→ cannot authorize a merge
\end{verbatim}

A multi-proposal semantic subject must therefore be connected by explicit \texttt{equivalent} comparison evidence. A \texttt{distinct} or \texttt{uncertain} comparison inside one semantic subject is an integrity failure. An \texttt{equivalent} comparison across two different subjects is also an integrity failure.

This rule allows safe degradation toward singleton semantic subjects without inventing equivalence.

\hypertarget{element-identity}{%
\subsubsection{Element identity}\label{element-identity}}

\texttt{element\_type} is explicitly \textbf{not} part of semantic-subject identity.

Different Personas may recognize the same engineering concept while proposing different modeling classifications. Such disagreement must become visible classification variance on one Review subject.

Example:

\begin{verbatim}
"separate client application"

Persona A → interface
Persona B → system
Persona C → system

Recognition:
3 / 3 Personas

Classification:
2 × system
1 × interface

→ one semantic subject
→ visible variance
→ explicit Human decision
\end{verbatim}

\hypertarget{relationship-identity}{%
\subsubsection{Relationship identity}\label{relationship-identity}}

Relationship consolidation is based on:

\begin{enumerate}
\def\labelenumi{\arabic{enumi}.}
\tightlist
\item
  already consolidated semantic source subject;
\item
  already consolidated semantic target subject; and
\item
  semantic meaning of the relationship predicate.
\end{enumerate}

Surface wording alone is not relationship identity. For example, \texttt{retain}, \texttt{retains}, and \texttt{should\ retain} may represent one semantic relationship when comparison evidence supports equivalence.

Relationship consolidation must not silently repair unresolved endpoints. Reviewable endpoint uncertainty remains eligible for Human Review under the existing Human-in-the-Loop boundary.

\hypertarget{persona-aware-consensus}{%
\subsubsection{Persona-aware consensus}\label{persona-aware-consensus}}

Consensus is counted by \textbf{Persona perspective}, not by raw Agent run.

Multiple runs of one Persona do not create multiple votes:

\begin{verbatim}
Persona A / Run 1 ─┐
Persona A / Run 2 ─┴→ one Persona perspective

Persona B / Run 1 ─┐
Persona B / Run 2 ─┴→ one Persona perspective
\end{verbatim}

Divergence between repeated runs of one Persona is retained as \textbf{intra-Persona instability}. It is evidence, not additional voting weight.

Full Persona agreement remains a processing observation only and never constitutes Human approval.

\hypertarget{semantic-comparator}{%
\subsubsection{Semantic comparator}\label{semantic-comparator}}

A later implementation slice may use a semantic model for comparisons. The comparison context must remain bounded to the minimum information needed for semantic equivalence assessment, for example:

\begin{itemize}
\tightlist
\item
  exact proposal reference;
\item
  candidate / relationship wording;
\item
  proposed classification where relevant;
\item
  concise proposal description;
\item
  exact source-evidence reference / statement.
\end{itemize}

The entire legacy document, unrelated Agent rationale, or unrelated project context must not be resent by default.

Deterministic pre-filtering may reduce the number of semantic comparisons, but it must not authorize a semantic merge unless the resulting consolidation contract contains explicit valid equivalence evidence.

\hypertarget{c1--immutable-consolidation-contract}{%
\subsubsection{C1 --- Immutable consolidation contract}\label{c1--immutable-consolidation-contract}}

The first implementation slice introduces only the artifact contract and its architecture tests. C1 performs no LLM call and no production clustering.

The artifact binds:

\begin{itemize}
\tightlist
\item
  exact Project ID;
\item
  exact Processing Run ID;
\item
  immutable upstream artifact references and fingerprints;
\item
  exact proposal references;
\item
  proposal kind;
\item
  Agent ID;
\item
  Persona ID;
\item
  run index;
\item
  exact source-evidence references;
\item
  semantic subjects and their exact member proposal references;
\item
  explicit pairwise comparison outcomes and comparison provenance;
\item
  an input-set fingerprint; and
\item
  an artifact fingerprint.
\end{itemize}

The C1 contract requires:

\begin{itemize}
\tightlist
\item
  deterministic ordering;
\item
  unique upstream artifact references;
\item
  unique proposal references;
\item
  every proposal to belong to exactly one semantic subject;
\item
  no proposal to occur in more than one semantic subject;
\item
  proposal kind to match semantic-subject kind;
\item
  unique unordered comparison pairs;
\item
  no self-comparison;
\item
  comparison references to resolve to known proposals of the same kind;
\item
  multi-member subjects to have a connected graph of explicit \texttt{equivalent} comparisons;
\item
  no \texttt{distinct} or \texttt{uncertain} comparison within one subject;
\item
  no \texttt{equivalent} comparison across different subjects;
\item
  exact fingerprint validation on deserialization.
\end{itemize}

Repeated proposals from the same Persona with different run indexes are valid input. C1 deliberately does not count them as independent votes.

\hypertarget{implementation-slices}{%
\subsubsection{Implementation slices}\label{implementation-slices}}

\begin{verbatim}
C1  SemanticConsolidationArtifact + fail-closed contracts
C2  Element semantic clustering + Persona-aware consensus
C3  Relationship semantic clustering + consensus
C4  Review projection/cards + empirical before/after test
\end{verbatim}

C2--C4 must be reviewed against the persisted C1 contract and existing Human Review authority before integration.

\hypertarget{acceptance-criteria-1}{%
\subsubsection{Acceptance criteria}\label{acceptance-criteria-1}}

The completed architecture is accepted when:

\begin{itemize}
\tightlist
\item
  semantically equivalent wordings become one Review subject;
\item
  different element classifications of the same concept remain one subject with visible variance;
\item
  repeated runs of one Persona do not add votes;
\item
  truly distinct engineering concepts remain separate;
\item
  uncertain equivalence is never silently merged;
\item
  relationships receive semantic consolidation and Persona-aware consensus;
\item
  every original proposal and its source evidence remain inspectable;
\item
  the Human can split an incorrect semantic cluster before approval; and
\item
  full semantic agreement never implies Human approval.
\end{itemize}

No fixed target number of Review decisions is an acceptance criterion. The solution must not be optimized to the current WP-12 synthetic dataset.

\hypertarget{wp-12-evaluation-use}{%
\subsubsection{WP-12 evaluation use}\label{wp-12-evaluation-use}}

The existing WP-12 retest state with 68 Human decisions is retained as immutable before-evidence. After C1--C4 are accepted, a new isolated Processing Run shall use the same source / Persona / run configuration for an empirical before/after comparison.

No additional live LLM test is required for C1.

\hypertarget{consequences-10}{%
\subsubsection{Consequences}\label{consequences-10}}

\hypertarget{positive-3}{%
\paragraph{Positive}\label{positive-3}}

\begin{itemize}
\tightlist
\item
  Human Review operates on engineering subjects rather than duplicated strings.
\item
  Classification variance becomes visible instead of fragmenting identity.
\item
  Persona consensus becomes independent of repeated-run count.
\item
  Relationship consensus gains a defined semantic basis.
\item
  Semantic grouping remains traceable, reversible, and subordinate to Human authority.
\item
  The architecture can degrade safely toward non-merge instead of inventing equivalence.
\end{itemize}

\hypertarget{negative--cost}{%
\paragraph{Negative / cost}\label{negative--cost}}

\begin{itemize}
\tightlist
\item
  Processing gains an additional derived artifact and later a semantic comparison step.
\item
  Semantic comparison can add latency and token cost when a model is used.
\item
  The UI must later expose grouping evidence and support a Human split action.
\item
  Stability of semantic grouping becomes an additional property that requires testing.
\end{itemize}

\hypertarget{rejected-alternatives-4}{%
\subsubsection{Rejected alternatives}\label{rejected-alternatives-4}}

\hypertarget{exact-normalized-name--element-type-as-identity}{%
\paragraph{Exact normalized name + element type as identity}\label{exact-normalized-name--element-type-as-identity}}

Rejected because it fragments one engineering concept when Personas disagree on classification and therefore hides useful variance.

\hypertarget{exact-normalized-name-only}{%
\paragraph{Exact normalized name only}\label{exact-normalized-name-only}}

Rejected because wording variation still creates duplicates and homonyms may be incorrectly merged.

\hypertarget{fuzzy-string-matching-as-merge-authority}{%
\paragraph{Fuzzy string matching as merge authority}\label{fuzzy-string-matching-as-merge-authority}}

Rejected because lexical similarity is not sufficient engineering evidence for semantic identity and may silently merge distinct concepts.

\hypertarget{make-the-agent-output-contract-perfect-before-human-review}{%
\paragraph{Make the Agent output contract-perfect before Human Review}\label{make-the-agent-output-contract-perfect-before-human-review}}

Rejected because semantic uncertainty is legitimate Human-in-the-Loop input. Only hard authority and integrity corruption should block review construction.

\hypertarget{run-semantic-consolidation-when-human-review-is-opened}{%
\paragraph{Run semantic consolidation when Human Review is opened}\label{run-semantic-consolidation-when-human-review-is-opened}}

Rejected because review opening would incur new model cost, could produce different grouping on repeated opens, and would make Human Review availability depend on a live external call.

\textbackslash newpage

\hypertarget{adr-026--source-anchored-multi-persona-interpretation-and-cross-unit-semantic-synthesis}{%
\subsection{ADR-026 --- Source-Anchored Multi-Persona Interpretation and Cross-Unit Semantic Synthesis}\label{adr-026--source-anchored-multi-persona-interpretation-and-cross-unit-semantic-synthesis}}

\hypertarget{status-8}{%
\subsubsection{Status}\label{status-8}}

Accepted

\hypertarget{date-8}{%
\subsubsection{Date}\label{date-8}}

2026-08-18

\hypertarget{context-11}{%
\subsubsection{Context}\label{context-11}}

WP12-E2E-DRY-001 exposed a structural limitation in the current Phase-F multi-agent interpretation and Human Review path.

The original BLK-003 finding showed that semantically equivalent Agent proposals could become separate Human Review subjects because identity and grouping were primarily derived after independent broad-context Agent execution.

ADR-025 introduced semantic proposal consolidation, persona-aware consensus, relationship consolidation, and semantic Human Review projection.

The subsequent BLK-003.1 corrections improved processing robustness:

\begin{itemize}
\tightlist
\item
  exact source-evidence statements no longer collide solely because they share one Agent-local source-information identifier,
\item
  unresolved or ambiguous relationship endpoints are preserved for Human Review instead of terminating Processing,
\item
  non-recoverable semantic processing errors are classified more precisely at the Project Ingestion service boundary.
\end{itemize}

These corrections resolved the observed hard Processing failure, but the controlled POST-BLK-003.1 empirical retest exposed a deeper architectural problem.

\hypertarget{wp12-empirical-evidence}{%
\paragraph{WP12 empirical evidence}\label{wp12-empirical-evidence}}

Controlled retest:

\begin{itemize}
\tightlist
\item
  Project: \texttt{887027}
\item
  Processing Run: \texttt{RUN-000001}
\item
  Source: \texttt{01\_product\_overview.md}
\item
  Personas: 3
\item
  Runs per Persona: 2
\item
  Processing result: \texttt{awaiting\_review}
\end{itemize}

Observed semantic consolidation result:

\begin{itemize}
\tightlist
\item
  66 raw Element proposals,
\item
  66 Element semantic subjects,
\item
  Element consolidation degraded to singletons,
\item
  warning: \texttt{semantic\_comparator\_unavailable},
\item
  46 raw Relationship proposals,
\item
  46 Relationship semantic subjects,
\item
  Relationship consolidation degraded to singletons,
\item
  warning: \texttt{relationship\_semantic\_comparator\_invalid},
\item
  112 resulting Human Review items.
\end{itemize}

The LLM calls themselves completed, but their global consolidation outputs were not safely consumable:

\begin{itemize}
\tightlist
\item
  the Element comparator response was syntactically invalid and could not be parsed under the strict contract,
\item
  the Relationship comparator response violated the required complete, non-overlapping proposal partition,
\item
  safe fallback behavior therefore retained every proposal as an independent singleton subject.
\end{itemize}

The fallback preserved evidence and authority boundaries, but the resulting Human Review workload was unacceptable.

\hypertarget{architectural-finding}{%
\paragraph{Architectural finding}\label{architectural-finding}}

The live Phase-F workflow currently gives a team-wide stage input to all Agent personas. Each persona independently derives candidates from broad source context, and semantic identity is established only afterwards.

Source references therefore provide provenance, but they do not necessarily provide an explicit orchestration contract that says:

\begin{quote}
These Persona outputs are independent interpretations of the same canonical source analysis unit.
\end{quote}

This turns semantic consolidation into a large global clustering and partitioning problem.

For the controlled retest, 66 Element proposals and 46 Relationship proposals had to be reconciled after broad-context derivation. A failure of Element consolidation also reduces the ability of downstream Relationship consolidation to identify common endpoints.

\hypertarget{existing-reusable-architecture}{%
\paragraph{Existing reusable architecture}\label{existing-reusable-architecture}}

The repository already contains several compatible upstream concepts.

\texttt{SourceProjection} provides deterministic source content, ordered persisted segments, source locators, segment identifiers, and exact offsets.

\texttt{InformationUnitSourceAnchor} provides exact source ranges within persisted Source Projection segments.

\texttt{semantic\_extraction} provides persona/run-bound candidate representations with exact source anchors and excerpts.

\texttt{semantic\_consensus} already models:

\begin{itemize}
\tightlist
\item
  required Personas,
\item
  repeated Persona runs,
\item
  intra-Persona stability,
\item
  distinct-Persona voting,
\item
  exact source-evidence grouping,
\item
  field variance,
\item
  Human Review requirements,
\item
  and proposed Information Unit drafts.
\end{itemize}

These concepts shall be reused where compatible.

The existing \texttt{InformationUnit} concept shall not be repurposed as the pre-interpretation orchestration anchor. An Information Unit is already an interpreted, independently reviewable semantic claim with classification, epistemic state, extraction provenance, confidence, and content fingerprint.

A separate source-level orchestration identity is therefore required.

\begin{center}\rule{0.5\linewidth}{0.5pt}\end{center}

\hypertarget{decision-11}{%
\subsubsection{Decision}\label{decision-11}}

\hypertarget{1-introduce-canonical-source-analysis-units}{%
\paragraph{1. Introduce canonical Source Analysis Units}\label{1-introduce-canonical-source-analysis-units}}

Phase F shall introduce an immutable source-level orchestration artifact:

\texttt{SourceAnalysisUnit}

Proposed identifier form:

\texttt{SAU-000001}

A Source Analysis Unit answers only:

\begin{quote}
Which exact portion of the original projected source are the Personas being asked to interpret together?
\end{quote}

It is not an engineering interpretation and does not itself assert a model element, requirement, function, relationship, or other semantic meaning.

Each Source Analysis Unit shall bind at least:

\begin{itemize}
\tightlist
\item
  project identity,
\item
  source identity,
\item
  source projection identity,
\item
  source analysis unit identity,
\item
  exact Source Projection segment/range anchors,
\item
  exact source excerpt,
\item
  deterministic source-order information,
\item
  segmentation-profile/version reference,
\item
  immutable content fingerprint.
\end{itemize}

Source Analysis Units shall remain persona-independent.

No Agent may create or redefine the identity of the Source Analysis Unit it is asked to assess.

\begin{center}\rule{0.5\linewidth}{0.5pt}\end{center}

\hypertarget{2-initial-source-analysis-unit-segmentation}{%
\paragraph{2. Initial Source Analysis Unit segmentation}\label{2-initial-source-analysis-unit-segmentation}}

For the first implementation, the deterministic persisted Source Projection segments are the default Source Analysis Unit boundaries.

For text and Markdown sources, the existing Source Projection adapter currently creates ordered non-empty blank-line blocks.

This default is intentionally conservative:

\begin{itemize}
\tightlist
\item
  it already has deterministic identity and source locations,
\item
  it avoids introducing another LLM-controlled segmentation authority,
\item
  it keeps all Personas on the same bounded source context,
\item
  and it is sufficient for the WP12 controlled source fixtures.
\end{itemize}

A later versioned segmentation profile may deterministically subdivide Source Projection segments where evidence shows that smaller analysis units materially improve comparison quality.

Such subdivision must preserve exact source anchors and must not depend on the interpretation of an individual Persona.

\begin{center}\rule{0.5\linewidth}{0.5pt}\end{center}

\hypertarget{3-source-anchored-persona-execution}{%
\paragraph{3. Source-anchored Persona execution}\label{3-source-anchored-persona-execution}}

All configured Personas for a semantic/derivation team shall assess the same Source Analysis Unit.

Conceptually:

\begin{verbatim}
SourceAnalysisUnit SAU-000003
        |
        +-- Persona A / Run 1
        +-- Persona A / Run 2
        +-- Persona B / Run 1
        +-- Persona B / Run 2
        +-- Persona C / Run 1
        +-- Persona C / Run 2
\end{verbatim}

Every produced proposal shall retain:

\begin{itemize}
\tightlist
\item
  \texttt{source\_analysis\_unit\_id},
\item
  exact Agent identity,
\item
  Persona identity,
\item
  Persona run index,
\item
  exact source evidence,
\item
  exact Agent wording,
\item
  confidence and rationale,
\item
  and existing immutable provenance.
\end{itemize}

The Source Analysis Unit identity is the comparison scope, not the semantic identity of the proposed engineering subject.

One Source Analysis Unit may legitimately produce zero, one, or many proposed engineering subjects.

\begin{center}\rule{0.5\linewidth}{0.5pt}\end{center}

\hypertarget{4-local-source-anchored-semantic-consolidation}{%
\paragraph{4. Local source-anchored semantic consolidation}\label{4-local-source-anchored-semantic-consolidation}}

Semantic consolidation shall first operate inside one Source Analysis Unit.

The local comparator receives only proposals produced for that Source Analysis Unit.

Its task is to identify which proposals represent the same local engineering subject and which remain distinct or uncertain.

Conceptually:

\begin{verbatim}
SAU-000003
    |
    +-- Persona A: "remote expert"
    +-- Persona B: "external expert"
    +-- Persona C: "remote specialist"
    |
    +--> Local Subject: Remote Expert
\end{verbatim}

The local result shall preserve:

\begin{itemize}
\tightlist
\item
  every exact Agent proposal,
\item
  Persona/run provenance,
\item
  recognition consensus,
\item
  classification variance,
\item
  intra-Persona instability,
\item
  evidence,
\item
  and unresolved semantic differences.
\end{itemize}

Semantic uncertainty shall remain reviewable.

It shall not be converted into a technical Processing failure merely because the system cannot safely merge proposals.

\begin{center}\rule{0.5\linewidth}{0.5pt}\end{center}

\hypertarget{5-persona-voting-and-repeated-runs}{%
\paragraph{5. Persona voting and repeated runs}\label{5-persona-voting-and-repeated-runs}}

A Persona contributes at most one recognition vote to one local engineering subject.

Repeated runs of the same Persona measure stability and shall not create additional independent votes.

The existing semantic-consensus principles for distinct-Persona voting and intra-Persona stability shall be reused where compatible.

\begin{center}\rule{0.5\linewidth}{0.5pt}\end{center}

\hypertarget{6-cross-unit-semantic-synthesis}{%
\paragraph{6. Cross-Unit Semantic Synthesis}\label{6-cross-unit-semantic-synthesis}}

After all Source Analysis Units of the applicable source scope have been processed, the locally consolidated engineering subjects shall be synthesized across Source Analysis Units.

The purpose is to establish project/source-level semantic continuity.

Conceptually:

\begin{verbatim}
SAU-000003 / "remote expert"
SAU-000007 / "expert"
SAU-000011 / "remote specialist"
              |
              +--> Engineering Subject: Remote Expert
\end{verbatim}

Cross-unit synthesis operates on locally consolidated subjects rather than on the complete pool of raw Agent proposals.

This substantially reduces comparison scope while preserving all contributing source and Agent provenance.

Cross-unit synthesis shall be conservative:

\begin{itemize}
\tightlist
\item
  explicit equivalence may consolidate subjects,
\item
  explicit distinction keeps subjects separate,
\item
  uncertainty does not authorize automatic merging,
\item
  unresolved cases remain available to Human Review.
\end{itemize}

\begin{center}\rule{0.5\linewidth}{0.5pt}\end{center}

\hypertarget{7-relationship-synthesis-and-endpoint-rebinding}{%
\paragraph{7. Relationship synthesis and endpoint rebinding}\label{7-relationship-synthesis-and-endpoint-rebinding}}

Relationships are first derived with exact source-local proposal provenance.

After cross-unit Element synthesis, Relationship endpoints shall be rebound to the resulting synthesized Element subjects where exactly resolvable.

Unresolved or ambiguous endpoint bindings remain explicit Human Review findings.

Relationship synthesis shall not invent an endpoint or relationship merely to complete a graph.

A degraded or uncertain Relationship interpretation shall remain reviewable provided source evidence and provenance remain trustworthy.

Actual provenance, identity, or artifact-integrity corruption remains a hard Processing failure.

\begin{center}\rule{0.5\linewidth}{0.5pt}\end{center}

\hypertarget{8-compiled-human-review}{%
\paragraph{8. Compiled Human Review}\label{8-compiled-human-review}}

Human Review shall operate on compiled Engineering Subjects rather than raw Agent proposals.

The primary review subject shall expose, as applicable:

\begin{itemize}
\tightlist
\item
  synthesized Engineering Subject,
\item
  all contributing Source Analysis Units,
\item
  exact source excerpts,
\item
  all contributing Persona interpretations,
\item
  Persona recognition consensus,
\item
  classification variance,
\item
  intra-Persona instability,
\item
  relationship findings,
\item
  open questions,
\item
  and exact immutable provenance.
\end{itemize}

The reviewer shall not be forced to accept or reject each independent Agent proposal when multiple proposals are merely alternative interpretations of the same engineering subject.

Human Review remains the engineering authorization boundary.

\begin{center}\rule{0.5\linewidth}{0.5pt}\end{center}

\hypertarget{9-separation-of-identities}{%
\paragraph{9. Separation of identities}\label{9-separation-of-identities}}

The following identities represent different concepts and shall not be collapsed:

\texttt{SourceAnalysisUnit}

\begin{itemize}
\tightlist
\item
  Which exact source context is being jointly assessed?
\end{itemize}

\texttt{EvidenceReference}

\begin{itemize}
\tightlist
\item
  Which exact source statement or evidence occurrence supports a proposal?
\end{itemize}

\texttt{AgentProposal}

\begin{itemize}
\tightlist
\item
  What did one exact Persona run propose?
\end{itemize}

\texttt{LocalSemanticSubject}

\begin{itemize}
\tightlist
\item
  Which engineering subject is recognized within one Source Analysis Unit?
\end{itemize}

\texttt{SynthesizedEngineeringSubject}

\begin{itemize}
\tightlist
\item
  Which engineering subject persists across Source Analysis Units?
\end{itemize}

\texttt{InformationUnit}

\begin{itemize}
\tightlist
\item
  Which reviewed semantic claim is eligible for downstream engineering use?
\end{itemize}

Review workspace identity remains separate from synthesized engineering-subject identity.

\texttt{SES-......} and \texttt{SRS-......} remain D4 synthesized-domain identifiers and shall be preserved in Cross-Unit Synthesis artifacts and exact evidence locators.

The existing Review Workspace and Approved Input \texttt{stable\_subject\_key} contract remains lower-case and namespaced. D5 therefore maps synthesized identities deterministically and one-to-one:

\begin{verbatim}
SES-000001 -> semantic:element:ses-000001
SRS-000001 -> semantic:relationship:srs-000001
\end{verbatim}

This mapping does not create a new semantic subject and does not replace the D4 synthesized identifier. It is the downstream Review/Approved-Input identity representation of the same synthesized authority subject.

This separation is required for traceability, comparability, and safe Human-in-the-Loop processing.

\begin{center}\rule{0.5\linewidth}{0.5pt}\end{center}

\hypertarget{10-comparator-failure-behavior}{%
\paragraph{10. Comparator failure behavior}\label{10-comparator-failure-behavior}}

Comparator execution remains advisory semantic processing.

A malformed, unavailable, incomplete, or contract-invalid comparator response shall never authorize an unsafe merge.

Fallback behavior must preserve proposals and evidence.

However, a comparator fallback that would create an unacceptably large Human Review workload shall be surfaced as a processing-quality finding rather than silently treated as successful semantic consolidation.

The workflow may still reach Human Review when integrity is preserved, but the result shall remain visibly degraded.

\begin{center}\rule{0.5\linewidth}{0.5pt}\end{center}

\hypertarget{11-implementation-slices}{%
\paragraph{11. Implementation slices}\label{11-implementation-slices}}

Implementation shall proceed in controlled slices.

\hypertarget{d1--source-analysis-unit-contract}{%
\subparagraph{D1 --- Source Analysis Unit contract}\label{d1--source-analysis-unit-contract}}

Introduce:

\begin{itemize}
\tightlist
\item
  immutable Source Analysis Unit type,
\item
  deterministic identifiers,
\item
  source-anchor validation,
\item
  Source Projection binding,
\item
  deterministic fingerprint,
\item
  read/persistence contract,
\item
  focused tests.
\end{itemize}

No Agent orchestration changes occur in D1.

\hypertarget{d2--unit-bound-persona-execution}{%
\subparagraph{D2 --- Unit-bound Persona execution}\label{d2--unit-bound-persona-execution}}

Change Phase-F orchestration so all configured Persona runs process one explicit Source Analysis Unit at a time.

Persist the Source Analysis Unit binding in every relevant Agent result.

\hypertarget{d3--source-anchored-local-semantic-consolidation}{%
\subparagraph{D3 --- Source-anchored local semantic consolidation}\label{d3--source-anchored-local-semantic-consolidation}}

Consolidate Element and Relationship proposals inside one Source Analysis Unit.

Reuse existing semantic-extraction and semantic-consensus contracts where compatible.

\hypertarget{d4--cross-unit-synthesis-and-relationship-rebinding}{%
\subparagraph{D4 --- Cross-unit synthesis and relationship rebinding}\label{d4--cross-unit-synthesis-and-relationship-rebinding}}

Synthesize local subjects across Source Analysis Units and rebind Relationship endpoints conservatively.

\hypertarget{d5--human-review-integration-and-empirical-retest}{%
\subparagraph{D5 --- Human Review integration and empirical retest}\label{d5--human-review-integration-and-empirical-retest}}

Project synthesized Engineering Subjects into the existing Human Review workflow.

Repeat the controlled WP12 Source-01 empirical test and compare:

\begin{itemize}
\tightlist
\item
  raw Agent proposals,
\item
  local subjects,
\item
  cross-unit synthesized subjects,
\item
  Human Review items,
\item
  Persona comparison coverage,
\item
  retained evidence/provenance.
\end{itemize}

\begin{center}\rule{0.5\linewidth}{0.5pt}\end{center}

\hypertarget{consequences-11}{%
\subsubsection{Consequences}\label{consequences-11}}

\hypertarget{positive-consequences-6}{%
\paragraph{Positive consequences}\label{positive-consequences-6}}

\begin{itemize}
\tightlist
\item
  Personas are explicitly aligned on the same source context.
\item
  Source provenance becomes an orchestration input rather than only a post-hoc traceability attribute.
\item
  Semantic comparison becomes bounded and local before project-wide synthesis.
\item
  Global LLM partition problems are materially reduced.
\item
  Persona consensus becomes easier to interpret.
\item
  Repeated Persona runs remain stability evidence rather than extra votes.
\item
  Human Review operates on engineering subjects rather than raw Agent outputs.
\item
  Exact Agent wording and source evidence remain fully traceable.
\item
  Existing Source Projection, semantic extraction, semantic consensus, and Human Review concepts can be reused.
\item
  Comparator degradation no longer automatically explodes the primary review model without an explicit quality signal.
\end{itemize}

\hypertarget{trade-offs-5}{%
\paragraph{Trade-offs}\label{trade-offs-5}}

\begin{itemize}
\tightlist
\item
  Phase-F orchestration becomes more granular.
\item
  More execution records are created because Agent runs are bound to Source Analysis Units.
\item
  Cross-unit synthesis introduces an additional explicit semantic layer.
\item
  Runtime and LLM call count may increase if every unit is processed independently.
\item
  Batching and caching may be required later for performance, but they must not weaken source-unit identity or Persona comparability.
\item
  Existing broad-context prompts and reports require adaptation.
\end{itemize}

\begin{center}\rule{0.5\linewidth}{0.5pt}\end{center}

\hypertarget{rejected-alternatives-5}{%
\subsubsection{Rejected alternatives}\label{rejected-alternatives-5}}

\hypertarget{only-repair-the-global-semantic-comparator}{%
\paragraph{Only repair the global semantic comparator}\label{only-repair-the-global-semantic-comparator}}

Rejected as the sole solution.

The WP12 retest showed concrete parser and partition failures, but even a technically valid global comparator would still have to reconcile a large pool of independently generated broad-context proposals after semantic alignment had already been lost.

Comparator robustness should still improve, but it is not sufficient as the primary architecture.

\hypertarget{use-informationunit-as-the-pre-agent-source-unit}{%
\paragraph{\texorpdfstring{Use \texttt{InformationUnit} as the pre-Agent source unit}{Use InformationUnit as the pre-Agent source unit}}\label{use-informationunit-as-the-pre-agent-source-unit}}

Rejected.

An Information Unit is already an interpreted semantic claim with engineering classification, epistemic state, extraction provenance, confidence, and review-oriented semantics.

Using it as the pre-interpretation orchestration anchor would conflate source identity with engineering interpretation.

\hypertarget{let-each-persona-create-its-own-source-unit-identifiers}{%
\paragraph{Let each Persona create its own source-unit identifiers}\label{let-each-persona-create-its-own-source-unit-identifiers}}

Rejected.

This recreates the observed problem because identical identifier strings from different Agent outputs do not establish shared source identity.

Source Analysis Unit identity must be created before Persona interpretation and must be persona-independent.

\hypertarget{human-review-after-every-source-analysis-unit}{%
\paragraph{Human Review after every Source Analysis Unit}\label{human-review-after-every-source-analysis-unit}}

Rejected as the default workflow.

It would create excessive interaction overhead and prevent a coherent compiled review of cross-unit semantic continuity.

Human Review occurs after local consolidation and cross-unit synthesis.

\hypertarget{pure-string-matching-as-semantic-identity}{%
\paragraph{Pure string matching as semantic identity}\label{pure-string-matching-as-semantic-identity}}

Rejected.

Lexical normalization is useful evidence but cannot establish semantic equivalence across legitimate wording variation.

\hypertarget{unrestricted-embedding-or-llm-similarity-as-automatic-authority}{%
\paragraph{Unrestricted embedding or LLM similarity as automatic authority}\label{unrestricted-embedding-or-llm-similarity-as-automatic-authority}}

Rejected.

Similarity may support semantic comparison, but uncertainty must not silently authorize merges or replace Human Review.

\begin{center}\rule{0.5\linewidth}{0.5pt}\end{center}

\hypertarget{related-decisions-2}{%
\subsubsection{Related decisions}\label{related-decisions-2}}

\begin{itemize}
\tightlist
\item
  ADR-016 --- Human Review Workspace and Approved Input Promotion Architecture
\item
  ADR-017 --- Simple-by-Default Interaction and Progressive Disclosure
\item
  ADR-025 --- Semantic Proposal Consolidation and Persona-Aware Consensus
\end{itemize}

\begin{center}\rule{0.5\linewidth}{0.5pt}\end{center}

\hypertarget{wp12-test-status-impact}{%
\subsubsection{WP12 test status impact}\label{wp12-test-status-impact}}

ADR-026 does not mark BLK-003 as closed.

Current interpretation:

\begin{verbatim}
BLK-003.1 processing robustness
PASS at focused test / controlled processing level

BLK-003 semantic consolidation effectiveness
OPEN

POST-BLK-003.1 empirical result
Processing reached Human Review
but semantic consolidation degraded to singletons
and produced 112 Review Items
\end{verbatim}

WP12-E2E-DRY-001 remains the same interrupted controlled dry run.

The run shall not be restarted solely to hide the observed findings.

After implementation and focused validation of D1-D5, the affected Source-01 processing path shall be empirically retested before WP12 continues through the remaining blocked stages.

\begin{center}\rule{0.5\linewidth}{0.5pt}\end{center}

\hypertarget{implementation-note-2}{%
\subsubsection{Implementation note}\label{implementation-note-2}}

No D1-D5 production implementation shall precede acceptance of this ADR.

The implementation shall reuse existing Source Projection, source-anchor, semantic extraction, semantic consensus, project isolation, fingerprint, Human Review, and immutable provenance contracts where compatible.

The architecture shall preserve the distinction between:

\begin{itemize}
\tightlist
\item
  deterministic source segmentation,
\item
  Agent interpretation,
\item
  semantic consolidation,
\item
  Human Review,
\item
  and downstream engineering authority.
\end{itemize}

\textbackslash newpage

\hypertarget{adr-027--source-grounded-evidence-detection-and-persona-interpretation-architecture}{%
\subsection{ADR-027 --- Source-Grounded Evidence Detection and Persona Interpretation Architecture}\label{adr-027--source-grounded-evidence-detection-and-persona-interpretation-architecture}}

Status

\textbf{Accepted --- responsibility-boundary recovery accepted after Source-to-Human-Review audit. Implementation remains incremental and regression-controlled.}

Accepted

2026-08-21

Date

2026-08-20

\begin{center}\rule{0.5\linewidth}{0.5pt}\end{center}

\hypertarget{context-12}{%
\subsubsection{Context}\label{context-12}}

WP-12 end-to-end formative testing showed that the technical authority chain can reach Human Review, but the semantic quality of the review input is not yet acceptable.

Representative real single-source run:

\begin{verbatim}
Project 877791 / RUN-000001
3 personas × 1 run

93 element proposals
41 relationship proposals
134 raw proposals

D3: 70 elements + 39 relationships = 109 subjects
D4: 70 elements + 39 relationships = 109 subjects
Human Review: 110 items = 70 Elements + 39 Relationships + 1 Open Question
\end{verbatim}

The D4-to-Review routing is technically operational. The formative result exposed:

\begin{enumerate}
\def\labelenumi{\arabic{enumi}.}
\tightlist
\item
  source purity regression --- processing/orchestration/task metadata can become engineering subjects,
\item
  model relevance regression --- candidate content can be generic rather than concrete model-relevant engineering information,
\item
  subject multiplication --- persona-generated candidate identities can create parallel subject sets which later stages attempt to consolidate.
\end{enumerate}

Review Item count is therefore a diagnostic symptom, not the optimization target.

A deliberately simple external benchmark on \texttt{legacy/demo/wp12/01\_product\_overview.md} showed that a general-purpose LLM can identify a small set of source-grounded, model-relevant passages with a direct task and strict source boundary. The benchmark returned eight major source-grounded findings. It is qualitative diagnostic evidence only; it is not a gold standard and shall not become model authority.

ADR-011 already established several still-valid foundations: deterministic Source Projection, source anchors/excerpts, source-traceable Information Units, \texttt{engineering\_source} vs \texttt{context\_only}, semantic extraction, persona perspectives, consensus/variance, repeated runs as stability evidence rather than votes, ontology/terminology boundaries, and Human Review separation.

The recovery shall therefore first simplify and realign the existing architecture rather than add another post-processing contract.

\begin{center}\rule{0.5\linewidth}{0.5pt}\end{center}

\hypertarget{relationship-to-adr-025-and-adr-026}{%
\subsubsection{Relationship to ADR-025 and ADR-026}\label{relationship-to-adr-025-and-adr-026}}

ADR-025 --- Semantic Proposal Consolidation and Persona-Aware Consensus and ADR-026 --- Source-Anchored Multi-Persona Interpretation and Cross-Unit Semantic Synthesis remain accepted historical architecture decisions and document the evolution that led to this recovery.

ADR-027 refines the upstream responsibility sequence exposed by the subsequent empirical WP-12 test. In particular, source anchoring alone is not sufficient if Personas can still independently create engineering-subject populations that must later be reconciled.

ADR-027 therefore moves the common source-grounded evidence space ahead of Persona interpretation. Existing ADR-025/ADR-026 components remain reusable only where their responsibility is compatible with this clarified sequencing.

\hypertarget{decision-12}{%
\subsubsection{Decision}\label{decision-12}}

\hypertarget{1-separate-evidence-detection-from-persona-interpretation}{%
\paragraph{1. Separate Evidence Detection from Persona Interpretation}\label{1-separate-evidence-detection-from-persona-interpretation}}

\begin{verbatim}
Evidence Detection
"Which passages of the Engineering Source contain
potentially model-relevant engineering information?"
        ↓
Persona Interpretation
"What does this already identified source-grounded
engineering information mean?"
\end{verbatim}

Detection establishes the common source-grounded evidence space. Persona interpretation operates on that common evidence space.

\hypertarget{1a-evidence-detection-is-a-specialized-persona-independent-agent-task}{%
\paragraph{1a. Evidence Detection Is a Specialized Persona-Independent Agent Task}\label{1a-evidence-detection-is-a-specialized-persona-independent-agent-task}}

Evidence Detection is an explicit LLM-backed preparation task with one narrowly bounded responsibility:

\begin{verbatim}
Engineering Source
+ project/source identity
+ source projection / processing scope
+ reference examples and detection guidance
        ↓
Specialized Evidence Detection Agent
        ↓
exact source-grounded evidence spans
\end{verbatim}

The detector may use curated repository examples, modeling guidance and other reference knowledge to learn what kinds of engineering information are worth marking. Those inputs remain \texttt{context\_only} guidance. They shall never become positive Project evidence.

The detector shall not:

\begin{itemize}
\tightlist
\item
  create requirements, actors, functions, interfaces or other model elements,
\item
  perform architecture derivation,
\item
  assign final engineering meaning on behalf of interpretation personas,
\item
  create multiple evidence identities merely because several personas or runs will later consume the evidence.
\end{itemize}

Detector output must be independently verifiable against the Source Projection. At minimum, every accepted evidence span is bound by exact source anchors and an exact source excerpt.

\hypertarget{1b-source-registration-and-source-preparation-are-architecturally-separate}{%
\paragraph{1b. Source Registration and Source Preparation Are Architecturally Separate}\label{1b-source-registration-and-source-preparation-are-architecturally-separate}}

The Source Registry remains the immutable authority for Project Source identity. LLM execution shall not be added to the Source Registry contract itself.

The system distinguishes:

\begin{verbatim}
REGISTER SOURCE
    ↓
PREPARE SOURCE
    ↓
Source Projection
    ↓
Evidence Detection
    ↓
SOURCE READY FOR INTERPRETATION
\end{verbatim}

The UI may present registration and preparation as one convenient user action. That is a UX choice, not a collapse of architectural responsibilities.

Evidence Detection is version/configuration dependent. Its persisted result must therefore remain traceable to the immutable Source version and to the detector configuration used to create it. Changing persona configuration alone shall not require Evidence Detection to be repeated.

\hypertarget{2-reference-knowledge-is-guidance-never-engineering-evidence}{%
\paragraph{2. Reference Knowledge Is Guidance, Never Engineering Evidence}\label{2-reference-knowledge-is-guidance-never-engineering-evidence}}

Reference knowledge may include Apollo 11 reference material, SysML v2/KerML references, modeling guidance, framework definitions, agent roles/personas, recipes, project principles, terminology and ontology references.

It may answer:

\begin{verbatim}
"What kinds of engineering information should I look for?"
\end{verbatim}

It shall never answer:

\begin{verbatim}
"What engineering information exists in this Project Source?"
\end{verbatim}

Only a registered \texttt{engineering\_source} may provide positive engineering evidence. Prompt text, task instructions, recipes, orchestration manifests, processing metadata and reference models shall never become positive Project engineering evidence merely because they were visible to an LLM.

\hypertarget{3-source-grounded-evidence-is-the-common-discussion-basis}{%
\paragraph{3. Source-Grounded Evidence Is the Common Discussion Basis}\label{3-source-grounded-evidence-is-the-common-discussion-basis}}

Every detected Evidence Unit shall remain bound to the Engineering Source through at least:

\begin{verbatim}
project_id
source_id
source_projection_id
source_anchor(s)
exact source_excerpt
\end{verbatim}

This is the digital equivalent of a text-marker highlight. Evidence identity is source-grounded. An agent-generated candidate name shall not be the primary identity used to determine whether personas are discussing the same source information.

\hypertarget{4-personas-interpret-the-same-evidence}{%
\paragraph{4. Personas Interpret the Same Evidence}\label{4-personas-interpret-the-same-evidence}}

\begin{verbatim}
Evidence E-xxx
      │
 ┌────┼────┐
 ▼    ▼    ▼
P1   P2   P3
      │
      ▼
Consensus / Variance
\end{verbatim}

Personas may disagree about relevance confidence, professional interpretation, information type, modality, epistemic class, terminology, model relevance, uncertainty or missing information.

Their disagreement is first-class engineering information. It shall not automatically create multiple independent Review Subjects.

\hypertarget{5-persona-and-run-counts-shall-not-multiply-engineering-subjects}{%
\paragraph{5. Persona and Run Counts Shall Not Multiply Engineering Subjects}\label{5-persona-and-run-counts-shall-not-multiply-engineering-subjects}}

The number of engineering subjects shall be driven by distinct source-grounded engineering information, not by persona/run count.

\begin{verbatim}
Review Subjects ≠ Source Subjects × Personas × Runs
\end{verbatim}

Repeated runs measure intra-persona stability and are not additional independent votes. Each persona contributes at most one effective vote to one evidence-centered consensus assessment.

\hypertarget{6-evidence-detection-may-express-uncertainty}{%
\paragraph{6. Evidence Detection May Express Uncertainty}\label{6-evidence-detection-may-express-uncertainty}}

Evidence detection may classify a passage as:

\begin{verbatim}
relevant
not_relevant
uncertain
\end{verbatim}

One passage may contain multiple independently reviewable claims. The key constraint is that their identities remain source-grounded and are not multiplied merely by persona execution.

\hypertarget{7-review-item-count-is-a-diagnostic-not-an-optimization-objective}{%
\paragraph{7. Review Item Count Is a Diagnostic, Not an Optimization Objective}\label{7-review-item-count-is-a-diagnostic-not-an-optimization-objective}}

Evaluation shall prioritize:

\begin{enumerate}
\def\labelenumi{\arabic{enumi}.}
\tightlist
\item
  source purity,
\item
  model relevance,
\item
  source grounding,
\item
  clear interpretation/classification,
\item
  useful consensus/variance,
\item
  Human Review usability.
\end{enumerate}

A high Review Item count can be correct if the Source contains many independent claims. A count increase caused only by adding personas/runs indicates an architectural or consolidation defect.

\hypertarget{8-terminology-and-ontology-mapping-are-supporting-services}{%
\paragraph{8. Terminology and Ontology Mapping Are Supporting Services}\label{8-terminology-and-ontology-mapping-are-supporting-services}}

Semantic normalization, terminology mapping and ontology alignment may improve comparability and cross-source consistency. They shall not create new positive Engineering Source evidence or new subjects without a source-grounded basis.

\hypertarget{9-human-engineering-review-precedes-architecture-derivation}{%
\paragraph{9. Human Engineering Review Precedes Architecture Derivation}\label{9-human-engineering-review-precedes-architecture-derivation}}

\begin{verbatim}
Engineering Source
→ Deterministic Source Projection
→ Source-Grounded Evidence Detection
→ Persona Interpretation
→ Consensus / Variance
→ optional Semantic Normalization / Ontology Alignment
→ Human Review
→ Approved Engineering Information
→ Architecture Derivation
→ SysML v2 Generation
\end{verbatim}

AI-generated interpretation is evidence for Human Review, not Approved Engineering Information.

\hypertarget{9a-architecture-derivation-may-again-use-multiple-personas}{%
\paragraph{9a. Architecture Derivation May Again Use Multiple Personas}\label{9a-architecture-derivation-may-again-use-multiple-personas}}

Multi-persona processing remains valuable after Human Engineering Review. However, its responsibility changes.

Before Human Review, personas answer:

\begin{verbatim}
"What does this same source-grounded Evidence mean?"
\end{verbatim}

After approval, model-derivation personas may answer:

\begin{verbatim}
"How can this Approved Engineering Information be represented
as a coherent system/model architecture?"
\end{verbatim}

Target sequence:

\begin{verbatim}
Approved Engineering Information
        ↓
 ┌──────┼──────┐
 ▼      ▼      ▼
Model  Model  Model
Persona Persona Persona
 └──────┼──────┘
        ↓
Architecture / Model Candidate Derivation
        ↓
Semantic Comparison / Candidate Consolidation
        ↓
Model Candidate Review
        ↓
Approved Internal Model
        ↓
SysML v2 Generation
\end{verbatim}

Whether this downstream derivation phase shall reuse existing personas or introduce new task-specific modeling personas is intentionally deferred until that stage is implemented and evaluated. The architecture requires the persona branch; it does not prematurely freeze its future persona set.

\hypertarget{10-existing-p9-responsibilities-have-no-grandfathering-protection}{%
\paragraph{10. Existing P9 Responsibilities Have No Grandfathering Protection}\label{10-existing-p9-responsibilities-have-no-grandfathering-protection}}

Every material P9/D3/D4 responsibility shall be classified as:

\begin{verbatim}
KEEP
MOVE DOWNSTREAM
REDUCE TO ADAPTER
BYPASS
RETIRE
\end{verbatim}

A responsibility shall not be retained merely because downstream code currently expects it.

\begin{center}\rule{0.5\linewidth}{0.5pt}\end{center}

\hypertarget{thesis-relevant-architecture-finding}{%
\subsubsection{Thesis-Relevant Architecture Finding}\label{thesis-relevant-architecture-finding}}

The WP-12 formative run exposed a responsibility-ordering defect rather than a simple clustering or prompt-quality problem.

\hypertarget{previous-responsibility-sequence}{%
\paragraph{Previous responsibility sequence}\label{previous-responsibility-sequence}}

In the previous active path, each interpretation persona could independently:

\begin{enumerate}
\def\labelenumi{\arabic{enumi}.}
\tightlist
\item
  select what it considered meaningful engineering information,
\item
  assign persona-local \texttt{source\_info\_id} identities,
\item
  interpret/classify that selected information,
\item
  feed a derivation step that created model-oriented candidate elements and relationships before the first Human Engineering Review.
\end{enumerate}

Downstream D3/D4 semantic consolidation then attempted to determine which of those independently created candidate populations referred to the same engineering subject.

This caused a fundamental comparability problem:

\begin{verbatim}
Persona A selection ≠ Persona B selection ≠ Persona C selection
\end{verbatim}

Therefore differences between persona outputs could not be interpreted cleanly as professional interpretation variance. A difference could instead be caused by different source selection, segmentation, abstraction, naming or early modeling decisions.

\hypertarget{observed-consequences}{%
\paragraph{Observed consequences}\label{observed-consequences}}

The architecture produced several characteristic symptoms:

\begin{itemize}
\tightlist
\item
  \textbf{persona-driven subject multiplication} --- additional personas could create additional subject populations rather than additional views on one subject;
\item
  \textbf{unstable engineering-subject identity} --- candidate names/types became part of the mechanism used to infer common subjects;
\item
  \textbf{false or ambiguous variance} --- source-selection differences and genuine interpretation differences were mixed;
\item
  \textbf{excessive downstream semantic repair} --- later LLM consolidation had to reconstruct commonality that had not been established upstream;
\item
  \textbf{premature model derivation} --- model structure was proposed from unreviewed interpretation;
\item
  \textbf{Human Review overload} --- the reviewer received a large model-oriented subject population rather than a compact set of source-grounded engineering information;
\item
  \textbf{technical success without engineering effectiveness} --- the processing chain could reach Human Review correctly while the review content remained unsuitable for engineering use.
\end{itemize}

The representative real single-source run made the effect visible:

\begin{verbatim}
3 personas × 1 run
93 element proposals
41 relationship proposals
134 raw proposals
109 D3/D4 semantic subjects
110 Human Review items
\end{verbatim}

These counts are evidence of the symptom, not a target to optimize directly.

\hypertarget{architectural-correction}{%
\paragraph{Architectural correction}\label{architectural-correction}}

ADR-027 introduces an explicit boundary:

\begin{verbatim}
DETECTION
Which source passages are potentially engineering-relevant?
        ↓
fixed source-grounded Evidence identity
        ↓
INTERPRETATION
What does this same Evidence mean from different professional perspectives?
\end{verbatim}

Only after Evidence identity is fixed do personas branch. This makes consensus/variance meaningful because each persona is now discussing the same source-grounded object.

After Human Engineering Review and approval, a second persona branch may be used for architecture/model derivation. The key change is therefore not the removal of personas or LLM reasoning, but the placement of those responsibilities on the correct side of the engineering approval boundary.

\hypertarget{thesis-interpretation}{%
\paragraph{Thesis interpretation}\label{thesis-interpretation}}

This is a first-class formative engineering result of the prototype:

\begin{quote}
Source grounding alone is insufficient for reliable multi-persona semantic processing when evidence detection and persona-specific interpretation share the same responsibility boundary. A common evidence identity must be established before persona variance can be interpreted as engineering variance.
\end{quote}

This finding should be carried into the thesis as an observed failure, root-cause analysis, architectural correction and subsequent evaluation target.

\begin{center}\rule{0.5\linewidth}{0.5pt}\end{center}

\hypertarget{relationship-to-adr-011}{%
\subsubsection{Relationship to ADR-011}\label{relationship-to-adr-011}}

ADR-027 does not replace ADR-011 wholesale. The following remain valid unless separately changed:

\begin{itemize}
\tightlist
\item
  deterministic Source Projection,
\item
  immutable Source identity,
\item
  source anchors/excerpts,
\item
  Information Unit traceability,
\item
  \texttt{engineering\_source} / \texttt{context\_only},
\item
  explicit epistemic classification,
\item
  independent persona perspectives,
\item
  deterministic consensus/variance,
\item
  repeated runs as stability evidence rather than votes,
\item
  explicit terminology/ontology mapping,
\item
  Human Review separation.
\end{itemize}

ADR-027 sharpens the sequencing: a common source-grounded evidence space is established before persona interpretations become downstream Review Subjects.

\begin{center}\rule{0.5\linewidth}{0.5pt}\end{center}

\hypertarget{qualitative-benchmark}{%
\subsubsection{Qualitative Benchmark}\label{qualitative-benchmark}}

The 2026-08-20 external diagnostic benchmark used exactly:

\begin{verbatim}
legacy/demo/wp12/01_product_overview.md
\end{verbatim}

A simple source-bounded prompt returned eight major findings around:

\begin{itemize}
\tightlist
\item
  remote microscope collaboration capability,
\item
  microscope operator and remote expert,
\item
  workstation and remote client context,
\item
  remote consultation purpose,
\item
  live-image observation,
\item
  temporary remote control subject to operator permission,
\item
  operator responsibility/controller transparency,
\item
  session-information retention,
\item
  explicitly unspecified protocol/deployment/performance/latency/regulatory information and open product questions.
\end{itemize}

The benchmark also showed that a general-purpose LLM can jump too early to modeling choices such as State Machine/Guard. Therefore Turing shall distinguish engineering meaning from later model-structure decisions.

The benchmark is diagnostic only and shall not become Project authority.

\begin{center}\rule{0.5\linewidth}{0.5pt}\end{center}

\hypertarget{recovery-audit}{%
\subsubsection{Recovery Audit}\label{recovery-audit}}

Before implementation changes, trace the active Source-to-Human-Review path end-to-end.

For every material component classify:

\begin{verbatim}
KEEP
MOVE
REDUCE
BYPASS
REMOVE / RETIRE
\end{verbatim}

The audit shall answer:

\begin{enumerate}
\def\labelenumi{\arabic{enumi}.}
\tightlist
\item
  What exact content is supplied to each LLM call?
\item
  Which inputs are Engineering Source vs reference/context/instruction?
\item
  Where are source-grounded anchors first established?
\item
  Where is a new subject identity first created?
\item
  Where do personas begin to create independent subject sets?
\item
  Which P4 semantic-extraction / semantic-consensus components are active in the live path?
\item
  Why does current P9 derivation appear before the first Engineering Human Review?
\item
  Which D3/D4 responsibilities remain necessary once evidence-centered subject identity is restored?
\item
  Which contracts are current authority, and which are legacy compatibility ballast?
\item
  Can the corrected path reuse existing Source Projection, Information Unit, consensus, Human Review and Approved Input infrastructure?
\end{enumerate}

No new end-to-end LLM run should be used as improvement evidence until this audit identifies and corrects the source-contamination / subject-multiplication path.

\begin{center}\rule{0.5\linewidth}{0.5pt}\end{center}

\hypertarget{catia--presentation-rule}{%
\subsubsection{CATIA / Presentation Rule}\label{catia--presentation-rule}}

Intended future high-level System Behavior:

\begin{verbatim}
REFERENCE KNOWLEDGE
        │ guidance only
        ▼
ENGINEERING SOURCE
        ↓
Register Source
        ↓
Prepare Source
        ↓
Deterministic Source Projection
        ↓
Specialized Evidence Detection Agent
        ↓
Source-Grounded Evidence
        ↓
 ┌──────┼──────┐
 ▼      ▼      ▼
P1     P2     P3
Interpret / Classify the SAME Evidence
 └──────┼──────┘
        ↓
Consensus / Variance
        ↓
optional Semantic Normalization / Ontology Alignment
        ↓
Human Engineering Review
        ↓
Approved Engineering Information
        ↓
 ┌──────┼──────┐
 ▼      ▼      ▼
Model  Model  Model
Persona Persona Persona
 └──────┼──────┘
        ↓
Architecture / Model Candidate Derivation
        ↓
Candidate Consolidation
        ↓
Model Candidate Review
        ↓
Approved Internal Model
        ↓
SysML v2
\end{verbatim}

Add this to the authoritative CATIA model only after the implementation is aligned and ADR-027 is accepted. Do not model it merely to improve the Monday presentation.

\begin{center}\rule{0.5\linewidth}{0.5pt}\end{center}

\hypertarget{status-and-next-decision}{%
\subsubsection{Status and Next Decision}\label{status-and-next-decision}}

ADR-027 is \texttt{Accepted} as the governing architecture-recovery direction.

The completed Source-to-first-Human-Review audit established that the primary BLK-003 root cause is the mixing of evidence detection, persona interpretation and pre-review model derivation.

Implementation proceeds incrementally:

\begin{verbatim}
R1  Accept/document ADR-027 and thesis finding
R2  Introduce source-grounded Evidence contract and persistence
R3  Add specialized Evidence Detection Agent and Source Preparation
R4  Make interpretation personas consume the same persisted Evidence
R5  Move model derivation behind Human Engineering Approval and evaluate
    the required downstream modeling personas
\end{verbatim}

After at least one material corrected processing path exists:

\begin{itemize}
\tightlist
\item
  run focused automated tests,
\item
  run one bounded real single-source LLM test,
\item
  evaluate source purity, model relevance, source grounding, persona behavior on shared Evidence, consensus/variance and Human Review usability,
\item
  retain a successful persisted run as the preferred Monday demo reference.
\end{itemize}

\textbackslash newpage

\hypertarget{adr-028a--model-derivation-mode-and-review-escalation-architecture}{%
\subsection{ADR-028a --- Model Derivation Mode and Review Escalation Architecture}\label{adr-028a--model-derivation-mode-and-review-escalation-architecture}}

\hypertarget{status-9}{%
\subsubsection{Status}\label{status-9}}

Accepted

\hypertarget{date-9}{%
\subsubsection{Date}\label{date-9}}

2026-08-21

\hypertarget{context-13}{%
\subsubsection{Context}\label{context-13}}

The corrected WP12 architecture establishes source-grounded Evidence before persona interpretation and introduces Human Engineering Review before target model derivation.

Phase H already consumes only active Approved Inputs and already supports both strict deterministic profile projection and LLM-assisted hybrid projection. The Human Model Candidate Review remains mandatory regardless of the selected derivation strategy.

The architecture includes an explicit energy-saving mode and a Human-controlled escalation path:

\begin{itemize}
\tightlist
\item
  model derivation may run without any LLM call,
\item
  LLM assistance is optional rather than mandatory,
\item
  the system may recommend a strategy based on deterministic projection coverage,
\item
  the Human remains free to select the strategy,
\item
  and a rejected deterministic Candidate may be regenerated with LLM assistance without repeating the upstream Engineering Review.
\end{itemize}

A further architectural issue follows from ADR-020 H9-05. The normal hybrid path does not send deterministically mapped Approved Inputs to the LLM. A Candidate that was deterministically mapped successfully but later rejected by the Human Model Review therefore requires an explicit escalation exception: the Approved Inputs supporting that rejected predecessor Candidate must become eligible for LLM re-proposal even if deterministic coverage still reports them as \texttt{mapped}.

\begin{center}\rule{0.5\linewidth}{0.5pt}\end{center}

\hypertarget{decision-13}{%
\subsubsection{Decision}\label{decision-13}}

\hypertarget{r5-01--two-explicit-derivation-modes}{%
\paragraph{R5-01 --- Two explicit derivation modes}\label{r5-01--two-explicit-derivation-modes}}

The supported Human-selectable derivation modes are:

\begin{itemize}
\tightlist
\item
  \texttt{eco\_deterministic}
\item
  \texttt{llm\_assisted}
\end{itemize}

\texttt{eco\_deterministic} performs no LLM request.

\texttt{llm\_assisted} remains deterministic-first but may invoke bounded modeling LLM reasoning for eligible Approved Inputs.

The selected mode does not remove the downstream Model Candidate Review.

\begin{center}\rule{0.5\linewidth}{0.5pt}\end{center}

\hypertarget{r5-02--strategy-recommendation-is-advisory}{%
\paragraph{R5-02 --- Strategy recommendation is advisory}\label{r5-02--strategy-recommendation-is-advisory}}

Before Candidate generation, deterministic projection coverage shall be assessed.

The system may recommend:

\begin{itemize}
\tightlist
\item
  \texttt{eco\_deterministic} when all projectable Approved Inputs are deterministically mapped,
\item
  \texttt{llm\_assisted} when projection remains ambiguous or unmapped,
\item
  or \texttt{llm\_assisted} when a predecessor Candidate has been rejected and explicit Human escalation is requested.
\end{itemize}

The recommendation is advisory. It shall not silently select or approve a derivation strategy on behalf of the Human.

\begin{center}\rule{0.5\linewidth}{0.5pt}\end{center}

\hypertarget{r5-03--eco-mode-remains-fail-closed}{%
\paragraph{R5-03 --- Eco mode remains fail-closed}\label{r5-03--eco-mode-remains-fail-closed}}

Eco mode shall not weaken deterministic profile safety.

If deterministic projection cannot establish the target mapping required to generate a complete Candidate Set, the strict deterministic deriver may fail closed.

The system shall not force a target mapping merely to keep the no-LLM path running.

\begin{center}\rule{0.5\linewidth}{0.5pt}\end{center}

\hypertarget{r5-04--human-model-review-remains-mandatory}{%
\paragraph{R5-04 --- Human Model Review remains mandatory}\label{r5-04--human-model-review-remains-mandatory}}

Candidate Sets generated in either mode shall pass the same Candidate Review authority boundary.

A deterministic result is not approved merely because it was reproducible.

An LLM-assisted result is not approved merely because multiple modeling perspectives agree.

\begin{center}\rule{0.5\linewidth}{0.5pt}\end{center}

\hypertarget{r5-05--review-rejection-may-escalate-to-llm-assisted-regeneration}{%
\paragraph{R5-05 --- Review rejection may escalate to LLM-assisted regeneration}\label{r5-05--review-rejection-may-escalate-to-llm-assisted-regeneration}}

A rejected predecessor Candidate may trigger a successor Candidate Set with:

\begin{itemize}
\tightlist
\item
  \texttt{predecessor\_candidate\_set\_id},
\item
  a non-empty \texttt{regeneration\_reason},
\item
  and derivation mode \texttt{llm\_assisted}.
\end{itemize}

The predecessor Candidate Set remains immutable and traceable.

\begin{center}\rule{0.5\linewidth}{0.5pt}\end{center}

\hypertarget{r5-06--explicit-escalation-eligibility}{%
\paragraph{R5-06 --- Explicit escalation eligibility}\label{r5-06--explicit-escalation-eligibility}}

For review-driven regeneration, the system shall resolve each currently rejected predecessor Candidate to its immutable Approved Input references.

Those Approved Input IDs become explicit escalation targets.

Explicit escalation targets may enter LLM-assisted modeling even when the deterministic resolver still classifies them as \texttt{mapped}.

This is the only accepted exception to the normal ADR-020 rule that uniquely mapped inputs are not sent to the LLM.

The exception requires an explicit Human Model Review rejection bound to the exact predecessor Candidate snapshot.

\begin{center}\rule{0.5\linewidth}{0.5pt}\end{center}

\hypertarget{r5-07--no-upstream-reinterpretation}{%
\paragraph{R5-07 --- No upstream reinterpretation}\label{r5-07--no-upstream-reinterpretation}}

Review escalation shall not return to Source Evidence Detection, pre-review persona interpretation, or Human Engineering Review unless the Human explicitly changes the approved engineering information itself.

The LLM escalation task remains target-model derivation only.

\begin{center}\rule{0.5\linewidth}{0.5pt}\end{center}

\hypertarget{r5-08--modeling-personas-are-llm-mode-only}{%
\paragraph{R5-08 --- Modeling personas are LLM-mode only}\label{r5-08--modeling-personas-are-llm-mode-only}}

The modeling persona team is not executed in Eco mode.

In \texttt{llm\_assisted} mode, eligible Approved Inputs may be assessed by the accepted modeling perspectives:

\begin{itemize}
\tightlist
\item
  rules-focused,
\item
  architecture-focused,
\item
  conservative review.
\end{itemize}

All modeling personas receive the same Approved Input identities and the same profile-controlled target options.

\begin{center}\rule{0.5\linewidth}{0.5pt}\end{center}

\hypertarget{r5-09--modeling-comparison-is-bounded-by-profile-rules}{%
\paragraph{R5-09 --- Modeling comparison is bounded by profile rules}\label{r5-09--modeling-comparison-is-bounded-by-profile-rules}}

Modeling-persona comparison shall operate on profile-controlled mapping results, not free-form semantic subjects.

Agreement may yield one proposed mapping.

Material disagreement shall remain explicit variance / ambiguity.

No comparison step may invent new Approved Engineering Information.

\begin{center}\rule{0.5\linewidth}{0.5pt}\end{center}

\hypertarget{target-lifecycle}{%
\subsubsection{Target lifecycle}\label{target-lifecycle}}

\begin{verbatim}
Approved Engineering Information
        |
        v
Deterministic Projection Coverage
        |
        +--> Recommendation: Eco / LLM-assisted
        |
        v
Human-selected Derivation Mode
   |                         |
   | Eco                     | LLM-assisted
   v                         v
Strict deterministic      deterministic-first
projection               + modeling personas
   |                         |
   +------------+------------+
                v
        Model Candidate Set
                |
                v
        Human Model Review
          |             |
       accepted       rejected
          |             |
          v             v
     Internal Model   Regenerate successor
                        |
                        +--> LLM-assisted escalation
                             of rejected Candidate's
                             Approved Input support
\end{verbatim}

\begin{center}\rule{0.5\linewidth}{0.5pt}\end{center}

\hypertarget{consequences-12}{%
\subsubsection{Consequences}\label{consequences-12}}

\hypertarget{positive-4}{%
\paragraph{Positive}\label{positive-4}}

\begin{itemize}
\tightlist
\item
  zero-LLM model derivation remains available,
\item
  deterministic results still receive Human authorization,
\item
  LLM use becomes targeted and explainable,
\item
  review rejection becomes a first-class feedback mechanism,
\item
  successful upstream Engineering Review is not repeated unnecessarily,
\item
  and predecessor/successor Candidate Sets preserve full traceability.
\end{itemize}

\hypertarget{trade-offs-6}{%
\paragraph{Trade-offs}\label{trade-offs-6}}

\begin{itemize}
\tightlist
\item
  the LLM-assisted executor must distinguish normal unresolved targets from explicit review-escalation targets,
\item
  the recommendation layer is not itself an authority decision,
\item
  and deterministic Eco mode may legitimately fail when the selected profile cannot resolve the approved information.
\end{itemize}

\begin{center}\rule{0.5\linewidth}{0.5pt}\end{center}

\hypertarget{implementation-decomposition-1}{%
\subsubsection{Implementation decomposition}\label{implementation-decomposition-1}}

\begin{verbatim}
R5a  Derivation mode, recommendation and review-escalation contract
R5b  Modeling-persona execution and bounded comparison for LLM-assisted mode
R5c  Guided Workflow integration and full regression / demo path
\end{verbatim}

R5a introduces no new LLM call.

\textbackslash newpage

\hypertarget{adr-028b--context-preserving-canonical-engineering-subject-discovery}{%
\subsection{ADR-028b --- Context-Preserving Canonical Engineering Subject Discovery}\label{adr-028b--context-preserving-canonical-engineering-subject-discovery}}

Status

\textbf{Accepted --- architecture direction approved on 2026-08-23; implementation and live validation pending.}

Date

2026-08-23

\begin{center}\rule{0.5\linewidth}{0.5pt}\end{center}

\hypertarget{context-14}{%
\subsubsection{Context}\label{context-14}}

WP-12 recovery under ADR-027 successfully restored a shared, source-grounded Evidence population before Persona interpretation. The corrected live path demonstrated:

\begin{itemize}
\tightlist
\item
  deterministic Source grounding,
\item
  stable \texttt{EVD-*} identity,
\item
  all Personas operating on the same Evidence population,
\item
  no multiplication of Review Subjects by Persona count,
\item
  successful propagation into Human Review.
\end{itemize}

However, formative Human Review testing exposed a second semantic defect.

The active interpretation contract effectively maps one \texttt{SourceEvidence} object to one interpreted statement and one information type. A single Evidence block can contain multiple distinct engineering subjects and assertions. The result is therefore often a paraphrase of the Evidence block rather than an MBSE-relevant decomposition.

Example:

\begin{verbatim}
The remote expert shall be able to observe the live microscope image.
During the session, the expert may also take temporary control of the microscope
when the operator permits it.
The operator remains responsible for the local session and must be able to
understand who currently controls the microscope.
\end{verbatim}

This source passage contains several independently meaningful engineering subjects, for example:

\begin{verbatim}
Remote Expert
Live Microscope Image Observation
Temporary Microscope Control
Operator Permission
Operator Responsibility
Current Controller Awareness
\end{verbatim}

Treating the whole passage as one semantic subject loses this structure.

A second issue is context starvation. Provenance units were made deliberately small to guarantee exact source grounding. Those small units are useful for identity and traceability, but they are not necessarily the correct context window for LLM interpretation.

The architecture therefore needs to distinguish:

\begin{verbatim}
small deterministic provenance units
from
larger context-preserving interpretation windows
\end{verbatim}

A third issue is repeated mention. The same engineering subject can occur several times in one source. Repeated mentions shall not be interpreted independently by every Persona.

\begin{center}\rule{0.5\linewidth}{0.5pt}\end{center}

\hypertarget{relationship-to-adr-027}{%
\subsubsection{Relationship to ADR-027}\label{relationship-to-adr-027}}

ADR-028 refines ADR-027; it does not reverse it.

ADR-027 remains authoritative that:

\begin{enumerate}
\def\labelenumi{\arabic{enumi}.}
\tightlist
\item
  source-grounded identity is established before Persona variance,
\item
  Personas shall not independently create incomparable subject populations,
\item
  Reference Knowledge is guidance only,
\item
  Human Review precedes model derivation,
\item
  Persona/run counts shall not multiply engineering subjects.
\end{enumerate}

ADR-028 adds one missing level between SourceEvidence and Persona interpretation:

\begin{verbatim}
SourceEvidence / Source Spans
→ Mention Discovery
→ Canonical Engineering Subjects
→ Shared Persona Interpretation
\end{verbatim}

ADR-027's phrase "Personas interpret the same Evidence" is therefore refined to:

\begin{quote}
Personas interpret the same canonical engineering subjects, with all of those subjects bound to exact source-grounded Evidence and adequate surrounding source context.
\end{quote}

\begin{center}\rule{0.5\linewidth}{0.5pt}\end{center}

\hypertarget{decision-14}{%
\subsubsection{Decision}\label{decision-14}}

\hypertarget{1-provenance-unit-and-interpretation-context-are-separate-concepts}{%
\paragraph{1. Provenance Unit and Interpretation Context Are Separate Concepts}\label{1-provenance-unit-and-interpretation-context-are-separate-concepts}}

The system shall not assume that the smallest deterministic source unit is also the best LLM interpretation context.

\begin{verbatim}
Provenance unit
→ small, exact, deterministic, source-bound

Interpretation context
→ larger, readable, context-preserving, still source-only
\end{verbatim}

For a small source document, the interpretation context may be the complete Engineering Source. For larger documents it may be a section or bounded multi-paragraph window.

Every discovered Mention and Canonical Subject shall still bind to exact deterministic source spans.

\hypertarget{2-one-sourceevidence-does-not-equal-one-engineering-subject}{%
\paragraph{2. One SourceEvidence Does Not Equal One Engineering Subject}\label{2-one-sourceevidence-does-not-equal-one-engineering-subject}}

The following cardinalities are valid:

\begin{verbatim}
1 SourceEvidence → 0..n Mentions
1 SourceEvidence → 0..n Canonical Engineering Subjects
1 Canonical Engineering Subject → 1..n Mentions
1 Canonical Engineering Subject → 1..n SourceEvidence references
\end{verbatim}

SourceEvidence establishes that source content is engineering-relevant and provides provenance.

It is not the semantic subject identity.

\hypertarget{3-mention-discovery-precedes-professional-interpretation}{%
\paragraph{3. Mention Discovery Precedes Professional Interpretation}\label{3-mention-discovery-precedes-professional-interpretation}}

A shared discovery stage shall identify source mentions of potentially MBSE-relevant engineering subjects without assigning the final professional interpretation.

Examples include mentions of:

\begin{itemize}
\tightlist
\item
  people or roles,
\item
  systems or applications,
\item
  capabilities,
\item
  functions or behaviors,
\item
  information or data,
\item
  states,
\item
  constraints,
\item
  requirements,
\item
  interfaces,
\item
  relationships,
\item
  unresolved engineering questions.
\end{itemize}

The discovery stage may propose a neutral canonical label but shall not make the final SysML v2 representation decision.

\hypertarget{4-repeated-mentions-shall-be-consolidated-before-persona-interpretation}{%
\paragraph{4. Repeated Mentions Shall Be Consolidated Before Persona Interpretation}\label{4-repeated-mentions-shall-be-consolidated-before-persona-interpretation}}

Multiple mentions of the same engineering subject shall resolve to one Canonical Subject whenever equivalence is sufficiently established.

Example:

\begin{verbatim}
"The remote expert"
"the expert"
"remote expert"

→ one Canonical Subject:
SUBJ-000003 Remote Expert
\end{verbatim}

The Canonical Subject retains all supporting Mention references.

Personas shall interpret \texttt{SUBJ-000003} once each, not once per Mention.

\hypertarget{5-ambiguous-identity-shall-not-be-silently-merged}{%
\paragraph{5. Ambiguous Identity Shall Not Be Silently Merged}\label{5-ambiguous-identity-shall-not-be-silently-merged}}

If two mentions may refer to the same subject but equivalence is uncertain, the system shall preserve that uncertainty.

Allowed:

\begin{verbatim}
SUBJ-000010
SUBJ-000011
possible_equivalence = uncertain
\end{verbatim}

Not allowed:

\begin{verbatim}
silent merge based only on lexical similarity
\end{verbatim}

Deterministic normalization may merge trivial representation variants such as case and surrounding whitespace, but semantic/coreference consolidation beyond trivial equivalence is explicit and traceable.

\hypertarget{6-canonical-subject-identity-is-established-before-persona-variance}{%
\paragraph{6. Canonical Subject Identity Is Established Before Persona Variance}\label{6-canonical-subject-identity-is-established-before-persona-variance}}

Canonical Subject identity shall not contain the Persona's semantic classification.

For example:

\begin{verbatim}
SUBJ-000003
canonical_label: Remote Expert
mentions: MNT-000004, MNT-000009, MNT-000014
\end{verbatim}

Personas then independently assess the same Subject:

\begin{verbatim}
P1 → Actor
P2 → External Actor
P3 → Actor
\end{verbatim}

\texttt{Actor} vs \texttt{External\ Actor} is interpretation variance, not Subject identity variance.

\hypertarget{7-personas-receive-the-same-subject-population-and-adequate-context}{%
\paragraph{7. Personas Receive the Same Subject Population and Adequate Context}\label{7-personas-receive-the-same-subject-population-and-adequate-context}}

Each configured Persona shall receive:

\begin{verbatim}
same Canonical Subject IDs
same Mention bindings
same source-only interpretation context
same Source authority
\end{verbatim}

Persona prompts shall ask for professional interpretation/classification of the supplied subjects.

They shall not independently rediscover the subject population.

\hypertarget{8-one-persona-produces-at-most-one-effective-interpretation-per-canonical-subject}{%
\paragraph{8. One Persona Produces At Most One Effective Interpretation per Canonical Subject}\label{8-one-persona-produces-at-most-one-effective-interpretation-per-canonical-subject}}

For one effective run:

\begin{verbatim}
(Persona ID, Canonical Subject ID) → at most one effective interpretation
\end{verbatim}

Repeated runs remain stability evidence and shall not multiply professional votes.

\hypertarget{9-interpretation-shall-be-mbse-relevant-but-pre-model}{%
\paragraph{9. Interpretation Shall Be MBSE-Relevant but Pre-Model}\label{9-interpretation-shall-be-mbse-relevant-but-pre-model}}

Persona interpretation may classify engineering meaning such as:

\begin{verbatim}
actor
stakeholder
system
external_system
capability
function
behavior
requirement
constraint
information
interface
state
relationship
open_question
other_engineering_subject
\end{verbatim}

The exact controlled vocabulary is an implementation contract and may evolve.

This stage shall not prematurely decide the final SysML v2 model representation where multiple valid representations remain possible.

For example:

\begin{verbatim}
engineering meaning: Actor
\end{verbatim}

is permitted, while an unnecessary early decision about the exact later SysML containment/package structure is not.

\hypertarget{10-assertionrelationship-meaning-may-be-separate-from-entity-subjects}{%
\paragraph{10. Assertion/Relationship Meaning May Be Separate from Entity Subjects}\label{10-assertionrelationship-meaning-may-be-separate-from-entity-subjects}}

The system shall support engineering semantics that are not simple noun-like entities.

Example source:

\begin{verbatim}
The expert may take temporary control when the operator permits it.
\end{verbatim}

Possible canonical subjects include:

\begin{verbatim}
Remote Expert
Temporary Microscope Control
Operator Permission
\end{verbatim}

and an assertion/relationship may express:

\begin{verbatim}
Temporary Microscope Control
is constrained by
Operator Permission
\end{verbatim}

The implementation may represent entity-like subjects and assertion-like subjects with one shared base contract or explicit subtypes, provided identity and provenance remain deterministic.

\hypertarget{11-human-review-operates-on-canonical-engineering-subjects}{%
\paragraph{11. Human Review Operates on Canonical Engineering Subjects}\label{11-human-review-operates-on-canonical-engineering-subjects}}

Human Review shall present professionally meaningful subjects, not raw Evidence blocks.

Target presentation:

\begin{verbatim}
Remote Expert
Proposed engineering meaning: Actor

P1: Actor
P2: External Actor
P3: Actor

Source mentions:
- SENT-002: "The remote expert ..."
- SENT-007: "... the expert ..."
\end{verbatim}

The Human may accept, modify, reject, defer, mark out of scope, split, merge or resolve uncertainty according to the Review contract.

\hypertarget{12-model-derivation-remains-downstream}{%
\paragraph{12. Model Derivation Remains Downstream}\label{12-model-derivation-remains-downstream}}

The corrected sequence is:

\begin{verbatim}
Engineering Source
→ Deterministic Source Projection / Source Spans
→ Source-Grounded Evidence Detection
→ Context-Preserving Mention Discovery
→ Cross-Mention Canonical Subject Consolidation
→ Shared Persona Interpretation
→ Field-Level Consensus / Variance
→ Human Engineering Review
→ Approved Engineering Information
→ Architecture / Model Derivation
→ SysML v2
\end{verbatim}

No Canonical Subject becomes Approved Engineering Information without Human Review.

\begin{center}\rule{0.5\linewidth}{0.5pt}\end{center}

\hypertarget{canonical-identity-principle}{%
\subsubsection{Canonical Identity Principle}\label{canonical-identity-principle}}

The architecture shall preserve the following invariant:

\begin{quote}
Mentions are many; canonical engineering subjects are unique within the resolved source scope; professional interpretations are many per subject only because Personas may legitimately disagree.
\end{quote}

Formally:

\begin{verbatim}
Source Span
    ↓
Mention
    ↓
Canonical Engineering Subject
    ↓
Persona Interpretation
    ↓
Consensus / Variance
    ↓
Human Review
\end{verbatim}

Identity shall be established before variance.

\begin{center}\rule{0.5\linewidth}{0.5pt}\end{center}

\hypertarget{source-context-principle}{%
\subsubsection{Source Context Principle}\label{source-context-principle}}

A key lesson from WP-12 is:

\begin{quote}
Small units are required for exact provenance; larger context is often required for correct interpretation.
\end{quote}

The architecture shall therefore never reduce LLM context merely to mirror persistence granularity.

Source context supplied to the LLM shall remain bounded to registered Engineering Source content plus clearly separated reference guidance. Reference guidance shall never become positive engineering evidence.

\begin{center}\rule{0.5\linewidth}{0.5pt}\end{center}

\hypertarget{expected-effect-on-the-wp-12-example}{%
\subsubsection{Expected Effect on the WP-12 Example}\label{expected-effect-on-the-wp-12-example}}

A context-preserving discovery of the Remote Microscope Collaboration source is expected to identify a small canonical population resembling, without prescribing as a gold standard:

\begin{verbatim}
Remote Microscope Collaboration
Microscope Operator
Remote Expert
Microscope Workstation
Client Application
Live Microscope View Sharing
Remote Consultation
Live Image Observation
Temporary Microscope Control
Operator Permission
Operator Responsibility
Current Controller Awareness
Session Information Retention
Retention Period
Connection Quality Limits
Remotely Controllable Microscope Functions
\end{verbatim}

The exact result remains LLM-generated and subject to Human Review.

Repeated appearances of e.g. \texttt{Remote\ Expert} shall contribute additional provenance Mentions to one Canonical Subject rather than create duplicate subjects.

\begin{center}\rule{0.5\linewidth}{0.5pt}\end{center}

\hypertarget{rejected-alternatives-6}{%
\subsubsection{Rejected Alternatives}\label{rejected-alternatives-6}}

\hypertarget{a-continue-one-evidence-to-one-interpretation}{%
\paragraph{A. Continue One-Evidence-to-One-Interpretation}\label{a-continue-one-evidence-to-one-interpretation}}

Rejected because one Evidence block can contain multiple independent engineering subjects and assertions.

\hypertarget{b-let-every-persona-rediscover-subjects-independently}{%
\paragraph{B. Let Every Persona Rediscover Subjects Independently}\label{b-let-every-persona-rediscover-subjects-independently}}

Rejected because it recreates the original subject-population mismatch and makes downstream consensus compare non-equivalent populations.

\hypertarget{c-match-independent-persona-outputs-only-afterward}{%
\paragraph{C. Match Independent Persona Outputs Only Afterward}\label{c-match-independent-persona-outputs-only-afterward}}

Rejected as the primary architecture because late semantic matching recreates the exact reconciliation problem ADR-027 was introduced to avoid.

\hypertarget{d-use-sentence-identity-as-engineering-subject-identity}{%
\paragraph{D. Use Sentence Identity as Engineering Subject Identity}\label{d-use-sentence-identity-as-engineering-subject-identity}}

Rejected because one sentence may contain multiple engineering subjects and one engineering subject may appear across multiple sentences.

\hypertarget{e-interpret-every-mention-independently}{%
\paragraph{E. Interpret Every Mention Independently}\label{e-interpret-every-mention-independently}}

Rejected because repeated mentions of the same engineering subject would multiply work and votes without adding engineering value.

\begin{center}\rule{0.5\linewidth}{0.5pt}\end{center}

\hypertarget{implementation-slices-1}{%
\subsubsection{Implementation Slices}\label{implementation-slices-1}}

\hypertarget{r4c1--contract-and-deterministic-identity}{%
\paragraph{R4c.1 --- Contract and deterministic identity}\label{r4c1--contract-and-deterministic-identity}}

Introduce immutable contracts for:

\begin{verbatim}
SourceSpanReference
EngineeringMention
CanonicalEngineeringSubject
CanonicalSubjectSet
\end{verbatim}

including fingerprints and fail-closed validation.

\hypertarget{r4c2--context-preserving-subject-discovery}{%
\paragraph{R4c.2 --- Context-preserving Subject Discovery}\label{r4c2--context-preserving-subject-discovery}}

Use a dedicated LLM task over a source-only context window to discover \texttt{0..n} Mentions and propose cross-mention subject groupings.

The system, not the LLM, binds returned span IDs to exact source text.

\hypertarget{r4c3--canonical-consolidation}{%
\paragraph{R4c.3 --- Canonical consolidation}\label{r4c3--canonical-consolidation}}

Resolve trivial deterministic equivalence directly and validate LLM-proposed cross-mention grouping.

Unknown spans, impossible mention ranges, duplicate identities, unsupported merges and ungrounded subjects fail closed.

\hypertarget{r4c4--shared-persona-interpretation}{%
\paragraph{R4c.4 --- Shared Persona Interpretation}\label{r4c4--shared-persona-interpretation}}

All Personas interpret the same Canonical Subject Set and return one professional interpretation per subject.

\hypertarget{r4c5--subject-centered-consensus-and-human-review}{%
\paragraph{R4c.5 --- Subject-centered Consensus and Human Review}\label{r4c5--subject-centered-consensus-and-human-review}}

Consensus compares fields for the same \texttt{SUBJ-*}. Human Review is generated from those subjects and their Persona interpretations.

\hypertarget{r4c6--live-validation-and-documentation-closeout}{%
\paragraph{R4c.6 --- Live validation and documentation closeout}\label{r4c6--live-validation-and-documentation-closeout}}

Validate on the WP-12 source, run focused/full regression, update recovery findings and SSOT, and only then stage/commit.

\begin{center}\rule{0.5\linewidth}{0.5pt}\end{center}

\hypertarget{acceptance-criteria-2}{%
\subsubsection{Acceptance Criteria}\label{acceptance-criteria-2}}

R4c is successful only if the real single-source test demonstrates all of the following:

\begin{enumerate}
\def\labelenumi{\arabic{enumi}.}
\tightlist
\item
  source context remains readable to the LLM,
\item
  every Canonical Subject is grounded in exact source spans,
\item
  one passage may yield multiple MBSE-relevant subjects,
\item
  repeated mentions of one subject are consolidated,
\item
  Personas interpret the same Subject population,
\item
  adding Personas does not add Subject identities,
\item
  Human Review displays meaningful engineering semantics,
\item
  no positive engineering subject is derived from instructions/reference knowledge,
\item
  no premature SysML v2 model structure is treated as approved fact,
\item
  the resulting Review is usable by a Systems Engineer.
\end{enumerate}

\begin{center}\rule{0.5\linewidth}{0.5pt}\end{center}

\hypertarget{thesis-relevant-finding}{%
\subsubsection{Thesis-Relevant Finding}\label{thesis-relevant-finding}}

The WP-12 recovery produced a second architectural finding:

\begin{quote}
Provenance granularity and interpretation granularity are different concerns. Exact traceability benefits from small deterministic source units, while semantic understanding benefits from larger context windows. A robust AI4MBSE pipeline should therefore establish canonical engineering-subject identity from source-bound mentions before professional interpretation, while supplying sufficient surrounding source context to the interpreting model.
\end{quote}

This finding complements ADR-027's earlier conclusion that source evidence detection and professional interpretation must be separated.

\textbackslash newpage

\hypertarget{adr-029--human-reviewed-model-placement-before-model-assembly}{%
\subsection{ADR-029 --- Human-Reviewed Model Placement Before Model Assembly}\label{adr-029--human-reviewed-model-placement-before-model-assembly}}

\hypertarget{status-10}{%
\subsubsection{Status}\label{status-10}}

Accepted

\hypertarget{date-10}{%
\subsubsection{Date}\label{date-10}}

2026-08-24

\hypertarget{context-15}{%
\subsubsection{Context}\label{context-15}}

The first live LLM-assisted Phase-H generation attempt for WP-12 Project \texttt{120412} exposed \texttt{BLK-006}.

The three modeling personas successfully produced and persisted bounded projection results for 14 unresolved Approved Inputs. The consolidated result contained:

\begin{verbatim}
proposed_mapping: 0
ambiguous:        10
unmapped:          4
\end{verbatim}

The current hybrid deriver treats every non-\texttt{proposed\_mapping} result as a generation failure and therefore aborts before a reviewable Candidate Set can be produced.

This behavior is inconsistent with ADR-028. Material modeling disagreement was intended to remain explicit variance while Human Model Review remains the authority boundary.

A second architectural issue became visible at the same time: the current Phase-H step tries to resolve local framework placement and assemble a complete model in one operation. This combines two materially different reasoning tasks.

The accepted recovery pattern from R4c is applicable again:

\begin{verbatim}
local identification / interpretation
-> Human authority
-> later synthesis
\end{verbatim}

For model creation this becomes:

\begin{verbatim}
placement
-> Human placement authority
-> model assembly
-> Final Model Review
\end{verbatim}

\hypertarget{decision-15}{%
\subsubsection{Decision}\label{decision-15}}

\hypertarget{mpa-01--model-placement-is-a-separate-stage}{%
\paragraph{MPA-01 --- Model placement is a separate stage}\label{mpa-01--model-placement-is-a-separate-stage}}

Before model assembly, every model-promotable Approved Engineering Subject shall receive an explicit placement decision against the pinned Model Structure Profile.

Placement answers:

\begin{quote}
Where does this approved engineering information belong in the model framework?
\end{quote}

The exact placement is expressed through a profile-controlled rule. The rule therefore carries the target framework location, model area, element type and, where encoded by the profile, the Stakeholder / System / Subsystem level.

\hypertarget{mpa-02--dedicated-modeling-personas}{%
\paragraph{MPA-02 --- Dedicated modeling personas}\label{mpa-02--dedicated-modeling-personas}}

LLM-assisted placement uses three dedicated modeling personas:

\begin{enumerate}
\def\labelenumi{\arabic{enumi}.}
\tightlist
\item
  SysML / Profile Modeler
\item
  System Architecture Modeler
\item
  Conservative Modeling Reviewer
\end{enumerate}

They belong to the model-placement responsibility and shall not reuse the legacy-processing / derivation-assessment persona definitions.

All three receive the same Approved Engineering Information and the same pinned placement options.

\hypertarget{mpa-03--persona-agreement-is-not-a-gate}{%
\paragraph{MPA-03 --- Persona agreement is not a gate}\label{mpa-03--persona-agreement-is-not-a-gate}}

The three personas reason independently.

Valid outcomes include:

\begin{verbatim}
3:0 agreement
2:1 variance
1:1:1 variance
ambiguous
unmapped
\end{verbatim}

No majority result and no unanimity requirement may silently create engineering or model authority.

Persona agreement is advisory evidence only.

\hypertarget{mpa-04--human-model-placement-review-is-the-authority-boundary}{%
\paragraph{MPA-04 --- Human Model Placement Review is the authority boundary}\label{mpa-04--human-model-placement-review-is-the-authority-boundary}}

Persona placement results are collected into a reviewable placement bundle.

The Human reviewer decides the authoritative placement for each required Approved Subject.

The UX shall intentionally follow the already established Human Engineering Review pattern where practical:

\begin{itemize}
\tightlist
\item
  Pending Decisions
\item
  Reviewed Decisions
\item
  All
\item
  Needs Attention
\item
  explicit persona proposals / variance
\item
  immutable decision history
\item
  explicit reopen / correction
\end{itemize}

A Human placement decision may select only a placement allowed by the pinned profile unless a separately governed exception path is introduced.

\hypertarget{mpa-05--approved-model-placement-set-precedes-assembly}{%
\paragraph{MPA-05 --- Approved Model Placement Set precedes assembly}\label{mpa-05--approved-model-placement-set-precedes-assembly}}

A complete, Human-resolved placement population forms the input authority for model assembly.

Model assembly answers a different question:

\begin{quote}
How do the already placed engineering building blocks form one coherent model?
\end{quote}

Assembly may construct hierarchy, topology and relationships and may create multiple coordinated model views. One source does not imply one diagram, and one approved engineering subject does not imply one complete diagram.

\hypertarget{mpa-06--relationships-remain-engineering-authority-for-later-assembly}{%
\paragraph{MPA-06 --- Relationships remain engineering authority for later assembly}\label{mpa-06--relationships-remain-engineering-authority-for-later-assembly}}

Accepted semantic Relationships from Approved Engineering Information are not forced through the per-Subject placement decision merely to obtain a target relationship type.

They remain authoritative engineering semantics and are consumed during the later assembly / target-representation stage.

\hypertarget{mpa-07--preview-belongs-primarily-after-placement-resolution}{%
\paragraph{MPA-07 --- Preview belongs primarily after placement resolution}\label{mpa-07--preview-belongs-primarily-after-placement-resolution}}

The placement stage may show a lightweight framework overview to communicate where Subjects are being sorted.

Rich model preview belongs after placement decisions have been resolved and assembly has produced a coherent model draft. That later review may expose multiple coordinated diagrams and a textual notation view.

\hypertarget{mpa-08--candidate--model-materialization-follows-human-placement}{%
\paragraph{MPA-08 --- Candidate / model materialization follows Human placement}\label{mpa-08--candidate--model-materialization-follows-human-placement}}

A complete Model Candidate Set or Internal Model shall not be created by silently resolving modeling-persona variance.

Human-resolved placement authority must precede deterministic model materialization / assembly.

\hypertarget{relationship-to-adr-028}{%
\subsubsection{Relationship to ADR-028}\label{relationship-to-adr-028}}

ADR-028 remains authoritative for:

\begin{itemize}
\tightlist
\item
  \texttt{eco\_deterministic} versus \texttt{llm\_assisted},
\item
  advisory strategy recommendation,
\item
  mandatory Human authority,
\item
  review-driven LLM escalation,
\item
  and no upstream reinterpretation.
\end{itemize}

This ADR supersedes any implementation interpretation of ADR-028 R5-09 that treats persona unanimity as a prerequisite for successful LLM-assisted processing. Modeling comparison preserves variance for Human resolution.

\hypertarget{target-lifecycle-1}{%
\subsubsection{Target lifecycle}\label{target-lifecycle-1}}

\begin{verbatim}
Approved Engineering Information
        |
        v
Deterministic placement coverage
        |
        v
Human-selected mode
   |                 |
   | Eco             | LLM-assisted
   v                 v
profile proposal   3 placement personas
   |                 |
   +--------+--------+
            v
Model Placement Review Bundle
            |
            v
Human Model Placement Review
            |
            v
Approved Model Placement Set
            |
            v
Model Assembly
            |
            +--> coordinated diagram view(s)
            +--> textual notation
            |
            v
Final Model Review
            |
            v
Approved downstream model
\end{verbatim}

\hypertarget{consequences-13}{%
\subsubsection{Consequences}\label{consequences-13}}

\hypertarget{positive-5}{%
\paragraph{Positive}\label{positive-5}}

\begin{itemize}
\tightlist
\item
  placement uncertainty becomes reviewable rather than fatal,
\item
  Stakeholder / System / Subsystem assignment becomes an explicit Human modeling decision,
\item
  persona disagreement becomes useful engineering evidence,
\item
  model assembly receives cleaner, already placed building blocks,
\item
  and the Final Model Review can focus on the assembled model rather than re-litigating local placement.
\end{itemize}

\hypertarget{trade-offs-7}{%
\paragraph{Trade-offs}\label{trade-offs-7}}

\begin{itemize}
\tightlist
\item
  Phase H gains an additional persisted Human authority boundary,
\item
  current direct Candidate generation must be decomposed,
\item
  and the existing Candidate Review / Final Model Review integration must be reconciled with the new placement stage without weakening traceability.
\end{itemize}

\hypertarget{implementation-slices-2}{%
\subsubsection{Implementation slices}\label{implementation-slices-2}}

\begin{verbatim}
BLK-006 C1  architecture + dedicated model-placement department
BLK-006 C2  immutable placement proposal / comparison contracts
BLK-006 C3  persisted Human Model Placement Review
BLK-006 C4  Guided Workflow UI analogous to Human Engineering Review
BLK-006 C5  placement-authoritative model assembly + Candidate materialization
BLK-006 C6  live same-project retest + downstream Final Model Review
\end{verbatim}

\textbackslash newpage

\hypertarget{adr-030--semantic-interpretation-and-controlled-classification-alignment}{%
\subsection{ADR-030 --- Semantic Interpretation and Controlled Classification Alignment}\label{adr-030--semantic-interpretation-and-controlled-classification-alignment}}

Semantic Interpretation and Controlled Classification Alignment

Status

Accepted

Date

2026-08-26

\hypertarget{context-16}{%
\subsubsection{Context}\label{context-16}}

Cross-source processing tests exposed a recurring semantic-contract failure that is distinct from source grounding and from system integrity.

An LLM may return a professionally meaningful classification expression such as \texttt{architecture}, \texttt{condition} or \texttt{classification} even when ADR-011 defines a closed internal \texttt{information\_type} vocabulary. Treating every such expression as a failed interpretation couples two responsibilities that should remain separate:

\begin{enumerate}
\def\labelenumi{\arabic{enumi}.}
\tightlist
\item
  understanding what the Engineering Source means; and
\item
  representing that meaning with the repository's controlled vocabulary.
\end{enumerate}

The observed terms are test vectors, not mapping rules. This ADR does not authorize hard-coded aliases such as \texttt{architecture\ -\textgreater{}\ logical\_element} or \texttt{condition\ -\textgreater{}\ constraint}; those mappings are context-dependent.

\hypertarget{decision-16}{%
\subsubsection{Decision}\label{decision-16}}

The semantic workflow shall explicitly separate professional interpretation from controlled representation:

\begin{verbatim}
Engineering Source
    -> LLM professional interpretation
    -> controlled classification alignment
    -> strict internal semantic contract
    -> consensus / Human Engineering Review
\end{verbatim}

An LLM-produced classification expression is a semantic proposal, not an authoritative internal enum value.

\hypertarget{controlled-classification-alignment}{%
\paragraph{Controlled classification alignment}\label{controlled-classification-alignment}}

Before an interpretation enters the strict internal semantic contract, any out-of-vocabulary required classification expression shall be aligned to the accepted controlled vocabulary.

The alignment input retains at least immutable item identity (\texttt{EVD-*} or \texttt{SUBJ-*}), field name, raw LLM expression, interpreted statement, relevant source-grounded context when available and the complete allowed target vocabulary.

The alignment may change only the explicitly identified classification field. It may not change item identity or population, interpreted statements, source grounding, rationale, uncertainty, missing-evidence content, relationships or already-valid fields.

A value already contained in the controlled vocabulary passes through unchanged. Pure case/whitespace normalization to one unique controlled token may be performed deterministically.

If semantic inference is required, exactly one bounded alignment LLM request may translate the raw expression to the controlled vocabulary using context. The mapper reports \texttt{mapped}, \texttt{ambiguous} or \texttt{unmapped}.

For \texttt{information\_type}, \texttt{ambiguous} and \texttt{unmapped} normalize to the already accepted neutral value \texttt{unclassified}. If the mapper itself violates its contract, \texttt{information\_type} may likewise fail safely to \texttt{unclassified} rather than inventing a more specific meaning. A classification dimension without an accepted neutral value remains fail-closed if alignment cannot produce a valid controlled target.

The original LLM output remains immutable. Every non-identity alignment is retained as an auditable record containing item identity, field name, raw value, normalized value, mapping status, concise rationale/fallback reason, mapper response identity when applicable and a deterministic fingerprint. Only normalized controlled values may enter downstream semantic processing.

\hypertarget{boundary-to-ontology-mapping}{%
\paragraph{Boundary to ontology mapping}\label{boundary-to-ontology-mapping}}

Controlled Classification Alignment is not ADR-011 P4.3 Terminology and Ontology Candidate Mapping. It maps an interpreted classification expression to the existing engineering classification vocabulary. Project Glossary, Turing Core, IOF Core and BFO mapping remain downstream and may not become accidental \texttt{information\_type} values.

\hypertarget{no-source-specific-aliases}{%
\paragraph{No source-specific aliases}\label{no-source-specific-aliases}}

No source-specific or example-specific semantic alias table is permitted. The system shall not contain hard-coded translations for the terms observed in the current test sources.

\hypertarget{validation-boundary}{%
\paragraph{Validation boundary}\label{validation-boundary}}

Strict deterministic parsing remains mandatory after alignment. Classification alignment shall not repair malformed JSON, missing or duplicate item identities, changed item population, invalid source references, malformed relationships or system-owned integrity corruption.

The existing bounded R4c classification repair is superseded as the primary handling of out-of-vocabulary classification expressions. A compatibility projection may temporarily remain, but the architectural authority is the explicit alignment artifact.

\hypertarget{consequences-14}{%
\subsubsection{Consequences}\label{consequences-14}}

Semantically reasonable wording no longer fails only because the LLM did not reproduce one exact internal token. Downstream stages receive normalized controlled values and unresolved precision becomes visible as \texttt{unclassified} instead of a technical processing failure.

An additional LLM call may be required for persona outputs containing out-of-vocabulary classifications. Raw and normalized semantics remain fully traceable.

\hypertarget{acceptance-criteria-3}{%
\subsubsection{Acceptance Criteria}\label{acceptance-criteria-3}}

\begin{enumerate}
\def\labelenumi{\arabic{enumi}.}
\tightlist
\item
  Valid controlled classifications pass without an alignment LLM call.
\item
  Lexical-only normalization is deterministic.
\item
  Out-of-vocabulary Information Types can be contextually mapped by one bounded alignment call.
\item
  Ambiguous/unmapped Information Types become \texttt{unclassified}.
\item
  An invalid mapper response cannot inject an uncontrolled value.
\item
  Raw LLM output remains unchanged and alignment provenance is auditable.
\item
  Structural/identity/integrity failures remain fail-closed.
\item
  The same alignment service is reusable by Evidence and Subject interpretation.
\item
  No source-specific alias for observed test expressions is introduced.
\end{enumerate}

\textbackslash newpage

\hypertarget{adr-031--semantic-field-consistency-alignment}{%
\subsection{ADR-031 --- Semantic Field Consistency Alignment}\label{adr-031--semantic-field-consistency-alignment}}

Semantic Field Consistency Alignment

Status

Accepted

Date

2026-08-26

\hypertarget{context-17}{%
\subsubsection{Context}\label{context-17}}

ADR-030 separates professional semantic interpretation from controlled classification vocabulary. Cross-source retry RUN-000004 then exposed a different failure class after classification alignment had succeeded.

The LLM returned a valid controlled \texttt{epistemic\_class="explicit"} together with a non-null object-valued \texttt{missing\_evidence}. The object contained explanatory semantic text, but the strict internal contract requires the coupled invariant:

\begin{itemize}
\tightlist
\item
  \texttt{assumption} -\textgreater{} non-empty textual \texttt{missing\_evidence}
\item
  \texttt{explicit} / \texttt{interpretation} -\textgreater{} \texttt{missing\_evidence=null}
\end{itemize}

This is not a classification-vocabulary problem. It is a semantic consistency problem between two coupled fields.

Blindly coercing every non-null value to null would be unsafe because a non-null value may represent a genuine evidence gap and therefore justify \texttt{epistemic\_class="assumption"} instead. The system must preserve meaning rather than merely force schema validity.

The same coupled invariant exists in both Evidence Interpretation and Subject Interpretation.

\hypertarget{decision-17}{%
\subsubsection{Decision}\label{decision-17}}

A reusable Semantic Field Consistency Alignment step shall run after Controlled Classification Alignment and before the strict interpretation parser:

\begin{verbatim}
LLM professional interpretation
    -> Controlled Classification Alignment
    -> Semantic Field Consistency Alignment
    -> strict deterministic interpretation contract
    -> downstream consensus / Human Engineering Review
\end{verbatim}

The first supported coupled field set is:

\begin{verbatim}
epistemic_class + missing_evidence
\end{verbatim}

\hypertarget{trigger}{%
\paragraph{Trigger}\label{trigger}}

No additional LLM request is made when the pair already satisfies the internal contract.

A consistency-alignment need is created only when:

\begin{itemize}
\tightlist
\item
  \texttt{epistemic\_class} is one of the accepted pre-review values; and
\item
  the coupled \texttt{missing\_evidence} value violates the deterministic pair invariant.
\end{itemize}

Malformed item identity, changed population, malformed JSON, invalid relationships and other structural failures remain outside this recovery mechanism.

\hypertarget{bounded-contextual-resolution}{%
\paragraph{Bounded contextual resolution}\label{bounded-contextual-resolution}}

All inconsistent pairs in one persona output are batched into exactly one bounded mapper request.

For each requested item the mapper sees immutable item identity, interpreted statement, raw epistemic class, raw \texttt{missing\_evidence}, and source-grounded context when available.

The mapper may change only:

\begin{itemize}
\tightlist
\item
  \texttt{epistemic\_class}
\item
  \texttt{missing\_evidence}
\end{itemize}

It may not change identity, population, interpreted statement, information type, modality, rationale, uncertainties, relationships or source grounding.

The normalized pair must satisfy exactly one of:

\begin{itemize}
\tightlist
\item
  \texttt{explicit} + \texttt{null}
\item
  \texttt{interpretation} + \texttt{null}
\item
  \texttt{assumption} + non-empty trimmed text
\end{itemize}

If the raw non-null content represents a genuine evidence gap, the mapper may normalize to \texttt{assumption} and preserve that meaning as concise textual \texttt{missing\_evidence}. If it does not represent a genuine evidence gap, the mapper may retain/choose \texttt{explicit} or \texttt{interpretation} and normalize \texttt{missing\_evidence} to null.

There is no automatic neutral fallback for an unresolved pair. Mapper execution or contract failure remains fail-closed because silently discarding possible missing-evidence meaning would reduce semantic quality.

\hypertarget{auditability}{%
\paragraph{Auditability}\label{auditability}}

The raw LLM output remains immutable.

Every non-identity consistency decision is persisted with raw and normalized pair values, rationale, mapper response identity and deterministic fingerprint. Only the normalized output enters the strict parser.

\hypertarget{boundary-to-adr-030}{%
\paragraph{Boundary to ADR-030}\label{boundary-to-adr-030}}

ADR-030 remains authoritative for controlled classification vocabulary and continues to prohibit Classification Alignment from modifying \texttt{missing\_evidence}.

ADR-031 is a separate downstream consistency layer for semantically coupled fields and does not weaken ADR-030's field authority.

\hypertarget{acceptance-criteria-4}{%
\subsubsection{Acceptance Criteria}\label{acceptance-criteria-4}}

\begin{enumerate}
\def\labelenumi{\arabic{enumi}.}
\tightlist
\item
  Already-consistent pairs pass without an additional LLM request.
\item
  Inconsistent pairs are batched into one bounded request per persona output.
\item
  The mapper may modify only \texttt{epistemic\_class} and \texttt{missing\_evidence}.
\item
  Every normalized pair satisfies the deterministic pair invariant.
\item
  Genuine evidence-gap meaning can be retained by normalizing to \texttt{assumption} plus textual \texttt{missing\_evidence}.
\item
  No automatic nulling or other lossy fallback is permitted on mapper execution/validation failure.
\item
  Raw outputs remain immutable and decisions are auditable.
\item
  The service is reusable by Evidence and Subject Interpretation.
\item
  Structural, identity and relationship failures remain fail-closed.
\end{enumerate}

\textbackslash newpage

\hypertarget{adr-032--project-level-multi-source-reconciliation-and-controlled-engineering-evolution}{%
\subsection{ADR-032 --- Project-Level Multi-Source Reconciliation and Controlled Engineering Evolution}\label{adr-032--project-level-multi-source-reconciliation-and-controlled-engineering-evolution}}

\hypertarget{status-11}{%
\subsubsection{Status}\label{status-11}}

Accepted

\hypertarget{date-11}{%
\subsubsection{Date}\label{date-11}}

2026-08-31

\hypertarget{context-18}{%
\subsubsection{Context}\label{context-18}}

The Turing Generator currently processes registered engineering sources through source-bound Processing Runs.

ADR-012 intentionally established this boundary:

\begin{itemize}
\tightlist
\item
  one Processing Run processes one primary registered source;
\item
  Information Units and downstream processing evidence remain source-traceable;
\item
  equivalent, overlapping, or contradictory information from different sources remains separate during source-level processing;
\item
  project-level comparison is a later responsibility.
\end{itemize}

ADR-025 and ADR-026 subsequently introduced source-anchored semantic consolidation and Persona-aware interpretation. These capabilities improve semantic consistency inside one source-processing context but do not provide project-level reconciliation between independent registered sources.

WP-12 Multi-Source validation exposed BLK-002:

\begin{quote}
Local Processing Artifact identities such as \texttt{IU-000001} can legitimately occur in different Processing Runs, while existing project-level lifecycle aggregation incorrectly treats \texttt{(artifact\_type,\ artifact\_id)} as a globally unique artifact identity.
\end{quote}

Resolving this collision is necessary but not sufficient for genuine Multi-Source processing.

A project may contain:

\begin{itemize}
\tightlist
\item
  several engineering sources describing the same product;
\item
  complementary documents covering different engineering viewpoints or abstraction levels;
\item
  contradictory or materially different statements;
\item
  context-only documents that establish product terminology and system understanding;
\item
  accidentally registered documents that belong to another project;
\item
  already Human-reviewed engineering information and an existing accepted model when additional sources are registered later.
\end{itemize}

The architecture must therefore distinguish:

\begin{enumerate}
\def\labelenumi{\arabic{enumi}.}
\tightlist
\item
  source-local processing;
\item
  project-fit assessment;
\item
  cross-source semantic reconciliation;
\item
  Human Engineering Authority;
\item
  model-impact reconciliation.
\end{enumerate}

The system must not force unrelated information into a common semantic subject merely because it belongs to the same Project.

Likewise, absence of semantic overlap does not imply that a source belongs to another Project. A BOM, interface specification, requirement specification and user-needs document may contain very different terminology while still describing the same product.

Context-only sources exist explicitly to provide product and domain understanding and shall be usable for this purpose without becoming Engineering Authority.

Finally, newly registered information must not automatically override already reviewed information merely because the source is newer. Existing Human-reviewed Approved Inputs remain Engineering Authority until a subsequent Human Review explicitly changes that authority.

\begin{center}\rule{0.5\linewidth}{0.5pt}\end{center}

\hypertarget{decision-18}{%
\subsubsection{Decision}\label{decision-18}}

\hypertarget{1-preserve-source-bound-processing}{%
\paragraph{1. Preserve source-bound Processing}\label{1-preserve-source-bound-processing}}

Existing source-level Processing remains source-bound.

The following conceptual flow remains unchanged:

\begin{verbatim}
Registered Source
→ Source Projection
→ Source Analysis Units / Evidence
→ Source-local Interpretation
→ Canonical Engineering Subjects
→ Persona Interpretation / Consensus
\end{verbatim}

No source-level Processing stage may silently merge evidence from independent registered engineering sources.

All source-derived artifacts retain exact Source, Processing Run and Attempt provenance.

\begin{center}\rule{0.5\linewidth}{0.5pt}\end{center}

\hypertarget{2-qualify-processing-artifact-identity-by-processing-run-at-project-wide-boundaries}{%
\paragraph{2. Qualify Processing Artifact identity by Processing Run at project-wide boundaries}\label{2-qualify-processing-artifact-identity-by-processing-run-at-project-wide-boundaries}}

A \texttt{ProcessingArtifactReference} remains a run-local artifact reference and shall not be extended with additional persisted identity fields solely to resolve BLK-002.

Within one Processing Run, local artifact identity remains:

\begin{verbatim}
(artifact_type, artifact_id)
\end{verbatim}

At boundaries that aggregate several Processing Runs, effective artifact identity becomes:

\begin{verbatim}
(processing_run_id, artifact_type, artifact_id)
\end{verbatim}

Therefore:

\begin{verbatim}
RUN-000001 / information_unit / IU-000001
RUN-000002 / information_unit / IU-000001
\end{verbatim}

are legal independent artifacts.

Within the same Processing Run, conflicting immutable references, duplicate publication and other existing integrity violations remain fail-closed.

No historical Processing Artifact manifest migration is required.

\begin{center}\rule{0.5\linewidth}{0.5pt}\end{center}

\hypertarget{3-introduce-a-project-fit-assessment-before-cross-source-reconciliation}{%
\paragraph{3. Introduce a Project Fit Assessment before cross-source reconciliation}\label{3-introduce-a-project-fit-assessment-before-cross-source-reconciliation}}

A processed source shall not automatically participate in project-level semantic reconciliation merely because it was registered inside the Project.

A new derived processing boundary shall assess:

\begin{quote}
Is this source plausibly part of the product/system represented by this Project?
\end{quote}

The assessment may use:

\begin{itemize}
\tightlist
\item
  explicit \texttt{context\_only} sources;
\item
  project/product contextual information;
\item
  terminology and concepts derived from existing project sources;
\item
  existing active Human-reviewed engineering information where available.
\end{itemize}

Context information is evidence for project understanding. It is not Engineering Authority merely because it is used by the Project Fit Assessment.

The Project Fit Assessment produces one of:

\begin{verbatim}
plausible_in_scope
uncertain
likely_out_of_scope
\end{verbatim}

The result is machine-generated processing evidence.

It does not modify:

\begin{itemize}
\tightlist
\item
  the Source Manifest;
\item
  the registered Source Role;
\item
  Engineering Authority;
\item
  Approved Inputs;
\item
  the engineering model.
\end{itemize}

\texttt{plausible\_in\_scope} permits the source to continue toward project-level reconciliation.

\texttt{uncertain} and \texttt{likely\_out\_of\_scope} prevent the source from entering project-level engineering reconciliation until a Human Processing Decision resolves its treatment.

The source remains registered and auditable in every case.

Existing source dispositions defined by ADR-012 remain authoritative for operational source treatment:

\begin{verbatim}
in_scope
context_only
out_of_scope
\end{verbatim}

The Project Fit Assessment may support such a Human decision but does not replace it.

\begin{center}\rule{0.5\linewidth}{0.5pt}\end{center}

\hypertarget{4-introduce-project-level-cross-source-semantic-reconciliation}{%
\paragraph{4. Introduce Project-Level Cross-Source Semantic Reconciliation}\label{4-introduce-project-level-cross-source-semantic-reconciliation}}

Only sources admitted to project-level engineering processing participate in Cross-Source Semantic Reconciliation.

This boundary operates after sufficiently stable source-local Engineering Subjects and Persona interpretation/consensus are available.

Its purpose is not to determine model placement.

It answers:

\begin{quote}
How are engineering statements from different sources semantically related?
\end{quote}

Supported semantic relations are:

\begin{verbatim}
equivalent
complementary
potential_conflict
distinct
uncertain
\end{verbatim}

\hypertarget{equivalent}{%
\subparagraph{Equivalent}\label{equivalent}}

Two or more source contributions express substantially the same engineering meaning.

They may be presented as one project-level review subject while retaining all individual source evidence.

\hypertarget{complementary}{%
\subparagraph{Complementary}\label{complementary}}

Contributions relate to the same engineering concern but provide different, non-competing information, abstraction levels, viewpoints or detail.

They shall not be collapsed into one statement merely because they overlap semantically.

\hypertarget{potential-conflict}{%
\subparagraph{Potential conflict}\label{potential-conflict}}

Contributions appear to make materially incompatible claims concerning the same engineering concern.

The system records the variance but does not select a winner.

\hypertarget{distinct}{%
\subparagraph{Distinct}\label{distinct}}

The contributions represent separate engineering subjects.

\hypertarget{uncertain}{%
\subparagraph{Uncertain}\label{uncertain}}

The semantic relationship cannot safely be determined.

Uncertain contributions remain separate and visible for Human reasoning.

\begin{center}\rule{0.5\linewidth}{0.5pt}\end{center}

\hypertarget{5-llm-assisted-semantic-comparison-has-no-engineering-authority}{%
\paragraph{5. LLM-assisted semantic comparison has no Engineering Authority}\label{5-llm-assisted-semantic-comparison-has-no-engineering-authority}}

Project Fit Assessment and Cross-Source Semantic Reconciliation may use bounded LLM calls.

LLMs may propose semantic relationships and supporting reasoning evidence.

They may not authorize:

\begin{itemize}
\tightlist
\item
  Engineering Approval;
\item
  source priority;
\item
  automatic truth selection;
\item
  automatic supersession;
\item
  model modification;
\item
  semantic merging without valid evidence.
\end{itemize}

The following never constitute Engineering Authority:

\begin{verbatim}
LLM confidence
number of agreeing sources
source order
source age
source type
Persona consensus
semantic equivalence
\end{verbatim}

Deterministic validation shall verify provenance, referenced identities and reconciliation consistency.

Missing, malformed or uncertain semantic evidence shall fail closed against automatic merge authority.

\begin{center}\rule{0.5\linewidth}{0.5pt}\end{center}

\hypertarget{6-preserve-human-engineering-authority-before-project-level-reconciliation}{%
\paragraph{6. Preserve Human Engineering Authority before project-level reconciliation}\label{6-preserve-human-engineering-authority-before-project-level-reconciliation}}

Cross-Source Semantic Reconciliation produces review evidence only and has no Engineering Authority.

The existing Human Review, Approved Input and Approved Engineering Information contracts remain the authority boundary for source-derived engineering information. Each participating engineering Source therefore retains its existing source-local path:

\begin{verbatim}
Source / Processing Evidence
        ↓
Human Engineering Review
        ↓
Approved Input / Approved Engineering Information
\end{verbatim}

Cross-source semantic evidence does not bypass, replace or synthesize the result of that Human Engineering Review.

After source-local Engineering Authority exists, the Multi-Source extension introduces one common project-level Human Engineering Authority boundary. At this boundary, Human reviewers reconcile the relationship between already reviewed source-local authority using the exact Cross-Source Semantic Reconciliation evidence.

The reviewer may determine that reviewed engineering information shall:

\begin{verbatim}
remain independent
coexist as valid information concerning the same engineering concern
supersede previously accepted project-level authority
remain unresolved
\end{verbatim}

Existing source provenance remains attached to every reviewed result and every project-level authority decision.

The common project-level Human Engineering Authority boundary is an authority boundary, not a requirement to represent all cross-source evidence inside one physical source-bound \texttt{ReviewDocument}.

\begin{center}\rule{0.5\linewidth}{0.5pt}\end{center}

\hypertarget{6a-project-level-engineering-authority-reconciliation}{%
\paragraph{6a. Project-Level Engineering Authority Reconciliation}\label{6a-project-level-engineering-authority-reconciliation}}

The existing Human Review, Approved Input and Approved Engineering Information contracts remain source-bound and shall not be migrated solely to support Multi-Source processing.

Cross-source semantic evidence shall therefore not be forced into a legacy source-bound \texttt{ReviewDocument}, and no artificial primary Source shall be assigned to engineering information derived from several independent Sources.

Existing \texttt{stable\_subject\_key} values shall not be interpreted as project-wide cross-source semantic identity unless such continuity has been explicitly established. Equality or inequality of source-local \texttt{stable\_subject\_key} values alone does not establish cross-source Engineering Authority continuity.

After source-local Human Review and Approved Input / Approved Engineering Information projection, a new project-level boundary is introduced:

\begin{verbatim}
Source-local Approved Engineering Authority
        +
Cross-Source Semantic Reconciliation Evidence
        ↓
Project Engineering Authority Reconciliation
        ↓
Explicit Human Authority Decision
        ↓
Project Engineering Authority State
\end{verbatim}

Project Engineering Authority Reconciliation references existing immutable source-local Approved Inputs / Approved Engineering Information and the exact Cross-Source Semantic Reconciliation evidence. It does not rewrite those artifacts.

The Human reviewer may establish whether reviewed engineering information shall:

\begin{verbatim}
remain independent
coexist as valid information concerning the same engineering concern
supersede previously accepted project-level authority
remain unresolved
\end{verbatim}

A project-level authority concern or continuity identity may only be established through this explicit Human decision boundary. Machine-generated semantic relationships such as \texttt{equivalent}, \texttt{complementary}, or \texttt{potential\_conflict} do not create that identity and do not authorize an authority transition.

Source-local Approved Input lifecycle state and project-level Engineering Authority state are separate concepts. Cross-source reconciliation may therefore determine that an otherwise valid source-local Approved Input is no longer active project-level authority without modifying or deleting its immutable source-local authority record.

Where the existing Approved Input supersession contract can represent an accepted successor without falsifying provenance or continuity, it may be reused. Where its source-local \texttt{stable\_subject\_key} contract is insufficient for genuine cross-source continuity, project-level authority reconciliation remains authoritative and shall not fabricate a matching \texttt{stable\_subject\_key}.

If Human Review determines that the accepted engineering meaning requires a newly synthesized or materially changed engineering statement that is not represented by an existing Approved Input, the reconciliation layer shall not author that statement automatically. The new statement must pass an explicit Human Engineering Review and become new reviewed Engineering Authority before it may affect the model.

The resulting Project Engineering Authority State is the authoritative input to Model Impact Reconciliation in S5.

The phrase "one common Human Engineering Review boundary" denotes a common Human-authority boundary at project level; it does not require all cross-source evidence to be represented by one physical \texttt{ReviewDocument}.

\hypertarget{7-existing-reviewed-engineering-information-remains-authority-when-new-sources-arrive}{%
\paragraph{7. Existing reviewed Engineering Information remains authority when new sources arrive}\label{7-existing-reviewed-engineering-information-remains-authority-when-new-sources-arrive}}

When additional sources are registered after engineering information has already been reviewed and modeled:

\begin{verbatim}
existing active Approved Inputs / AEI
= current Engineering Authority

new source
= new evidence / change candidate
\end{verbatim}

A newer source does not automatically supersede existing authority.

Cross-source reconciliation may identify semantic overlap or material variance between new source-derived information and existing active reviewed engineering information.

Human Engineering Review determines whether the accepted authority shall:

\begin{verbatim}
remain unchanged
be extended
be replaced/superseded
coexist as separately valid information
remain unresolved
\end{verbatim}

Existing Approved Input successor and supersession mechanisms shall be reused wherever compatible.

A changed reviewed successor creates new authority; previous authority is retained immutably and may become \texttt{superseded}.

No existing reviewed evidence is rewritten.

\begin{center}\rule{0.5\linewidth}{0.5pt}\end{center}

\hypertarget{8-separate-engineering-authority-from-model-authority}{%
\paragraph{8. Separate Engineering Authority from Model Authority}\label{8-separate-engineering-authority-from-model-authority}}

The existing engineering model is not the primary authority for deciding whether new source information is correct.

Authority order remains:

\begin{verbatim}
Human-reviewed Engineering Information
        ↓
Approved Input / AEI
        ↓
Accepted Engineering Model
        ↓
SysML v2 Representation
\end{verbatim}

The model represents accepted Engineering Authority.

It does not override new evidence merely because the new evidence differs from the current model.

\begin{center}\rule{0.5\linewidth}{0.5pt}\end{center}

\hypertarget{9-introduce-model-impact-reconciliation-after-engineering-review}{%
\paragraph{9. Introduce Model Impact Reconciliation after Engineering Review}\label{9-introduce-model-impact-reconciliation-after-engineering-review}}

Only after new Engineering Authority has been established may the system compare it with the currently accepted model.

A separate Model Impact Reconciliation boundary shall evaluate the impact of accepted engineering changes on existing model elements.

It may propose:

\begin{verbatim}
retain
extend
modify
new
supersede
unresolved
\end{verbatim}

Potential model impact shall be derived using:

\begin{itemize}
\tightlist
\item
  active Engineering Authority;
\item
  existing model traceability;
\item
  model classification and placement information;
\item
  existing accepted model structure.
\end{itemize}

Model Impact Reconciliation is advisory evidence.

It shall not directly mutate the accepted model.

Human Model Review remains authoritative for model changes.

This separation prevents semantically related statements on different engineering abstraction levels from being falsely treated as competing model elements during source processing.

\begin{center}\rule{0.5\linewidth}{0.5pt}\end{center}

\hypertarget{10-preserve-explicit-product-evolution-and-change-control}{%
\paragraph{10. Preserve explicit product evolution and Change Control}\label{10-preserve-explicit-product-evolution-and-change-control}}

When Human Review establishes that newly processed information replaces previously accepted engineering information:

\begin{verbatim}
old reviewed authority
        ↓
superseded by
        ↓
new reviewed authority
\end{verbatim}

the change shall remain explicitly traceable.

Corresponding model changes shall be represented as a subsequent accepted model state rather than rewriting historical authority or historical model evidence.

The architecture therefore supports product evolution without using document age as automatic truth.

\begin{center}\rule{0.5\linewidth}{0.5pt}\end{center}

\hypertarget{consequences-15}{%
\subsubsection{Consequences}\label{consequences-15}}

\hypertarget{positive-6}{%
\paragraph{Positive}\label{positive-6}}

\begin{itemize}
\tightlist
\item
  Genuine Multi-Source processing becomes possible without rewriting source-level Processing.
\item
  Sources accidentally registered in the wrong Project can be detected before project-level semantic integration.
\item
  Context-only documents obtain a clear architectural role in establishing product understanding.
\item
  Complementary information is preserved instead of being falsely treated as conflict.
\item
  Cross-source conflicts remain explicit and provenance-preserving.
\item
  Existing reviewed information remains authoritative until Human-controlled change.
\item
  Existing Approved Input supersession mechanisms can support auditable Change Control.
\item
  Model updates become consequences of reviewed engineering change rather than direct consequences of LLM output.
\item
  Existing Single-Source behavior remains valid.
\item
  No migration of persisted \texttt{ProcessingArtifactReference} contracts is required.
\end{itemize}

\hypertarget{negative--cost-1}{%
\paragraph{Negative / Cost}\label{negative--cost-1}}

\begin{itemize}
\tightlist
\item
  Multi-Source processing requires additional LLM-assisted reasoning boundaries.
\item
  A project-level reconciliation contract and Human Review bridge must be introduced.
\item
  Incremental source ingestion must consider existing active Engineering Authority.
\item
  Model impact analysis adds an additional reasoning step before model update.
\item
  More explicit provenance and lifecycle state must be retained across boundaries.
\end{itemize}

These costs are accepted because collapsing Source Fit, semantic reconciliation, Engineering Authority and Model Change into one step would create unsafe hidden authority and unreliable model evolution.

\begin{center}\rule{0.5\linewidth}{0.5pt}\end{center}

\hypertarget{affected-components}{%
\subsubsection{Affected Components}\label{affected-components}}

Existing components retained:

\begin{verbatim}
modules/project_sources/
modules/project_processing/
modules/source_evidence/
modules/engineering_subjects/
modules/subject_interpretation/
modules/subject_consensus/
modules/human_review/
modules/review_workspace/
modules/approved_input/
Approved Engineering Information projection
Model Candidate / Placement processing
Internal Engineering Model
SysML v2 generation
SYSIDE validation
Final Human Release
\end{verbatim}

New or extended responsibilities are expected around:

\begin{verbatim}
project-fit assessment
project-level cross-source semantic reconciliation
multi-source Human Review projection / bridge
engineering-authority continuity
model-impact reconciliation
\end{verbatim}

Existing immutable Project \texttt{000116} Gate-3 evidence shall not be modified or regenerated by this work.

\begin{center}\rule{0.5\linewidth}{0.5pt}\end{center}

\hypertarget{supersedes}{%
\subsubsection{Supersedes}\label{supersedes}}

None.

This ADR extends, but does not supersede:

\begin{verbatim}
ADR-012 Processing State and Artifact Organization
ADR-016 Human Review Workspace and Approved Input Promotion Architecture
ADR-025 Semantic Proposal Consolidation and Persona-Aware Consensus
ADR-026 Source-Anchored Multi-Persona Interpretation and Cross-Unit Semantic Synthesis
ADR-029 Human-Reviewed Model Placement Before Model Assembly
ADR-023 Final Model Review and Output Publication Architecture
\end{verbatim}

Where this ADR introduces project-level behavior, the existing source-level contracts remain valid.

\begin{center}\rule{0.5\linewidth}{0.5pt}\end{center}

\hypertarget{related-roadmap-phase}{%
\subsubsection{Related Roadmap Phase}\label{related-roadmap-phase}}

\begin{verbatim}
BLK-002 — Cross-Source Processing Artifact Identity Collision
WP-12 Multi-Source Acceptance
30-Day Thesis Completion Track A
Multi-Source → Safe Demo
\end{verbatim}

\begin{center}\rule{0.5\linewidth}{0.5pt}\end{center}

\hypertarget{related-implementation}{%
\subsubsection{Related Implementation}\label{related-implementation}}

Implementation is divided into bounded, independently testable slices.

\hypertarget{s1--contextual-processing-artifact-identity}{%
\paragraph{S1 --- Contextual Processing Artifact Identity}\label{s1--contextual-processing-artifact-identity}}

At project-wide lifecycle and aggregation boundaries:

\begin{verbatim}
(processing_run_id, artifact_type, artifact_id)
\end{verbatim}

becomes effective identity.

No persisted Processing Artifact schema change.

\hypertarget{s2--project-fit--source-admissibility}{%
\paragraph{S2 --- Project Fit / Source Admissibility}\label{s2--project-fit--source-admissibility}}

Introduce the Project Fit Assessment and its Human escalation boundary.

\hypertarget{s3--project-level-cross-source-semantic-reconciliation}{%
\paragraph{S3 --- Project-Level Cross-Source Semantic Reconciliation}\label{s3--project-level-cross-source-semantic-reconciliation}}

Introduce provenance-preserving:

\begin{verbatim}
equivalent
complementary
potential_conflict
distinct
uncertain
\end{verbatim}

classification.

\hypertarget{s4--project-level-engineering-authority-reconciliation}{%
\paragraph{S4 --- Project-Level Engineering Authority Reconciliation}\label{s4--project-level-engineering-authority-reconciliation}}

Connect source-local Approved Input / AEI authority and Cross-Source Semantic Reconciliation evidence to the explicit project-level Human Authority Decision defined by Sections 6 and 6a, without migrating or falsifying source-bound Review / Approved Input / AEI provenance.

\hypertarget{s5--model-impact-reconciliation}{%
\paragraph{S5 --- Model Impact Reconciliation}\label{s5--model-impact-reconciliation}}

Compare newly accepted Engineering Authority with the existing accepted model and produce reviewable model-impact proposals.

Each implementation slice requires focused tests.

Before final acceptance:

\begin{verbatim}
focused slice tests
full repository regression
git diff --check
real heterogeneous Multi-Source E2E
downstream deterministic SysML v2 generation
real SYSIDE validation
Final Human Release
SSOT update
\end{verbatim}

\textbackslash newpage

\hypertarget{adr-033--concern-centric-project-reconciliation-and-coherent-model-handoff}{%
\subsection{ADR-033 --- Concern-Centric Project Reconciliation and Coherent Model Handoff}\label{adr-033--concern-centric-project-reconciliation-and-coherent-model-handoff}}

\begin{itemize}
\tightlist
\item
  \textbf{Status:} Accepted
\item
  \textbf{Date:} 2026-08-31
\item
  \textbf{Decision scope:} BLK-002 project-level multi-source reconciliation
\item
  \textbf{Refines:} ADR-032
\item
  \textbf{Replaces:} pair-centric S3 orchestration introduced experimentally by I2D.4
\item
  \textbf{Does not replace:} source-local Human Review / Approved Input authority, S4 Human Project Authority, S5 Model Impact authority boundaries, or Human Final Model Review
\end{itemize}

\hypertarget{context-19}{%
\subsubsection{Context}\label{context-19}}

BLK-002 must reconcile multiple independent Engineering Sources without collapsing their provenance or allowing an LLM to become Engineering Authority.

The first S3 implementation asked one LLM call to reason across all supplied Subjects and to construct arbitrary cross-source relations. A later experiment constrained each call to one Source pair. Both approaches remained relation-centric: the LLM still had to choose which Subjects should be compared.

Real E2E execution showed that this is the wrong abstraction. The engineering question is not primarily "which pair of Subjects relates to which other pair?". It is:

\begin{enumerate}
\def\labelenumi{\arabic{enumi}.}
\tightlist
\item
  Do the Sources belong to the same Project/System context?
\item
  Which source-local Subjects concern the same engineering topic?
\item
  For each such topic, are the statements compatible, complementary, conflicting, distinct, or uncertain?
\item
  Which engineering meaning is accepted by Human Project Authority?
\item
  How does the complete accepted project-level meaning become one coherent Model Candidate Set?
\end{enumerate}

The useful analogy is a collection of stickers: Project Fit determines whether a sticker belongs to the same album/tournament; semantic indexing groups stickers by the number/topic they concern; only then are all stickers of that topic assessed together.

\hypertarget{decision-19}{%
\subsubsection{Decision}\label{decision-19}}

Project reconciliation is \textbf{concern-centric, not pair-centric}.

\hypertarget{s2--project-fit}{%
\paragraph{S2 --- Project Fit}\label{s2--project-fit}}

S2 remains unchanged.

Its question is whether each Engineering Source is plausibly in scope for the current Project/System. Only Sources with an admitted Project Fit gate state may enter project-level semantic reconciliation.

S2 answers the "same tournament/album?" question. It does not establish semantic identity between Subjects and does not create Engineering Authority.

\hypertarget{s3a--global-project-semantic-index}{%
\paragraph{S3A --- Global Project Semantic Index}\label{s3a--global-project-semantic-index}}

S3A receives all admitted, source-bound Project semantic Subjects in one bounded semantic-indexing task.

Its single task is:

\begin{quote}
Group supplied Subjects by shared engineering concern.
\end{quote}

S3A does \textbf{not}:

\begin{itemize}
\tightlist
\item
  decide which Source is correct;
\item
  decide conflict or compatibility;
\item
  merge Subjects;
\item
  create model elements;
\item
  create Engineering Authority;
\item
  rank Sources;
\item
  create persisted Case identity.
\end{itemize}

The LLM sees only transient opaque Subject transport identifiers such as \texttt{SUBJ-0001}. Real \texttt{project\_subject:...} identity remains application-controlled.

The deterministic application layer validates:

\begin{itemize}
\tightlist
\item
  every supplied Subject appears exactly once;
\item
  no unknown Subject reference is accepted;
\item
  no Subject appears in more than one group;
\item
  no Subject is silently dropped;
\item
  a multi-member Case contains Subjects from at least two different Sources;
\item
  Subjects without a cross-source counterpart become Singleton Cases.
\end{itemize}

After successful validation, application code deterministically orders Cases and assigns \texttt{CASE-000001}, \texttt{CASE-000002}, ... . LLM-generated labels are descriptive only and never Case identity.

Case identity is fingerprint-bound to:

\begin{itemize}
\tightlist
\item
  Project ID;
\item
  exact global S3 input fingerprint;
\item
  sorted real member \texttt{project\_subject:...} references.
\end{itemize}

\hypertarget{s3b--reconciliation-case-assessment}{%
\paragraph{S3B --- Reconciliation Case Assessment}\label{s3b--reconciliation-case-assessment}}

S3B processes one non-singleton Reconciliation Case per LLM call.

Its single task is:

\begin{quote}
Assess how all statements belonging to this one engineering concern fit together.
\end{quote}

Allowed outcomes are:

\begin{itemize}
\tightlist
\item
  \texttt{equivalent}
\item
  \texttt{complementary}
\item
  \texttt{potential\_conflict}
\item
  \texttt{distinct}
\item
  \texttt{uncertain}
\end{itemize}

\texttt{unique} is reserved for Singleton Cases and is derived deterministically without an LLM call.

S3B assesses the Case \textbf{as a whole}. It does not expand the Case into a graph of pairwise relations.

For example, if three Sources state viewer limits of 2, 2 and 5, the result is one Case with two claim groups and \texttt{potential\_conflict}, not three pairwise relations.

Claim groups are non-authoritative evidence describing materially different variants within one Case. For \texttt{potential\_conflict}, claim groups must partition all Case members into at least two explicit competing groups.

S3B remains relationship/meaning evidence only. Human Project Authority remains required for every non-singleton Case.

\hypertarget{deterministic-project-reconciliation-summary}{%
\paragraph{Deterministic Project Reconciliation Summary}\label{deterministic-project-reconciliation-summary}}

After every Case has an assessment, application code derives the global Project Reconciliation summary without an additional LLM call.

The summary may state:

\begin{itemize}
\tightlist
\item
  number of Cases by outcome;
\item
  whether potential cross-source conflicts were detected;
\item
  whether uncertainty exists;
\item
  whether regrouping is required because a Case was assessed as \texttt{distinct};
\item
  whether Human Project Authority is required.
\end{itemize}

The system must say "no cross-source conflict detected" rather than claiming the Project is objectively contradiction-free.

\hypertarget{s4--human-project-authority}{%
\paragraph{S4 --- Human Project Authority}\label{s4--human-project-authority}}

S4 operates on Reconciliation Cases, not synthetic pairwise relations.

S4 may establish that source-bound engineering meaning:

\begin{itemize}
\tightlist
\item
  remains independently active;
\item
  coexists with other source-bound meaning;
\item
  is superseded by explicitly selected source-bound meaning;
\item
  remains unresolved.
\end{itemize}

No synthetic merged Approved Engineering Input is created.

Singleton Cases retain their existing source-local authority and do not require an additional semantic LLM assessment.

\hypertarget{project-engineering-state-and-coherent-downstream-model-handoff}{%
\paragraph{Project Engineering State and coherent downstream model handoff}\label{project-engineering-state-and-coherent-downstream-model-handoff}}

The end product of project reconciliation is \textbf{not a list of isolated Cases}.

Resolved Cases, unique retained Subjects, exact source-local authority bindings, and explicit Human Project decisions together form the Project Engineering State.

That state is handed to the downstream model derivation pipeline as one coherent project-level engineering handoff.

The downstream model derivation outcome is one coherent \textbf{Model Candidate Set} which may contain, together:

\begin{itemize}
\tightlist
\item
  model elements;
\item
  relationships;
\item
  properties / attributes;
\item
  constraints;
\item
  allocations / usages;
\item
  other required structural model constructs.
\end{itemize}

Cases are not written into the model one by one.

S3/S4 establish approved engineering meaning. The model-candidate / Phase-H layer decides how that meaning is represented structurally in SysML v2.

The complete Model Candidate Set is then compared with the accepted model state, processed through Model Impact / model proposal logic, reviewed by Human Model Authority, and only then released to the model.

\hypertarget{llm-design-principle}{%
\paragraph{LLM design principle}\label{llm-design-principle}}

For project reconciliation:

\begin{quote}
LLMs perform bounded semantic judgments. Deterministic application logic owns identity, coverage, orchestration, aggregation, and authority boundaries.
\end{quote}

And:

\begin{quote}
A singular LLM task does not necessarily mean a tiny input. S3A may see the complete bounded Subject set because its task is only semantic indexing. S3B sees one Case because its task is only Case assessment.
\end{quote}

If the bounded S3A input is ever too large, the system must fail closed until an explicit indexing reduction/hierarchy architecture is accepted. It must not silently chunk and infer global grouping equivalence.

\hypertarget{consequences-16}{%
\subsubsection{Consequences}\label{consequences-16}}

\hypertarget{positive-7}{%
\paragraph{Positive}\label{positive-7}}

\begin{itemize}
\tightlist
\item
  Reconciliation matches the actual engineering question rather than creating O(n²) relation evidence.
\item
  Conflicts are represented once per engineering concern.
\item
  Human review sees coherent topics and variants rather than redundant pairwise edges.
\item
  Source provenance remains intact.
\item
  The LLM no longer controls identity or pair construction.
\item
  The downstream model receives one coherent approved project meaning rather than Case-by-Case model mutations.
\end{itemize}

\hypertarget{trade-offs-8}{%
\paragraph{Trade-offs}\label{trade-offs-8}}

\begin{itemize}
\tightlist
\item
  S3A is a global semantic indexing task and therefore requires an explicit bounded-input contract.
\item
  Existing S4 relation-oriented implementation must be adapted to Case-oriented authority in a later slice.
\item
  Existing pairwise I2D.4 runtime code must be removed/replaced before the next live Project Reconciliation execution.
\item
  Persisted project-reconciliation schema changes, if required, must be additive and separately accepted.
\end{itemize}

\hypertarget{migration--implementation-plan}{%
\subsubsection{Migration / implementation plan}\label{migration--implementation-plan}}

\begin{enumerate}
\def\labelenumi{\arabic{enumi}.}
\tightlist
\item
  \textbf{I2D.5A} --- freeze ADR-033 and additive in-memory Case contracts.
\item
  \textbf{I2D.5B} --- implement S3A global semantic indexing with transient Subject refs.
\item
  \textbf{I2D.5C} --- implement S3B one-Case assessment and deterministic project summary.
\item
  \textbf{I2D.5D} --- replace I2D.4 runtime orchestration and define additive persistence.
\item
  \textbf{I2D.5E} --- adapt S4 Human Project Authority from relation decisions to Case decisions.
\item
  Reconnect S5 / I1 project-level handoff so all resolved/unique meaning enters one coherent Model Candidate Set.
\item
  Resume Project 308131 E2E only after the pairwise runtime path has been removed and the Case path is green.
\end{enumerate}

No immutable accepted evidence is rewritten by this ADR.

\textbackslash newpage

\hypertarget{adr-034--source-provenance-does-not-constrain-concern-grouping}{%
\subsection{ADR-034 --- Source Provenance Does Not Constrain Concern Grouping}\label{adr-034--source-provenance-does-not-constrain-concern-grouping}}

\textbf{Status:} Accepted refinement of ADR-033 \textbf{Scope:} BLK-002 / I2D.5D4

\hypertarget{context-20}{%
\subsubsection{Context}\label{context-20}}

ADR-033 replaced pair-centric cross-source reconciliation with concern-centric Project Reconciliation Cases. During the first live Project 308131 run, S3A returned a multi-member semantic group whose Subjects all originated from one Source. The existing Case contract rejected that group because a multi-member Case was required to span at least two Sources.

That restriction carries forward source-pair semantics into a concern-centric architecture. It is especially inappropriate for well-separated engineering documents, where several Subjects from one document may legitimately address the same engineering concern.

The source of a Subject remains important for provenance, traceability, source- local Human Authority, and later project-level authority. It is not semantic evidence that determines whether Subjects concern the same topic.

\hypertarget{decision-20}{%
\subsubsection{Decision}\label{decision-20}}

Reconciliation Case membership is determined solely by shared engineering concern.

\texttt{source\_id} is provenance only and SHALL NOT constrain semantic concern grouping.

Therefore:

\begin{itemize}
\tightlist
\item
  A one-member group becomes a Singleton Case.
\item
  A multi-member Case MAY contain Subjects from one Source or multiple Sources.
\item
  \texttt{unique} is defined by one Subject, not by one Source.
\item
  Every real source binding remains preserved on every Subject and Case.
\item
  Same-source multi-member Cases remain non-singleton Cases and therefore receive the normal S3B bounded semantic assessment.
\item
  No automatic cross-group merging, source balancing, or source-based splitting is permitted.
\item
  Exact Subject coverage, no duplicates, known references, deterministic Case identity, and immutable provenance remain fail-closed constraints.
\end{itemize}

\hypertarget{rationale}{%
\subsubsection{Rationale}\label{rationale}}

The concern-centric architecture separates semantics from provenance:

\begin{quote}
Semantic meaning decides grouping. Source provenance records where evidence came from.
\end{quote}

This is analogous to grouping Panini stickers by motif regardless of the kiosk where a packet was purchased. The purchase location is traceability metadata, not part of the motif identity.

\hypertarget{consequences-17}{%
\subsubsection{Consequences}\label{consequences-17}}

The former invariant

\texttt{multi-member\ Case\ =\textgreater{}\ at\ least\ two\ Sources}

is removed from the S3A Case contract and persistence validation.

Case identity remains unchanged and continues to bind:

\begin{itemize}
\tightlist
\item
  project identity,
\item
  exact semantic input fingerprint,
\item
  exact sorted Subject references.
\end{itemize}

No persisted PRC cycle schema changes. Existing semantic-index schema 1.0.0 remains readable, and provenance-bearing semantic-index schema 1.1.0 remains the live S3A format.

Human Project Authority remains Case-aware and is still a later gate. This ADR does not create authority and does not alter source-local Approved Engineering Input.

\textbackslash newpage

\hypertarget{adr-035--safe-demo-review-compression-and-final-review-routing}{%
\subsection{ADR-035 --- Safe-Demo Review Compression and Final Review Routing}\label{adr-035--safe-demo-review-compression-and-final-review-routing}}

\hypertarget{status-12}{%
\subsubsection{Status}\label{status-12}}

Accepted for v0.5.0 feature implementation. Integration into \texttt{main} remains gated by Safe-Demo acceptance.

\hypertarget{date-12}{%
\subsubsection{Date}\label{date-12}}

2026-09-10

\hypertarget{context-21}{%
\subsubsection{Context}\label{context-21}}

The v0.4.0 application exposes several technically distinct authority and transformation steps in one long Model Proposal / Internal Model working surface.

A 2026-09-10 informal professor demo showed that the late-stage workflow is functionally usable but unnecessarily time-consuming in Focused View.

The most visible problem was Target-Model Formulation: deterministic, single-candidate formulation items were presented as many individual Human review cards with individual rationale entry.

At the same time, the system already has a separate top-level Phase-L \texttt{Final\ Model\ Review} workspace after generated SysML v2 and validation.

The earlier whole-assembly gate inside \texttt{app/model\_final\_review\_ui.py} was also user-facing as \texttt{Final\ Model\ Review}, creating naming ambiguity already tracked by \texttt{OBS-032}.

ADR-024 requires the engineer-facing UI to be task-oriented, simple by default and fully traceable underneath. This ADR refines that UX direction for v0.5.0 without redesigning the underlying authority schemas.

\hypertarget{decision-21}{%
\subsubsection{Decision}\label{decision-21}}

\hypertarget{d1--keep-the-real-phase-l-final-model-review}{%
\paragraph{D1 --- Keep the real Phase-L Final Model Review}\label{d1--keep-the-real-phase-l-final-model-review}}

The top-level Final Model Review after generated SysML v2 and validation remains the final whole-model engineering review before release/publication authority.

It is not removed or replaced.

\hypertarget{d2--rename-the-pre-generation-assembly-gate}{%
\paragraph{D2 --- Rename the pre-generation assembly gate}\label{d2--rename-the-pre-generation-assembly-gate}}

The earlier whole-assembly review is user-facing as:

\begin{verbatim}
Model Assembly Review
\end{verbatim}

It remains responsible for:

\begin{itemize}
\tightlist
\item
  reviewing the assembled model;
\item
  resolving remaining target Relationship representation where necessary;
\item
  approving materialization of the authority-backed Internal Model.
\end{itemize}

\hypertarget{d3--compress-deterministic-target-model-formulation-in-focused-view}{%
\paragraph{D3 --- Compress deterministic Target-Model Formulation in Focused View}\label{d3--compress-deterministic-target-model-formulation-in-focused-view}}

The bounded current Target-Model Formulation bridge produces one candidate per review item.

Focused View shall not require an engineer to open and accept every deterministic single-candidate mapping separately.

One action may batch-confirm all pending single-candidate items that are not:

\begin{verbatim}
unresolved_human_review
\end{verbatim}

Existing persisted decisions are resumed.

\texttt{unresolved\_human\_review} remains explicit Human attention and blocks authority finalization.

\texttt{intentionally\_not\_materialized} may be included in the batch only when the Focused UI explicitly tells the engineer that the approved engineering information remains authoritative but is not formally materialized under the current bounded notation.

\hypertarget{d4--preserve-existing-formulation-authority-contracts-for-v050}{%
\paragraph{D4 --- Preserve existing formulation authority contracts for v0.5.0}\label{d4--preserve-existing-formulation-authority-contracts-for-v050}}

v0.5.0 does not remove TFA/TFD persistence or redesign authority schemas.

Batch confirmation delegates to the existing exact write service and persists one exact decision per item with the same immutable binding and fingerprints.

This is an interaction compression, not hidden weakening of authority.

A future architecture revision may reassess whether deterministic mappings require Human authority at all, but that is outside the Safe-Demo slice.

\hypertarget{d5--batch-only-clean-llm-model-quality-proposals}{%
\paragraph{D5 --- Batch only clean LLM Model Quality proposals}\label{d5--batch-only-clean-llm-model-quality-proposals}}

Model Quality refinement remains LLM-assisted and therefore retains Human authority.

Focused View may offer:

\begin{verbatim}
Accept all clean refinements
\end{verbatim}

only for pending proposals satisfying all of:

\begin{verbatim}
meaning_preserved == true
unsupported_information_added == false
requires_human_attention == false
\end{verbatim}

All other proposals remain individual Human review.

Modify and reject remain explicit rationale-backed actions.

Technical View continues to expose the complete proposal and decision evidence.

\hypertarget{d6--route-directly-to-the-real-final-model-review}{%
\paragraph{D6 --- Route directly to the real Final Model Review}\label{d6--route-directly-to-the-real-final-model-review}}

After authority-backed SysML v2 is generated, Focused View presents a compact artifact summary and an explicit:

\begin{verbatim}
Go to Final Model Review
\end{verbatim}

action.

That action routes to the existing top-level Final Model Review workspace.

Generated code remains fully inspectable in Technical View and again in Final Model Review.

\hypertarget{d7--safe-demo-integration-gate}{%
\paragraph{D7 --- Safe-Demo integration gate}\label{d7--safe-demo-integration-gate}}

v0.4.0 remains the functional fallback.

v0.5.0 may be merged into \texttt{main} only after:

\begin{verbatim}
focused tests
→ appropriate broader regression
→ complete repository regression
→ fresh E2E project
→ generated SysML v2
→ SYSIDE validation
→ Phase-L Final Model Review
→ release/publication path
→ exact Safe-Demo rehearsal
→ explicit Human acceptance
\end{verbatim}

No new unrelated feature work is bundled into this slice.

\hypertarget{consequences-18}{%
\subsubsection{Consequences}\label{consequences-18}}

Positive:

\begin{itemize}
\tightlist
\item
  substantially fewer low-value clicks in the demo and normal Focused workflow;
\item
  clearer separation of deterministic transformation and Human engineering review;
\item
  clear naming between Model Assembly Review and Phase-L Final Model Review;
\item
  Human attention remains concentrated on ambiguity, LLM-risk and whole-model authority;
\item
  existing traceability and authority contracts remain intact;
\item
  v0.4.0 rollback remains available.
\end{itemize}

Trade-offs:

\begin{itemize}
\tightlist
\item
  the current persisted authority model still records item-level formulation decisions even though the UI performs one batch confirmation;
\item
  batch acceptance can persist several decisions before a later item fails, so the operation is resumable rather than transactionally all-or-nothing;
\item
  the long-term question whether deterministic mappings need Human authority is deliberately deferred;
\item
  real-data robustness and synthesis-relevance findings are not solved by this ADR.
\end{itemize}
